% ============================================================
%  يُترجَم بـ XeLaTeX  (في Overleaf: Menu → Compiler → XeLaTeX)
%  الخط: Amiri (متوفر في Overleaf، أو حمّله من github.com/aliftype/amiri)
% ============================================================
\documentclass[12pt,a4paper]{article}
\usepackage[a4paper,margin=1.7cm,top=2.3cm,bottom=2.9cm,headheight=15pt,footskip=1.5cm]{geometry}
\usepackage{amsmath,amssymb}
\usepackage{icomma}
\usepackage[english,bidi=default,layout=lists]{babel}
\babelprovide[import,main]{arabic}
\babelfont{rm}[Scale=1.08]{Amiri}
\babelfont{sf}[Scale=1.08]{Amiri}
\usepackage[most]{tcolorbox}
\usepackage{tikz}
\usetikzlibrary{arrows.meta,calc,shadings}
\usepackage{array}
\usepackage{colortbl}
\usepackage{fancyhdr}
\usepackage{titlesec}
\usepackage{enumitem}
\usepackage{collcell}
\usepackage{needspace}
\usepackage{etoolbox}
\newcolumntype{C}{>{\collectcell\ar}c<{\endcollectcell}}
\pretocmd{\section}{\Needspace{15\baselineskip}}{}{}
\pretocmd{\subsection}{\Needspace{11\baselineskip}}{}{}

\definecolor{cDef}{HTML}{1F4E9C}
\definecolor{cThm}{HTML}{B3261E}
\definecolor{cRem}{HTML}{2E7D32}
\definecolor{cEx}{HTML}{C76B00}
\definecolor{cMeth}{HTML}{6A1B9A}
\definecolor{cSol}{HTML}{00695C}
\definecolor{cSec}{HTML}{1F3A68}

\titleformat{\section}{\Large\bfseries\color{cSec}}{\thesection.}{0.6em}{}[{\color{cSec!60}\titlerule[1pt]}]
\titleformat{\subsection}{\large\bfseries\color{cThm!85!black}}{\thesubsection.}{0.5em}{}
\titlespacing*{\section}{0pt}{14pt}{8pt}
\titlespacing*{\subsection}{0pt}{10pt}{5pt}
\setlength{\parindent}{0pt}
\setlength{\parskip}{4pt}
\setlist{itemsep=2pt,topsep=3pt}

\newtcolorbox{arbox}[3][]{enhanced,
  colback=#2!5,colframe=#2!80!black,arc=2mm,boxrule=0.8pt,
  left=6pt,right=6pt,top=8pt,bottom=5pt,
  fonttitle=\bfseries,coltitle=white,
  title={\foreignlanguage{arabic}{#3}},
  attach boxed title to top right={xshift=-5mm,yshift=-3mm},
  boxed title style={colback=#2!80!black,arc=1.5mm,boxrule=0pt},
  before upper={\selectlanguage{arabic}},
  before skip=8pt,after skip=8pt,breakable,#1}
\newenvironment{BOX}[2]{\begin{otherlanguage}{english}\begin{arbox}{#1}{#2}}{\end{arbox}\end{otherlanguage}}
\newenvironment{tarif}[1][تعريف]{\begin{BOX}{cDef}{#1}}{\end{BOX}}
\newenvironment{khassia}[1][خاصية]{\begin{BOX}{cThm}{#1}}{\end{BOX}}
\newenvironment{molahada}[1][ملاحظة]{\begin{BOX}{cRem}{#1}}{\end{BOX}}
\newenvironment{mithal}[1][مثال]{\begin{BOX}{cEx}{#1}}{\end{BOX}}
\newenvironment{tariqa}[1][طريقة]{\begin{BOX}{cMeth}{#1}}{\end{BOX}}

\newcounter{exo}
\newenvironment{tamrin}[1]{\stepcounter{exo}\begin{BOX}{cSec}{التمرين \theexo: #1}}{\end{BOX}}
\newenvironment{hal}{\Needspace{12\baselineskip}\begin{otherlanguage}{english}\begin{arbox}[breakable]{cSol}{حل التمرين \theexo}}{\end{arbox}\end{otherlanguage}}

\newenvironment{LTR}{\begin{otherlanguage}{english}}{\end{otherlanguage}}
\newcommand{\ar}[1]{\foreignlanguage{arabic}{#1}}

\renewcommand{\dots}{\mathinner{\ldotp\ldotp\ldotp}}
\renewcommand{\ldots}{\mathinner{\ldotp\ldotp\ldotp}}
\renewcommand{\cdots}{\mathinner{\cdotp\cdotp\cdotp}}
\newcommand{\N}{\mathbb{N}}
\newcommand{\Z}{\mathbb{Z}}
\newcommand{\D}{\mathbb{D}}
\newcommand{\Q}{\mathbb{Q}}
\newcommand{\R}{\mathbb{R}}
\newcommand{\pgcd}{\operatorname{PGCD}}
\newcommand{\ppcm}{\operatorname{PPCM}}

% \NLseg{a}{b}{[}{]} ; \NLright{a}{[} ; \NLleft{b}{]}  (interval on a number line)
\newcommand{\NLseg}[4]{\begin{LTR}\begin{tikzpicture}[baseline=-0.6ex]
 \draw[-{Stealth[length=2mm]},gray!70](0,0)--(3.4,0);
 \draw[line width=2pt,cThm!80!black](0.7,0)--(2.7,0);
 \node[cThm!80!black,font=\LARGE] at (0.7,0){$#3$};
 \node[cThm!80!black,font=\LARGE] at (2.7,0){$#4$};
 \node[below=3pt,font=\footnotesize] at (0.7,0){$#1$};
 \node[below=3pt,font=\footnotesize] at (2.7,0){$#2$};
\end{tikzpicture}\end{LTR}}
\newcommand{\NLright}[2]{\begin{LTR}\begin{tikzpicture}[baseline=-0.6ex]
 \draw[-{Stealth[length=2mm]},gray!70](0,0)--(3.4,0);
 \draw[line width=2pt,cThm!80!black](1.3,0)--(3.3,0);
 \node[cThm!80!black,font=\LARGE] at (1.3,0){$#2$};
 \node[below=3pt,font=\footnotesize] at (1.3,0){$#1$};
\end{tikzpicture}\end{LTR}}
\newcommand{\NLleft}[2]{\begin{LTR}\begin{tikzpicture}[baseline=-0.6ex]
 \draw[-{Stealth[length=2mm]},gray!70](0,0)--(3.4,0);
 \draw[line width=2pt,cThm!80!black](0.05,0)--(2.1,0);
 \node[cThm!80!black,font=\LARGE] at (2.1,0){$#2$};
 \node[below=3pt,font=\footnotesize] at (2.1,0){$#1$};
\end{tikzpicture}\end{LTR}}

% ---------- social-media icons (pure TikZ) ----------
\definecolor{cFB}{HTML}{1877F2}
\definecolor{cTG}{HTML}{229ED9}
\definecolor{cIG1}{HTML}{FEDA75}
\definecolor{cIG2}{HTML}{D62976}
\definecolor{cIG3}{HTML}{4F5BD5}
\newcommand{\fbicon}{\begin{LTR}\begin{tikzpicture}[baseline=-0.75ex,x=1em,y=1em]
 \fill[cFB] (0,0) circle (0.55);
 \node[white,inner sep=0,font=\usefont{TU}{lmss}{bx}{n}\small] at (0.04,-0.02) {f};
\end{tikzpicture}\end{LTR}}
\newcommand{\igicon}{\begin{LTR}\begin{tikzpicture}[baseline=-0.75ex,x=1em,y=1em]
 \shade[lower left=cIG1,lower right=cIG2,upper left=cIG2,upper right=cIG3,rounded corners=0.28em] (-0.55,-0.55) rectangle (0.55,0.55);
 \draw[white,line width=0.09em,rounded corners=0.18em] (-0.36,-0.36) rectangle (0.36,0.36);
 \draw[white,line width=0.09em] (0,0) circle (0.17);
 \fill[white] (0.22,0.22) circle (0.05);
\end{tikzpicture}\end{LTR}}
\newcommand{\tgicon}{\begin{LTR}\begin{tikzpicture}[baseline=-0.75ex,x=1em,y=1em]
 \fill[cTG] (0,0) circle (0.55);
 \fill[white] (-0.33,0.0) -- (0.32,0.27) -- (0.2,-0.3) -- (0.0,-0.15) -- (-0.08,-0.28) -- (-0.08,-0.08) -- cycle;
 \fill[cTG!70!black] (-0.08,-0.08) -- (0.25,0.2) -- (0.0,-0.15) -- (-0.08,-0.28) -- cycle;
\end{tikzpicture}\end{LTR}}
% page number: black box, white two-digit number, at the right of the social row
\newcommand{\PageBox}{\begin{LTR}\tikz[baseline=(p.base)]\node[fill=black,text=white,inner sep=2.5pt,minimum width=1.6em,minimum height=1.3em,font=\usefont{TU}{lmss}{bx}{n}\small](p){\ifnum\value{page}<10 0\fi\arabic{page}};\end{LTR}}
% visual left-to-right order: Telegram | Instagram | Facebook | page box
\newcommand{\SocialRow}{\begin{otherlanguage}{english}%
 \ar{\textbf{الأستاذ سوايسية هشام للرياضيات}}\ \tgicon\hfill
 \textbf{prof\_hicham\_math}\ \igicon\hfill
 \ar{\textbf{الأستاذ سوايسية هشام للرياضيات}}\ \fbicon\hspace{2.2em}\PageBox\end{otherlanguage}}

\pagestyle{fancy}
\fancyhf{}
\renewcommand{\headrulewidth}{0.6pt}
\renewcommand{\footrulewidth}{0.4pt}
% \SetHeader{right head: document title}{left head: level -- stream}
\newcommand{\SetHeader}[2]{%
 \fancyhead[C]{\parbox{\textwidth}{\ar{#1\hfill #2}}}%
 \fancyfoot[C]{\parbox{\textwidth}{\small\SocialRow}}}
\usepackage{tkz-tab}
\usepackage{multicol}
% variation / sign table (LTR)
\newenvironment{TAB}{\par\begin{center}\begin{LTR}\begin{tikzpicture}[scale=0.95,every node/.style={font=\small}]}{\end{tikzpicture}\end{LTR}\end{center}}
% grid + axes, clipped: \begin{GR}{xmin}{xmax}{ymin}{ymax}{unit cm}
\newenvironment{GR}[5]{\par\begin{center}\begin{LTR}\begin{tikzpicture}[x=#5cm,y=#5cm,font=\footnotesize]
 \draw[gray!25,very thin] (#1,#3) grid (#2,#4);
 \draw[-{Stealth[length=2mm]}] (#1,0)--(#2+0.3,0) node[below left]{$x$};
 \draw[-{Stealth[length=2mm]}] (0,#3)--(0,#4+0.3) node[below left]{$y$};
 \node[below left] at (0,0){$O$};
 \begin{scope}\clip (#1,#3) rectangle (#2,#4);}
{\end{scope}\end{tikzpicture}\end{LTR}\end{center}}
\newcommand{\Ticks}[2]{\foreach \t in {#1} \node[below,font=\scriptsize] at (\t,0){$\t$}; \foreach \t in {#2} \node[left,font=\scriptsize] at (0,\t){$\t$};}
\newcommand{\bx}[1]{\ensuremath{\colorbox{cDef!12}{$\displaystyle #1$}}}
\newcommand{\TwoCol}[4]{\par\begin{LTR}\noindent\begin{minipage}[t]{#1\textwidth}\selectlanguage{arabic}#2\end{minipage}\hfill\begin{minipage}[t]{#3\textwidth}\selectlanguage{arabic}#4\end{minipage}\end{LTR}\par}
\newcommand{\RL}[4]{\TwoCol{#3}{#4}{#1}{#2}}  % \RL{w1}{right content}{w2}{left content}
\newcommand{\ok}{$\textcolor{cRem}{\checkmark}$}
\newcommand{\limn}{\displaystyle\lim_{n\to+\infty}}
% sub-questions, each on its own line
\newenvironment{sq}{\begin{itemize}[leftmargin=2.4em,itemsep=3pt]}{\end{itemize}}
\newcommand{\qa}{\item[(أ)]}\newcommand{\qb}{\item[(ب)]}\newcommand{\qc}{\item[(جـ)]}\newcommand{\qd}{\item[(د)]}\newcommand{\qe}{\item[(هـ)]}
\newcommand{\pa}{\par\smallskip\textbf{(أ)} }\newcommand{\pb}{\par\smallskip\textbf{(ب)} }\newcommand{\pc}{\par\smallskip\textbf{(جـ)} }\newcommand{\pd}{\par\smallskip\textbf{(د)} }\newcommand{\pe}{\par\smallskip\textbf{(هـ)} }
\newcommand{\QQ}[1]{\par\smallskip\textbf{#1.} }
\newenvironment{tmr}{\stepcounter{exo}\begin{BOX}{cSec}{تمرين محلول \theexo}}{\end{BOX}}
\newcommand{\ABC}[3]{\\[2pt]\begin{tabular}{p{0.28\textwidth}p{0.28\textwidth}p{0.28\textwidth}}(جـ) \ #3 & (ب) \ #2 & (أ) \ #1\end{tabular}}
\pretocmd{\section}{\Needspace{24\baselineskip}}{}{}
\newtcbox{\bxtbox}{on line,colback=cDef!12,colframe=cDef!12,boxrule=0pt,arc=1pt,boxsep=0pt,left=2pt,right=2pt,top=1pt,bottom=1pt}
\newcommand{\bxt}[1]{\foreignlanguage{english}{\bxtbox{\ar{#1}}}}
\SetHeader{\small\bfseries\color{cSec}تمارين البكالوريا: الدالة اللوغاريتمية}{\small\color{cSec}السنة الثالثة ثانوي -- الشعب العلمية}
\begin{document}
\begin{LTR}
\begin{tcolorbox}[enhanced,colback=cSec,colframe=cSec,arc=3mm,halign=center,boxrule=0pt,top=12pt,bottom=12pt,before upper={\selectlanguage{arabic}}]
{\color{white}\fontsize{26}{32}\selectfont\bfseries تمارين البكالوريا في الدالة اللوغاريتمية}\\[8pt]
{\color{white}\large مع الحلول المفصلة خطوة بخطوة}\\[5pt]
{\color{white}\large علوم تجريبية -- تقني رياضي -- رياضيات}
\end{tcolorbox}
\end{LTR}

\begin{molahada}
حُذفت من هذه التمارين الأسئلة الخاصة بالدوال الأصلية وحساب المساحات (لم تُدرَس بعد). مسائل شعبة الرياضيات في آخر الوثيقة مسائل بمستوى البكالوريا.
\end{molahada}
\section{تمارين البكالوريا -- شعبة علوم تجريبية}

% ===================== التمرين 1 =====================
\begin{tamrin}{بكالوريا 2009 -- الموضوع الأول}
\textbf{I)} $h$ دالة عددية معرفة على $]-1;+\infty[$ كما يلي: $h(x)=x^2+2x+\ln(x+1)$.
\QQ{1} أحسب $\lim\limits_{x\to -1^{+}}h(x)$ و $\lim\limits_{x\to +\infty}h(x)$.
\QQ{2} بيّن أنه من أجل كل عدد حقيقي $x$ من المجال $]-1;+\infty[$: $h'(x)=\dfrac{1+2(x+1)^2}{x+1}$، واستنتج اتجاه تغيّر الدالة $h$ ثم أنجز جدول تغيّراتها.
\QQ{3} أحسب $h(0)$ واستنتج إشارة $h(x)$ حسب قيم $x$.

\textbf{II)} لتكن $f$ دالة عددية معرفة على المجال $]-1;+\infty[$ كما يلي: $f(x)=x-1-\dfrac{\ln(x+1)}{x+1}$.
نرمز بـ $(C_f)$ لتمثيلها البياني في المستوي المنسوب إلى المعلم المتعامد والمتجانس $(O;\vec i,\vec j)$.
\QQ{1}
\begin{sq}
\qa أحسب $\lim\limits_{x\to -1^{+}}f(x)$، ثم فسّر هذه النتيجة بيانيا.
\qb باستخدام النتيجة $\lim\limits_{t\to +\infty}\dfrac{e^{t}}{t}=+\infty$، برهن أن $\lim\limits_{u\to +\infty}\dfrac{\ln u}{u}=0$.
\qc استنتج $\lim\limits_{x\to +\infty}f(x)$.
\qd أحسب $\lim\limits_{x\to +\infty}\big[f(x)-(x-1)\big]$ واستنتج وجود مستقيم مقارب مائل للمنحنى $(C_f)$.
\qe أدرس وضعية المنحنى $(C_f)$ بالنسبة إلى المستقيم المقارب المائل.
\end{sq}
\QQ{2} بيّن أنه من أجل كل عدد حقيقي $x$ من المجال $]-1;+\infty[$: $f'(x)=\dfrac{h(x)}{(x+1)^2}$، ثم شكّل جدول تغيّرات الدالة $f$.
\QQ{3} بيّن أن المنحنى $(C_f)$ يقطع المستقيم ذا المعادلة $y=2$ عند النقطة التي فاصلتها محصورة بين $3{,}3$ و $3{,}4$.
\QQ{4} أرسم $(C_f)$.
\end{tamrin}

\begin{hal}
\textbf{الجزء الأول:}

\QQ{1} \textbf{حساب النهايتين:}
\begin{itemize}
\item عند $-1$ من اليمين: $\lim\limits_{x\to -1^{+}}(x^2+2x)=1-2=-1$، و $\lim\limits_{x\to -1^{+}}(x+1)=0^{+}$ إذن $\lim\limits_{x\to -1^{+}}\ln(x+1)=-\infty$، ومنه بالجمع:
\[\bx{\lim_{x\to -1^{+}}h(x)=-\infty}\]
\item عند $+\infty$: $\lim\limits_{x\to +\infty}(x^2+2x)=+\infty$ و $\lim\limits_{x\to +\infty}\ln(x+1)=+\infty$، ومنه:
\[\bx{\lim_{x\to +\infty}h(x)=+\infty}\]
\end{itemize}

\QQ{2} الدالة $h$ قابلة للاشتقاق على $]-1;+\infty[$ (مجموع دالة كثير حدود ودالة $x\mapsto \ln(x+1)$ حيث $x+1>0$)، ومن أجل كل $x>-1$:
\[h'(x)=2x+2+\frac{1}{x+1}=2(x+1)+\frac{1}{x+1}=\frac{2(x+1)^2+1}{x+1}\]
إذن: \bx{h'(x)=\dfrac{1+2(x+1)^2}{x+1}}.

\textbf{اتجاه التغيّر:} من أجل كل $x>-1$ لدينا $x+1>0$ و $1+2(x+1)^2>0$، إذن $h'(x)>0$، ومنه الدالة $h$ \textbf{متزايدة تماما} على $]-1;+\infty[$.

\textbf{جدول التغيّرات:}
\begin{TAB}
\tkzTabInit[lgt=2,espcl=4]{$x$/1,$h'(x)$/1,$h(x)$/2}{$-1$,$+\infty$}
\tkzTabLine{d,+,}
\tkzTabVar{D-/$-\infty$,+/$+\infty$}
\end{TAB}

\QQ{3} $h(0)=0^2+2\times 0+\ln 1=0$، إذن \bx{h(0)=0}.

بما أن $h$ متزايدة تماما على $]-1;+\infty[$ و $h(0)=0$ فإن:
\begin{itemize}
\item من أجل $-1<x<0$: $h(x)<h(0)$ أي $h(x)<0$.
\item من أجل $x>0$: $h(x)>h(0)$ أي $h(x)>0$.
\end{itemize}
\begin{TAB}
\tkzTabInit[lgt=2,espcl=3]{$x$/1,$h(x)$/1}{$-1$,$0$,$+\infty$}
\tkzTabLine{d,-,z,+,}
\end{TAB}

\textbf{الجزء الثاني:}

\QQ{1}
\pa لدينا $\lim\limits_{x\to -1^{+}}\ln(x+1)=-\infty$ و $\lim\limits_{x\to -1^{+}}(x+1)=0^{+}$، إذن $\lim\limits_{x\to -1^{+}}\dfrac{\ln(x+1)}{x+1}=-\infty$، ومنه $\lim\limits_{x\to -1^{+}}\left[-\dfrac{\ln(x+1)}{x+1}\right]=+\infty$، وبما أن $\lim\limits_{x\to -1^{+}}(x-1)=-2$:
\[\bx{\lim_{x\to -1^{+}}f(x)=+\infty}\]
\textbf{التفسير البياني:} المستقيم ذو المعادلة $x=-1$ مقارب عمودي للمنحنى $(C_f)$.

\pb نضع $u=e^{t}$ أي $t=\ln u$؛ عندما $u\to +\infty$ فإن $t\to +\infty$، ومنه:
\[\frac{\ln u}{u}=\frac{t}{e^{t}}=\frac{1}{\;\dfrac{e^{t}}{t}\;}\]
وبما أن $\lim\limits_{t\to +\infty}\dfrac{e^{t}}{t}=+\infty$ فإن $\lim\limits_{t\to +\infty}\dfrac{1}{e^{t}/t}=0$، إذن: \bx{\lim\limits_{u\to +\infty}\dfrac{\ln u}{u}=0}.

\pc نضع $u=x+1$؛ عندما $x\to +\infty$ فإن $u\to +\infty$، ومنه حسب (ب): $\lim\limits_{x\to +\infty}\dfrac{\ln(x+1)}{x+1}=\lim\limits_{u\to +\infty}\dfrac{\ln u}{u}=0$.
وبما أن $\lim\limits_{x\to +\infty}(x-1)=+\infty$ فإن:
\[\bx{\lim_{x\to +\infty}f(x)=+\infty}\]

\pd $f(x)-(x-1)=-\dfrac{\ln(x+1)}{x+1}$، ومنه حسب (جـ):
\[\bx{\lim_{x\to +\infty}\big[f(x)-(x-1)\big]=0}\]
إذن المستقيم $(\Delta)$ ذو المعادلة \bx{y=x-1} مقارب مائل للمنحنى $(C_f)$ عند $+\infty$.

\pe ندرس إشارة الفرق $f(x)-(x-1)=-\dfrac{\ln(x+1)}{x+1}$. بما أن $x+1>0$ فإن الإشارة هي عكس إشارة $\ln(x+1)$:
\begin{itemize}
\item $\ln(x+1)=0 \iff x+1=1\iff x=0$.
\item $\ln(x+1)>0\iff x+1>1\iff x>0$.
\end{itemize}
\begin{TAB}
\tkzTabInit[lgt=3,espcl=3]{$x$/1,$f(x)-(x-1)$/1}{$-1$,$0$,$+\infty$}
\tkzTabLine{d,+,z,-,}
\end{TAB}
\begin{itemize}
\item على $]-1;0[$: $(C_f)$ \textbf{فوق} المستقيم $(\Delta)$.
\item على $]0;+\infty[$: $(C_f)$ \textbf{تحت} المستقيم $(\Delta)$.
\item $(C_f)$ يقطع $(\Delta)$ في النقطة $(0;-1)$.
\end{itemize}

\QQ{2} الدالة $f$ قابلة للاشتقاق على $]-1;+\infty[$، ومن أجل كل $x>-1$:
\[\left(\frac{\ln(x+1)}{x+1}\right)'=\frac{\frac{1}{x+1}\times(x+1)-\ln(x+1)\times 1}{(x+1)^2}=\frac{1-\ln(x+1)}{(x+1)^2}\]
ومنه:
\begin{align*}
f'(x)&=1-\frac{1-\ln(x+1)}{(x+1)^2}=\frac{(x+1)^2-1+\ln(x+1)}{(x+1)^2}\\
&=\frac{x^2+2x+1-1+\ln(x+1)}{(x+1)^2}=\frac{x^2+2x+\ln(x+1)}{(x+1)^2}
\end{align*}
إذن: \bx{f'(x)=\dfrac{h(x)}{(x+1)^2}}.

بما أن $(x+1)^2>0$ فإن إشارة $f'(x)$ هي إشارة $h(x)$ (الجزء الأول، السؤال 3): $f$ متناقصة تماما على $]-1;0]$ ومتزايدة تماما على $[0;+\infty[$.
و $f(0)=0-1-\dfrac{\ln 1}{1}=-1$.
\begin{TAB}
\tkzTabInit[lgt=2,espcl=3]{$x$/1,$f'(x)$/1,$f(x)$/2}{$-1$,$0$,$+\infty$}
\tkzTabLine{d,-,z,+,}
\tkzTabVar{D+/$+\infty$,-/$-1$,+/$+\infty$}
\end{TAB}

\QQ{3} الدالة $f$ مستمرة ومتزايدة تماما على المجال $[3{,}3;3{,}4]$ (لأن $[3{,}3;3{,}4]\subset[0;+\infty[$)، ولدينا:
\[f(3{,}3)=2{,}3-\frac{\ln 4{,}3}{4{,}3}\approx 1{,}961<2\qquad;\qquad f(3{,}4)=2{,}4-\frac{\ln 4{,}4}{4{,}4}\approx 2{,}063>2\]
إذن $2$ محصور بين $f(3{,}3)$ و $f(3{,}4)$، وحسب مبرهنة القيم المتوسطة توجد قيمة وحيدة $x_0\in]3{,}3;3{,}4[$ تحقق $f(x_0)=2$؛ أي أن $(C_f)$ يقطع المستقيم ذا المعادلة $y=2$ عند نقطة فاصلتها \bx{3{,}3<x_0<3{,}4} ($x_0\approx 3{,}34$).

\textbf{ملاحظة:} يقطع $(C_f)$ المستقيم $y=2$ أيضا في نقطة ثانية فاصلتها في $]-1;0[$ (تقريبا $-0{,}69$) لأن $f$ تتناقص من $+\infty$ إلى $-1$ على هذا المجال.

\QQ{4} \textbf{الرسم:} المستقيمان المقاربان $x=-1$ و $(\Delta):y=x-1$، والنقطة $(0;-1)$ (ذروة صغرى بمماس أفقي وهي أيضا نقطة تقاطع $(C_f)$ مع $(\Delta)$).
\begin{GR}{-2}{6}{-2}{6}{0.9}
 \Ticks{-1,1,2,3,4,5}{-1,1,2,3,4,5}
 \draw[cThm,dashed,thick] (-1,-2)--(-1,6);
 \draw[cRem,thick] (-2,-3)--(6,5) node[pos=0.93,above left]{$(\Delta)$};
 \draw[gray,dashed] (-2,2)--(6,2);
 \draw[cDef,very thick,domain=-0.82:6,samples=200,smooth] plot (\x,{\x-1-ln(\x+1)/(\x+1)});
 \fill[cDef] (0,-1) circle (1.6pt);
 \fill[cEx] (3.338,2) circle (1.6pt);
 \node[cDef] at (5.3,3.1){$(C_f)$};
 \node[cEx,above] at (3.338,2){$x_0$};
\end{GR}
\end{hal}

% ===================== التمرين 2 =====================
\begin{tamrin}{بكالوريا 2010 -- الموضوع الأول}
\textbf{I)} $f$ الدالة العددية المعرفة على المجال $I=\left]\frac12;+\infty\right[$ بـ: $f(x)=1+\ln(2x-1)$.
وليكن $(C_f)$ تمثيلها البياني في المستوي المنسوب إلى المعلم المتعامد والمتجانس $(O;\vec i,\vec j)$.
\QQ{1} أحسب $\lim\limits_{x\to +\infty}f(x)$ و $\lim\limits_{x\to \frac12^{+}}f(x)$.
\QQ{2} بيّن أن الدالة $f$ متزايدة تماما على المجال $I$، ثم شكّل جدول تغيّراتها.
\QQ{3} عيّن فاصلة النقطة من $(C_f)$ التي تكون فيها المماس موازيا للمستقيم $(d)$ ذي المعادلة $y=x$.
\QQ{4}
\begin{sq}
\qa أثبت أنه من أجل كل $x$ من $I$ يمكن كتابة $f(x)$ على الشكل: $f(x)=\ln(x+a)+b$ حيث $a$، $b$ عددان حقيقيان يطلب تعيينهما.
\qb استنتج أنه يمكن رسم $(C_f)$ انطلاقا من $(C)$ منحنى الدالة اللوغاريتمية النيبيرية $\ln$، ثم أرسم $(C)$ و $(C_f)$.
\end{sq}

\textbf{II)} نعتبر الدالة العددية $g$ المعرفة على المجال $I$ بـ: $g(x)=f(x)-x$.
\QQ{1} أحسب $\lim\limits_{x\to \frac12^{+}}g(x)$ ثم بيّن أن $\lim\limits_{x\to +\infty}g(x)=-\infty$.
\QQ{2} أدرس اتجاه تغيّر الدالة $g$ على $I$ ثم شكّل جدول تغيّراتها.
\QQ{3}
\begin{sq}
\qa أحسب $g(1)$ ثم بيّن أن المعادلة $g(x)=0$ تقبل في المجال $\left]\frac32;+\infty\right[$ حلا وحيدا $\alpha$. تحقق أن $2<\alpha<3$.
\qb أرسم $(C_g)$ منحنى الدالة $g$ على المجال $\left]\frac12;5\right]$ في المعلم السابق.
\qc استنتج إشارة $g(x)$ على المجال $I$ ثم حدّد وضعية المنحنى $(C_f)$ بالنسبة إلى $(d)$.
\end{sq}
\QQ{4} برهن أنه من أجل كل عدد حقيقي $x$ من المجال $]1;\alpha[$ فإن: $f(x)$ ينتمي إلى المجال $]1;\alpha[$.

\textbf{III)} نسمي $(u_n)$ المتتالية العددية المعرفة على $\N^{*}$ كما يأتي: $u_n=f\left(1+\dfrac{1}{2n}\right)$.
\QQ{1} عيّن قيمة العدد الطبيعي $n$ التي من أجلها يكون: $u_n=1+2\ln 3-3\ln 2$.
\QQ{2} أحسب بدلالة $n$ المجموع $S_n$ حيث: $S_n=u_1+u_2+\cdots+u_n$.
\end{tamrin}

\begin{hal}
\textbf{الجزء الأول:}

\QQ{1}
\begin{itemize}
\item عند $+\infty$: $\lim\limits_{x\to+\infty}(2x-1)=+\infty$ و $\lim\limits_{X\to+\infty}\ln X=+\infty$، إذن بالتركيب \bx{\lim\limits_{x\to +\infty}f(x)=+\infty}.
\item عند $\frac12$ من اليمين: $\lim\limits_{x\to \frac12^{+}}(2x-1)=0^{+}$ و $\lim\limits_{X\to 0^{+}}\ln X=-\infty$، إذن \bx{\lim\limits_{x\to \frac12^{+}}f(x)=-\infty}.

ومنه المستقيم ذو المعادلة $x=\frac12$ مقارب عمودي لـ $(C_f)$.
\end{itemize}

\QQ{2} الدالة $f$ قابلة للاشتقاق على $I$، ومن أجل كل $x\in I$:
\[f'(x)=\frac{(2x-1)'}{2x-1}=\frac{2}{2x-1}\]
ومن أجل كل $x>\frac12$ لدينا $2x-1>0$ إذن $f'(x)>0$، ومنه $f$ \textbf{متزايدة تماما} على $I$.
\begin{TAB}
\tkzTabInit[lgt=2,espcl=4]{$x$/1,$f'(x)$/1,$f(x)$/2}{$\frac12$,$+\infty$}
\tkzTabLine{d,+,}
\tkzTabVar{D-/$-\infty$,+/$+\infty$}
\end{TAB}

\QQ{3} المماس عند النقطة ذات الفاصلة $x_0$ يوازي $(d):y=x$ معناه أن لهما نفس معامل التوجيه، أي $f'(x_0)=1$:
\[\frac{2}{2x_0-1}=1\iff 2x_0-1=2\iff x_0=\frac32\]
و $\frac32\in I$، إذن فاصلة النقطة المطلوبة هي \bx{x_0=\dfrac32}، وترتيبها $f\left(\frac32\right)=1+\ln 2$.

\QQ{4}
\pa من أجل كل $x\in I$: $2x-1=2\left(x-\frac12\right)$ مع $x-\frac12>0$، إذن:
\[f(x)=1+\ln\left[2\left(x-\tfrac12\right)\right]=1+\ln 2+\ln\left(x-\tfrac12\right)\]
ومنه \bx{f(x)=\ln\left(x-\tfrac12\right)+1+\ln 2} أي \bx{a=-\tfrac12} و \bx{b=1+\ln 2}.

\pb لدينا $f(x)=\ln\left(x-\frac12\right)+(1+\ln 2)$، إذن $(C_f)$ هو صورة $(C)$ بالانسحاب الذي شعاعه
\[\bx{\vec v=\tfrac12\,\vec i+(1+\ln 2)\,\vec j}\]
(كل نقطة $M(x;\ln x)$ من $(C)$ تنسحب إلى $M'\left(x+\frac12;\ln x+1+\ln 2\right)$ من $(C_f)$). الرسم مع $(C_g)$ في السؤال \textbf{II-3(ب)}.

\textbf{الجزء الثاني:} $g(x)=1+\ln(2x-1)-x$.

\QQ{1} $\lim\limits_{x\to \frac12^{+}}f(x)=-\infty$ و $\lim\limits_{x\to \frac12^{+}}x=\frac12$، إذن \bx{\lim\limits_{x\to \frac12^{+}}g(x)=-\infty}.

عند $+\infty$ (حالة عدم تعيين $+\infty-\infty$) نكتب من أجل $x>\frac12$:
\[g(x)=x\left[\frac{1}{x}+\frac{\ln(2x-1)}{2x-1}\times\frac{2x-1}{x}-1\right]\]
لدينا $\lim\limits_{x\to+\infty}\frac1x=0$، و $\lim\limits_{x\to+\infty}\frac{\ln(2x-1)}{2x-1}=\lim\limits_{X\to+\infty}\frac{\ln X}{X}=0$، و $\lim\limits_{x\to+\infty}\frac{2x-1}{x}=2$؛ إذن ما بين القوسين يؤول إلى $-1$، ومنه:
\[\bx{\lim_{x\to +\infty}g(x)=-\infty}\]

\QQ{2} $g$ قابلة للاشتقاق على $I$، ومن أجل كل $x\in I$:
\[g'(x)=\frac{2}{2x-1}-1=\frac{2-(2x-1)}{2x-1}=\frac{3-2x}{2x-1}\]
بما أن $2x-1>0$ فإن إشارة $g'(x)$ هي إشارة $3-2x$: $g'(x)>0$ على $\left]\frac12;\frac32\right[$، $g'(x)<0$ على $\left]\frac32;+\infty\right[$، و $g'\left(\frac32\right)=0$.

$g$ متزايدة تماما على $\left]\frac12;\frac32\right]$ ومتناقصة تماما على $\left[\frac32;+\infty\right[$، و $g\left(\frac32\right)=1+\ln2-\frac32=\ln 2-\frac12\approx 0{,}19$.
\begin{TAB}
\tkzTabInit[lgt=2,espcl=3]{$x$/1,$g'(x)$/1,$g(x)$/2}{$\frac12$,$\frac32$,$+\infty$}
\tkzTabLine{d,+,z,-,}
\tkzTabVar{D-/$-\infty$,+/$\ln2-\frac12$,-/$-\infty$}
\end{TAB}

\QQ{3}
\pa $g(1)=1+\ln(2-1)-1=\ln 1=0$، إذن \bx{g(1)=0}.

على المجال $\left[\frac32;+\infty\right[$: $g$ مستمرة ومتناقصة تماما، و $g\left(\frac32\right)=\ln2-\frac12>0$ و $\lim\limits_{x\to+\infty}g(x)=-\infty$، إذن حسب مبرهنة القيم المتوسطة المعادلة $g(x)=0$ تقبل حلا وحيدا $\alpha$ في $\left]\frac32;+\infty\right[$.

\textbf{التحقق:} $g(2)=\ln 3-1\approx 0{,}099>0$ و $g(3)=\ln5-2\approx -0{,}391<0$، إذن $g(2)\times g(3)<0$ ومنه \bx{2<\alpha<3} ($\alpha\approx 2{,}26$).

\pb \textbf{الرسم} (المنحنى $(C)$ بالأخضر، $(C_f)$ بالأزرق، $(C_g)$ بالبرتقالي، و $(d):y=x$):
\begin{GR}{-1}{5.5}{-3}{4.5}{1.15}
 \Ticks{-1,1,2,3,4,5}{-2,-1,1,2,3,4}
 \draw[cThm,dashed] (0.5,-3)--(0.5,4.5);
 \draw[gray!70!black,thick] (-1,-1)--(4.5,4.5) node[pos=0.97,left]{$(d)$};
 \draw[cRem,thick,domain=0.05:5.5,samples=150,smooth] plot (\x,{ln(\x)});
 \draw[cDef,very thick,domain=0.525:5.5,samples=150,smooth] plot (\x,{1+ln(2*\x-1)});
 \draw[cEx,very thick,domain=0.52:5,samples=150,smooth] plot (\x,{1+ln(2*\x-1)-\x});
 \draw[cDef,dashed] (0.5,1.193)--(2.5,3.193);
 \fill[cDef] (1,1) circle (1.6pt) (2.256,2.256) circle (1.6pt) (1.5,1.693) circle (1.6pt);
 \fill[cEx] (1,0) circle (1.6pt) (2.256,0) circle (1.6pt);
 \node[cEx,below right] at (2.256,0){$\alpha$};
 \node[cRem] at (5.0,1.25){$(C)$};
 \node[cDef] at (4.9,3.6){$(C_f)$};
 \node[cEx] at (4.3,-1.75){$(C_g)$};
\end{GR}

\pc من جدول تغيّرات $g$ ومن $g(1)=g(\alpha)=0$ (والحل $1$ وحيد في $\left]\frac12;\frac32\right]$ لأن $g$ متزايدة تماما عليه):
\begin{TAB}
\tkzTabInit[lgt=2,espcl=2.5]{$x$/1,$g(x)$/1}{$\frac12$,$1$,$\alpha$,$+\infty$}
\tkzTabLine{d,-,z,+,z,-,}
\end{TAB}
بما أن $g(x)=f(x)-x$:
\begin{itemize}
\item على $\left]\frac12;1\right[\cup]\alpha;+\infty[$: $(C_f)$ \textbf{تحت} $(d)$.
\item على $]1;\alpha[$: $(C_f)$ \textbf{فوق} $(d)$.
\item $(C_f)$ يقطع $(d)$ في النقطتين $(1;1)$ و $(\alpha;\alpha)$.
\end{itemize}

\QQ{4} ليكن $x\in]1;\alpha[$، أي $1<x<\alpha$. بما أن $f$ متزايدة تماما على $I$ فإن: $f(1)<f(x)<f(\alpha)$.
لدينا $f(1)=1+\ln1=1$، ومن $g(\alpha)=0$ نجد $f(\alpha)-\alpha=0$ أي $f(\alpha)=\alpha$. إذن:
\bx{1<f(x)<\alpha} أي \bx{f(x)\in]1;\alpha[}.

\textbf{الجزء الثالث:} من أجل كل $n\in\N^{*}$: $2\left(1+\frac{1}{2n}\right)-1=1+\frac1n=\frac{n+1}{n}>0$، إذن:
\[\bx{u_n=1+\ln\left(\frac{n+1}{n}\right)}\]

\QQ{1} $1+2\ln3-3\ln2=1+\ln 9-\ln 8=1+\ln\frac98$. إذن:
\[u_n=1+\ln\frac98\iff \ln\frac{n+1}{n}=\ln\frac98\iff \frac{n+1}{n}=\frac98\iff 8n+8=9n\iff n=8\]
(لأن الدالة $\ln$ متزايدة تماما على $]0;+\infty[$). إذن \bx{n=8}.

\QQ{2} $u_k=1+\ln(k+1)-\ln k$، ومنه:
\begin{align*}
S_n&=\sum_{k=1}^{n}\big[1+\ln(k+1)-\ln k\big]\\
&=n+\big[\ln2-\ln1\big]+\big[\ln3-\ln2\big]+\cdots+\big[\ln(n+1)-\ln n\big]
\end{align*}
تُختزل الحدود (مجموع تلسكوبي) فيبقى $\ln(n+1)-\ln 1$، إذن:
\[\bx{S_n=n+\ln(n+1)}\]
\end{hal}

% ===================== التمرين 3 =====================
\begin{tamrin}{بكالوريا 2011 -- الموضوع الأول}
\textbf{I)} نعتبر الدالة $g$ المعرفة على $\R-\{-1\}$ بـ: $g(x)=\dfrac{x-1}{x+1}$،
و $(C_g)$ تمثيلها البياني في المستوي المنسوب إلى المعلم المتعامد والمتجانس $(O;\vec i,\vec j)$ (الشكل المقابل).
\begin{GR}{-5}{5}{-3}{4}{0.6}
 \Ticks{1}{1}
 \draw[gray,dashed] (-5,1)--(5,1);
 \draw[gray,dashed] (-1,-3)--(-1,4);
 \draw[cDef,very thick,domain=-5:-1.45,samples=100,smooth] plot (\x,{(\x-1)/(\x+1)});
 \draw[cDef,very thick,domain=-0.5:5,samples=100,smooth] plot (\x,{(\x-1)/(\x+1)});
 \node[cDef] at (-3.6,3.2){$(C_g)$};
\end{GR}
\textbf{بقراءة بيانية:}
\begin{sq}
\qa شكّل جدول تغيّرات الدالة $g$.
\qb حلّ بيانيا المتراجحة $g(x)>0$.
\qc عيّن بيانيا قيم $x$ التي يكون من أجلها $0<g(x)<1$.
\end{sq}

\textbf{II)} لتكن الدالة $f$ المعرفة على المجال $]1;+\infty[$ بـ: $f(x)=\dfrac{x-1}{x+1}+\ln\left(\dfrac{x-1}{x+1}\right)$،
و $(C_f)$ تمثيلها البياني في المستوي المنسوب إلى المعلم المتعامد والمتجانس $(O;\vec i,\vec j)$.
\QQ{1} أحسب $\lim\limits_{x\to 1^{+}}f(x)$ و $\lim\limits_{x\to +\infty}f(x)$ ثم فسّر النتيجتين هندسيا.
\QQ{2}
\begin{sq}
\qa بيّن أنه من أجل كل عدد حقيقي $x$ من المجال $]1;+\infty[$: $g'(x)=\dfrac{2}{(x+1)^2}$.
\qb أحسب $f'(x)$ وأدرس إشارتها ثم شكّل جدول تغيّرات الدالة $f$.
\end{sq}
\QQ{3}
\begin{sq}
\qa باستعمال الجزء (I) السؤال (جـ)، عيّن إشارة العبارة $\ln\left(\dfrac{x-1}{x+1}\right)$ على المجال $]1;+\infty[$.
\qb $\alpha$ عدد حقيقي. بيّن أن الدالة $x\mapsto (x-\alpha)\ln(x-\alpha)-x$ هي دالة أصلية للدالة $x\mapsto \ln(x-\alpha)$ على المجال $]\alpha;+\infty[$.
\qc تحقق أنه من أجل كل عدد حقيقي $x$ من المجال $]1;+\infty[$: $g(x)=1-\dfrac{2}{x+1}$، ثم عيّن دالة أصلية للدالة $f$ على المجال $]1;+\infty[$.
\end{sq}
\end{tamrin}

\begin{hal}
\textbf{الجزء الأول (قراءة بيانية):}

\pa من الشكل: $g$ متزايدة تماما على كل من المجالين $]-\infty;-1[$ و $]-1;+\infty[$، والمستقيمان $x=-1$ و $y=1$ مقاربان لـ $(C_g)$:
\begin{TAB}
\tkzTabInit[lgt=2,espcl=3]{$x$/1,$g(x)$/2}{$-\infty$,$-1$,$+\infty$}
\tkzTabVar{-/$1$,+D-/$+\infty$/$-\infty$,+/$1$}
\end{TAB}

\pb $g(x)>0$ معناه أن $(C_g)$ يقع فوق محور الفواصل، ويحدث ذلك بيانيا لما $x<-1$ أو $x>1$ (المنحنى يقطع محور الفواصل في النقطة $(1;0)$). إذن مجموعة الحلول:
\[\bx{S=]-\infty;-1[\,\cup\,]1;+\infty[}\]

\pc $0<g(x)<1$ معناه أن $(C_g)$ محصور بين محور الفواصل والمستقيم $y=1$: على $]-\infty;-1[$ يكون $g(x)>1$، وعلى $]-1;1]$ يكون $g(x)\le 0$، وعلى $]1;+\infty[$ يكون $0<g(x)<1$. إذن:
\[\bx{0<g(x)<1\iff x\in]1;+\infty[}\]

\textbf{الجزء الثاني:} نلاحظ أن $f(x)=g(x)+\ln\big(g(x)\big)$، وهي معرفة على $]1;+\infty[$ لأن $g(x)>0$ عليه.

\QQ{1}
\begin{itemize}
\item $\lim\limits_{x\to1^{+}}g(x)=\dfrac{0^{+}}{2}=0^{+}$، و $\lim\limits_{X\to 0^{+}}\ln X=-\infty$، إذن \bx{\lim\limits_{x\to 1^{+}}f(x)=-\infty}.

\textbf{التفسير:} المستقيم ذو المعادلة $x=1$ مقارب عمودي لـ $(C_f)$.
\item $\lim\limits_{x\to+\infty}g(x)=\lim\limits_{x\to+\infty}\dfrac{x}{x}=1$، و $\lim\limits_{X\to 1}\ln X=0$، إذن \bx{\lim\limits_{x\to +\infty}f(x)=1+0=1}.

\textbf{التفسير:} المستقيم ذو المعادلة $y=1$ مقارب أفقي لـ $(C_f)$ عند $+\infty$.
\end{itemize}

\QQ{2}
\pa $g$ دالة ناطقة قابلة للاشتقاق على $]1;+\infty[$، ومن أجل كل $x>1$:
\[g'(x)=\frac{1\times(x+1)-(x-1)\times 1}{(x+1)^2}=\frac{x+1-x+1}{(x+1)^2}=\bx{\frac{2}{(x+1)^2}}\]

\pb من أجل $x>1$: $f(x)=g(x)+\ln\big(g(x)\big)$ مع $g(x)>0$، إذن:
\[f'(x)=g'(x)+\frac{g'(x)}{g(x)}=g'(x)\left(1+\frac{1}{g(x)}\right)=\frac{2}{(x+1)^2}\left(1+\frac{x+1}{x-1}\right)=\frac{2}{(x+1)^2}\times\frac{2x}{x-1}\]
\[\bx{f'(x)=\frac{4x}{(x-1)(x+1)^2}}\]
من أجل $x>1$: $4x>0$، $x-1>0$ و $(x+1)^2>0$، إذن $f'(x)>0$، ومنه $f$ \textbf{متزايدة تماما} على $]1;+\infty[$.
\begin{TAB}
\tkzTabInit[lgt=2,espcl=4]{$x$/1,$f'(x)$/1,$f(x)$/2}{$1$,$+\infty$}
\tkzTabLine{d,+,}
\tkzTabVar{D-/$-\infty$,+/$1$}
\end{TAB}

\QQ{3}
\pa حسب الجزء (I) السؤال (جـ): من أجل كل $x\in]1;+\infty[$ لدينا $0<\dfrac{x-1}{x+1}<1$، وبما أن الدالة $\ln$ متزايدة تماما على $]0;+\infty[$ و $\ln 1=0$ فإن:
من أجل كل $x\in]1;+\infty[$: \bx{\ln\left(\dfrac{x-1}{x+1}\right)<0}.

\pb نضع $F_\alpha(x)=(x-\alpha)\ln(x-\alpha)-x$ على $]\alpha;+\infty[$ (حيث $x-\alpha>0$). $F_\alpha$ قابلة للاشتقاق، و:
\[F_\alpha'(x)=1\times\ln(x-\alpha)+(x-\alpha)\times\frac{1}{x-\alpha}-1=\ln(x-\alpha)+1-1=\ln(x-\alpha)\]
إذن \bx{F_\alpha'(x)=\ln(x-\alpha)}، أي $F_\alpha$ دالة أصلية للدالة $x\mapsto\ln(x-\alpha)$ على $]\alpha;+\infty[$.

\pc $1-\dfrac{2}{x+1}=\dfrac{x+1-2}{x+1}=\dfrac{x-1}{x+1}=g(x)$ \ok.

ومن أجل $x>1$ لدينا $x-1>0$ و $x+1>0$، إذن $\ln\left(\dfrac{x-1}{x+1}\right)=\ln(x-1)-\ln(x+1)$، ومنه:
\[f(x)=1-\frac{2}{x+1}+\ln(x-1)-\ln(x+1)\]
نأخذ دالة أصلية لكل حد:
\begin{itemize}
\item $x\mapsto 1$: دالتها الأصلية $x$.
\item $x\mapsto -\dfrac{2}{x+1}$: دالتها الأصلية $-2\ln(x+1)$ (لأن $x+1>0$).
\item $x\mapsto \ln(x-1)$: حسب (ب) مع $\alpha=1$: $(x-1)\ln(x-1)-x$.
\item $x\mapsto \ln(x+1)$: حسب (ب) مع $\alpha=-1$ (و $]1;+\infty[\subset]-1;+\infty[$): $(x+1)\ln(x+1)-x$.
\end{itemize}
إذن دالة أصلية $F$ للدالة $f$ على $]1;+\infty[$:
\begin{align*}
F(x)&=x-2\ln(x+1)+(x-1)\ln(x-1)-x-\big[(x+1)\ln(x+1)-x\big]\\
&=x+(x-1)\ln(x-1)-(x+3)\ln(x+1)
\end{align*}
\[\bx{F(x)=x+(x-1)\ln(x-1)-(x+3)\ln(x+1)}\]
\end{hal}

% ===================== التمرين 4 =====================
\begin{tamrin}{بكالوريا 2012 -- الموضوع الأول}
نعتبر الدالة $f$ المعرفة على المجال $]-\infty;0[$ كما يلي: $f(x)=x+5+6\ln\left(\dfrac{x}{x-1}\right)$.
$(C_f)$ تمثيلها البياني في المستوي المنسوب إلى المعلم المتعامد والمتجانس $(O;\vec i,\vec j)$.
\QQ{1}
\begin{sq}
\qa أحسب $\lim\limits_{x\to 0^{-}}f(x)$، ثم فسّر النتيجة هندسيا.
\qb أحسب $\lim\limits_{x\to -\infty}f(x)$.
\end{sq}
\QQ{2} بيّن أنه من أجل كل $x$ من $]-\infty;0[$: $f'(x)=\dfrac{x^2-x-6}{x(x-1)}$.
استنتج اتجاه تغيّر الدالة $f$، ثم شكّل جدول تغيّراتها.
\QQ{3}
\begin{sq}
\qa بيّن أن المستقيم $(\Delta)$ الذي معادلة له: $y=x+5$ هو مستقيم مقارب مائل للمنحنى $(C_f)$ عند $-\infty$.
\qb أدرس وضع المنحنى $(C_f)$ بالنسبة للمستقيم $(\Delta)$.
\end{sq}
\QQ{4} بيّن أن المعادلة $f(x)=0$ تقبل حلين $\alpha$ و $\beta$ حيث $-3{,}5<\alpha<-3{,}4$ و $-1{,}1<\beta<-1$.
\QQ{5} أنشئ المنحنى $(C_f)$ والمستقيم $(\Delta)$.
\QQ{6}
\begin{sq}
\qa نعتبر النقطتين $A\left(-1;3+6\ln\frac34\right)$ و $B\left(-2;\frac52+6\ln\frac34\right)$.
بيّن أن $y=\dfrac12x+\dfrac72+6\ln\dfrac34$ معادلة ديكارتية للمستقيم $(AB)$.
\qb بيّن أن المستقيم $(AB)$ يمس المنحنى $(C_f)$ في نقطة $M_0$ يطلب تعيين إحداثيتيها.
\end{sq}
\end{tamrin}

\begin{hal}
\QQ{1}
\pa لما $x\to 0^{-}$: $x\to 0^{-}$ و $x-1\to -1$، إذن $\dfrac{x}{x-1}\to 0^{+}$، ومنه $\ln\left(\dfrac{x}{x-1}\right)\to -\infty$. وبما أن $x+5\to 5$:
\[\bx{\lim_{x\to 0^{-}}f(x)=-\infty}\]
\textbf{التفسير الهندسي:} المستقيم ذو المعادلة $x=0$ (محور التراتيب) مقارب عمودي لـ $(C_f)$.

\pb $\lim\limits_{x\to-\infty}\dfrac{x}{x-1}=\lim\limits_{x\to-\infty}\dfrac{x}{x}=1$، إذن $\lim\limits_{x\to-\infty}\ln\left(\dfrac{x}{x-1}\right)=\ln 1=0$، و $\lim\limits_{x\to-\infty}(x+5)=-\infty$، ومنه:
\[\bx{\lim_{x\to -\infty}f(x)=-\infty}\]

\QQ{2} من أجل $x<0$ لدينا $\dfrac{x}{x-1}>0$. نضع $u(x)=\dfrac{x}{x-1}$، إذن $u'(x)=\dfrac{(x-1)-x}{(x-1)^2}=\dfrac{-1}{(x-1)^2}$، و:
\[\frac{u'(x)}{u(x)}=\frac{-1}{(x-1)^2}\times\frac{x-1}{x}=\frac{-1}{x(x-1)}\]
ومنه:
\[f'(x)=1+6\times\frac{-1}{x(x-1)}=\frac{x(x-1)-6}{x(x-1)}=\bx{\frac{x^2-x-6}{x(x-1)}}\]

\textbf{الإشارة:} من أجل $x<0$: $x<0$ و $x-1<0$ إذن $x(x-1)>0$؛ فإشارة $f'(x)$ هي إشارة $x^2-x-6=(x+2)(x-3)$ ، وجذراه $-2$ و $3$ (و $3\notin]-\infty;0[$).
على $]-\infty;0[$: $x^2-x-6>0$ لما $x<-2$، و $x^2-x-6<0$ لما $-2<x<0$.

إذن $f$ متزايدة تماما على $]-\infty;-2]$ ومتناقصة تماما على $[-2;0[$، و:
\[f(-2)=-2+5+6\ln\frac{-2}{-3}=3+6\ln\frac23\approx 0{,}57\]
\begin{TAB}
\tkzTabInit[lgt=2,espcl=3]{$x$/1,$f'(x)$/1,$f(x)$/2}{$-\infty$,$-2$,$0$}
\tkzTabLine{,+,z,-,d}
\tkzTabVar{-/$-\infty$,+/$3+6\ln\frac23$,-D/$-\infty$}
\end{TAB}

\QQ{3}
\pa $f(x)-(x+5)=6\ln\left(\dfrac{x}{x-1}\right)$، وحسب 1(ب): $\lim\limits_{x\to-\infty}6\ln\left(\dfrac{x}{x-1}\right)=0$، إذن \bx{\lim\limits_{x\to-\infty}\big[f(x)-(x+5)\big]=0}، ومنه $(\Delta):y=x+5$ مستقيم مقارب مائل لـ $(C_f)$ عند $-\infty$.

\pb ندرس إشارة $f(x)-(x+5)=6\ln\left(\dfrac{x}{x-1}\right)$. لدينا:
\[\frac{x}{x-1}-1=\frac{x-(x-1)}{x-1}=\frac{1}{x-1}<0\quad(x<0)\]
إذن $0<\dfrac{x}{x-1}<1$، ومنه $\ln\left(\dfrac{x}{x-1}\right)<0$، أي $f(x)-(x+5)<0$ من أجل كل $x<0$.
\textbf{النتيجة:} المنحنى $(C_f)$ يقع \textbf{تحت} المستقيم $(\Delta)$ على $]-\infty;0[$.

\QQ{4}
\begin{itemize}
\item \textbf{على $]-\infty;-2]$:} $f$ مستمرة ومتزايدة تماما، و $f(-3{,}5)\approx-0{,}008<0$ و $f(-3{,}4)\approx 0{,}053>0$؛ إذن حسب مبرهنة القيم المتوسطة توجد قيمة وحيدة $\alpha\in]-3{,}5;-3{,}4[$ تحقق $f(\alpha)=0$ (وهي وحيدة على $]-\infty;-2]$ لأن $f$ تقابل منه نحو $]-\infty;f(-2)]$).
\item \textbf{على $[-2;0[$:} $f$ مستمرة ومتناقصة تماما، و $f(-1{,}1)\approx 0{,}020>0$ و $f(-1)=4+6\ln\frac12\approx -0{,}159<0$؛ إذن توجد قيمة وحيدة $\beta\in]-1{,}1;-1[$ تحقق $f(\beta)=0$.
\end{itemize}
إذن المعادلة $f(x)=0$ تقبل حلين بالضبط: \bx{-3{,}5<\alpha<-3{,}4} و \bx{-1{,}1<\beta<-1}.

\QQ{5} \textbf{الرسم:} (المقاربان: محور التراتيب و $(\Delta)$؛ الذروة $(-2;0{,}57)$؛ $\alpha\approx-3{,}49$ و $\beta\approx-1{,}09$؛ والمماس $(AB)$ من السؤال 6).
\begin{GR}{-7}{1}{-4}{3}{0.95}
 \Ticks{-6,-5,-4,-3,-2,-1}{-3,-2,-1,1,2}
 \draw[cRem,thick] (-7,-2)--(-2,3) node[pos=0.5,above left]{$(\Delta)$};
 \draw[cEx,thick,dashed] (-7,-1.726)--(1,2.274) node[pos=0.8,below right]{$(AB)$};
 \draw[cDef,very thick,domain=-7:-0.08,samples=200,smooth] plot (\x,{\x+5+6*ln(\x/(\x-1))});
 \fill[cDef] (-2,0.567) circle (1.6pt) (-3,0.274) circle (1.6pt) (-3.487,0) circle (1.6pt) (-1.088,0) circle (1.6pt);
 \node[cDef,above] at (-3,0.3){$M_0$};
 \node[cDef,above left] at (-3.487,0){$\alpha$};
 \node[cDef,above left] at (-1.088,0){$\beta$};
 \node[cDef] at (-5.6,-2.4){$(C_f)$};
\end{GR}

\QQ{6}
\pa معامل توجيه $(AB)$:
\[a=\frac{y_A-y_B}{x_A-x_B}=\frac{\left(3+6\ln\frac34\right)-\left(\frac52+6\ln\frac34\right)}{-1-(-2)}=\frac{\frac12}{1}=\frac12\]
معادلة $(AB)$ من الشكل $y=\frac12x+b$، وبما أن $A\in(AB)$: $3+6\ln\frac34=-\frac12+b$ إذن $b=\frac72+6\ln\frac34$. ومنه:
\[\bx{(AB):\ y=\frac12x+\frac72+6\ln\frac34}\]

\pb يكون $(AB)$ مماسا لـ $(C_f)$ في نقطة فاصلتها $x_0$ إذا كان $f'(x_0)=\frac12$ والنقطة $\big(x_0;f(x_0)\big)$ تنتمي إلى $(AB)$.
\begin{align*}
f'(x)=\frac12&\iff \frac{x^2-x-6}{x(x-1)}=\frac12\iff 2x^2-2x-12=x^2-x\\
&\iff x^2-x-12=0\iff (x+3)(x-4)=0
\end{align*}
و $4\notin]-\infty;0[$، إذن $x_0=-3$، و $f(-3)=-3+5+6\ln\dfrac{-3}{-4}=2+6\ln\dfrac34$.

\textbf{التحقق:} $\frac12(-3)+\frac72+6\ln\frac34=2+6\ln\frac34=f(-3)$، إذن النقطة $\big(-3;f(-3)\big)$ تنتمي إلى $(AB)$، والمماس عندها معامل توجيهه $\frac12$ مثل $(AB)$؛ فالمماس هو $(AB)$ نفسه.
\[\bx{M_0\left(-3\,;\,2+6\ln\tfrac34\right)}\qquad(M_0\approx(-3;0{,}27))\]
\end{hal}

% ===================== التمرين 5 =====================
\begin{tamrin}{بكالوريا 2013 -- الموضوع الثاني}
\textbf{I)} $g$ الدالة المعرفة على المجال $]-1;+\infty[$ بـ: $g(x)=x^2+2x+4-2\ln(x+1)$.
\QQ{1} أدرس تغيّرات الدالة $g$، ثم شكّل جدول تغيّراتها.
\QQ{2} استنتج أنه من أجل كل $x$ من المجال $]-1;+\infty[$: $g(x)>0$.

\textbf{II)} $f$ الدالة المعرفة على المجال $]-1;+\infty[$ بـ: $f(x)=x-\dfrac{1-2\ln(x+1)}{x+1}$،
و $(C_f)$ تمثيلها البياني في المستوي المنسوب إلى المعلم المتعامد والمتجانس $(O;\vec i,\vec j)$ (وحدة الطول $2\,cm$).
\QQ{1}
\begin{sq}
\qa أحسب $\lim\limits_{x\to -1^{+}}f(x)$. فسّر النتيجة بيانيا.
\qb أحسب $\lim\limits_{x\to +\infty}f(x)$.
\end{sq}
\QQ{2}
\begin{sq}
\qa بيّن أنه من أجل كل $x$ من $]-1;+\infty[$: $f'(x)=\dfrac{g(x)}{(x+1)^2}$.
\qb أدرس اتجاه تغيّر الدالة $f$ على المجال $]-1;+\infty[$، ثم شكّل جدول تغيّراتها.
\qc بيّن أن المعادلة $f(x)=0$ تقبل حلا وحيدا $\alpha$ في المجال $]-1;+\infty[$، ثم تحقق أن $0<\alpha<0{,}5$.
\end{sq}
\QQ{3}
\begin{sq}
\qa بيّن أن المستقيم $(\Delta)$ ذا المعادلة $y=x$ مقارب مائل للمنحنى $(C_f)$ عند $+\infty$.
\qb أدرس وضعية المنحنى $(C_f)$ بالنسبة إلى المستقيم $(\Delta)$.
\end{sq}
\QQ{4} نقبل أن المستقيم $(T)$ ذا المعادلة $y=x+\dfrac{2}{\sqrt{e^3}}$ مماس للمنحنى $(C_f)$ في نقطة فاصلتها $x_0$.
\begin{sq}
\qa أحسب $x_0$.
\qb أرسم المستقيمين المقاربين والمماس $(T)$ ثم المنحنى $(C_f)$.
\qc عيّن بيانيا قيم الوسيط الحقيقي $m$ بحيث تقبل المعادلة $f(x)=x+m$ حلين متمايزين.
\end{sq}
\end{tamrin}

\begin{hal}
\textbf{الجزء الأول:}

\QQ{1} \textbf{النهايات:}
\begin{itemize}
\item عند $-1^{+}$: $\lim\limits_{x\to-1^{+}}(x^2+2x+4)=3$ و $\lim\limits_{x\to-1^{+}}\ln(x+1)=-\infty$ إذن $-2\ln(x+1)\to+\infty$، ومنه \bx{\lim\limits_{x\to -1^{+}}g(x)=+\infty}.
\item عند $+\infty$: لدينا $x^2+2x+4=(x+1)^2+3$، إذن:
\[g(x)=(x+1)^2\left[1-2\,\frac{\ln(x+1)}{(x+1)^2}\right]+3\]
وبما أن $\lim\limits_{X\to+\infty}\dfrac{\ln X}{X^2}=\lim\limits_{X\to+\infty}\dfrac{\ln X}{X}\times\dfrac1X=0$ فإن ما بين القوسين يؤول إلى $1$، ومنه \bx{\lim\limits_{x\to +\infty}g(x)=+\infty}.
\end{itemize}
\textbf{المشتقة:} $g$ قابلة للاشتقاق على $]-1;+\infty[$ و:
\[g'(x)=2x+2-\frac{2}{x+1}=\frac{2(x+1)^2-2}{x+1}=\frac{2\big[(x+1)-1\big]\big[(x+1)+1\big]}{x+1}=\bx{\frac{2x(x+2)}{x+1}}\]
من أجل $x>-1$: $x+1>0$ و $x+2>0$، إذن إشارة $g'(x)$ هي إشارة $x$: $g$ متناقصة تماما على $]-1;0]$ ومتزايدة تماما على $[0;+\infty[$، و $g(0)=4-2\ln1=4$.
\begin{TAB}
\tkzTabInit[lgt=2,espcl=3]{$x$/1,$g'(x)$/1,$g(x)$/2}{$-1$,$0$,$+\infty$}
\tkzTabLine{d,-,z,+,}
\tkzTabVar{D+/$+\infty$,-/$4$,+/$+\infty$}
\end{TAB}

\QQ{2} من جدول التغيّرات، $g$ تقبل قيمة حدية صغرى مطلقة على $]-1;+\infty[$ هي $g(0)=4$، إذن من أجل كل $x>-1$: $g(x)\ge 4>0$، ومنه \bx{g(x)>0}.

\textbf{الجزء الثاني:}

\QQ{1}
\pa $\lim\limits_{x\to-1^{+}}\big(1-2\ln(x+1)\big)=+\infty$ و $\lim\limits_{x\to-1^{+}}(x+1)=0^{+}$، إذن $\lim\limits_{x\to-1^{+}}\dfrac{1-2\ln(x+1)}{x+1}=+\infty$، ومنه:
\[\bx{\lim_{x\to -1^{+}}f(x)=-1-(+\infty)=-\infty}\]
\textbf{التفسير:} المستقيم ذو المعادلة $x=-1$ مقارب عمودي لـ $(C_f)$.

\pb $\dfrac{1-2\ln(x+1)}{x+1}=\dfrac{1}{x+1}-2\,\dfrac{\ln(x+1)}{x+1}$، ولدينا $\lim\limits_{x\to+\infty}\dfrac{1}{x+1}=0$ و $\lim\limits_{x\to+\infty}\dfrac{\ln(x+1)}{x+1}=\lim\limits_{X\to+\infty}\dfrac{\ln X}{X}=0$؛ إذن هذا الكسر يؤول إلى $0$، و $\lim\limits_{x\to+\infty}x=+\infty$، ومنه:
\[\bx{\lim_{x\to +\infty}f(x)=+\infty}\]

\QQ{2}
\pa نضع $u(x)=1-2\ln(x+1)$ و $v(x)=x+1$، إذن $u'(x)=\dfrac{-2}{x+1}$ و $v'(x)=1$:
\[\left(\frac uv\right)'(x)=\frac{\frac{-2}{x+1}(x+1)-\big(1-2\ln(x+1)\big)}{(x+1)^2}=\frac{-3+2\ln(x+1)}{(x+1)^2}\]
ومنه:
\begin{align*}
f'(x)&=1-\frac{2\ln(x+1)-3}{(x+1)^2}=\frac{(x+1)^2+3-2\ln(x+1)}{(x+1)^2}\\
&=\frac{x^2+2x+4-2\ln(x+1)}{(x+1)^2}=\bx{\frac{g(x)}{(x+1)^2}}
\end{align*}

\pb حسب الجزء الأول $g(x)>0$، و $(x+1)^2>0$، إذن $f'(x)>0$: الدالة $f$ \textbf{متزايدة تماما} على $]-1;+\infty[$.
\begin{TAB}
\tkzTabInit[lgt=2,espcl=4]{$x$/1,$f'(x)$/1,$f(x)$/2}{$-1$,$+\infty$}
\tkzTabLine{d,+,}
\tkzTabVar{D-/$-\infty$,+/$+\infty$}
\end{TAB}

\pc $f$ مستمرة ومتزايدة تماما على $]-1;+\infty[$ وتأخذ قيمها في $]-\infty;+\infty[$ (من النهايتين)، و $0\in\R$؛ إذن المعادلة $f(x)=0$ تقبل حلا وحيدا $\alpha$ في $]-1;+\infty[$.
\[f(0)=0-\frac{1-0}{1}=-1<0\qquad;\qquad f(0{,}5)=0{,}5-\frac{1-2\ln1{,}5}{1{,}5}\approx 0{,}374>0\]
إذن $f(0)\times f(0{,}5)<0$، ومنه \bx{0<\alpha<0{,}5} ($\alpha\approx 0{,}33$).

\QQ{3}
\pa $f(x)-x=-\dfrac{1-2\ln(x+1)}{x+1}$، وحسب 1(ب) هذا الكسر يؤول إلى $0$، إذن \bx{\lim\limits_{x\to+\infty}\big[f(x)-x\big]=0}، ومنه $(\Delta):y=x$ مقارب مائل لـ $(C_f)$ عند $+\infty$.

\pb $f(x)-x=\dfrac{2\ln(x+1)-1}{x+1}$، وبما أن $x+1>0$ فإشارته هي إشارة $2\ln(x+1)-1$:
\[2\ln(x+1)-1>0\iff \ln(x+1)>\frac12\iff x+1>\sqrt e\iff x>\sqrt e-1\]
\begin{TAB}
\tkzTabInit[lgt=2.5,espcl=3]{$x$/1,$f(x)-x$/1}{$-1$,$\sqrt e-1$,$+\infty$}
\tkzTabLine{d,-,z,+,}
\end{TAB}
\begin{itemize}
\item على $]-1;\sqrt e-1[$: $(C_f)$ \textbf{تحت} $(\Delta)$.
\item على $]\sqrt e-1;+\infty[$: $(C_f)$ \textbf{فوق} $(\Delta)$.
\item $(C_f)$ يقطع $(\Delta)$ في النقطة $\left(\sqrt e-1;\sqrt e-1\right)$ ($\sqrt e-1\approx 0{,}65$).
\end{itemize}

\QQ{4}
\pa $(T)$ يوازي $(\Delta)$ (معامل توجيهه $1$)، إذن $f'(x_0)=1$:
\begin{align*}
\frac{(x_0+1)^2+3-2\ln(x_0+1)}{(x_0+1)^2}=1&\iff 3-2\ln(x_0+1)=0\\
&\iff \ln(x_0+1)=\frac32\iff x_0=e^{\frac32}-1
\end{align*}
\textbf{التحقق:} $f(x_0)=x_0-\dfrac{1-3}{e^{3/2}}=x_0+\dfrac{2}{\sqrt{e^3}}$، فمعادلة المماس: $y=1\cdot(x-x_0)+x_0+\dfrac{2}{\sqrt{e^3}}=x+\dfrac{2}{\sqrt{e^3}}$ \ok.
\[\bx{x_0=e^{\frac32}-1=\sqrt{e^3}-1\approx 3{,}48}\qquad\left(\frac{2}{\sqrt{e^3}}\approx 0{,}45\right)\]

\pb \textbf{الرسم} (رُسم بوحدة مصغّرة لضيق المساحة):
\begin{GR}{-1.5}{4.5}{-3}{5}{1.25}
 \Ticks{1,2,3,4}{-2,-1,1,2,3,4}
 \draw[cThm,dashed,thick] (-1,-3)--(-1,5);
 \draw[cRem,thick] (-1.5,-1.5)--(4.5,4.5) node[pos=0.9,below right]{$(\Delta)$};
 \draw[cEx,thick] (-1.5,-1.054)--(4.5,4.946) node[pos=0.72,above left]{$(T)$};
 \draw[cDef,very thick,domain=-0.24:4.5,samples=200,smooth] plot (\x,{\x-(1-2*ln(\x+1))/(\x+1)});
 \fill[cDef] (3.482,3.928) circle (1.6pt) (0.6487,0.6487) circle (1.6pt) (0.327,0) circle (1.6pt);
 \node[cDef,below right] at (0.327,0){$\alpha$};
 \node[cDef] at (0.9,-1.6){$(C_f)$};
\end{GR}

\pc $f(x)=x+m$: حلولها هي فواصل نقط تقاطع $(C_f)$ مع المستقيم $(D_m):y=x+m$ الموازي لـ $(\Delta)$.
\begin{itemize}
\item إذا كان $m\le 0$: $(D_m)$ يقطع $(C_f)$ في نقطة واحدة (حل وحيد).
\item إذا كان $0<m<\dfrac{2}{\sqrt{e^3}}$: $(D_m)$ محصور بين $(\Delta)$ و $(T)$ فيقطع $(C_f)$ في نقطتين (حلان متمايزان).
\item إذا كان $m=\dfrac{2}{\sqrt{e^3}}$: $(D_m)=(T)$ يمس $(C_f)$ (حل وحيد $x_0$).
\item إذا كان $m>\dfrac{2}{\sqrt{e^3}}$: لا توجد نقط تقاطع (لا حلول).
\end{itemize}
إذن المعادلة تقبل حلين متمايزين من أجل \bx{m\in\left]0;\dfrac{2}{\sqrt{e^3}}\right[}.
\end{hal}

% ===================== التمرين 6 =====================
\begin{tamrin}{بكالوريا 2014 -- الموضوع الأول}
نعتبر الدالة العددية $f$ المعرفة على المجال $]0;+\infty[$ كما يلي: $f(x)=1+\dfrac{2\ln x}{x}$،
و $(C_f)$ تمثيلها البياني في المستوي المنسوب إلى المعلم المتعامد والمتجانس $(O;\vec i,\vec j)$.
\QQ{1}
\begin{sq}
\qa أحسب $\lim\limits_{x\to 0^{+}}f(x)$ و $\lim\limits_{x\to +\infty}f(x)$، فسّر النتيجتين هندسيا.
\qb أدرس اتجاه تغيّر الدالة $f$ على المجال $]0;+\infty[$ ثم شكّل جدول تغيّراتها.
\end{sq}
\QQ{2}
\begin{sq}
\qa أدرس وضعية المنحنى $(C_f)$ بالنسبة إلى المستقيم $(\Delta)$ الذي معادلته: $y=1$.
\qb أكتب معادلة المماس $(T)$ للمنحنى $(C_f)$ في النقطة ذات الفاصلة $1$.
\qc بيّن أن المعادلة $f(x)=0$ تقبل في المجال $]0;1[$ حلا وحيدا $\alpha$، حيث $e^{-0{,}4}<\alpha<e^{-0{,}3}$.
\end{sq}
\QQ{3} أنشئ $(T)$ و $(C_f)$.
\QQ{4} لتكن الدالة $h$ المعرفة على $\R-\{0\}$ كما يلي: $h(x)=1+\dfrac{2\ln|x|}{|x|}$،
وليكن $(C_h)$ تمثيلها البياني في نفس المعلم السابق.
\begin{sq}
\qa بيّن أنه من أجل كل عدد حقيقي $x$ غير معدوم: $h(x)-h(-x)=0$. ماذا تستنتج؟
\qb أنشئ المنحنى $(C_h)$ اعتمادا على المنحنى $(C_f)$.
\qc ناقش بيانيا، حسب قيم الوسيط الحقيقي $m$، عدد حلول المعادلة: $\ln x^2=(m-1)|x|$.
\end{sq}
\end{tamrin}

\begin{hal}
\QQ{1}
\pa
\begin{itemize}
\item عند $0^{+}$: $\lim\limits_{x\to0^{+}}\ln x=-\infty$ و $\lim\limits_{x\to0^{+}}\dfrac1x=+\infty$، إذن $\lim\limits_{x\to0^{+}}\dfrac{2\ln x}{x}=\lim\limits_{x\to0^{+}}2\ln x\times\dfrac1x=-\infty$، ومنه \bx{\lim\limits_{x\to 0^{+}}f(x)=-\infty}.

\textbf{التفسير:} المستقيم ذو المعادلة $x=0$ (محور التراتيب) مقارب عمودي لـ $(C_f)$.
\item عند $+\infty$: $\lim\limits_{x\to+\infty}\dfrac{\ln x}{x}=0$ (نهاية شهيرة)، ومنه \bx{\lim\limits_{x\to +\infty}f(x)=1}.

\textbf{التفسير:} المستقيم ذو المعادلة $y=1$ مقارب أفقي لـ $(C_f)$ عند $+\infty$.
\end{itemize}

\pb $f$ قابلة للاشتقاق على $]0;+\infty[$، و:
\[f'(x)=2\times\frac{\frac1x\times x-\ln x\times1}{x^2}=\bx{\frac{2(1-\ln x)}{x^2}}\]
إشارة $f'(x)$ هي إشارة $1-\ln x$ (لأن $x^2>0$): $1-\ln x>0\iff \ln x<1\iff x<e$.
إذن $f$ متزايدة تماما على $]0;e]$ ومتناقصة تماما على $[e;+\infty[$، و $f(e)=1+\dfrac2e\approx 1{,}74$.
\begin{TAB}
\tkzTabInit[lgt=2,espcl=3]{$x$/1,$f'(x)$/1,$f(x)$/2}{$0$,$e$,$+\infty$}
\tkzTabLine{d,+,z,-,}
\tkzTabVar{D-/$-\infty$,+/$1+\frac2e$,-/$1$}
\end{TAB}

\QQ{2}
\pa $f(x)-1=\dfrac{2\ln x}{x}$، و $x>0$، إذن إشارة الفرق هي إشارة $\ln x$:
\begin{TAB}
\tkzTabInit[lgt=2,espcl=3]{$x$/1,$f(x)-1$/1}{$0$,$1$,$+\infty$}
\tkzTabLine{d,-,z,+,}
\end{TAB}
$(C_f)$ \textbf{تحت} $(\Delta)$ على $]0;1[$، و\textbf{فوق} $(\Delta)$ على $]1;+\infty[$، ويقطعه في النقطة $(1;1)$.

\pb $f(1)=1+0=1$ و $f'(1)=\dfrac{2(1-0)}{1}=2$، إذن:
\[(T):\ y=f'(1)(x-1)+f(1)=2(x-1)+1\quad\Longrightarrow\quad \bx{(T):\ y=2x-1}\]

\pc $f$ مستمرة ومتزايدة تماما على $]0;1]$ (لأن $]0;1]\subset]0;e]$)، و $\lim\limits_{x\to0^{+}}f(x)=-\infty$ و $f(1)=1>0$؛ إذن المعادلة $f(x)=0$ تقبل حلا وحيدا $\alpha$ في $]0;1[$.
\[f\left(e^{-0{,}4}\right)=1+\frac{2(-0{,}4)}{e^{-0{,}4}}=1-0{,}8\,e^{0{,}4}\approx-0{,}19<0\]
\[f\left(e^{-0{,}3}\right)=1+\frac{2(-0{,}3)}{e^{-0{,}3}}=1-0{,}6\,e^{0{,}3}\approx 0{,}19>0\]
إذن \bx{e^{-0{,}4}<\alpha<e^{-0{,}3}} ($\alpha\approx 0{,}70$).

\QQ{3} و \textbf{4(ب)} \textbf{الرسم:} $(C_f)$ بالأزرق، ونظيره بالنسبة لمحور التراتيب (بالأحمر المتقطع) يكمل $(C_h)$:
\begin{GR}{-8}{8}{-3}{3}{0.8}
 \Ticks{-6,-4,-2,1,2,4,6}{-2,-1,1,2}
 \draw[cRem,thick] (-8,1)--(8,1) node[pos=0.88,above]{$(\Delta)$};
 \draw[cEx,thick] (-1,-3)--(2,3) node[pos=0.93,left]{$(T)$};
 \draw[cDef,very thick,domain=0.42:8,samples=200,smooth] plot (\x,{1+2*ln(\x)/\x});
 \draw[cThm,very thick,dashed,domain=0.42:8,samples=200,smooth] plot (-\x,{1+2*ln(\x)/\x});
 \fill[cDef] (1,1) circle (1.8pt) (2.718,1.736) circle (1.8pt) (0.703,0) circle (1.8pt);
 \draw[gray,densely dotted] (2.718,0)--(2.718,1.736) (-2.718,0)--(-2.718,1.736);
 \node[below,font=\scriptsize] at (2.718,0){$e$};
 \node[below,font=\scriptsize] at (-2.718,0){$-e$};
 \node[cDef] at (6,2.2){$(C_f)$};
 \node[cThm] at (-6,2.2){$(C_h)$};
\end{GR}

\QQ{4}
\pa من أجل كل $x\ne0$: $|-x|=|x|$، إذن $h(-x)=1+\dfrac{2\ln|-x|}{|-x|}=1+\dfrac{2\ln|x|}{|x|}=h(x)$، أي \bx{h(x)-h(-x)=0}.

\textbf{الاستنتاج:} مجموعة التعريف $\R^{*}$ متناظرة بالنسبة إلى $0$ و $h(-x)=h(x)$، إذن $h$ \textbf{دالة زوجية}، ومنه $(C_h)$ متناظر بالنسبة إلى محور التراتيب.

\pb من أجل $x>0$: $|x|=x$ إذن $h(x)=f(x)$؛ فالجزء من $(C_h)$ على $]0;+\infty[$ هو $(C_f)$ نفسه، والجزء على $]-\infty;0[$ هو نظير $(C_f)$ بالنسبة إلى محور التراتيب (انظر الرسم أعلاه).

\pc من أجل $x\ne0$: $\ln x^2=\ln|x|^2=2\ln|x|$، إذن:
\begin{align*}
\ln x^2=(m-1)|x|&\iff 2\ln|x|=(m-1)|x|\iff \frac{2\ln|x|}{|x|}=m-1\\
&\iff 1+\frac{2\ln|x|}{|x|}=m\iff h(x)=m
\end{align*}
(قسمنا على $|x|>0$). عدد الحلول هو عدد نقط تقاطع $(C_h)$ مع المستقيم الأفقي $y=m$:
\begin{itemize}
\item $m\le 1$: حلان متعاكسان (أحدهما في $]0;1]$ والآخر نظيره).
\item $1<m<1+\dfrac2e$: أربعة حلول (حلان موجبان وحلان سالبان).
\item $m=1+\dfrac2e$: حلان هما $-e$ و $e$.
\item $m>1+\dfrac2e$: لا توجد حلول.
\end{itemize}
\end{hal}

% ===================== التمرين 7 =====================
\begin{tamrin}{بكالوريا 2015 -- الموضوع الأول}
المستوي منسوب إلى المعلم المتعامد والمتجانس $(O;\vec i,\vec j)$.

\textbf{I)} $(\gamma)$ التمثيل البياني للدالة $x\mapsto \ln x$ و $(\Delta)$ المستقيم ذو المعادلة $y=-x+3$.
$\alpha$ هي فاصلة نقطة تقاطع $(\gamma)$ و $(\Delta)$ (الشكل).
\begin{GR}{-0.5}{5}{-2}{3}{0.8}
 \Ticks{1}{1}
 \draw[cThm,thick] (-0.5,3.5)--(5,-2) node[pos=0.12,right]{$(\Delta)$};
 \draw[cDef,very thick,domain=0.14:5,samples=150,smooth] plot (\x,{ln(\x)});
 \draw[dashed] (2.208,0)--(2.208,0.792)--(0,0.792);
 \node[below,font=\scriptsize] at (2.208,0){$\alpha$};
 \node[cDef] at (4.3,1.9){$(\gamma)$};
\end{GR}
\QQ{1} بقراءة بيانية حدّد وضعية $(\gamma)$ بالنسبة إلى $(\Delta)$ على $]0;+\infty[$.
\QQ{2} $g$ الدالة المعرفة على المجال $]0;+\infty[$ بـ: $g(x)=x-3+\ln x$.
استنتج حسب قيم $x$ إشارة $g(x)$.
\QQ{3} تحقق أن: $2{,}2<\alpha<2{,}3$.

\textbf{II)} $f$ الدالة المعرفة على المجال $]0;+\infty[$ بـ: $f(x)=\left(1-\dfrac1x\right)(\ln x-2)$، و $(C_f)$ تمثيلها البياني.
\QQ{1} أحسب $\lim\limits_{x\to 0^{+}}f(x)$ و $\lim\limits_{x\to +\infty}f(x)$.
\QQ{2} أثبت أنه من أجل كل $x$ من $]0;+\infty[$: $f'(x)=\dfrac{g(x)}{x^2}$، ثم شكّل جدول تغيّرات الدالة $f$.
\QQ{3} بيّن أن: $f(\alpha)=\dfrac{-(\alpha-1)^2}{\alpha}$، ثم استنتج حصرا للعدد $f(\alpha)$.
\QQ{4} أدرس وضعية $(C_f)$ بالنسبة إلى حامل محور الفواصل، ثم أنشئ $(C_f)$ على المجال $]0;e^2]$.
\end{tamrin}

\begin{hal}
\textbf{الجزء الأول:}

\QQ{1} من الشكل:
\begin{itemize}
\item على $]0;\alpha[$: $(\gamma)$ يقع \textbf{تحت} $(\Delta)$.
\item على $]\alpha;+\infty[$: $(\gamma)$ يقع \textbf{فوق} $(\Delta)$.
\item $(\gamma)$ و $(\Delta)$ يتقاطعان في النقطة ذات الفاصلة $\alpha$.
\end{itemize}

\QQ{2} لدينا $g(x)=\ln x-(-x+3)$، أي $g(x)$ هو الفرق بين ترتيب نقطة من $(\gamma)$ وترتيب نقطة من $(\Delta)$ لهما نفس الفاصلة $x$. إذن حسب السؤال 1:
\begin{TAB}
\tkzTabInit[lgt=2,espcl=3]{$x$/1,$g(x)$/1}{$0$,$\alpha$,$+\infty$}
\tkzTabLine{d,-,z,+,}
\end{TAB}
$g(x)<0$ على $]0;\alpha[$، $g(\alpha)=0$، و $g(x)>0$ على $]\alpha;+\infty[$.

\QQ{3} $g$ مستمرة على $[2{,}2;2{,}3]$ ومتزايدة تماما عليه (لأن $g'(x)=1+\frac1x>0$)، و:
\[g(2{,}2)=-0{,}8+\ln 2{,}2\approx-0{,}012<0\qquad;\qquad g(2{,}3)=-0{,}7+\ln2{,}3\approx 0{,}133>0\]
إذن حسب مبرهنة القيم المتوسطة الحل الوحيد $\alpha$ للمعادلة $g(x)=0$ يحقق \bx{2{,}2<\alpha<2{,}3}.

\textbf{الجزء الثاني:}

\QQ{1}
\begin{itemize}
\item عند $0^{+}$: $\lim\limits_{x\to0^{+}}\left(1-\frac1x\right)=-\infty$ و $\lim\limits_{x\to0^{+}}(\ln x-2)=-\infty$، إذن بالجداء \bx{\lim\limits_{x\to 0^{+}}f(x)=+\infty} (المستقيم $x=0$ مقارب عمودي).
\item عند $+\infty$: $\lim\limits_{x\to+\infty}\left(1-\frac1x\right)=1$ و $\lim\limits_{x\to+\infty}(\ln x-2)=+\infty$، إذن \bx{\lim\limits_{x\to +\infty}f(x)=+\infty}.
\end{itemize}

\QQ{2} $f$ قابلة للاشتقاق على $]0;+\infty[$ (جداء دالتين قابلتين للاشتقاق)، و $(uv)'=u'v+uv'$:
\[f'(x)=\frac{1}{x^2}(\ln x-2)+\left(1-\frac1x\right)\frac1x=\frac{\ln x-2}{x^2}+\frac{x-1}{x^2}=\frac{x-3+\ln x}{x^2}\]
إذن \bx{f'(x)=\dfrac{g(x)}{x^2}}، وإشارتها هي إشارة $g(x)$ (لأن $x^2>0$): $f$ متناقصة تماما على $]0;\alpha]$ ومتزايدة تماما على $[\alpha;+\infty[$.
\begin{TAB}
\tkzTabInit[lgt=2,espcl=3]{$x$/1,$f'(x)$/1,$f(x)$/2}{$0$,$\alpha$,$+\infty$}
\tkzTabLine{d,-,z,+,}
\tkzTabVar{D+/$+\infty$,-/$f(\alpha)$,+/$+\infty$}
\end{TAB}

\QQ{3} بما أن $g(\alpha)=0$ فإن $\alpha-3+\ln\alpha=0$ أي $\ln\alpha=3-\alpha$، ومنه $\ln\alpha-2=1-\alpha$. إذن:
\[f(\alpha)=\left(1-\frac1\alpha\right)(1-\alpha)=\frac{\alpha-1}{\alpha}\times\big(-(\alpha-1)\big)=\bx{\frac{-(\alpha-1)^2}{\alpha}}\]
\textbf{الحصر:} نعتبر الدالة $\varphi(t)=\dfrac{-(t-1)^2}{t}$ على $[2{,}2;2{,}3]$:
\[\varphi'(t)=-\frac{2(t-1)\,t-(t-1)^2}{t^2}=-\frac{(t-1)(t+1)}{t^2}<0\]
إذن $\varphi$ متناقصة تماما على $[2{,}2;2{,}3]$، ومن $2{,}2<\alpha<2{,}3$ نجد $\varphi(2{,}3)<\varphi(\alpha)<\varphi(2{,}2)$:
\[-\frac{1{,}69}{2{,}3}<f(\alpha)<-\frac{1{,}44}{2{,}2}\quad\Longrightarrow\quad \bx{-0{,}74<f(\alpha)<-0{,}65}\qquad(f(\alpha)\approx-0{,}66)\]

\QQ{4} $f(x)=\dfrac{x-1}{x}\,(\ln x-2)$. ندرس إشارة كل عامل:
\begin{itemize}
\item $\dfrac{x-1}{x}$ (مع $x>0$) له إشارة $x-1$: ينعدم عند $1$.
\item $\ln x-2=0\iff x=e^2$، و $\ln x-2>0\iff x>e^2$.
\end{itemize}
\begin{TAB}
\tkzTabInit[lgt=2.5,espcl=2.5]{$x$/1,$\frac{x-1}{x}$/1,$\ln x-2$/1,$f(x)$/1}{$0$,$1$,$e^2$,$+\infty$}
\tkzTabLine{d,-,z,+,t,+,}
\tkzTabLine{d,-,t,-,z,+,}
\tkzTabLine{d,+,z,-,z,+,}
\end{TAB}
\begin{itemize}
\item على $]0;1[\cup]e^2;+\infty[$: $(C_f)$ \textbf{فوق} محور الفواصل.
\item على $]1;e^2[$: $(C_f)$ \textbf{تحت} محور الفواصل.
\item $(C_f)$ يقطع محور الفواصل في النقطتين $(1;0)$ و $(e^2;0)$ ($e^2\approx 7{,}39$).
\end{itemize}
\textbf{الرسم على $]0;e^2]$:}
\begin{GR}{-0.5}{8}{-1.5}{4}{1}
 \Ticks{1,2,3,4,5,6,7}{-1,1,2,3}
 \draw[cDef,very thick,domain=0.42:7.389,samples=200,smooth] plot (\x,{(1-1/\x)*(ln(\x)-2)});
 \fill[cDef] (1,0) circle (1.8pt) (7.389,0) circle (1.8pt) (2.208,-0.661) circle (1.8pt);
 \draw[dashed] (2.208,0)--(2.208,-0.661);
 \node[above,font=\scriptsize] at (2.208,0){$\alpha$};
 \node[above,font=\scriptsize] at (7.389,0){$e^2$};
 \node[cDef] at (1.3,2.8){$(C_f)$};
\end{GR}
\end{hal}

% ===================== التمرين 8 =====================
\begin{tamrin}{بكالوريا 2016 -- الموضوع الأول -- الدورة الأولى}
\textbf{I)} $g$ الدالة المعرفة على المجال $]0;+\infty[$ بـ: $g(x)=x^2+1-\ln x$.
\QQ{1} أدرس اتجاه تغيّر الدالة $g$.
\QQ{2} أحسب $g\left(\dfrac{\sqrt2}{2}\right)$ ثم بيّن أنه من أجل كل عدد حقيقي $x$ من المجال $]0;+\infty[$: $g(x)>0$.

\textbf{II)} $f$ الدالة المعرفة على المجال $]0;+\infty[$ بـ: $f(x)=\dfrac{\ln x}{x}+x-1$،
و $(C)$ تمثيلها البياني في المستوي المنسوب إلى المعلم المتعامد والمتجانس $(O;\vec i,\vec j)$.
\QQ{1} أحسب $\lim\limits_{x\to 0^{+}}f(x)$ و $\lim\limits_{x\to +\infty}f(x)$.
\QQ{2}
\begin{sq}
\qa بيّن أنه من أجل كل عدد حقيقي $x$ من المجال $]0;+\infty[$: $f'(x)=\dfrac{g(x)}{x^2}$.
\qb شكّل جدول تغيّرات الدالة $f$.
\end{sq}
\QQ{3} أكتب معادلة للمماس $(T)$ للمنحنى $(C)$ في النقطة التي فاصلتها $1$.
\QQ{4}
\begin{sq}
\qa بيّن أن $(C)$ يقبل مستقيما مقاربا مائلا $(\Delta)$ حيث $y=x-1$ معادلة له.
\qb أدرس الوضع النسبي لـ $(C)$ و $(\Delta)$.
\end{sq}
\QQ{5} أرسم المستقيمين $(T)$ و $(\Delta)$ ثم المنحنى $(C)$.
\QQ{6} $m$ عدد حقيقي. $(\Delta_m)$ المستقيم حيث $y=mx-m$ معادلة له.
\begin{sq}
\qa تحقق أنه من أجل كل عدد حقيقي $m$، النقطة $A(1;0)$ تنتمي إلى المستقيم $(\Delta_m)$.
\qb ناقش بيانيا وحسب قيم الوسيط الحقيقي $m$ عدد حلول المعادلة $f(x)=mx-m$.
\end{sq}
\end{tamrin}

\begin{hal}
\textbf{الجزء الأول:}

\QQ{1} $g$ قابلة للاشتقاق على $]0;+\infty[$، و:
\[g'(x)=2x-\frac1x=\bx{\frac{2x^2-1}{x}}\]
بما أن $x>0$ فإن إشارة $g'(x)$ هي إشارة $2x^2-1$، و $2x^2-1=0\iff x=\dfrac{\sqrt2}{2}$ (لأن $x>0$).
إذن $g$ متناقصة تماما على $\left]0;\frac{\sqrt2}{2}\right]$ ومتزايدة تماما على $\left[\frac{\sqrt2}{2};+\infty\right[$.
\begin{TAB}
\tkzTabInit[lgt=2,espcl=3]{$x$/1,$g'(x)$/1,$g(x)$/2}{$0$,$\frac{\sqrt2}{2}$,$+\infty$}
\tkzTabLine{d,-,z,+,}
\tkzTabVar{D+/$+\infty$,-/$\frac{3+\ln2}{2}$,+/$+\infty$}
\end{TAB}
(النهايتان: عند $0^{+}$: $-\ln x\to+\infty$؛ وعند $+\infty$: $g(x)=x^2\left(1+\frac1{x^2}-\frac{\ln x}{x}\cdot\frac1x\right)\to+\infty$.)

\QQ{2} $\left(\frac{\sqrt2}{2}\right)^2=\frac12$ و $\ln\frac{\sqrt2}{2}=\ln\frac{1}{\sqrt2}=-\frac12\ln2$، إذن:
\[g\left(\frac{\sqrt2}{2}\right)=\frac12+1+\frac12\ln2=\bx{\frac{3+\ln 2}{2}}\approx 1{,}85\]
وهي القيمة الحدية الصغرى للدالة $g$ على $]0;+\infty[$، وهي موجبة تماما؛ إذن من أجل كل $x>0$: $g(x)\ge\frac{3+\ln2}{2}>0$، ومنه \bx{g(x)>0}.

\textbf{الجزء الثاني:}

\QQ{1}
\begin{itemize}
\item عند $0^{+}$: $\dfrac{\ln x}{x}=\ln x\times\dfrac1x$ مع $\ln x\to-\infty$ و $\dfrac1x\to+\infty$، إذن $\dfrac{\ln x}{x}\to-\infty$، و $x-1\to-1$، ومنه \bx{\lim\limits_{x\to 0^{+}}f(x)=-\infty} (محور التراتيب مقارب عمودي).
\item عند $+\infty$: $\lim\limits_{x\to+\infty}\dfrac{\ln x}{x}=0$ و $\lim\limits_{x\to+\infty}(x-1)=+\infty$، ومنه \bx{\lim\limits_{x\to +\infty}f(x)=+\infty}.
\end{itemize}

\QQ{2}
\pa $f$ قابلة للاشتقاق على $]0;+\infty[$، و:
\[f'(x)=\frac{\frac1x\cdot x-\ln x\cdot 1}{x^2}+1=\frac{1-\ln x}{x^2}+\frac{x^2}{x^2}=\frac{x^2+1-\ln x}{x^2}=\bx{\frac{g(x)}{x^2}}\]

\pb $g(x)>0$ و $x^2>0$، إذن $f'(x)>0$: $f$ \textbf{متزايدة تماما} على $]0;+\infty[$.
\begin{TAB}
\tkzTabInit[lgt=2,espcl=4]{$x$/1,$f'(x)$/1,$f(x)$/2}{$0$,$+\infty$}
\tkzTabLine{d,+,}
\tkzTabVar{D-/$-\infty$,+/$+\infty$}
\end{TAB}

\QQ{3} $f(1)=\frac{0}{1}+1-1=0$ و $f'(1)=\dfrac{g(1)}{1}=1+1-0=2$، إذن:
\[(T):\ y=f'(1)(x-1)+f(1)=2(x-1)\quad\Longrightarrow\quad\bx{(T):\ y=2x-2}\]

\QQ{4}
\pa $f(x)-(x-1)=\dfrac{\ln x}{x}$ و $\lim\limits_{x\to+\infty}\dfrac{\ln x}{x}=0$، إذن \bx{\lim\limits_{x\to+\infty}\big[f(x)-(x-1)\big]=0}، ومنه $(\Delta):y=x-1$ مستقيم مقارب مائل لـ $(C)$ عند $+\infty$.

\pb إشارة $f(x)-(x-1)=\dfrac{\ln x}{x}$ هي إشارة $\ln x$ (لأن $x>0$):
\begin{TAB}
\tkzTabInit[lgt=3,espcl=3]{$x$/1,$f(x)-(x-1)$/1}{$0$,$1$,$+\infty$}
\tkzTabLine{d,-,z,+,}
\end{TAB}
$(C)$ \textbf{تحت} $(\Delta)$ على $]0;1[$، و\textbf{فوق} $(\Delta)$ على $]1;+\infty[$، ويقطعه في النقطة $A(1;0)$.

\QQ{5} \textbf{الرسم:}
\begin{GR}{-0.5}{5}{-3}{4.5}{1}
 \Ticks{1,2,3,4}{-2,-1,1,2,3,4}
 \draw[cRem,thick] (-0.5,-1.5)--(5,4) node[pos=0.8,below right]{$(\Delta)$};
 \draw[cEx,thick] (-0.25,-2.5)--(3.25,4.5) node[pos=0.9,left]{$(T)$};
 \draw[cDef,very thick,domain=0.4:5,samples=200,smooth] plot (\x,{ln(\x)/\x+\x-1});
 \fill[cDef] (1,0) circle (1.8pt);
 \node[below right] at (1,0){$A$};
 \node[cDef] at (3.3,3.9){$(C)$};
\end{GR}

\QQ{6}
\pa من أجل كل $m\in\R$: $m\times1-m=0=y_A$، إذن \bx{A(1;0)\in(\Delta_m)}: كل المستقيمات $(\Delta_m)$ تشمل النقطة $A$، و $m$ معامل توجيهها.

\pb حلول المعادلة $f(x)=mx-m$ هي فواصل نقط تقاطع $(C)$ مع $(\Delta_m)$. بما أن $A\in(C)$ (لأن $f(1)=0$) فإن $x=1$ حل دائما.
ندوّر المستقيم $(\Delta_m)$ حول $A$ مع مراعاة أن ميل $(\Delta)$ هو $1$ وميل $(T)$ هو $2$:
\begin{itemize}
\item $m\le 1$: $(\Delta_m)$ يقطع $(C)$ في $A$ فقط: \textbf{حل وحيد} $x=1$.
\item $1<m<2$: $(\Delta_m)$ يقطع $(C)$ في $A$ وفي نقطة أخرى فاصلتها أكبر من $1$: \textbf{حلان}.
\item $m=2$: $(\Delta_m)=(T)$ مماس لـ $(C)$ في $A$: \textbf{حل وحيد} $x=1$.
\item $m>2$: $(\Delta_m)$ يقطع $(C)$ في $A$ وفي نقطة أخرى فاصلتها في $]0;1[$: \textbf{حلان}.
\end{itemize}
\end{hal}

% ===================== التمرين 9 =====================
\begin{tamrin}{بكالوريا 2016 -- الموضوع الأول -- الدورة الثانية}
\textbf{I)} $g$ الدالة المعرفة على $]-1;+\infty[$ بـ: $g(x)=-1+(x+1)e+2\ln(x+1)$ (العدد $e$ هو أساس اللوغاريتم النيبيري).
\QQ{1} أدرس تغيّرات الدالة $g$، ثم شكّل جدول تغيّراتها.
\QQ{2} بيّن أن للمعادلة $g(x)=0$ حلا وحيدا $\alpha$ حيث $-0{,}34<\alpha<-0{,}33$.
\QQ{3} استنتج إشارة $g(x)$ حسب قيم العدد الحقيقي $x$ على المجال $]-1;+\infty[$.

\textbf{II)} $f$ الدالة المعرفة على المجال $]-1;+\infty[$ بـ: $f(x)=\dfrac{e}{x+1}+\dfrac{\ln(x+1)}{(x+1)^2}$،
و $(C_f)$ تمثيلها البياني في المستوي المنسوب إلى المعلم المتعامد والمتجانس $(O;\vec i,\vec j)$.
\QQ{1}
\begin{sq}
\qa بيّن أن $\lim\limits_{x\to -1^{+}}f(x)=-\infty$ وأحسب $\lim\limits_{x\to +\infty}f(x)$، فسّر النتيجتين هندسيا.
\qb بيّن أنه من أجل كل $x$ من المجال $]-1;+\infty[$: $f'(x)=\dfrac{-g(x)}{(x+1)^3}$.
\qc أدرس اتجاه تغيّر الدالة $f$ على المجال $]-1;+\infty[$، ثم شكّل جدول تغيّراتها.
\qd أرسم المنحنى $(C_f)$ (نقبل أن: $f(\alpha)\approx 3{,}16$).
\end{sq}
\QQ{2} نعتبر الدالة العددية $k$ المعرفة على $]-1;1[$ بـ: $k(x)=f(-|x|)$، و $(C_k)$ تمثيلها البياني في المعلم السابق.
\begin{sq}
\qa بيّن أن الدالة $k$ زوجية.
\qb بيّن كيف يمكن استنتاج المنحنى $(C_k)$ انطلاقا من المنحنى $(C_f)$ ثم أرسمه (دون دراسة تغيّرات الدالة $k$).
\qc ناقش بيانيا وحسب قيم الوسيط الحقيقي $m$ عدد وإشارة حلول المعادلة $k(x)=m$.
\end{sq}
\end{tamrin}

\begin{hal}
\textbf{الجزء الأول:}

\QQ{1} \textbf{النهايات:}
\begin{itemize}
\item عند $-1^{+}$: $(x+1)e\to 0$ و $\ln(x+1)\to-\infty$، إذن \bx{\lim\limits_{x\to -1^{+}}g(x)=-\infty}.
\item عند $+\infty$: $(x+1)e\to+\infty$ و $\ln(x+1)\to+\infty$، إذن \bx{\lim\limits_{x\to +\infty}g(x)=+\infty}.
\end{itemize}
\textbf{المشتقة:} $g'(x)=e+\dfrac{2}{x+1}$، ومن أجل $x>-1$: $\dfrac{2}{x+1}>0$ و $e>0$، إذن $g'(x)>0$: $g$ \textbf{متزايدة تماما} على $]-1;+\infty[$.
\begin{TAB}
\tkzTabInit[lgt=2,espcl=4]{$x$/1,$g'(x)$/1,$g(x)$/2}{$-1$,$+\infty$}
\tkzTabLine{d,+,}
\tkzTabVar{D-/$-\infty$,+/$+\infty$}
\end{TAB}

\QQ{2} $g$ مستمرة ومتزايدة تماما على $]-1;+\infty[$ وتأخذ قيمها في $\R$، إذن المعادلة $g(x)=0$ تقبل حلا وحيدا $\alpha$. ولدينا:
\[g(-0{,}34)=-1+0{,}66\,e+2\ln0{,}66\approx-0{,}037<0\]
\[g(-0{,}33)=-1+0{,}67\,e+2\ln0{,}67\approx 0{,}020>0\]
إذن \bx{-0{,}34<\alpha<-0{,}33}.

\QQ{3} بما أن $g$ متزايدة تماما و $g(\alpha)=0$:
\begin{TAB}
\tkzTabInit[lgt=2,espcl=3]{$x$/1,$g(x)$/1}{$-1$,$\alpha$,$+\infty$}
\tkzTabLine{d,-,z,+,}
\end{TAB}

\textbf{الجزء الثاني:}

\QQ{1}
\pa من أجل $x>-1$: $f(x)=\dfrac{e(x+1)+\ln(x+1)}{(x+1)^2}$.
لما $x\to-1^{+}$: البسط $e(x+1)+\ln(x+1)\to 0+(-\infty)=-\infty$ والمقام $(x+1)^2\to 0^{+}$، إذن:
\[\bx{\lim_{x\to -1^{+}}f(x)=-\infty}\]
وعند $+\infty$: $\dfrac{e}{x+1}\to 0$ و $\dfrac{\ln(x+1)}{(x+1)^2}=\dfrac{\ln(x+1)}{x+1}\times\dfrac{1}{x+1}\to 0\times0=0$، إذن:
\[\bx{\lim_{x\to +\infty}f(x)=0}\]
\textbf{التفسير الهندسي:} المستقيم $x=-1$ مقارب عمودي لـ $(C_f)$، والمستقيم $y=0$ (محور الفواصل) مقارب أفقي لـ $(C_f)$ عند $+\infty$.

\pb
\[\left(\frac{e}{x+1}\right)'=\frac{-e}{(x+1)^2}\]
\[\left(\frac{\ln(x+1)}{(x+1)^2}\right)'=\frac{\frac{1}{x+1}(x+1)^2-\ln(x+1)\cdot2(x+1)}{(x+1)^4}=\frac{1-2\ln(x+1)}{(x+1)^3}\]
ومنه:
\[f'(x)=\frac{-e(x+1)}{(x+1)^3}+\frac{1-2\ln(x+1)}{(x+1)^3}=\frac{-\big[-1+(x+1)e+2\ln(x+1)\big]}{(x+1)^3}=\bx{\frac{-g(x)}{(x+1)^3}}\]

\pc من أجل $x>-1$: $(x+1)^3>0$، إذن إشارة $f'(x)$ عكس إشارة $g(x)$: $f'(x)>0$ على $]-1;\alpha[$ و $f'(x)<0$ على $]\alpha;+\infty[$.
$f$ متزايدة تماما على $]-1;\alpha]$ ومتناقصة تماما على $[\alpha;+\infty[$.
\begin{TAB}
\tkzTabInit[lgt=2,espcl=3]{$x$/1,$f'(x)$/1,$f(x)$/2}{$-1$,$\alpha$,$+\infty$}
\tkzTabLine{d,+,z,-,}
\tkzTabVar{D-/$-\infty$,+/$f(\alpha)$,-/$0$}
\end{TAB}

\pd \textbf{الرسم} ($f(\alpha)\approx3{,}16$، $f(0)=e\approx2{,}72$، و $(C_f)$ يقطع محور الفواصل عند $x\approx-0{,}63$). المنحنى $(C_f)$ بالأزرق، و $(C_k)$ (السؤال 2) يتكوّن من الجزء الأزرق على $]-1;0]$ ونظيره بالأحمر:
\begin{GR}{-1.5}{5}{-2}{4}{1.25}
 \Ticks{-1,1,2,3,4}{-1,1,2,3}
 \draw[cThm,dashed,thick] (-1,-2)--(-1,4);
 \draw[cThm,dashed,thick] (1,-2)--(1,4);
 \draw[cDef,very thick,domain=-0.67:5,samples=250,smooth] plot (\x,{exp(1)/(\x+1)+ln(\x+1)/((\x+1)^2)});
 \draw[cThm,very thick,dashed,domain=-0.67:0,samples=150,smooth] plot (-\x,{exp(1)/(\x+1)+ln(\x+1)/((\x+1)^2)});
 \fill[cDef] (-0.3336,3.165) circle (1.8pt) (0,2.718) circle (1.8pt);
 \fill[cThm] (0.3336,3.165) circle (1.8pt);
 \draw[gray,densely dotted] (-1.5,3.165)--(1.5,3.165);
 \draw[gray,densely dotted] (-0.3336,0)--(-0.3336,3.165);
 \node[below,font=\scriptsize] at (-0.3336,0){$\alpha$};
 \node[cDef] at (3.2,1.3){$(C_f)$};
 \node[cThm] at (1.5,-1.2){$(C_k)$};
\end{GR}

\QQ{2}
\pa من أجل كل $x\in]-1;1[$: $-x\in]-1;1[$ (المجال متناظر بالنسبة إلى $0$)، و:
\[k(-x)=f(-|-x|)=f(-|x|)=k(x)\]
إذن الدالة $k$ \textbf{زوجية}، و $(C_k)$ متناظر بالنسبة إلى محور التراتيب.

\pb من أجل $x\in]-1;0]$: $|x|=-x$ إذن $-|x|=x$ ومنه $k(x)=f(x)$.
\begin{itemize}
\item على $]-1;0]$: $(C_k)$ ينطبق على جزء $(C_f)$ الموافق لـ $]-1;0]$.
\item على $[0;1[$: $(C_k)$ هو نظير الجزء السابق بالنسبة إلى محور التراتيب (لأن $k$ زوجية).
\end{itemize}
(انظر الرسم أعلاه.)

\pc حلول $k(x)=m$ هي فواصل نقط تقاطع $(C_k)$ مع المستقيم الأفقي $y=m$. على $]-1;0]$ تتزايد $k$ من $-\infty$ إلى $f(\alpha)$ ثم تتناقص إلى $k(0)=e$:
\begin{itemize}
\item $m<e$: حلان متعاكسان (أحدهما سالب في $]-1;\alpha[$ والآخر موجب).
\item $m=e$: ثلاثة حلول: $0$، وحل سالب، وحل موجب (متعاكسان).
\item $e<m<f(\alpha)$: أربعة حلول: حلان سالبان وحلان موجبان.
\item $m=f(\alpha)$: حلان متعاكسان هما $\alpha$ و $-\alpha$.
\item $m>f(\alpha)$: لا توجد حلول.
\end{itemize}
\end{hal}

% ===================== التمرين 10 =====================
\begin{tamrin}{بكالوريا 2017 -- الموضوع الأول}
نعتبر الدالة $f$ المعرفة على $D$ حيث $D=]-\infty;-1[\,\cup\,]1;+\infty[$ بـ: $f(x)=\dfrac23x+\ln\left(\dfrac{x-1}{x+1}\right)$،
و $(C_f)$ تمثيلها البياني في المستوي المنسوب إلى المعلم المتعامد والمتجانس $(O;\vec i,\vec j)$.
\QQ{1} بيّن أن الدالة $f$ فردية ثم فسّر ذلك بيانيا.
\QQ{2} أحسب النهايات التالية: $\lim\limits_{x\to 1^{+}}f(x)$، $\lim\limits_{x\to -1^{-}}f(x)$، $\lim\limits_{x\to -\infty}f(x)$ و $\lim\limits_{x\to +\infty}f(x)$.
استنتج أن $(C_f)$ يقبل مستقيمين مقاربين موازيين لحامل محور التراتيب.
\QQ{3} بيّن أنه من أجل كل عدد حقيقي $x$ من $D$: $f'(x)=\dfrac23\left(\dfrac{x^2+2}{x^2-1}\right)$.
استنتج اتجاه تغيّر الدالة $f$، ثم شكّل جدول تغيّراتها.
\QQ{4} بيّن أن المعادلة $f(x)=0$ تقبل في المجال $]1;+\infty[$ حلا وحيدا $\alpha$ حيث $1{,}8<\alpha<1{,}9$.
\QQ{5} بيّن أن المستقيم $(\Delta)$ ذا المعادلة: $y=\dfrac23x$ هو مستقيم مقارب مائل للمنحنى $(C_f)$ ثم أدرس وضعية المنحنى $(C_f)$ بالنسبة إلى المستقيم $(\Delta)$.
\QQ{6} أنشئ المنحنى $(C_f)$ والمستقيم $(\Delta)$.
\QQ{7} $m$ وسيط حقيقي، ناقش بيانيا حسب قيم الوسيط الحقيقي $m$ عدد حلول المعادلة:
\[(2-3|m|)x+3\ln\left(\frac{x-1}{x+1}\right)=0\]
\end{tamrin}

\begin{molahada}
في نص المجلة: «بيّن أن المعادلة $f(x)=0$ تقبل حلا وحيدا $\alpha$» دون ذكر المجال، بينما للمعادلة على $D$ حلان هما $\alpha$ و $-\alpha$ (لأن $f$ فردية). صحّحنا السؤال 4 بإضافة «في المجال $]1;+\infty[$».
\end{molahada}

\begin{hal}
\QQ{1} \textbf{المجال $D$ متناظر بالنسبة إلى $0$:} إذا كان $x\in D$ فإن $|x|>1$ ومنه $|-x|>1$ أي $-x\in D$.
ومن أجل كل $x\in D$:
\begin{align*}
f(-x)&=-\frac23x+\ln\left(\frac{-x-1}{-x+1}\right)=-\frac23x+\ln\left(\frac{x+1}{x-1}\right)\\
&=-\frac23x-\ln\left(\frac{x-1}{x+1}\right)=-f(x)
\end{align*}
(لأن $\ln\frac1a=-\ln a$). إذن الدالة $f$ \textbf{فردية}، \textbf{بيانيا:} $(C_f)$ متناظر بالنسبة إلى المبدأ $O$.

\QQ{2}
\begin{itemize}
\item لما $x\to1^{+}$: $x-1\to0^{+}$ و $x+1\to2$، إذن $\dfrac{x-1}{x+1}\to0^{+}$ و $\ln\left(\dfrac{x-1}{x+1}\right)\to-\infty$، ومنه \bx{\lim\limits_{x\to 1^{+}}f(x)=-\infty}.
\item لما $x\to-1^{-}$: $x-1\to-2$ و $x+1\to0^{-}$، إذن $\dfrac{x-1}{x+1}\to+\infty$ و $\ln\left(\dfrac{x-1}{x+1}\right)\to+\infty$، ومنه \bx{\lim\limits_{x\to -1^{-}}f(x)=+\infty}.
\item لما $x\to\pm\infty$: $\dfrac{x-1}{x+1}\to1$ إذن $\ln\left(\dfrac{x-1}{x+1}\right)\to0$، ومنه:
\[\bx{\lim_{x\to +\infty}f(x)=+\infty}\qquad;\qquad\bx{\lim_{x\to -\infty}f(x)=-\infty}\]
\end{itemize}
\textbf{الاستنتاج:} المستقيمان ذوا المعادلتين $x=1$ و $x=-1$ مقاربان لـ $(C_f)$ موازيان لحامل محور التراتيب.

\QQ{3} نضع $u(x)=\dfrac{x-1}{x+1}>0$ على $D$، فـ $u'(x)=\dfrac{(x+1)-(x-1)}{(x+1)^2}=\dfrac{2}{(x+1)^2}$، و:
\[\frac{u'(x)}{u(x)}=\frac{2}{(x+1)^2}\times\frac{x+1}{x-1}=\frac{2}{(x-1)(x+1)}=\frac{2}{x^2-1}\]
ومنه:
\[f'(x)=\frac23+\frac{2}{x^2-1}=\frac{2(x^2-1)+6}{3(x^2-1)}=\frac{2x^2+4}{3(x^2-1)}=\bx{\frac23\left(\frac{x^2+2}{x^2-1}\right)}\]
من أجل $x\in D$: $x^2>1$ إذن $x^2-1>0$، و $x^2+2>0$؛ ومنه $f'(x)>0$: $f$ \textbf{متزايدة تماما} على كل من $]-\infty;-1[$ و $]1;+\infty[$.
\begin{TAB}
\tkzTabInit[lgt=2,espcl=3]{$x$/1,$f'(x)$/1,$f(x)$/2}{$-\infty$,$-1$,$1$,$+\infty$}
\tkzTabLine{,+,d,h,d,+,}
\tkzTabVar{-/$-\infty$,+DH/$+\infty$,D-/$-\infty$,+/$+\infty$}
\end{TAB}

\QQ{4} على $]1;+\infty[$: $f$ مستمرة ومتزايدة تماما وتأخذ قيمها في $]-\infty;+\infty[$، إذن المعادلة $f(x)=0$ تقبل حلا وحيدا $\alpha$ في هذا المجال. و:
\[f(1{,}8)=1{,}2+\ln\frac{0{,}8}{2{,}8}\approx-0{,}053<0\qquad;\qquad f(1{,}9)=\frac{3{,}8}{3}+\ln\frac{0{,}9}{2{,}9}\approx0{,}097>0\]
إذن \bx{1{,}8<\alpha<1{,}9} ($\alpha\approx1{,}83$). (وبما أن $f$ فردية فالعدد $-\alpha$ هو الحل الوحيد على $]-\infty;-1[$.)

\QQ{5} $f(x)-\dfrac23x=\ln\left(\dfrac{x-1}{x+1}\right)$ و $\lim\limits_{x\to\pm\infty}\ln\left(\dfrac{x-1}{x+1}\right)=\ln1=0$، إذن \bx{\lim\limits_{x\to\pm\infty}\left[f(x)-\tfrac23x\right]=0}؛ ومنه $(\Delta):y=\frac23x$ مقارب مائل لـ $(C_f)$ عند $+\infty$ وعند $-\infty$.

\textbf{الوضعية:} $\dfrac{x-1}{x+1}-1=\dfrac{-2}{x+1}$.
\begin{itemize}
\item على $]1;+\infty[$: $x+1>0$ إذن $\dfrac{x-1}{x+1}<1$ ومنه $\ln\left(\dfrac{x-1}{x+1}\right)<0$: $(C_f)$ \textbf{تحت} $(\Delta)$.
\item على $]-\infty;-1[$: $x+1<0$ إذن $\dfrac{x-1}{x+1}>1$ ومنه $\ln\left(\dfrac{x-1}{x+1}\right)>0$: $(C_f)$ \textbf{فوق} $(\Delta)$.
\end{itemize}

\QQ{6} \textbf{الرسم:}
\begin{GR}{-5}{5}{-4}{4}{0.85}
 \Ticks{-4,-3,-2,-1,1,2,3,4}{-3,-2,-1,1,2,3}
 \draw[cThm,dashed,thick] (-1,-4)--(-1,4) (1,-4)--(1,4);
 \draw[cRem,thick] (-5,-3.333)--(5,3.333) node[pos=0.95,above left]{$(\Delta)$};
 \draw[cDef,very thick,domain=1.012:5,samples=200,smooth] plot (\x,{2*\x/3+ln((\x-1)/(\x+1))});
 \draw[cDef,very thick,domain=-5:-1.012,samples=200,smooth] plot (\x,{2*\x/3+ln((\x-1)/(\x+1))});
 \fill[cDef] (1.834,0) circle (1.8pt) (-1.834,0) circle (1.8pt);
 \node[cDef,above left,font=\scriptsize] at (1.834,0){$\alpha$};
 \node[cDef,below right,font=\scriptsize] at (-1.834,0){$-\alpha$};
 \node[cDef] at (3.6,1.2){$(C_f)$};
 \node[cDef] at (-3.6,-1.2){$(C_f)$};
\end{GR}

\QQ{7} من أجل $x\in D$:
\begin{align*}
(2-3|m|)x+3\ln\left(\frac{x-1}{x+1}\right)=0&\iff 2x+3\ln\left(\frac{x-1}{x+1}\right)=3|m|x\\
&\iff \frac23x+\ln\left(\frac{x-1}{x+1}\right)=|m|x\iff f(x)=|m|x
\end{align*}
حلول المعادلة هي فواصل نقط تقاطع $(C_f)$ مع المستقيم $(D_m):y=|m|x$ الذي يشمل المبدأ $O$ وميله $|m|\ge0$.
\begin{itemize}
\item إذا كان $0\le|m|<\dfrac23$، أي $-\dfrac23<m<\dfrac23$: ميل $(D_m)$ أصغر من ميل $(\Delta)$، فيقطع $(D_m)$ الفرع الأيمن من $(C_f)$ في نقطة واحدة، ويقطع الفرع الأيسر في نقطة نظيرتها بالنسبة إلى $O$: \textbf{حلان متعاكسان}.
\item إذا كان $|m|\ge\dfrac23$، أي $m\le-\dfrac23$ أو $m\ge\dfrac23$: من أجل $x>1$ لدينا $f(x)<\frac23x\le|m|x$، ومن أجل $x<-1$ لدينا $f(x)>\frac23x\ge|m|x$؛ إذن \textbf{لا توجد حلول}.
\end{itemize}
\end{hal}

% ===================== التمرين: بكالوريا 2017 م1 استثنائية =====================
\begin{tamrin}{بكالوريا 2017 -- الموضوع الأول -- الدورة الاستثنائية}
\textbf{I)} نعتبر $f$ الدالة العددية المعرفة على المجال $\left]-\frac12;+\infty\right[$ كما يلي:
\[f(x)=\frac{1+2\ln(2x+1)}{(2x+1)^2}.\]
و $(C_f)$ تمثيلها البياني في المستوي المنسوب إلى المعلم المتعامد والمتجانس $(O;\vec i,\vec j)$.
\QQ{1} أحسب النهايتين: $\lim\limits_{x\to+\infty}f(x)$ ، $\lim\limits_{x\to-\frac12}f(x)$ ثم فسّر النتيجتين بيانيا.
\QQ{2}
\begin{sq}
\qa بيّن أنه من أجل كل عدد حقيقي $x$ من المجال $\left]-\frac12;+\infty\right[$: \ $f'(x)=\dfrac{-8\ln(2x+1)}{(2x+1)^3}$.
\qb أدرس اتجاه تغيّر الدالة $f$ ثم شكّل جدول تغيّراتها.
\end{sq}
\QQ{3} حلّ في المجال $\left]-\frac12;+\infty\right[$ المعادلة $f(x)=0$ ، ثم استنتج إشارة $f(x)$.
\QQ{4} بيّن أن المنحنى $(C_f)$ يقبل نقطة انعطاف $\omega$ يُطلب تعيين إحداثييها، ثم أنشئ $(C_f)$.

\textbf{II)} لتكن الدالة $g$ المعرفة على $\left]-\frac12;+\infty\right[$ كما يلي: \ $g(x)=2\big[-x+\ln(2x+1)\big]$.
\QQ{1}
\begin{sq}
\qa أدرس اتجاه تغيّر الدالة $g$.
\qb بيّن أن المعادلة $g(x)=0$ تقبل حلّين أحدهما معدوم والآخر $\alpha$ حيث: $1{,}2<\alpha<1{,}3$.
\qc استنتج إشارة $g(x)$.
\end{sq}
\end{tamrin}

\begin{hal}
\textbf{الجزء I}

\QQ{1} \textbf{حساب النهايتين.}

عند $+\infty$: نضع $t=2x+1$ ، لما $x\to+\infty$ فإن $t\to+\infty$ و
\[f(x)=\frac{1}{t^2}+2\cdot\frac{\ln t}{t}\cdot\frac{1}{t}.\]
لأن $\lim\limits_{t\to+\infty}\frac{1}{t^2}=0$ و $\lim\limits_{t\to+\infty}\frac{\ln t}{t}=0$ (التزايد المقارن)، إذن: \ \bx{\lim_{x\to+\infty}f(x)=0}.

عند $-\frac12$ (من اليمين): $2x+1\to0^+$ إذن $\ln(2x+1)\to-\infty$ ومنه $1+2\ln(2x+1)\to-\infty$ ، والمقام $(2x+1)^2\to0^+$. إذن: \ \bx{\lim_{x\to-\frac12}f(x)=-\infty}.

\textbf{التفسير البياني:} المستقيم ذو المعادلة $x=-\frac12$ مقارب عمودي لـ $(C_f)$ ، والمستقيم ذو المعادلة $y=0$ (حامل محور الفواصل) مقارب أفقي لـ $(C_f)$ بجوار $+\infty$.

\QQ{2}
\pa الدالة $f$ قابلة للاشتقاق على $\left]-\frac12;+\infty\right[$ ، وهي من الشكل $\frac uv$ مع
$u(x)=1+2\ln(2x+1)$ ، $u'(x)=\frac{4}{2x+1}$ ، $v(x)=(2x+1)^2$ ، $v'(x)=4(2x+1)$.
\begin{align*}
f'(x)&=\frac{\frac{4}{2x+1}(2x+1)^2-4(2x+1)\big[1+2\ln(2x+1)\big]}{(2x+1)^4}
=\frac{4(2x+1)\big[1-1-2\ln(2x+1)\big]}{(2x+1)^4}\\
&=\frac{-8\ln(2x+1)}{(2x+1)^3}.
\end{align*}
\pb لأن $2x+1>0$ فإن $(2x+1)^3>0$ ، إذن إشارة $f'(x)$ هي إشارة $-\ln(2x+1)$:
\begin{itemize}
\item إذا كان $-\frac12<x<0$ فإن $0<2x+1<1$ ومنه $\ln(2x+1)<0$ إذن $f'(x)>0$.
\item إذا كان $x>0$ فإن $2x+1>1$ ومنه $\ln(2x+1)>0$ إذن $f'(x)<0$ ، و $f'(0)=0$.
\end{itemize}
إذن: \bxt{$f$ متزايدة تماما على $\left]-\tfrac12;0\right]$ ومتناقصة تماما على $[0;+\infty[$}. مع $f(0)=\frac{1+0}{1}=1$.
\begin{TAB}
\tkzTabInit[lgt=2,espcl=3]{$x$/1,$f'(x)$/1,$f(x)$/2}{$-\frac12$,$0$,$+\infty$}
\tkzTabLine{d,+,z,-,}
\tkzTabVar{D-/$-\infty$,+/$1$,-/$0$}
\end{TAB}

\QQ{3} $f(x)=0 \iff 1+2\ln(2x+1)=0 \iff \ln(2x+1)=-\frac12 \iff 2x+1=e^{-\frac12}\iff x=\dfrac{e^{-\frac12}-1}{2}$.

وهو ينتمي إلى المجال لأن $e^{-\frac12}>0$. إذن: \bx{x_0=\dfrac{e^{-\frac12}-1}{2}\approx-0{,}2}.

\textbf{إشارة $f(x)$:} المقام موجب تماما، فإشارة $f(x)$ هي إشارة $1+2\ln(2x+1)$ وهي دالة متزايدة تماما تنعدم عند $x_0$:
\begin{TAB}
\tkzTabInit[lgt=2,espcl=3]{$x$/1,$f(x)$/1}{$-\frac12$,$x_0$,$+\infty$}
\tkzTabLine{d,-,z,+,}
\end{TAB}

\QQ{4} نحسب المشتقة الثانية. بوضع $u=2x+1$ ( $u'=2$ ) لدينا $f'(x)=\frac{-8\ln u}{u^3}$:
\[f''(x)=-8\cdot\frac{\frac{2}{u}\,u^3-\ln u\cdot 3u^2\cdot 2}{u^6}=-8\cdot\frac{2u^2(1-3\ln u)}{u^6}=\frac{16\big(3\ln(2x+1)-1\big)}{(2x+1)^4}.\]
إشارة $f''(x)$ هي إشارة $3\ln(2x+1)-1$ ، وهو ينعدم عند $\ln(2x+1)=\frac13$ أي $x=\frac{e^{\frac13}-1}{2}$ ، سالب قبلها وموجب بعدها.

إذن $f''$ تنعدم مغيّرة إشارتها، و $(C_f)$ يقبل نقطة انعطاف:
\[f\Big(\tfrac{e^{1/3}-1}{2}\Big)=\frac{1+\frac23}{e^{2/3}}=\frac53e^{-\frac23}\qquad\Longrightarrow\qquad
\bx{\omega\Big(\frac{e^{\frac13}-1}{2};\ \frac53e^{-\frac23}\Big)\approx(0{,}2;\ 0{,}86)}\]

\textbf{الإنشاء:}
\begin{GR}{-1}{3.2}{-3}{1.6}{2}
\draw[dashed,cThm] (-0.5,-3)--(-0.5,1.6);
\draw[cDef,thick] plot coordinates {(-0.2798,-3.2978) (-0.2710,-2.6778) (-0.2622,-2.1516) (-0.2535,-1.7033) (-0.2447,-1.3202) (-0.2359,-0.9919) (-0.2271,-0.7097) (-0.2184,-0.4668) (-0.2096,-0.2572) (-0.2008,-0.0761) (-0.1921,0.0805) (-0.1833,0.2161) (-0.1745,0.3335) (-0.1658,0.4352) (-0.1570,0.5233) (-0.1482,0.5996) (-0.1395,0.6655) (-0.1307,0.7223) (-0.1219,0.7713) (-0.1131,0.8133) (-0.1044,0.8493) (-0.0956,0.8799) (-0.0868,0.9058) (-0.0781,0.9275) (-0.0693,0.9456) (-0.0605,0.9604) (-0.0518,0.9723) (-0.0430,0.9818) (-0.0342,0.9889) (-0.0255,0.9941) (-0.0167,0.9976) (-0.0079,0.9995) (0.0009,1.0000) (0.0096,0.9993) (0.0184,0.9975) (0.0272,0.9948) (0.0359,0.9912) (0.0447,0.9869) (0.0535,0.9819) (0.0622,0.9764) (0.0710,0.9704) (0.0798,0.9639) (0.0886,0.9571) (0.0973,0.9499) (0.1061,0.9425) (0.1149,0.9348) (0.1236,0.9269) (0.1324,0.9188) (0.1412,0.9106) (0.1499,0.9023) (0.1587,0.8938) (0.1675,0.8854) (0.1762,0.8768) (0.1850,0.8682) (0.1938,0.8596) (0.2026,0.8510) (0.2113,0.8424) (0.2201,0.8339) (0.2289,0.8253) (0.2376,0.8168) (0.2464,0.8083) (0.2552,0.7999) (0.2639,0.7915) (0.2727,0.7832) (0.2815,0.7750) (0.2902,0.7668) (0.2990,0.7587) (0.3078,0.7507) (0.3166,0.7428) (0.3253,0.7349) (0.3341,0.7271) (0.3429,0.7194) (0.3516,0.7118) (0.3604,0.7043) (0.3692,0.6969) (0.3779,0.6895) (0.3867,0.6823) (0.3955,0.6751) (0.4043,0.6681) (0.4130,0.6611) (0.4218,0.6542) (0.4306,0.6474) (0.4393,0.6407) (0.4481,0.6340) (0.4569,0.6275) (0.4656,0.6210) (0.4744,0.6147) (0.4832,0.6084) (0.4919,0.6022) (0.5007,0.5961) (0.5095,0.5900) (0.5183,0.5841) (0.5270,0.5782) (0.5358,0.5724) (0.5446,0.5667) (0.5533,0.5611) (0.5621,0.5556) (0.5709,0.5501) (0.5796,0.5447) (0.5884,0.5394) (0.5972,0.5341) (0.6059,0.5289) (0.6147,0.5238) (0.6235,0.5188) (0.6323,0.5138) (0.6410,0.5089) (0.6498,0.5040) (0.6586,0.4993) (0.6673,0.4946) (0.6761,0.4899) (0.6849,0.4853) (0.6936,0.4808) (0.7024,0.4764) (0.7112,0.4720) (0.7199,0.4676) (0.7287,0.4634) (0.7375,0.4591) (0.7463,0.4550) (0.7550,0.4509) (0.7638,0.4468) (0.7726,0.4428) (0.7813,0.4389) (0.7901,0.4350) (0.7989,0.4311) (0.8076,0.4273) (0.8164,0.4236) (0.8252,0.4199) (0.8340,0.4162) (0.8427,0.4126) (0.8515,0.4091) (0.8603,0.4056) (0.8690,0.4021) (0.8778,0.3987) (0.8866,0.3953) (0.8953,0.3920) (0.9041,0.3887) (0.9129,0.3854) (0.9216,0.3822) (0.9304,0.3790) (0.9392,0.3759) (0.9480,0.3728) (0.9567,0.3698) (0.9655,0.3668) (0.9743,0.3638) (0.9830,0.3608) (0.9918,0.3579) (1.0006,0.3551) (1.0093,0.3522) (1.0181,0.3494) (1.0269,0.3467) (1.0356,0.3439) (1.0444,0.3412) (1.0532,0.3386) (1.0620,0.3359) (1.0707,0.3333) (1.0795,0.3307) (1.0883,0.3282) (1.0970,0.3257) (1.1058,0.3232) (1.1146,0.3207) (1.1233,0.3183) (1.1321,0.3159) (1.1409,0.3135) (1.1497,0.3112) (1.1584,0.3089) (1.1672,0.3066) (1.1760,0.3043) (1.1847,0.3021) (1.1935,0.2999) (1.2023,0.2977) (1.2110,0.2955) (1.2198,0.2934) (1.2286,0.2912) (1.2373,0.2891) (1.2461,0.2871) (1.2549,0.2850) (1.2637,0.2830) (1.2724,0.2810) (1.2812,0.2790) (1.2900,0.2771) (1.2987,0.2751) (1.3075,0.2732) (1.3163,0.2713) (1.3250,0.2694) (1.3338,0.2676) (1.3426,0.2657) (1.3513,0.2639) (1.3601,0.2621) (1.3689,0.2603) (1.3777,0.2586) (1.3864,0.2568) (1.3952,0.2551) (1.4040,0.2534) (1.4127,0.2517) (1.4215,0.2500) (1.4303,0.2484) (1.4390,0.2467) (1.4478,0.2451) (1.4566,0.2435) (1.4654,0.2419) (1.4741,0.2403) (1.4829,0.2388) (1.4917,0.2372) (1.5004,0.2357) (1.5092,0.2342) (1.5180,0.2327) (1.5267,0.2312) (1.5355,0.2298) (1.5443,0.2283) (1.5530,0.2269) (1.5618,0.2254) (1.5706,0.2240) (1.5794,0.2226) (1.5881,0.2212) (1.5969,0.2199) (1.6057,0.2185) (1.6144,0.2172) (1.6232,0.2158) (1.6320,0.2145) (1.6407,0.2132) (1.6495,0.2119) (1.6583,0.2106) (1.6670,0.2094) (1.6758,0.2081) (1.6846,0.2069) (1.6934,0.2056) (1.7021,0.2044) (1.7109,0.2032) (1.7197,0.2020) (1.7284,0.2008) (1.7372,0.1996) (1.7460,0.1985) (1.7547,0.1973) (1.7635,0.1962) (1.7723,0.1950) (1.7811,0.1939) (1.7898,0.1928) (1.7986,0.1917) (1.8074,0.1906) (1.8161,0.1895) (1.8249,0.1884) (1.8337,0.1873) (1.8424,0.1863) (1.8512,0.1852) (1.8600,0.1842) (1.8687,0.1832) (1.8775,0.1821) (1.8863,0.1811) (1.8951,0.1801) (1.9038,0.1791) (1.9126,0.1781) (1.9214,0.1772) (1.9301,0.1762) (1.9389,0.1752) (1.9477,0.1743) (1.9564,0.1733) (1.9652,0.1724) (1.9740,0.1715) (1.9827,0.1705) (1.9915,0.1696) (2.0003,0.1687) (2.0091,0.1678) (2.0178,0.1669) (2.0266,0.1661) (2.0354,0.1652) (2.0441,0.1643) (2.0529,0.1634) (2.0617,0.1626) (2.0704,0.1617) (2.0792,0.1609) (2.0880,0.1601) (2.0967,0.1592) (2.1055,0.1584) (2.1143,0.1576) (2.1231,0.1568) (2.1318,0.1560) (2.1406,0.1552) (2.1494,0.1544) (2.1581,0.1536) (2.1669,0.1528) (2.1757,0.1521) (2.1844,0.1513) (2.1932,0.1505) (2.2020,0.1498) (2.2108,0.1490) (2.2195,0.1483) (2.2283,0.1476) (2.2371,0.1468) (2.2458,0.1461) (2.2546,0.1454) (2.2634,0.1447) (2.2721,0.1440) (2.2809,0.1433) (2.2897,0.1426) (2.2984,0.1419) (2.3072,0.1412) (2.3160,0.1405) (2.3248,0.1398) (2.3335,0.1392) (2.3423,0.1385) (2.3511,0.1378) (2.3598,0.1372) (2.3686,0.1365) (2.3774,0.1359) (2.3861,0.1352) (2.3949,0.1346) (2.4037,0.1340) (2.4124,0.1333) (2.4212,0.1327) (2.4300,0.1321) (2.4388,0.1315) (2.4475,0.1309) (2.4563,0.1303) (2.4651,0.1297) (2.4738,0.1291) (2.4826,0.1285) (2.4914,0.1279) (2.5001,0.1273) (2.5089,0.1267) (2.5177,0.1262) (2.5265,0.1256) (2.5352,0.1250) (2.5440,0.1245) (2.5528,0.1239) (2.5615,0.1233) (2.5703,0.1228) (2.5791,0.1222) (2.5878,0.1217) (2.5966,0.1212) (2.6054,0.1206) (2.6141,0.1201) (2.6229,0.1196) (2.6317,0.1190) (2.6405,0.1185) (2.6492,0.1180) (2.6580,0.1175) (2.6668,0.1170) (2.6755,0.1165) (2.6843,0.1159) (2.6931,0.1154) (2.7018,0.1149) (2.7106,0.1145) (2.7194,0.1140) (2.7281,0.1135) (2.7369,0.1130) (2.7457,0.1125) (2.7545,0.1120) (2.7632,0.1116) (2.7720,0.1111) (2.7808,0.1106) (2.7895,0.1102) (2.7983,0.1097) (2.8071,0.1092) (2.8158,0.1088) (2.8246,0.1083) (2.8334,0.1079) (2.8422,0.1074) (2.8509,0.1070) (2.8597,0.1065) (2.8685,0.1061) (2.8772,0.1057) (2.8860,0.1052) (2.8948,0.1048) (2.9035,0.1044) (2.9123,0.1039) (2.9211,0.1035) (2.9298,0.1031) (2.9386,0.1027) (2.9474,0.1023) (2.9562,0.1019) (2.9649,0.1014) (2.9737,0.1010) (2.9825,0.1006) (2.9912,0.1002) (3.0000,0.0998)};
\fill (0,1) circle (1.5pt) node[above right]{$1$};
\fill[cThm] (0.1978,0.8557) circle (1.5pt) node[above right]{$\omega$};
\fill (-0.1967,0) circle (1.5pt);
\draw[cRem,thick] (-0.4,1)--(0.6,1);
\Ticks{-1,1,2,3}{-3,-2,-1,1}
\node[cDef] at (2.2,0.5){$(C_f)$};
\node[cThm] at (-0.75,-2.6){$x=-\frac12$};
\end{GR}

\textbf{الجزء II}

\QQ{1}
\pa $g$ قابلة للاشتقاق على $\left]-\frac12;+\infty\right[$ و:
\[g'(x)=2\Big[-1+\frac{2}{2x+1}\Big]=2\cdot\frac{-2x-1+2}{2x+1}=\frac{2(1-2x)}{2x+1}.\]
لأن $2x+1>0$ فإن إشارة $g'(x)$ هي إشارة $1-2x$.
النهايات: $\lim\limits_{x\to-\frac12}g(x)=-\infty$ (لأن $\ln(2x+1)\to-\infty$)، و
$g(x)=2x\Big[-1+\frac{\ln(2x+1)}{2x+1}\cdot\frac{2x+1}{x}\Big]$ مع $\frac{\ln(2x+1)}{2x+1}\to0$ و $\frac{2x+1}{x}\to2$ ، إذن $\lim\limits_{x\to+\infty}g(x)=-\infty$.

إذن \bxt{$g$ متزايدة تماما على $\left]-\tfrac12;\tfrac12\right]$ ومتناقصة تماما على $\left[\tfrac12;+\infty\right[$} ، مع $g\big(\frac12\big)=2\ln2-1\approx0{,}39$.
\begin{TAB}
\tkzTabInit[lgt=2,espcl=3]{$x$/1,$g'(x)$/1,$g(x)$/2}{$-\frac12$,$\frac12$,$+\infty$}
\tkzTabLine{d,+,z,-,}
\tkzTabVar{D-/$-\infty$,+/$2\ln2-1$,-/$-\infty$}
\end{TAB}

\pb $g(0)=2\big[0+\ln1\big]=0$ ، إذن $0$ حلّ للمعادلة، وهو الحل الوحيد في $\left]-\frac12;\frac12\right]$ لأن $g$ متزايدة تماما عليه.

على $\left[\frac12;+\infty\right[$: $g$ مستمرة ومتناقصة تماما، و
$g(1{,}2)\approx0{,}048>0$ ، $g(1{,}3)\approx-0{,}038<0$ ، إذن $g(1{,}2)\times g(1{,}3)<0$ ،
وحسب مبرهنة القيم المتوسطة المعادلة تقبل حلا وحيدا $\alpha$ حيث \bx{1{,}2<\alpha<1{,}3}.

\pc من جدول التغيرات:
\begin{TAB}
\tkzTabInit[lgt=2,espcl=2.5]{$x$/1,$g(x)$/1}{$-\frac12$,$0$,$\alpha$,$+\infty$}
\tkzTabLine{d,-,z,+,z,-,}
\end{TAB}
\end{hal}

% ===================== بكالوريا 2018 م2 =====================
\begin{tamrin}{بكالوريا 2018 -- الموضوع الثاني}
\textbf{I)} $g$ الدالة العددية ذات المتغير الحقيقي $x$ المعرفة على $]0;+\infty[$ بـ: \ $g(x)=\dfrac1x-(\ln x)^2-\ln x-1$.

$(C_g)$ المنحنى البياني الممثل لها كما هو مبيّن في الشكل المقابل:
\begin{GR}{-0.5}{3.5}{-2.6}{4.6}{0.75}
\draw[cDef,thick] plot coordinates {(0.1374,4.3218) (0.1459,4.0728) (0.1544,3.8538) (0.1629,3.6598) (0.1714,3.4867) (0.1799,3.3311) (0.1884,3.1905) (0.1969,3.0627) (0.2054,2.9459) (0.2139,2.8387) (0.2224,2.7398) (0.2309,2.6481) (0.2394,2.5629) (0.2479,2.4834) (0.2564,2.4089) (0.2649,2.3389) (0.2734,2.2728) (0.2819,2.2104) (0.2904,2.1513) (0.2989,2.0950) (0.3074,2.0414) (0.3159,1.9902) (0.3244,1.9412) (0.3329,1.8942) (0.3414,1.8491) (0.3498,1.8056) (0.3583,1.7637) (0.3668,1.7231) (0.3753,1.6839) (0.3838,1.6459) (0.3923,1.6091) (0.4008,1.5733) (0.4093,1.5384) (0.4178,1.5045) (0.4263,1.4714) (0.4348,1.4391) (0.4433,1.4075) (0.4518,1.3766) (0.4603,1.3464) (0.4688,1.3168) (0.4773,1.2877) (0.4858,1.2592) (0.4943,1.2312) (0.5028,1.2037) (0.5113,1.1767) (0.5198,1.1501) (0.5283,1.1239) (0.5368,1.0981) (0.5453,1.0726) (0.5538,1.0476) (0.5623,1.0228) (0.5708,0.9984) (0.5792,0.9743) (0.5877,0.9504) (0.5962,0.9269) (0.6047,0.9036) (0.6132,0.8806) (0.6217,0.8578) (0.6302,0.8353) (0.6387,0.8130) (0.6472,0.7909) (0.6557,0.7690) (0.6642,0.7473) (0.6727,0.7258) (0.6812,0.7045) (0.6897,0.6834) (0.6982,0.6625) (0.7067,0.6417) (0.7152,0.6211) (0.7237,0.6006) (0.7322,0.5803) (0.7407,0.5602) (0.7492,0.5402) (0.7577,0.5203) (0.7662,0.5006) (0.7747,0.4810) (0.7832,0.4616) (0.7917,0.4422) (0.8002,0.4230) (0.8086,0.4039) (0.8171,0.3849) (0.8256,0.3661) (0.8341,0.3473) (0.8426,0.3287) (0.8511,0.3101) (0.8596,0.2917) (0.8681,0.2733) (0.8766,0.2551) (0.8851,0.2369) (0.8936,0.2189) (0.9021,0.2009) (0.9106,0.1831) (0.9191,0.1653) (0.9276,0.1476) (0.9361,0.1300) (0.9446,0.1124) (0.9531,0.0950) (0.9616,0.0776) (0.9701,0.0603) (0.9786,0.0431) (0.9871,0.0259) (0.9956,0.0089) (1.0041,-0.0081) (1.0126,-0.0250) (1.0211,-0.0419) (1.0295,-0.0587) (1.0380,-0.0754) (1.0465,-0.0920) (1.0550,-0.1086) (1.0635,-0.1251) (1.0720,-0.1416) (1.0805,-0.1580) (1.0890,-0.1743) (1.0975,-0.1906) (1.1060,-0.2068) (1.1145,-0.2229) (1.1230,-0.2390) (1.1315,-0.2550) (1.1400,-0.2710) (1.1485,-0.2869) (1.1570,-0.3028) (1.1655,-0.3186) (1.1740,-0.3343) (1.1825,-0.3500) (1.1910,-0.3657) (1.1995,-0.3813) (1.2080,-0.3968) (1.2165,-0.4123) (1.2250,-0.4277) (1.2335,-0.4431) (1.2420,-0.4585) (1.2505,-0.4737) (1.2589,-0.4890) (1.2674,-0.5042) (1.2759,-0.5193) (1.2844,-0.5344) (1.2929,-0.5495) (1.3014,-0.5645) (1.3099,-0.5795) (1.3184,-0.5944) (1.3269,-0.6092) (1.3354,-0.6241) (1.3439,-0.6389) (1.3524,-0.6536) (1.3609,-0.6683) (1.3694,-0.6830) (1.3779,-0.6976) (1.3864,-0.7121) (1.3949,-0.7267) (1.4034,-0.7412) (1.4119,-0.7556) (1.4204,-0.7700) (1.4289,-0.7844) (1.4374,-0.7987) (1.4459,-0.8130) (1.4544,-0.8273) (1.4629,-0.8415) (1.4714,-0.8557) (1.4798,-0.8698) (1.4883,-0.8839) (1.4968,-0.8980) (1.5053,-0.9120) (1.5138,-0.9260) (1.5223,-0.9400) (1.5308,-0.9539) (1.5393,-0.9678) (1.5478,-0.9816) (1.5563,-0.9954) (1.5648,-1.0092) (1.5733,-1.0229) (1.5818,-1.0367) (1.5903,-1.0503) (1.5988,-1.0640) (1.6073,-1.0776) (1.6158,-1.0912) (1.6243,-1.1047) (1.6328,-1.1182) (1.6413,-1.1317) (1.6498,-1.1451) (1.6583,-1.1585) (1.6668,-1.1719) (1.6753,-1.1853) (1.6838,-1.1986) (1.6923,-1.2119) (1.7008,-1.2251) (1.7092,-1.2384) (1.7177,-1.2515) (1.7262,-1.2647) (1.7347,-1.2778) (1.7432,-1.2909) (1.7517,-1.3040) (1.7602,-1.3171) (1.7687,-1.3301) (1.7772,-1.3431) (1.7857,-1.3560) (1.7942,-1.3689) (1.8027,-1.3818) (1.8112,-1.3947) (1.8197,-1.4075) (1.8282,-1.4203) (1.8367,-1.4331) (1.8452,-1.4459) (1.8537,-1.4586) (1.8622,-1.4713) (1.8707,-1.4840) (1.8792,-1.4966) (1.8877,-1.5092) (1.8962,-1.5218) (1.9047,-1.5344) (1.9132,-1.5469) (1.9217,-1.5595) (1.9302,-1.5719) (1.9386,-1.5844) (1.9471,-1.5968) (1.9556,-1.6092) (1.9641,-1.6216) (1.9726,-1.6340) (1.9811,-1.6463) (1.9896,-1.6586) (1.9981,-1.6709) (2.0066,-1.6831) (2.0151,-1.6954) (2.0236,-1.7076) (2.0321,-1.7198) (2.0406,-1.7319) (2.0491,-1.7440) (2.0576,-1.7561) (2.0661,-1.7682) (2.0746,-1.7803) (2.0831,-1.7923) (2.0916,-1.8043) (2.1001,-1.8163) (2.1086,-1.8283) (2.1171,-1.8402) (2.1256,-1.8521) (2.1341,-1.8640) (2.1426,-1.8759) (2.1511,-1.8878) (2.1595,-1.8996) (2.1680,-1.9114) (2.1765,-1.9232) (2.1850,-1.9349) (2.1935,-1.9467) (2.2020,-1.9584) (2.2105,-1.9701) (2.2190,-1.9817) (2.2275,-1.9934) (2.2360,-2.0050) (2.2445,-2.0166) (2.2530,-2.0282) (2.2615,-2.0398) (2.2700,-2.0513) (2.2785,-2.0628) (2.2870,-2.0743) (2.2955,-2.0858) (2.3040,-2.0972) (2.3125,-2.1087) (2.3210,-2.1201) (2.3295,-2.1315) (2.3380,-2.1428) (2.3465,-2.1542) (2.3550,-2.1655) (2.3635,-2.1768) (2.3720,-2.1881) (2.3805,-2.1994) (2.3889,-2.2106) (2.3974,-2.2219) (2.4059,-2.2331) (2.4144,-2.2443) (2.4229,-2.2554) (2.4314,-2.2666) (2.4399,-2.2777) (2.4484,-2.2888) (2.4569,-2.2999) (2.4654,-2.3110) (2.4739,-2.3221) (2.4824,-2.3331) (2.4909,-2.3441) (2.4994,-2.3551) (2.5079,-2.3661) (2.5164,-2.3770) (2.5249,-2.3880) (2.5334,-2.3989) (2.5419,-2.4098) (2.5504,-2.4207) (2.5589,-2.4316) (2.5674,-2.4424) (2.5759,-2.4532) (2.5844,-2.4640) (2.5929,-2.4748) (2.6014,-2.4856) (2.6098,-2.4964) (2.6183,-2.5071) (2.6268,-2.5178) (2.6353,-2.5285) (2.6438,-2.5392) (2.6523,-2.5499) (2.6608,-2.5605) (2.6693,-2.5712) (2.6778,-2.5818) (2.6863,-2.5924)};
\Ticks{1,2,3}{-2,-1,1,2,3,4}
\node[cDef] at (0.9,3.6){$(C_g)$};
\end{GR}
\par أحسب $g(1)$ ثم استنتج بيانيا إشارة $g(x)$.

\textbf{II)} $f$ الدالة العددية ذات المتغير الحقيقي $x$ المعرفة على $]0;+\infty[$ بـ: \ $f(x)=\dfrac{1+\ln x}{1+x\ln x}$ ،
و $(C_f)$ تمثيلها البياني في المستوي المنسوب إلى المعلم المتعامد والمتجانس $(O;\vec i,\vec j)$.
\QQ{1} أحسب $\lim\limits_{x\to0}f(x)$ وبيّن أن: $\lim\limits_{x\to+\infty}f(x)=0$ ، ثم فسّر النتيجتين بيانيا.
\QQ{2}
\begin{sq}
\qa بيّن أنه من أجل كل عدد حقيقي $x$ من المجال $]0;+\infty[$: \ $f'(x)=\dfrac{g(x)}{(1+x\ln x)^2}$.
\qb استنتج اتجاه تغيّر الدالة $f$ وشكّل جدول تغيّراتها.
\end{sq}
\QQ{3} بيّن أن $y=\Big(\dfrac{e^2}{e-1}\Big)x-\dfrac{e}{e-1}$ هي معادلة لـ $(T)$ مماس المنحنى $(C_f)$ في نقطة تقاطعه مع حامل محور الفواصل، ثم أرسم المماس $(T)$ والمنحنى $(C_f)$.
\QQ{4} عيّن بيانيا قيم الوسيط الحقيقي $m$ بحيث تقبل المعادلة $(e-1)f(x)=e^2x-me$ حلّين متمايزين.
\end{tamrin}

\begin{hal}
\textbf{الجزء I:} $g(1)=\frac11-(\ln1)^2-\ln1-1=1-0-0-1$ ، إذن \bx{g(1)=0}.

من الشكل: $(C_g)$ يقطع محور الفواصل في النقطة ذات الفاصلة $1$ ، ويقع فوقه على $]0;1[$ وتحته على $]1;+\infty[$:
\begin{TAB}
\tkzTabInit[lgt=2,espcl=3]{$x$/1,$g(x)$/1}{$0$,$1$,$+\infty$}
\tkzTabLine{d,+,z,-,}
\end{TAB}

\textbf{الجزء II}

\textbf{ملاحظة تمهيدية:} المقام لا ينعدم: نضع $\varphi(x)=x\ln x$ ، $\varphi'(x)=\ln x+1$ تنعدم عند $\frac1e$ ، فـ $\varphi$ تبلغ قيمتها الحدية الصغرى $\varphi\big(\frac1e\big)=-\frac1e$ ، إذن $1+x\ln x\geqslant1-\frac1e>0$.

\QQ{1} عند $0$: $\lim\limits_{x\to0}(1+\ln x)=-\infty$ و $\lim\limits_{x\to0}x\ln x=0$ ومنه $\lim\limits_{x\to0}(1+x\ln x)=1$ ، إذن \bx{\lim_{x\to0}f(x)=-\infty}.

عند $+\infty$: نقسم البسط والمقام على $x\ln x$ (حيث $x>1$):
\[f(x)=\frac{\dfrac{1}{x\ln x}+\dfrac1x}{\dfrac{1}{x\ln x}+1}\ \xrightarrow[x\to+\infty]{}\ \frac{0+0}{0+1}=0 ,\qquad\text{إذن}\quad \bx{\lim_{x\to+\infty}f(x)=0}.\]
\textbf{التفسير البياني:} المستقيم ذو المعادلة $x=0$ مقارب عمودي لـ $(C_f)$ ، والمستقيم $y=0$ مقارب أفقي لـ $(C_f)$ بجوار $+\infty$.

\QQ{2}
\pa $f=\frac uv$ مع $u=1+\ln x$ ، $u'=\frac1x$ ، $v=1+x\ln x$ ، $v'=\ln x+1$:
\begin{align*}
f'(x)&=\frac{\frac1x(1+x\ln x)-(1+\ln x)(1+\ln x)}{(1+x\ln x)^2}
=\frac{\frac1x+\ln x-1-2\ln x-(\ln x)^2}{(1+x\ln x)^2}\\
&=\frac{\frac1x-(\ln x)^2-\ln x-1}{(1+x\ln x)^2}=\frac{g(x)}{(1+x\ln x)^2}.
\end{align*}
\pb المقام موجب تماما، إذن إشارة $f'(x)$ هي إشارة $g(x)$: \bxt{$f$ متزايدة تماما على $]0;1]$ ومتناقصة تماما على $[1;+\infty[$} ، و $f(1)=\frac{1+0}{1+0}=1$.
\begin{TAB}
\tkzTabInit[lgt=2,espcl=3]{$x$/1,$f'(x)$/1,$f(x)$/2}{$0$,$1$,$+\infty$}
\tkzTabLine{d,+,z,-,}
\tkzTabVar{D-/$-\infty$,+/$1$,-/$0$}
\end{TAB}

\QQ{3} نقطة التقاطع مع محور الفواصل: $f(x)=0\iff1+\ln x=0\iff x=e^{-1}=\frac1e$.

$g\big(\frac1e\big)=e-(-1)^2-(-1)-1=e-1$ و $1+\frac1e\ln\frac1e=1-\frac1e=\frac{e-1}{e}$ ، إذن:
\[f'\Big(\frac1e\Big)=\frac{e-1}{\big(\frac{e-1}{e}\big)^2}=\frac{e^2}{e-1}.\]
معادلة المماس: $y=f'\big(\frac1e\big)\big(x-\frac1e\big)+f\big(\frac1e\big)=\frac{e^2}{e-1}\big(x-\frac1e\big)+0$ ، أي:
\[\bx{(T):\ y=\Big(\frac{e^2}{e-1}\Big)x-\frac{e}{e-1}}\qquad\big(\text{ميله}\approx4{,}3\big).\]
\begin{GR}{-0.5}{4.2}{-3.2}{1.8}{2}
\draw[cThm,thick] (0,-1.582)--(1.1,3.15);
\draw[cDef,thick] plot coordinates {(0.0142,-3.4667) (0.0148,-3.4282) (0.0154,-3.3896) (0.0161,-3.3510) (0.0168,-3.3123) (0.0176,-3.2736) (0.0184,-3.2348) (0.0192,-3.1959) (0.0200,-3.1570) (0.0210,-3.1181) (0.0210,-3.1141) (0.0219,-3.0791) (0.0229,-3.0400) (0.0239,-3.0009) (0.0250,-2.9617) (0.0261,-2.9225) (0.0273,-2.8832) (0.0286,-2.8439) (0.0299,-2.8044) (0.0311,-2.7704) (0.0313,-2.7649) (0.0327,-2.7254) (0.0342,-2.6857) (0.0358,-2.6459) (0.0374,-2.6061) (0.0391,-2.5661) (0.0409,-2.5260) (0.0411,-2.5229) (0.0428,-2.4858) (0.0448,-2.4454) (0.0469,-2.4049) (0.0491,-2.3642) (0.0511,-2.3275) (0.0513,-2.3234) (0.0537,-2.2823) (0.0562,-2.2410) (0.0588,-2.1995) (0.0611,-2.1645) (0.0616,-2.1577) (0.0645,-2.1156) (0.0675,-2.0733) (0.0706,-2.0306) (0.0712,-2.0234) (0.0739,-1.9876) (0.0774,-1.9442) (0.0810,-1.9004) (0.0812,-1.8980) (0.0912,-1.7844) (0.1012,-1.6799) (0.1112,-1.5826) (0.1213,-1.4913) (0.1313,-1.4048) (0.1413,-1.3224) (0.1513,-1.2436) (0.1614,-1.1679) (0.1714,-1.0948) (0.1814,-1.0242) (0.1914,-0.9557) (0.2015,-0.8892) (0.2115,-0.8246) (0.2215,-0.7616) (0.2315,-0.7003) (0.2415,-0.6405) (0.2516,-0.5822) (0.2616,-0.5252) (0.2716,-0.4697) (0.2816,-0.4154) (0.2917,-0.3624) (0.3017,-0.3107) (0.3117,-0.2603) (0.3217,-0.2111) (0.3317,-0.1631) (0.3418,-0.1163) (0.3518,-0.0707) (0.3618,-0.0263) (0.3718,0.0169) (0.3819,0.0590) (0.3919,0.0999) (0.4019,0.1396) (0.4119,0.1782) (0.4219,0.2156) (0.4320,0.2520) (0.4420,0.2872) (0.4520,0.3213) (0.4620,0.3543) (0.4721,0.3862) (0.4821,0.4171) (0.4921,0.4469) (0.5021,0.4756) (0.5122,0.5034) (0.5222,0.5301) (0.5322,0.5558) (0.5422,0.5806) (0.5522,0.6044) (0.5623,0.6273) (0.5723,0.6493) (0.5823,0.6703) (0.5923,0.6905) (0.6024,0.7098) (0.6124,0.7283) (0.6224,0.7460) (0.6324,0.7629) (0.6424,0.7790) (0.6525,0.7943) (0.6625,0.8089) (0.6725,0.8228) (0.6825,0.8360) (0.6926,0.8485) (0.7026,0.8604) (0.7126,0.8716) (0.7226,0.8822) (0.7326,0.8923) (0.7427,0.9017) (0.7527,0.9106) (0.7627,0.9190) (0.7727,0.9268) (0.7828,0.9342) (0.7928,0.9410) (0.8028,0.9474) (0.8128,0.9534) (0.8228,0.9589) (0.8329,0.9639) (0.8429,0.9686) (0.8529,0.9729) (0.8629,0.9769) (0.8730,0.9804) (0.8830,0.9836) (0.8930,0.9865) (0.9030,0.9891) (0.9131,0.9914) (0.9231,0.9934) (0.9331,0.9950) (0.9431,0.9965) (0.9531,0.9976) (0.9632,0.9986) (0.9732,0.9993) (0.9832,0.9997) (0.9932,1.0000) (1.0033,1.0000) (1.0133,0.9998) (1.0233,0.9995) (1.0333,0.9989) (1.0433,0.9982) (1.0534,0.9974) (1.0634,0.9963) (1.0734,0.9952) (1.0834,0.9938) (1.0935,0.9924) (1.1035,0.9908) (1.1135,0.9891) (1.1235,0.9873) (1.1335,0.9853) (1.1436,0.9833) (1.1536,0.9812) (1.1636,0.9789) (1.1736,0.9766) (1.1837,0.9742) (1.1937,0.9717) (1.2037,0.9691) (1.2137,0.9665) (1.2238,0.9638) (1.2338,0.9610) (1.2438,0.9582) (1.2538,0.9553) (1.2638,0.9523) (1.2739,0.9493) (1.2839,0.9463) (1.2939,0.9432) (1.3039,0.9401) (1.3140,0.9369) (1.3240,0.9337) (1.3340,0.9305) (1.3440,0.9272) (1.3540,0.9239) (1.3641,0.9206) (1.3741,0.9173) (1.3841,0.9139) (1.3941,0.9105) (1.4042,0.9071) (1.4142,0.9037) (1.4242,0.9002) (1.4342,0.8968) (1.4442,0.8933) (1.4543,0.8899) (1.4643,0.8864) (1.4743,0.8829) (1.4843,0.8794) (1.4944,0.8759) (1.5044,0.8724) (1.5144,0.8689) (1.5244,0.8654) (1.5345,0.8619) (1.5445,0.8584) (1.5545,0.8549) (1.5645,0.8514) (1.5745,0.8479) (1.5846,0.8444) (1.5946,0.8409) (1.6046,0.8374) (1.6146,0.8340) (1.6247,0.8305) (1.6347,0.8270) (1.6447,0.8236) (1.6547,0.8201) (1.6647,0.8167) (1.6748,0.8133) (1.6848,0.8099) (1.6948,0.8065) (1.7048,0.8031) (1.7149,0.7997) (1.7249,0.7963) (1.7349,0.7930) (1.7449,0.7896) (1.7549,0.7863) (1.7650,0.7830) (1.7750,0.7797) (1.7850,0.7764) (1.7950,0.7731) (1.8051,0.7699) (1.8151,0.7666) (1.8251,0.7634) (1.8351,0.7602) (1.8452,0.7570) (1.8552,0.7538) (1.8652,0.7506) (1.8752,0.7475) (1.8852,0.7443) (1.8953,0.7412) (1.9053,0.7381) (1.9153,0.7350) (1.9253,0.7319) (1.9354,0.7289) (1.9454,0.7258) (1.9554,0.7228) (1.9654,0.7198) (1.9754,0.7168) (1.9855,0.7138) (1.9955,0.7109) (2.0055,0.7079) (2.0155,0.7050) (2.0256,0.7021) (2.0356,0.6992) (2.0456,0.6963) (2.0556,0.6934) (2.0656,0.6906) (2.0757,0.6878) (2.0857,0.6849) (2.0957,0.6821) (2.1057,0.6794) (2.1158,0.6766) (2.1258,0.6739) (2.1358,0.6711) (2.1458,0.6684) (2.1558,0.6657) (2.1659,0.6630) (2.1759,0.6604) (2.1859,0.6577) (2.1959,0.6551) (2.2060,0.6525) (2.2160,0.6498) (2.2260,0.6473) (2.2360,0.6447) (2.2461,0.6421) (2.2561,0.6396) (2.2661,0.6371) (2.2761,0.6346) (2.2861,0.6321) (2.2962,0.6296) (2.3062,0.6271) (2.3162,0.6247) (2.3262,0.6222) (2.3363,0.6198) (2.3463,0.6174) (2.3563,0.6150) (2.3663,0.6126) (2.3763,0.6103) (2.3864,0.6079) (2.3964,0.6056) (2.4064,0.6033) (2.4164,0.6010) (2.4265,0.5987) (2.4365,0.5964) (2.4465,0.5942) (2.4565,0.5919) (2.4665,0.5897) (2.4766,0.5875) (2.4866,0.5853) (2.4966,0.5831) (2.5066,0.5809) (2.5167,0.5787) (2.5267,0.5766) (2.5367,0.5744) (2.5467,0.5723) (2.5568,0.5702) (2.5668,0.5681) (2.5768,0.5660) (2.5868,0.5639) (2.5968,0.5619) (2.6069,0.5598) (2.6169,0.5578) (2.6269,0.5558) (2.6369,0.5538) (2.6470,0.5518) (2.6570,0.5498) (2.6670,0.5478) (2.6770,0.5458) (2.6870,0.5439) (2.6971,0.5419) (2.7071,0.5400) (2.7171,0.5381) (2.7271,0.5362) (2.7372,0.5343) (2.7472,0.5324) (2.7572,0.5306) (2.7672,0.5287) (2.7772,0.5269) (2.7873,0.5250) (2.7973,0.5232) (2.8073,0.5214) (2.8173,0.5196) (2.8274,0.5178) (2.8374,0.5160) (2.8474,0.5142) (2.8574,0.5125) (2.8675,0.5107) (2.8775,0.5090) (2.8875,0.5073) (2.8975,0.5055) (2.9075,0.5038) (2.9176,0.5021) (2.9276,0.5004) (2.9376,0.4988) (2.9476,0.4971) (2.9577,0.4954) (2.9677,0.4938) (2.9777,0.4921) (2.9877,0.4905) (2.9977,0.4889) (3.0078,0.4873) (3.0178,0.4857) (3.0278,0.4841) (3.0378,0.4825) (3.0479,0.4809) (3.0579,0.4794) (3.0679,0.4778) (3.0779,0.4763) (3.0879,0.4747) (3.0980,0.4732) (3.1080,0.4717) (3.1180,0.4701) (3.1280,0.4686) (3.1381,0.4671) (3.1481,0.4657) (3.1581,0.4642) (3.1681,0.4627) (3.1782,0.4613) (3.1882,0.4598) (3.1982,0.4584) (3.2082,0.4569) (3.2182,0.4555) (3.2283,0.4541) (3.2383,0.4527) (3.2483,0.4512) (3.2583,0.4498) (3.2684,0.4485) (3.2784,0.4471) (3.2884,0.4457) (3.2984,0.4443) (3.3084,0.4430) (3.3185,0.4416) (3.3285,0.4403) (3.3385,0.4389) (3.3485,0.4376) (3.3586,0.4363) (3.3686,0.4350) (3.3786,0.4337) (3.3886,0.4324) (3.3986,0.4311) (3.4087,0.4298) (3.4187,0.4285) (3.4287,0.4272) (3.4387,0.4260) (3.4488,0.4247) (3.4588,0.4234) (3.4688,0.4222) (3.4788,0.4210) (3.4888,0.4197) (3.4989,0.4185) (3.5089,0.4173) (3.5189,0.4161) (3.5289,0.4149) (3.5390,0.4137) (3.5490,0.4125) (3.5590,0.4113) (3.5690,0.4101) (3.5791,0.4089) (3.5891,0.4078) (3.5991,0.4066) (3.6091,0.4054) (3.6191,0.4043) (3.6292,0.4031) (3.6392,0.4020) (3.6492,0.4009) (3.6592,0.3997) (3.6693,0.3986) (3.6793,0.3975) (3.6893,0.3964) (3.6993,0.3953) (3.7093,0.3942) (3.7194,0.3931) (3.7294,0.3920) (3.7394,0.3909) (3.7494,0.3898) (3.7595,0.3888) (3.7695,0.3877) (3.7795,0.3866) (3.7895,0.3856) (3.7995,0.3845) (3.8096,0.3835) (3.8196,0.3825) (3.8296,0.3814) (3.8396,0.3804) (3.8497,0.3794) (3.8597,0.3783) (3.8697,0.3773) (3.8797,0.3763) (3.8898,0.3753) (3.8998,0.3743) (3.9098,0.3733) (3.9198,0.3723) (3.9298,0.3713) (3.9399,0.3704) (3.9499,0.3694) (3.9599,0.3684) (3.9699,0.3675) (3.9800,0.3665) (3.9900,0.3655) (4.0000,0.3646)};
\fill (0.3679,0) circle (1.5pt) node[below right]{$\frac1e$};
\fill (1,1) circle (1.5pt);
\draw[cRem,thick] (0.6,1)--(1.4,1);
\Ticks{1,2,3,4}{-3,-2,-1,1}
\node[cDef] at (2.6,0.6){$(C_f)$};
\node[cThm] at (0.62,-0.9){$(T)$};
\end{GR}

\QQ{4} $(e-1)f(x)=e^2x-me\iff f(x)=\Big(\dfrac{e^2}{e-1}\Big)x-\dfrac{me}{e-1}$ (لأن $e-1\neq0$).

حلول المعادلة هي فواصل نقط تقاطع $(C_f)$ مع المستقيم $(D_m):\ y=\frac{e^2}{e-1}x-\frac{me}{e-1}$ الموازي لـ $(T)$ ، وترتيبه إلى المبدأ $-\frac{me}{e-1}$ ؛ ومن أجل $m=1$ يكون $(D_1)=(T)$.

بيانيا $(C_f)$ يقع تحت مماسه $(T)$ ويمسه في نقطة واحدة، ولما يُسحب $(T)$ نحو الأسفل يقطع $(C_f)$ في نقطتين:
\begin{itemize}
\item إذا كان $m<1$: $(D_m)$ فوق $(T)$ ، لا توجد حلول.
\item إذا كان $m=1$: حلّ مضاعف $x=\frac1e$.
\item إذا كان $m>1$: $(D_m)$ تحت $(T)$ فيقطع $(C_f)$ في نقطتين: حلّان متمايزان.
\end{itemize}
إذن: \bx{m\in\,]1;+\infty[}.
\end{hal}

% ===================== بكالوريا 2019 م1 =====================
\begin{tamrin}{بكالوريا 2019 -- الموضوع الأول}
$f$ الدالة العددية المعرفة على $]0;2[\,\cup\,]2;+\infty[$ بـ: \ $f(x)=\dfrac{1}{x-2}+\ln x$.

$(C_f)$ تمثيلها البياني في المستوي المنسوب إلى المعلم المتعامد والمتجانس $(O;\vec i,\vec j)$.
\QQ{1}
\begin{sq}
\qa أحسب $\lim\limits_{x\to0}f(x)$ ، $\lim\limits_{x\to2}f(x)$ (من اليسار) و $\lim\limits_{x\to2}f(x)$ (من اليمين) ثم فسّر النتائج بيانيا.
\qb أحسب $\lim\limits_{x\to+\infty}f(x)$.
\end{sq}
\QQ{2} أدرس اتجاه تغيّر الدالة $f$ على $]0;2[\,\cup\,]2;+\infty[$ ، ثم شكّل جدول تغيّراتها.
\QQ{3} نسمي $(\Gamma)$ المنحنى البياني للدالة اللوغاريتمية النيبيرية $\ln$ في المعلم السابق.
\begin{sq}
\qa أحسب $\lim\limits_{x\to+\infty}\big(f(x)-\ln x\big)$ ثم فسّر النتيجة بيانيا.
\qb أدرس وضعية $(C_f)$ بالنسبة للمنحنى $(\Gamma)$.
\end{sq}
\QQ{4} أرسم بعناية المنحنى $(\Gamma)$ ثم المنحنى $(C_f)$.
\QQ{5} $g$ الدالة المعرفة على $]-\infty;-1[\,\cup\,]-1;0[$ بـ: $g(x)=f(-2x)$. دون حساب عبارة $g(x)$ حدّد اتجاه تغيّر الدالة $g$.
\end{tamrin}

\begin{hal}
\QQ{1}
\pa
\begin{itemize}
\item $\lim\limits_{x\to0}\frac{1}{x-2}=-\frac12$ و $\lim\limits_{x\to0}\ln x=-\infty$ ، إذن \bx{\lim_{x\to0}f(x)=-\infty}.
\item لما $x\to2$ ، $x<2$: $x-2\to0^-$ ومنه $\frac{1}{x-2}\to-\infty$ و $\ln x\to\ln2$ ، إذن \bx{\lim_{x\to2^-}f(x)=-\infty}.
\item لما $x\to2$ ، $x>2$: $x-2\to0^+$ ومنه $\frac{1}{x-2}\to+\infty$ ، إذن \bx{\lim_{x\to2^+}f(x)=+\infty}.
\end{itemize}
\textbf{التفسير البياني:} المستقيمان ذوا المعادلتين $x=0$ و $x=2$ مقاربان عموديان لـ $(C_f)$.

\pb $\lim\limits_{x\to+\infty}\frac{1}{x-2}=0$ و $\lim\limits_{x\to+\infty}\ln x=+\infty$ ، إذن \bx{\lim_{x\to+\infty}f(x)=+\infty}.

\QQ{2} $f$ قابلة للاشتقاق على مجموعة تعريفها و:
\[f'(x)=-\frac{1}{(x-2)^2}+\frac1x=\frac{(x-2)^2-x}{x(x-2)^2}=\frac{x^2-5x+4}{x(x-2)^2}=\frac{(x-1)(x-4)}{x(x-2)^2}.\]
لأن $x(x-2)^2>0$ فإن إشارة $f'(x)$ هي إشارة $(x-1)(x-4)$ (موجبة خارج الجذرين وسالبة بينهما).

\bxt{$f$ متزايدة تماما على $]0;1]$ و $[4;+\infty[$ ، ومتناقصة تماما على $[1;2[$ و $]2;4]$}

$f(1)=\frac{1}{-1}+0=-1$ ، $f(4)=\frac12+\ln4\approx1{,}89$.
\begin{TAB}
\tkzTabInit[lgt=2,espcl=2.3]{$x$/1,$f'(x)$/1,$f(x)$/2.2}{$0$,$1$,$2$,$4$,$+\infty$}
\tkzTabLine{d,+,z,-,d,-,z,+,}
\tkzTabVar{D-/$-\infty$,+/$-1$,-D+/$-\infty$/$+\infty$,-/$\frac12+\ln4$,+/$+\infty$}
\end{TAB}

\QQ{3}
\pa $f(x)-\ln x=\dfrac{1}{x-2}$ ، إذن \bx{\lim_{x\to+\infty}\big(f(x)-\ln x\big)=0}.

\textbf{التفسير:} المنحنيان $(C_f)$ و $(\Gamma)$ متقاربان بجوار $+\infty$.

\pb إشارة $f(x)-\ln x=\frac{1}{x-2}$ هي إشارة $x-2$:
\begin{itemize}
\item على $]0;2[$: $f(x)-\ln x<0$ ، إذن $(C_f)$ يقع \textbf{تحت} $(\Gamma)$.
\item على $]2;+\infty[$: $f(x)-\ln x>0$ ، إذن $(C_f)$ يقع \textbf{فوق} $(\Gamma)$.
\end{itemize}

\QQ{4} الرسم:
\begin{GR}{-0.5}{8}{-5}{5}{0.85}
\draw[dashed,cThm] (2,-5)--(2,5);
\draw[cRem,thick] plot coordinates {(0.0041,-5.4895) (0.0043,-5.4542) (0.0044,-5.4186) (0.0046,-5.3826) (0.0048,-5.3462) (0.0049,-5.3094) (0.0051,-5.2723) (0.0053,-5.2349) (0.0055,-5.1971) (0.0057,-5.1591) (0.0060,-5.1207) (0.0062,-5.0820) (0.0065,-5.0430) (0.0067,-5.0038) (0.0070,-4.9643) (0.0073,-4.9245) (0.0076,-4.8844) (0.0079,-4.8441) (0.0082,-4.8036) (0.0085,-4.7629) (0.0089,-4.7219) (0.0093,-4.6807) (0.0097,-4.6393) (0.0101,-4.5977) (0.0105,-4.5559) (0.0110,-4.5140) (0.0114,-4.4718) (0.0119,-4.4295) (0.0124,-4.3870) (0.0130,-4.3444) (0.0135,-4.3016) (0.0141,-4.2587) (0.0148,-4.2156) (0.0154,-4.1724) (0.0161,-4.1290) (0.0168,-4.0855) (0.0176,-4.0420) (0.0183,-3.9983) (0.0192,-3.9544) (0.0200,-3.9105) (0.0209,-3.8665) (0.0210,-3.8610) (0.0219,-3.8224) (0.0229,-3.7782) (0.0239,-3.7339) (0.0250,-3.6895) (0.0261,-3.6450) (0.0273,-3.6004) (0.0286,-3.5558) (0.0299,-3.5111) (0.0312,-3.4664) (0.0327,-3.4215) (0.0342,-3.3766) (0.0357,-3.3317) (0.0374,-3.2866) (0.0391,-3.2416) (0.0409,-3.1965) (0.0411,-3.1919) (0.0428,-3.1513) (0.0448,-3.1060) (0.0469,-3.0608) (0.0490,-3.0155) (0.0513,-2.9701) (0.0537,-2.9247) (0.0562,-2.8793) (0.0588,-2.8338) (0.0611,-2.7945) (0.0615,-2.7883) (0.0644,-2.7427) (0.0674,-2.6971) (0.0705,-2.6515) (0.0738,-2.6059) (0.0773,-2.5602) (0.0809,-2.5145) (0.0812,-2.5110) (0.0847,-2.4688) (0.0887,-2.4230) (0.0928,-2.3773) (0.0972,-2.3315) (0.1012,-2.2903) (0.1017,-2.2857) (0.1065,-2.2398) (0.1115,-2.1940) (0.1167,-2.1481) (0.1213,-2.1096) (0.1222,-2.1022) (0.1279,-2.0563) (0.1339,-2.0103) (0.1402,-1.9644) (0.1413,-1.9566) (0.1468,-1.9184) (0.1537,-1.8725) (0.1610,-1.8265) (0.1614,-1.8240) (0.1814,-1.7069) (0.2015,-1.6021) (0.2215,-1.5072) (0.2416,-1.4206) (0.2616,-1.3409) (0.2817,-1.2670) (0.3017,-1.1983) (0.3218,-1.1339) (0.3418,-1.0735) (0.3619,-1.0165) (0.3819,-0.9626) (0.4020,-0.9114) (0.4220,-0.8627) (0.4420,-0.8163) (0.4621,-0.7720) (0.4821,-0.7295) (0.5022,-0.6888) (0.5222,-0.6496) (0.5423,-0.6120) (0.5623,-0.5757) (0.5824,-0.5406) (0.6024,-0.5068) (0.6225,-0.4740) (0.6425,-0.4424) (0.6626,-0.4116) (0.6826,-0.3818) (0.7027,-0.3529) (0.7227,-0.3247) (0.7428,-0.2974) (0.7628,-0.2707) (0.7829,-0.2448) (0.8029,-0.2195) (0.8230,-0.1949) (0.8430,-0.1708) (0.8630,-0.1473) (0.8831,-0.1243) (0.9031,-0.1019) (0.9232,-0.0799) (0.9432,-0.0584) (0.9633,-0.0374) (0.9833,-0.0168) (1.0034,0.0034) (1.0234,0.0232) (1.0435,0.0426) (1.0635,0.0616) (1.0836,0.0803) (1.1036,0.0986) (1.1237,0.1166) (1.1437,0.1343) (1.1638,0.1517) (1.1838,0.1687) (1.2039,0.1855) (1.2239,0.2020) (1.2440,0.2183) (1.2640,0.2343) (1.2840,0.2500) (1.3041,0.2655) (1.3241,0.2808) (1.3442,0.2958) (1.3642,0.3106) (1.3843,0.3252) (1.4043,0.3396) (1.4244,0.3537) (1.4444,0.3677) (1.4645,0.3815) (1.4845,0.3951) (1.5046,0.4085) (1.5246,0.4217) (1.5447,0.4348) (1.5647,0.4477) (1.5848,0.4604) (1.6048,0.4730) (1.6249,0.4854) (1.6449,0.4977) (1.6650,0.5098) (1.6850,0.5218) (1.7050,0.5336) (1.7251,0.5453) (1.7451,0.5568) (1.7652,0.5683) (1.7852,0.5796) (1.8053,0.5907) (1.8253,0.6018) (1.8454,0.6127) (1.8654,0.6235) (1.8855,0.6342) (1.9055,0.6448) (1.9256,0.6552) (1.9456,0.6656) (1.9657,0.6758) (1.9857,0.6860) (2.0058,0.6960) (2.0258,0.7060) (2.0459,0.7158) (2.0659,0.7256) (2.0860,0.7352) (2.1060,0.7448) (2.1260,0.7543) (2.1461,0.7637) (2.1661,0.7729) (2.1862,0.7822) (2.2062,0.7913) (2.2263,0.8003) (2.2463,0.8093) (2.2664,0.8182) (2.2864,0.8270) (2.3065,0.8357) (2.3265,0.8444) (2.3466,0.8530) (2.3666,0.8615) (2.3867,0.8699) (2.4067,0.8783) (2.4268,0.8866) (2.4468,0.8948) (2.4669,0.9029) (2.4869,0.9110) (2.5070,0.9191) (2.5270,0.9270) (2.5470,0.9349) (2.5671,0.9428) (2.5871,0.9506) (2.6072,0.9583) (2.6272,0.9659) (2.6473,0.9735) (2.6673,0.9811) (2.6874,0.9886) (2.7074,0.9960) (2.7275,1.0034) (2.7475,1.0107) (2.7676,1.0180) (2.7876,1.0252) (2.8077,1.0324) (2.8277,1.0395) (2.8478,1.0465) (2.8678,1.0535) (2.8879,1.0605) (2.9079,1.0674) (2.9280,1.0743) (2.9480,1.0811) (2.9680,1.0879) (2.9881,1.0946) (3.0081,1.1013) (3.0282,1.1080) (3.0482,1.1146) (3.0683,1.1211) (3.0883,1.1276) (3.1084,1.1341) (3.1284,1.1405) (3.1485,1.1469) (3.1685,1.1533) (3.1886,1.1596) (3.2086,1.1658) (3.2287,1.1721) (3.2487,1.1783) (3.2688,1.1844) (3.2888,1.1905) (3.3089,1.1966) (3.3289,1.2026) (3.3490,1.2086) (3.3690,1.2146) (3.3890,1.2205) (3.4091,1.2264) (3.4291,1.2323) (3.4492,1.2381) (3.4692,1.2439) (3.4893,1.2497) (3.5093,1.2554) (3.5294,1.2611) (3.5494,1.2668) (3.5695,1.2724) (3.5895,1.2780) (3.6096,1.2836) (3.6296,1.2891) (3.6497,1.2946) (3.6697,1.3001) (3.6898,1.3056) (3.7098,1.3110) (3.7299,1.3164) (3.7499,1.3217) (3.7700,1.3271) (3.7900,1.3324) (3.8100,1.3376) (3.8301,1.3429) (3.8501,1.3481) (3.8702,1.3533) (3.8902,1.3585) (3.9103,1.3636) (3.9303,1.3687) (3.9504,1.3738) (3.9704,1.3789) (3.9905,1.3839) (4.0105,1.3889) (4.0306,1.3939) (4.0506,1.3989) (4.0707,1.4038) (4.0907,1.4087) (4.1108,1.4136) (4.1308,1.4185) (4.1509,1.4233) (4.1709,1.4281) (4.1910,1.4329) (4.2110,1.4377) (4.2310,1.4424) (4.2511,1.4472) (4.2711,1.4519) (4.2912,1.4566) (4.3112,1.4612) (4.3313,1.4659) (4.3513,1.4705) (4.3714,1.4751) (4.3914,1.4797) (4.4115,1.4842) (4.4315,1.4887) (4.4516,1.4933) (4.4716,1.4978) (4.4917,1.5022) (4.5117,1.5067) (4.5318,1.5111) (4.5518,1.5155) (4.5719,1.5199) (4.5919,1.5243) (4.6120,1.5287) (4.6320,1.5330) (4.6520,1.5373) (4.6721,1.5416) (4.6921,1.5459) (4.7122,1.5502) (4.7322,1.5544) (4.7523,1.5586) (4.7723,1.5628) (4.7924,1.5670) (4.8124,1.5712) (4.8325,1.5754) (4.8525,1.5795) (4.8726,1.5836) (4.8926,1.5877) (4.9127,1.5918) (4.9327,1.5959) (4.9528,1.5999) (4.9728,1.6040) (4.9929,1.6080) (5.0129,1.6120) (5.0330,1.6160) (5.0530,1.6200) (5.0730,1.6239) (5.0931,1.6279) (5.1131,1.6318) (5.1332,1.6357) (5.1532,1.6396) (5.1733,1.6435) (5.1933,1.6474) (5.2134,1.6512) (5.2334,1.6551) (5.2535,1.6589) (5.2735,1.6627) (5.2936,1.6665) (5.3136,1.6703) (5.3337,1.6740) (5.3537,1.6778) (5.3738,1.6815) (5.3938,1.6853) (5.4139,1.6890) (5.4339,1.6927) (5.4540,1.6963) (5.4740,1.7000) (5.4940,1.7037) (5.5141,1.7073) (5.5341,1.7109) (5.5542,1.7146) (5.5742,1.7182) (5.5943,1.7217) (5.6143,1.7253) (5.6344,1.7289) (5.6544,1.7324) (5.6745,1.7360) (5.6945,1.7395) (5.7146,1.7430) (5.7346,1.7465) (5.7547,1.7500) (5.7747,1.7535) (5.7948,1.7570) (5.8148,1.7604) (5.8349,1.7638) (5.8549,1.7673) (5.8750,1.7707) (5.8950,1.7741) (5.9150,1.7775) (5.9351,1.7809) (5.9551,1.7843) (5.9752,1.7876) (5.9952,1.7910) (6.0153,1.7943) (6.0353,1.7976) (6.0554,1.8009) (6.0754,1.8043) (6.0955,1.8075) (6.1155,1.8108) (6.1356,1.8141) (6.1556,1.8174) (6.1757,1.8206) (6.1957,1.8239) (6.2158,1.8271) (6.2358,1.8303) (6.2559,1.8335) (6.2759,1.8367) (6.2960,1.8399) (6.3160,1.8431) (6.3360,1.8463) (6.3561,1.8494) (6.3761,1.8526) (6.3962,1.8557) (6.4162,1.8588) (6.4363,1.8620) (6.4563,1.8651) (6.4764,1.8682) (6.4964,1.8713) (6.5165,1.8743) (6.5365,1.8774) (6.5566,1.8805) (6.5766,1.8835) (6.5967,1.8866) (6.6167,1.8896) (6.6368,1.8926) (6.6568,1.8956) (6.6769,1.8986) (6.6969,1.9016) (6.7170,1.9046) (6.7370,1.9076) (6.7570,1.9106) (6.7771,1.9135) (6.7971,1.9165) (6.8172,1.9194) (6.8372,1.9224) (6.8573,1.9253) (6.8773,1.9282) (6.8974,1.9311) (6.9174,1.9340) (6.9375,1.9369) (6.9575,1.9398) (6.9776,1.9427) (6.9976,1.9456) (7.0177,1.9484) (7.0377,1.9513) (7.0578,1.9541) (7.0778,1.9570) (7.0979,1.9598) (7.1179,1.9626) (7.1380,1.9654) (7.1580,1.9682) (7.1780,1.9710) (7.1981,1.9738) (7.2181,1.9766) (7.2382,1.9794) (7.2582,1.9821) (7.2783,1.9849) (7.2983,1.9876) (7.3184,1.9904) (7.3384,1.9931) (7.3585,1.9959) (7.3785,1.9986) (7.3986,2.0013) (7.4186,2.0040) (7.4387,2.0067) (7.4587,2.0094) (7.4788,2.0121) (7.4988,2.0147) (7.5189,2.0174) (7.5389,2.0201) (7.5590,2.0227) (7.5790,2.0254) (7.5990,2.0280) (7.6191,2.0307) (7.6391,2.0333) (7.6592,2.0359) (7.6792,2.0385) (7.6993,2.0411) (7.7193,2.0437) (7.7394,2.0463) (7.7594,2.0489) (7.7795,2.0515) (7.7995,2.0541) (7.8196,2.0566) (7.8396,2.0592) (7.8597,2.0617) (7.8797,2.0643) (7.8998,2.0668) (7.9198,2.0694) (7.9399,2.0719) (7.9599,2.0744) (7.9800,2.0769) (8.0000,2.0794)};
\draw[cDef,thick] plot coordinates {(0.0070,-5.4649) (0.0073,-5.4252) (0.0076,-5.3852) (0.0079,-5.3450) (0.0082,-5.3045) (0.0086,-5.2638) (0.0089,-5.2230) (0.0093,-5.1819) (0.0097,-5.1406) (0.0101,-5.0991) (0.0105,-5.0574) (0.0110,-5.0155) (0.0110,-5.0113) (0.0114,-4.9735) (0.0119,-4.9313) (0.0125,-4.8889) (0.0130,-4.8464) (0.0136,-4.8038) (0.0142,-4.7610) (0.0148,-4.7181) (0.0154,-4.6750) (0.0160,-4.6378) (0.0161,-4.6318) (0.0168,-4.5885) (0.0176,-4.5451) (0.0184,-4.5016) (0.0192,-4.4580) (0.0201,-4.4143) (0.0210,-4.3705) (0.0210,-4.3671) (0.0219,-4.3266) (0.0229,-4.2827) (0.0239,-4.2386) (0.0250,-4.1945) (0.0260,-4.1548) (0.0262,-4.1503) (0.0273,-4.1061) (0.0286,-4.0618) (0.0299,-4.0174) (0.0310,-3.9802) (0.0313,-3.9730) (0.0327,-3.9286) (0.0342,-3.8840) (0.0358,-3.8395) (0.0361,-3.8320) (0.0374,-3.7949) (0.0392,-3.7503) (0.0410,-3.7056) (0.0411,-3.7032) (0.0461,-3.5894) (0.0511,-3.4876) (0.0561,-3.3954) (0.0611,-3.3112) (0.0661,-3.2337) (0.0711,-3.1620) (0.0761,-3.0953) (0.0811,-3.0330) (0.0861,-2.9744) (0.0911,-2.9193) (0.0961,-2.8672) (0.1012,-2.8178) (0.1062,-2.7709) (0.1112,-2.7262) (0.1162,-2.6835) (0.1212,-2.6427) (0.1262,-2.6037) (0.1312,-2.5662) (0.1362,-2.5301) (0.1412,-2.4955) (0.1462,-2.4621) (0.1512,-2.4299) (0.1562,-2.3988) (0.1612,-2.3687) (0.1662,-2.3396) (0.1713,-2.3114) (0.1763,-2.2841) (0.1813,-2.2576) (0.1863,-2.2319) (0.1913,-2.2069) (0.1963,-2.1826) (0.2013,-2.1589) (0.2063,-2.1359) (0.2113,-2.1135) (0.2163,-2.0916) (0.2213,-2.0703) (0.2263,-2.0495) (0.2313,-2.0292) (0.2364,-2.0094) (0.2414,-1.9901) (0.2464,-1.9712) (0.2514,-1.9527) (0.2564,-1.9346) (0.2614,-1.9169) (0.2664,-1.8996) (0.2714,-1.8826) (0.2764,-1.8660) (0.2814,-1.8498) (0.2864,-1.8338) (0.2914,-1.8182) (0.2964,-1.8029) (0.3015,-1.7879) (0.3065,-1.7732) (0.3115,-1.7587) (0.3165,-1.7445) (0.3215,-1.7306) (0.3265,-1.7169) (0.3315,-1.7035) (0.3365,-1.6903) (0.3415,-1.6773) (0.3465,-1.6646) (0.3515,-1.6521) (0.3565,-1.6398) (0.3615,-1.6277) (0.3665,-1.6158) (0.3716,-1.6041) (0.3766,-1.5926) (0.3816,-1.5813) (0.3866,-1.5702) (0.3916,-1.5593) (0.3966,-1.5485) (0.4016,-1.5379) (0.4066,-1.5275) (0.4116,-1.5172) (0.4166,-1.5071) (0.4216,-1.4972) (0.4266,-1.4874) (0.4316,-1.4778) (0.4367,-1.4683) (0.4417,-1.4589) (0.4467,-1.4497) (0.4517,-1.4406) (0.4567,-1.4317) (0.4617,-1.4229) (0.4667,-1.4143) (0.4717,-1.4057) (0.4767,-1.3973) (0.4817,-1.3890) (0.4867,-1.3809) (0.4917,-1.3728) (0.4967,-1.3649) (0.5018,-1.3571) (0.5068,-1.3494) (0.5118,-1.3418) (0.5168,-1.3344) (0.5218,-1.3270) (0.5268,-1.3197) (0.5318,-1.3126) (0.5368,-1.3056) (0.5418,-1.2986) (0.5468,-1.2918) (0.5518,-1.2850) (0.5568,-1.2784) (0.5618,-1.2719) (0.5668,-1.2654) (0.5719,-1.2591) (0.5769,-1.2528) (0.5819,-1.2467) (0.5869,-1.2406) (0.5919,-1.2346) (0.5969,-1.2287) (0.6019,-1.2229) (0.6069,-1.2172) (0.6119,-1.2116) (0.6169,-1.2060) (0.6219,-1.2006) (0.6269,-1.1952) (0.6319,-1.1899) (0.6370,-1.1847) (0.6420,-1.1796) (0.6470,-1.1745) (0.6520,-1.1696) (0.6570,-1.1647) (0.6620,-1.1599) (0.6670,-1.1552) (0.6720,-1.1505) (0.6770,-1.1459) (0.6820,-1.1414) (0.6870,-1.1370) (0.6920,-1.1327) (0.6970,-1.1284) (0.7021,-1.1242) (0.7071,-1.1201) (0.7121,-1.1160) (0.7171,-1.1120) (0.7221,-1.1081) (0.7271,-1.1043) (0.7321,-1.1005) (0.7371,-1.0969) (0.7421,-1.0932) (0.7471,-1.0897) (0.7521,-1.0862) (0.7571,-1.0828) (0.7621,-1.0795) (0.7672,-1.0762) (0.7722,-1.0730) (0.7772,-1.0699) (0.7822,-1.0668) (0.7872,-1.0638) (0.7922,-1.0609) (0.7972,-1.0580) (0.8022,-1.0553) (0.8072,-1.0525) (0.8122,-1.0499) (0.8172,-1.0473) (0.8222,-1.0448) (0.8272,-1.0423) (0.8322,-1.0400) (0.8373,-1.0377) (0.8423,-1.0354) (0.8473,-1.0332) (0.8523,-1.0311) (0.8573,-1.0291) (0.8623,-1.0271) (0.8673,-1.0252) (0.8723,-1.0234) (0.8773,-1.0216) (0.8823,-1.0199) (0.8873,-1.0183) (0.8923,-1.0167) (0.8973,-1.0152) (0.9024,-1.0138) (0.9074,-1.0124) (0.9124,-1.0111) (0.9174,-1.0099) (0.9224,-1.0088) (0.9274,-1.0077) (0.9324,-1.0067) (0.9374,-1.0057) (0.9424,-1.0049) (0.9474,-1.0041) (0.9524,-1.0033) (0.9574,-1.0027) (0.9624,-1.0021) (0.9675,-1.0016) (0.9725,-1.0011) (0.9775,-1.0008) (0.9825,-1.0005) (0.9875,-1.0002) (0.9925,-1.0001) (0.9975,-1.0000) (1.0025,-1.0000) (1.0075,-1.0001) (1.0125,-1.0002) (1.0175,-1.0005) (1.0225,-1.0008) (1.0275,-1.0012) (1.0325,-1.0016) (1.0376,-1.0022) (1.0426,-1.0028) (1.0476,-1.0035) (1.0526,-1.0043) (1.0576,-1.0051) (1.0626,-1.0061) (1.0676,-1.0071) (1.0726,-1.0082) (1.0776,-1.0094) (1.0826,-1.0107) (1.0876,-1.0120) (1.0926,-1.0135) (1.0976,-1.0150) (1.1027,-1.0167) (1.1077,-1.0184) (1.1127,-1.0202) (1.1177,-1.0221) (1.1227,-1.0241) (1.1277,-1.0262) (1.1327,-1.0284) (1.1377,-1.0307) (1.1427,-1.0331) (1.1477,-1.0355) (1.1527,-1.0381) (1.1577,-1.0408) (1.1627,-1.0436) (1.1678,-1.0465) (1.1728,-1.0495) (1.1778,-1.0526) (1.1828,-1.0558) (1.1878,-1.0591) (1.1928,-1.0625) (1.1978,-1.0661) (1.2028,-1.0697) (1.2078,-1.0735) (1.2128,-1.0774) (1.2178,-1.0814) (1.2228,-1.0856) (1.2278,-1.0898) (1.2328,-1.0942) (1.2379,-1.0987) (1.2429,-1.1033) (1.2479,-1.1081) (1.2529,-1.1130) (1.2579,-1.1181) (1.2629,-1.1233) (1.2679,-1.1286) (1.2729,-1.1340) (1.2779,-1.1397) (1.2829,-1.1454) (1.2879,-1.1513) (1.2929,-1.1574) (1.2979,-1.1636) (1.3030,-1.1700) (1.3080,-1.1765) (1.3130,-1.1832) (1.3180,-1.1901) (1.3230,-1.1972) (1.3280,-1.2044) (1.3330,-1.2118) (1.3380,-1.2194) (1.3430,-1.2272) (1.3480,-1.2352) (1.3530,-1.2433) (1.3580,-1.2517) (1.3630,-1.2602) (1.3681,-1.2690) (1.3731,-1.2780) (1.3781,-1.2872) (1.3831,-1.2966) (1.3881,-1.3063) (1.3931,-1.3162) (1.3981,-1.3263) (1.4031,-1.3366) (1.4081,-1.3473) (1.4131,-1.3581) (1.4181,-1.3693) (1.4231,-1.3806) (1.4281,-1.3923) (1.4332,-1.4043) (1.4382,-1.4165) (1.4432,-1.4290) (1.4482,-1.4419) (1.4532,-1.4550) (1.4582,-1.4685) (1.4632,-1.4823) (1.4682,-1.4964) (1.4732,-1.5108) (1.4782,-1.5257) (1.4832,-1.5409) (1.4882,-1.5564) (1.4932,-1.5724) (1.4982,-1.5887) (1.5033,-1.6055) (1.5083,-1.6226) (1.5133,-1.6403) (1.5183,-1.6583) (1.5233,-1.6768) (1.5283,-1.6958) (1.5333,-1.7153) (1.5383,-1.7353) (1.5433,-1.7558) (1.5483,-1.7768) (1.5533,-1.7984) (1.5583,-1.8206) (1.5633,-1.8433) (1.5684,-1.8667) (1.5734,-1.8907) (1.5784,-1.9153) (1.5834,-1.9407) (1.5884,-1.9667) (1.5934,-1.9935) (1.5984,-2.0210) (1.6034,-2.0493) (1.6084,-2.0785) (1.6134,-2.1084) (1.6184,-2.1393) (1.6234,-2.1710) (1.6284,-2.2038) (1.6335,-2.2375) (1.6385,-2.2722) (1.6435,-2.3080) (1.6485,-2.3449) (1.6535,-2.3830) (1.6585,-2.4223) (1.6635,-2.4628) (1.6685,-2.5047) (1.6735,-2.5480) (1.6785,-2.5927) (1.6835,-2.6389) (1.6885,-2.6868) (1.6935,-2.7363) (1.6985,-2.7875) (1.7036,-2.8406) (1.7086,-2.8956) (1.7136,-2.9527) (1.7186,-3.0119) (1.7236,-3.0734) (1.7286,-3.1372) (1.7336,-3.2036) (1.7386,-3.2726) (1.7436,-3.3444) (1.7486,-3.4193) (1.7536,-3.4973) (1.7586,-3.5786) (1.7636,-3.6636) (1.7687,-3.7523) (1.7737,-3.8451) (1.7787,-3.9423) (1.7837,-4.0440) (1.7887,-4.1508) (1.7937,-4.2628) (1.7987,-4.3806) (1.8037,-4.5046) (1.8087,-4.6352) (1.8137,-4.7729) (1.8187,-4.9185) (1.8237,-5.0724) (1.8287,-5.2356) (1.8338,-5.4087)};
\draw[cDef,thick] plot coordinates {(2.2265,5.2149) (2.2416,4.9469) (2.2566,4.7110) (2.2716,4.5020) (2.2867,4.3155) (2.3017,4.1482) (2.3167,3.9974) (2.3318,3.8607) (2.3468,3.7365) (2.3618,3.6231) (2.3769,3.5192) (2.3919,3.4237) (2.4069,3.3357) (2.4220,3.2544) (2.4370,3.1790) (2.4521,3.1091) (2.4671,3.0440) (2.4821,2.9833) (2.4972,2.9266) (2.5122,2.8735) (2.5272,2.8238) (2.5423,2.7772) (2.5573,2.7333) (2.5723,2.6920) (2.5874,2.6532) (2.6024,2.6165) (2.6174,2.5818) (2.6325,2.5490) (2.6475,2.5180) (2.6625,2.4886) (2.6776,2.4608) (2.6926,2.4343) (2.7076,2.4092) (2.7227,2.3853) (2.7377,2.3627) (2.7528,2.3411) (2.7678,2.3205) (2.7828,2.3009) (2.7979,2.2822) (2.8129,2.2644) (2.8279,2.2474) (2.8430,2.2311) (2.8580,2.2156) (2.8730,2.2008) (2.8881,2.1866) (2.9031,2.1731) (2.9181,2.1601) (2.9332,2.1477) (2.9482,2.1358) (2.9632,2.1244) (2.9783,2.1135) (2.9933,2.1031) (3.0084,2.0931) (3.0234,2.0835) (3.0384,2.0743) (3.0535,2.0655) (3.0685,2.0571) (3.0835,2.0490) (3.0986,2.0412) (3.1136,2.0338) (3.1286,2.0266) (3.1437,2.0198) (3.1587,2.0132) (3.1737,2.0069) (3.1888,2.0008) (3.2038,1.9950) (3.2188,1.9895) (3.2339,1.9841) (3.2489,1.9790) (3.2639,1.9741) (3.2790,1.9694) (3.2940,1.9649) (3.3091,1.9606) (3.3241,1.9564) (3.3391,1.9525) (3.3542,1.9487) (3.3692,1.9450) (3.3842,1.9416) (3.3993,1.9382) (3.4143,1.9350) (3.4293,1.9320) (3.4444,1.9291) (3.4594,1.9263) (3.4744,1.9237) (3.4895,1.9211) (3.5045,1.9187) (3.5195,1.9164) (3.5346,1.9142) (3.5496,1.9122) (3.5646,1.9102) (3.5797,1.9083) (3.5947,1.9065) (3.6098,1.9049) (3.6248,1.9033) (3.6398,1.9018) (3.6549,1.9003) (3.6699,1.8990) (3.6849,1.8977) (3.7000,1.8966) (3.7150,1.8955) (3.7300,1.8944) (3.7451,1.8935) (3.7601,1.8926) (3.7751,1.8918) (3.7902,1.8910) (3.8052,1.8903) (3.8202,1.8897) (3.8353,1.8891) (3.8503,1.8886) (3.8654,1.8881) (3.8804,1.8877) (3.8954,1.8874) (3.9105,1.8871) (3.9255,1.8868) (3.9405,1.8866) (3.9556,1.8865) (3.9706,1.8864) (3.9856,1.8863) (4.0007,1.8863) (4.0157,1.8863) (4.0307,1.8864) (4.0458,1.8865) (4.0608,1.8866) (4.0758,1.8868) (4.0909,1.8870) (4.1059,1.8873) (4.1209,1.8876) (4.1360,1.8879) (4.1510,1.8882) (4.1661,1.8886) (4.1811,1.8891) (4.1961,1.8895) (4.2112,1.8900) (4.2262,1.8905) (4.2412,1.8910) (4.2563,1.8916) (4.2713,1.8922) (4.2863,1.8928) (4.3014,1.8935) (4.3164,1.8941) (4.3314,1.8948) (4.3465,1.8955) (4.3615,1.8963) (4.3765,1.8970) (4.3916,1.8978) (4.4066,1.8986) (4.4216,1.8995) (4.4367,1.9003) (4.4517,1.9012) (4.4668,1.9021) (4.4818,1.9030) (4.4968,1.9039) (4.5119,1.9048) (4.5269,1.9058) (4.5419,1.9068) (4.5570,1.9077) (4.5720,1.9088) (4.5870,1.9098) (4.6021,1.9108) (4.6171,1.9119) (4.6321,1.9129) (4.6472,1.9140) (4.6622,1.9151) (4.6772,1.9162) (4.6923,1.9174) (4.7073,1.9185) (4.7224,1.9196) (4.7374,1.9208) (4.7524,1.9220) (4.7675,1.9232) (4.7825,1.9244) (4.7975,1.9256) (4.8126,1.9268) (4.8276,1.9280) (4.8426,1.9292) (4.8577,1.9305) (4.8727,1.9318) (4.8877,1.9330) (4.9028,1.9343) (4.9178,1.9356) (4.9328,1.9369) (4.9479,1.9382) (4.9629,1.9395) (4.9779,1.9408) (4.9930,1.9421) (5.0080,1.9435) (5.0231,1.9448) (5.0381,1.9462) (5.0531,1.9475) (5.0682,1.9489) (5.0832,1.9503) (5.0982,1.9517) (5.1133,1.9530) (5.1283,1.9544) (5.1433,1.9558) (5.1584,1.9572) (5.1734,1.9586) (5.1884,1.9601) (5.2035,1.9615) (5.2185,1.9629) (5.2335,1.9643) (5.2486,1.9658) (5.2636,1.9672) (5.2786,1.9687) (5.2937,1.9701) (5.3087,1.9716) (5.3238,1.9730) (5.3388,1.9745) (5.3538,1.9760) (5.3689,1.9775) (5.3839,1.9789) (5.3989,1.9804) (5.4140,1.9819) (5.4290,1.9834) (5.4440,1.9849) (5.4591,1.9864) (5.4741,1.9879) (5.4891,1.9894) (5.5042,1.9909) (5.5192,1.9924) (5.5342,1.9939) (5.5493,1.9954) (5.5643,1.9969) (5.5794,1.9985) (5.5944,2.0000) (5.6094,2.0015) (5.6245,2.0030) (5.6395,2.0046) (5.6545,2.0061) (5.6696,2.0076) (5.6846,2.0092) (5.6996,2.0107) (5.7147,2.0122) (5.7297,2.0138) (5.7447,2.0153) (5.7598,2.0169) (5.7748,2.0184) (5.7898,2.0200) (5.8049,2.0215) (5.8199,2.0231) (5.8349,2.0246) (5.8500,2.0262) (5.8650,2.0277) (5.8801,2.0293) (5.8951,2.0309) (5.9101,2.0324) (5.9252,2.0340) (5.9402,2.0355) (5.9552,2.0371) (5.9703,2.0387) (5.9853,2.0402) (6.0003,2.0418) (6.0154,2.0434) (6.0304,2.0449) (6.0454,2.0465) (6.0605,2.0481) (6.0755,2.0496) (6.0905,2.0512) (6.1056,2.0528) (6.1206,2.0543) (6.1356,2.0559) (6.1507,2.0575) (6.1657,2.0591) (6.1808,2.0606) (6.1958,2.0622) (6.2108,2.0638) (6.2259,2.0653) (6.2409,2.0669) (6.2559,2.0685) (6.2710,2.0701) (6.2860,2.0716) (6.3010,2.0732) (6.3161,2.0748) (6.3311,2.0764) (6.3461,2.0779) (6.3612,2.0795) (6.3762,2.0811) (6.3912,2.0827) (6.4063,2.0842) (6.4213,2.0858) (6.4364,2.0874) (6.4514,2.0889) (6.4664,2.0905) (6.4815,2.0921) (6.4965,2.0937) (6.5115,2.0952) (6.5266,2.0968) (6.5416,2.0984) (6.5566,2.0999) (6.5717,2.1015) (6.5867,2.1031) (6.6017,2.1046) (6.6168,2.1062) (6.6318,2.1078) (6.6468,2.1093) (6.6619,2.1109) (6.6769,2.1125) (6.6919,2.1140) (6.7070,2.1156) (6.7220,2.1172) (6.7371,2.1187) (6.7521,2.1203) (6.7671,2.1218) (6.7822,2.1234) (6.7972,2.1250) (6.8122,2.1265) (6.8273,2.1281) (6.8423,2.1296) (6.8573,2.1312) (6.8724,2.1327) (6.8874,2.1343) (6.9024,2.1359) (6.9175,2.1374) (6.9325,2.1390) (6.9475,2.1405) (6.9626,2.1421) (6.9776,2.1436) (6.9926,2.1452) (7.0077,2.1467) (7.0227,2.1482) (7.0378,2.1498) (7.0528,2.1513) (7.0678,2.1529) (7.0829,2.1544) (7.0979,2.1560) (7.1129,2.1575) (7.1280,2.1590) (7.1430,2.1606) (7.1580,2.1621) (7.1731,2.1636) (7.1881,2.1652) (7.2031,2.1667) (7.2182,2.1682) (7.2332,2.1698) (7.2482,2.1713) (7.2633,2.1728) (7.2783,2.1744) (7.2934,2.1759) (7.3084,2.1774) (7.3234,2.1789) (7.3385,2.1804) (7.3535,2.1820) (7.3685,2.1835) (7.3836,2.1850) (7.3986,2.1865) (7.4136,2.1880) (7.4287,2.1896) (7.4437,2.1911) (7.4587,2.1926) (7.4738,2.1941) (7.4888,2.1956) (7.5038,2.1971) (7.5189,2.1986) (7.5339,2.2001) (7.5489,2.2016) (7.5640,2.2031) (7.5790,2.2046) (7.5941,2.2061) (7.6091,2.2076) (7.6241,2.2091) (7.6392,2.2106) (7.6542,2.2121) (7.6692,2.2136) (7.6843,2.2151) (7.6993,2.2166) (7.7143,2.2181) (7.7294,2.2196) (7.7444,2.2211) (7.7594,2.2225) (7.7745,2.2240) (7.7895,2.2255) (7.8045,2.2270) (7.8196,2.2285) (7.8346,2.2299) (7.8496,2.2314) (7.8647,2.2329) (7.8797,2.2344) (7.8948,2.2358) (7.9098,2.2373) (7.9248,2.2388) (7.9399,2.2402) (7.9549,2.2417) (7.9699,2.2432) (7.9850,2.2446) (8.0000,2.2461)};
\fill (1,-1) circle (2pt); \fill (4,1.886) circle (2pt);
\Ticks{1,2,3,4,5,6,7}{-4,-3,-2,-1,1,2,3,4}
\node[cDef] at (3,4){$(C_f)$};
\node[cRem] at (7,1.4){$(\Gamma)$};
\node[cThm] at (2.6,-4.5){$x=2$};
\end{GR}

\QQ{5} $g$ مركّب الدالة $x\mapsto-2x$ متبوعة بـ $f$ ، إذن: $g'(x)=-2f'(-2x)$ ، فإشارة $g'(x)$ عكس إشارة $f'(-2x)$.

من أجل $x<0$ نضع $u=-2x>0$:
\begin{itemize}
\item $x\in\,]-\infty;-2[\iff u\in\,]4;+\infty[$: $f'(u)>0$ إذن $g'(x)<0$.
\item $x\in\,]-2;-1[\iff u\in\,]2;4[$: $f'(u)<0$ إذن $g'(x)>0$.
\item $x\in\,]-1;-\frac12[\iff u\in\,]1;2[$: $f'(u)<0$ إذن $g'(x)>0$.
\item $x\in\,]-\frac12;0[\iff u\in\,]0;1[$: $f'(u)>0$ إذن $g'(x)<0$.
\end{itemize}
\bxt{$g$ متناقصة تماما على $]-\infty;-2]$ و $\left[-\tfrac12;0\right[$ ، ومتزايدة تماما على $[-2;-1[$ و $\left]-1;-\tfrac12\right]$}
\end{hal}

% ===================== بكالوريا 2020 م1 =====================
\begin{tamrin}{بكالوريا 2020 -- الموضوع الأول}
$f$ الدالة العددية المعرفة على المجال $]0;+\infty[$ بـ: \ $f(x)=x-1-\dfrac{\ln x}{x^2}$.

$(C_f)$ تمثيلها البياني في المستوي المنسوب إلى المعلم المتعامد والمتجانس $(O;\vec i,\vec j)$ ، الوحدة $2\text{cm}$.
\QQ{1}
\begin{sq}
\qa أحسب $\lim\limits_{x\to0}f(x)$ وفسّر النتيجة بيانيا، ثم بيّن أن: $\lim\limits_{x\to+\infty}f(x)=+\infty$.
\qb بيّن أن المستقيم $(\Delta)$ ذا المعادلة $y=x-1$ مقارب مائل للمنحنى $(C_f)$ عند $+\infty$.
\qc أدرس وضعية $(C_f)$ بالنسبة للمستقيم $(\Delta)$.
\end{sq}
\QQ{2} $g$ الدالة العددية المعرفة على $]0;+\infty[$ بـ: $g(x)=x^3-1+2\ln x$.
\begin{sq}
\qa بيّن أن $g$ متزايدة تماما على $]0;+\infty[$.
\qb أحسب $g(1)$ ، ثم استنتج إشارة $g(x)$ حسب قيم $x$ من المجال $]0;+\infty[$.
\end{sq}
\QQ{3}
\begin{sq}
\qa بيّن أنه من أجل كل عدد حقيقي $x$ من المجال $]0;+\infty[$: \ $f'(x)=\dfrac{g(x)}{x^3}$.
\qb استنتج اتجاه تغيّر الدالة $f$ ثم شكّل جدول تغيّراتها.
\end{sq}
\QQ{4} بيّن أن المنحنى $(C_f)$ يقبل مماسا $(T)$ موازيا للمستقيم $(\Delta)$ ، يُطلب تعيين معادلة له.
\QQ{5} أنشئ $(\Delta)$ ، $(T)$ و $(C_f)$.
\QQ{6} $h$ الدالة العددية المعرفة على $]-\infty;0[\,\cup\,]0;+\infty[$ بـ: \ $h(x)=-|x|+1+\dfrac{\ln|x|}{x^2}$.
\begin{sq}
\qa بيّن أن $h$ دالة زوجية.
\qb اشرح كيف يتم إنشاء المنحنى $(C_h)$ الممثل للدالة $h$ انطلاقا من $(C_f)$. (لا يُطلب إنشاء $(C_h)$).
\end{sq}
\end{tamrin}

\begin{hal}
\QQ{1}
\pa $\lim\limits_{x\to0}\ln x=-\infty$ و $\lim\limits_{x\to0}\frac{1}{x^2}=+\infty$ ومنه $\lim\limits_{x\to0}\frac{\ln x}{x^2}=-\infty$ ، و $x-1\to-1$ ، إذن \bx{\lim_{x\to0}f(x)=+\infty}.

\textbf{التفسير:} المستقيم ذو المعادلة $x=0$ (محور التراتيب) مقارب عمودي لـ $(C_f)$.

عند $+\infty$: $\frac{\ln x}{x^2}=\frac{\ln x}{x}\cdot\frac1x\to0\times0=0$ و $x-1\to+\infty$ ، إذن \bx{\lim_{x\to+\infty}f(x)=+\infty}.

\pb $f(x)-(x-1)=-\dfrac{\ln x}{x^2}$ و $\lim\limits_{x\to+\infty}\Big(-\dfrac{\ln x}{x^2}\Big)=0$ ، إذن \bxt{$(\Delta): y=x-1$ مقارب مائل لـ $(C_f)$ عند $+\infty$}.

\pc إشارة $f(x)-(x-1)=-\frac{\ln x}{x^2}$ هي إشارة $-\ln x$ (لأن $x^2>0$):
\begin{itemize}
\item على $]0;1[$: $\ln x<0$ إذن $(C_f)$ فوق $(\Delta)$.
\item على $]1;+\infty[$: $\ln x>0$ إذن $(C_f)$ تحت $(\Delta)$.
\item $(C_f)$ يقطع $(\Delta)$ في النقطة $(1;0)$.
\end{itemize}

\QQ{2}
\pa $g'(x)=3x^2+\frac2x$ ، ومن أجل $x>0$: $3x^2>0$ و $\frac2x>0$ ، إذن $g'(x)>0$ و \bxt{$g$ متزايدة تماما على $]0;+\infty[$}.

\pb $g(1)=1-1+2\ln1=0$. لأن $g$ متزايدة تماما:
\begin{TAB}
\tkzTabInit[lgt=2,espcl=3]{$x$/1,$g(x)$/1}{$0$,$1$,$+\infty$}
\tkzTabLine{d,-,z,+,}
\end{TAB}

\QQ{3}
\pa $\Big(\dfrac{\ln x}{x^2}\Big)'=\dfrac{\frac1x\cdot x^2-2x\ln x}{x^4}=\dfrac{1-2\ln x}{x^3}$ ، إذن:
\[f'(x)=1-\frac{1-2\ln x}{x^3}=\frac{x^3-1+2\ln x}{x^3}=\frac{g(x)}{x^3}.\]
\pb $x^3>0$ ، إذن إشارة $f'(x)$ هي إشارة $g(x)$: \bxt{$f$ متناقصة تماما على $]0;1]$ ومتزايدة تماما على $[1;+\infty[$} ، و $f(1)=0$.
\begin{TAB}
\tkzTabInit[lgt=2,espcl=3]{$x$/1,$f'(x)$/1,$f(x)$/2}{$0$,$1$,$+\infty$}
\tkzTabLine{d,-,z,+,}
\tkzTabVar{D+/$+\infty$,-/$0$,+/$+\infty$}
\end{TAB}

\QQ{4} $(T)\parallel(\Delta)$ معناه أن ميل $(T)$ يساوي $1$:
\[f'(x)=1\iff\frac{x^3-1+2\ln x}{x^3}=1\iff -1+2\ln x=0\iff\ln x=\frac12\iff x=\sqrt e.\]
$f(\sqrt e)=\sqrt e-1-\frac{\frac12}{e}=\sqrt e-1-\frac{1}{2e}$ ، ومنه:
$y=1\cdot(x-\sqrt e)+\sqrt e-1-\frac1{2e}$ أي \bx{(T):\ y=x-1-\frac{1}{2e}}.

\QQ{5} الرسم (الوحدة $2\text{cm}$):
\begin{GR}{-0.5}{4}{-1.2}{4}{2}
\draw[cRem,thick] (0,-1)--(4,3);
\draw[cThm,thick] (0,-1.184)--(4,2.816);
\draw[cDef,thick] plot coordinates {(0.4163,4.4733) (0.4262,4.1216) (0.4361,3.8000) (0.4460,3.5055) (0.4559,3.2353) (0.4658,2.9873) (0.4757,2.7592) (0.4856,2.5492) (0.4955,2.3557) (0.5054,2.1772) (0.5153,2.0124) (0.5252,1.8600) (0.5351,1.7191) (0.5450,1.5887) (0.5549,1.4678) (0.5648,1.3558) (0.5747,1.2519) (0.5846,1.1555) (0.5945,1.0660) (0.6044,0.9829) (0.6143,0.9057) (0.6242,0.8339) (0.6341,0.7672) (0.6440,0.7051) (0.6539,0.6475) (0.6638,0.5939) (0.6737,0.5440) (0.6836,0.4977) (0.6935,0.4546) (0.7034,0.4146) (0.7133,0.3774) (0.7232,0.3429) (0.7331,0.3108) (0.7430,0.2812) (0.7529,0.2536) (0.7628,0.2282) (0.7727,0.2046) (0.7826,0.1829) (0.7925,0.1628) (0.8024,0.1444) (0.8123,0.1274) (0.8222,0.1118) (0.8321,0.0976) (0.8420,0.0846) (0.8519,0.0728) (0.8618,0.0621) (0.8717,0.0524) (0.8816,0.0438) (0.8915,0.0360) (0.9014,0.0292) (0.9113,0.0232) (0.9212,0.0179) (0.9311,0.0135) (0.9410,0.0097) (0.9509,0.0066) (0.9608,0.0041) (0.9707,0.0023) (0.9806,0.0010) (0.9905,0.0002) (1.0004,0.0000) (1.0103,0.0003) (1.0202,0.0010) (1.0301,0.0021) (1.0400,0.0037) (1.0499,0.0057) (1.0598,0.0081) (1.0697,0.0108) (1.0796,0.0139) (1.0895,0.0173) (1.0994,0.0210) (1.1093,0.0250) (1.1192,0.0293) (1.1291,0.0338) (1.1390,0.0387) (1.1489,0.0437) (1.1588,0.0490) (1.1687,0.0545) (1.1786,0.0603) (1.1885,0.0662) (1.1984,0.0724) (1.2083,0.0787) (1.2182,0.0852) (1.2281,0.0918) (1.2380,0.0987) (1.2479,0.1057) (1.2578,0.1128) (1.2677,0.1201) (1.2776,0.1275) (1.2875,0.1350) (1.2974,0.1427) (1.3073,0.1505) (1.3172,0.1584) (1.3271,0.1664) (1.3370,0.1745) (1.3469,0.1827) (1.3568,0.1910) (1.3667,0.1994) (1.3766,0.2079) (1.3865,0.2165) (1.3964,0.2251) (1.4063,0.2339) (1.4162,0.2427) (1.4261,0.2515) (1.4360,0.2605) (1.4459,0.2695) (1.4558,0.2786) (1.4657,0.2877) (1.4756,0.2969) (1.4855,0.3061) (1.4954,0.3154) (1.5053,0.3248) (1.5152,0.3342) (1.5251,0.3436) (1.5350,0.3531) (1.5449,0.3626) (1.5548,0.3722) (1.5647,0.3818) (1.5746,0.3915) (1.5845,0.4011) (1.5944,0.4109) (1.6043,0.4206) (1.6142,0.4304) (1.6241,0.4402) (1.6340,0.4501) (1.6439,0.4599) (1.6538,0.4698) (1.6637,0.4797) (1.6736,0.4897) (1.6835,0.4997) (1.6934,0.5097) (1.7033,0.5197) (1.7132,0.5297) (1.7231,0.5398) (1.7330,0.5499) (1.7429,0.5600) (1.7528,0.5701) (1.7627,0.5802) (1.7726,0.5904) (1.7825,0.6005) (1.7924,0.6107) (1.8023,0.6209) (1.8122,0.6311) (1.8221,0.6413) (1.8320,0.6516) (1.8419,0.6618) (1.8518,0.6721) (1.8617,0.6823) (1.8716,0.6926) (1.8815,0.7029) (1.8914,0.7132) (1.9013,0.7235) (1.9112,0.7338) (1.9211,0.7441) (1.9310,0.7545) (1.9409,0.7648) (1.9508,0.7752) (1.9607,0.7855) (1.9706,0.7959) (1.9805,0.8062) (1.9904,0.8166) (2.0003,0.8270) (2.0102,0.8374) (2.0201,0.8477) (2.0299,0.8581) (2.0398,0.8685) (2.0497,0.8789) (2.0596,0.8893) (2.0695,0.8997) (2.0794,0.9101) (2.0893,0.9206) (2.0992,0.9310) (2.1091,0.9414) (2.1190,0.9518) (2.1289,0.9622) (2.1388,0.9727) (2.1487,0.9831) (2.1586,0.9935) (2.1685,1.0039) (2.1784,1.0144) (2.1883,1.0248) (2.1982,1.0352) (2.2081,1.0457) (2.2180,1.0561) (2.2279,1.0666) (2.2378,1.0770) (2.2477,1.0874) (2.2576,1.0979) (2.2675,1.1083) (2.2774,1.1188) (2.2873,1.1292) (2.2972,1.1396) (2.3071,1.1501) (2.3170,1.1605) (2.3269,1.1710) (2.3368,1.1814) (2.3467,1.1918) (2.3566,1.2023) (2.3665,1.2127) (2.3764,1.2232) (2.3863,1.2336) (2.3962,1.2440) (2.4061,1.2545) (2.4160,1.2649) (2.4259,1.2754) (2.4358,1.2858) (2.4457,1.2962) (2.4556,1.3067) (2.4655,1.3171) (2.4754,1.3275) (2.4853,1.3379) (2.4952,1.3484) (2.5051,1.3588) (2.5150,1.3692) (2.5249,1.3797) (2.5348,1.3901) (2.5447,1.4005) (2.5546,1.4109) (2.5645,1.4213) (2.5744,1.4318) (2.5843,1.4422) (2.5942,1.4526) (2.6041,1.4630) (2.6140,1.4734) (2.6239,1.4838) (2.6338,1.4942) (2.6437,1.5046) (2.6536,1.5150) (2.6635,1.5254) (2.6734,1.5358) (2.6833,1.5462) (2.6932,1.5566) (2.7031,1.5670) (2.7130,1.5774) (2.7229,1.5878) (2.7328,1.5982) (2.7427,1.6086) (2.7526,1.6190) (2.7625,1.6294) (2.7724,1.6398) (2.7823,1.6501) (2.7922,1.6605) (2.8021,1.6709) (2.8120,1.6813) (2.8219,1.6917) (2.8318,1.7020) (2.8417,1.7124) (2.8516,1.7228) (2.8615,1.7331) (2.8714,1.7435) (2.8813,1.7539) (2.8912,1.7642) (2.9011,1.7746) (2.9110,1.7849) (2.9209,1.7953) (2.9308,1.8056) (2.9407,1.8160) (2.9506,1.8263) (2.9605,1.8367) (2.9704,1.8470) (2.9803,1.8574) (2.9902,1.8677) (3.0001,1.8781) (3.0100,1.8884) (3.0199,1.8987) (3.0298,1.9091) (3.0397,1.9194) (3.0496,1.9297) (3.0595,1.9401) (3.0694,1.9504) (3.0793,1.9607) (3.0892,1.9710) (3.0991,1.9814) (3.1090,1.9917) (3.1189,2.0020) (3.1288,2.0123) (3.1387,2.0226) (3.1486,2.0329) (3.1585,2.0432) (3.1684,2.0535) (3.1783,2.0638) (3.1882,2.0742) (3.1981,2.0845) (3.2080,2.0948) (3.2179,2.1051) (3.2278,2.1153) (3.2377,2.1256) (3.2476,2.1359) (3.2575,2.1462) (3.2674,2.1565) (3.2773,2.1668) (3.2872,2.1771) (3.2971,2.1874) (3.3070,2.1977) (3.3169,2.2079) (3.3268,2.2182) (3.3367,2.2285) (3.3466,2.2388) (3.3565,2.2490) (3.3664,2.2593) (3.3763,2.2696) (3.3862,2.2798) (3.3961,2.2901) (3.4060,2.3004) (3.4159,2.3106) (3.4258,2.3209) (3.4357,2.3312) (3.4456,2.3414) (3.4555,2.3517) (3.4654,2.3619) (3.4753,2.3722) (3.4852,2.3824) (3.4951,2.3927) (3.5050,2.4029) (3.5149,2.4132) (3.5248,2.4234) (3.5347,2.4337) (3.5446,2.4439) (3.5545,2.4541) (3.5644,2.4644) (3.5743,2.4746) (3.5842,2.4848) (3.5941,2.4951) (3.6040,2.5053) (3.6139,2.5155) (3.6238,2.5258) (3.6337,2.5360) (3.6436,2.5462) (3.6535,2.5564) (3.6634,2.5667) (3.6733,2.5769) (3.6832,2.5871) (3.6931,2.5973) (3.7030,2.6075) (3.7129,2.6177) (3.7228,2.6280) (3.7327,2.6382) (3.7426,2.6484) (3.7525,2.6586) (3.7624,2.6688) (3.7723,2.6790) (3.7822,2.6892) (3.7921,2.6994) (3.8020,2.7096) (3.8119,2.7198) (3.8218,2.7300) (3.8317,2.7402) (3.8416,2.7504) (3.8515,2.7606) (3.8614,2.7708) (3.8713,2.7810) (3.8812,2.7912) (3.8911,2.8014) (3.9010,2.8116) (3.9109,2.8217) (3.9208,2.8319) (3.9307,2.8421) (3.9406,2.8523) (3.9505,2.8625) (3.9604,2.8727) (3.9703,2.8828) (3.9802,2.8930) (3.9901,2.9032) (4.0000,2.9134)};
\fill (1,0) circle (1.5pt);
\fill (1.6487,0.4648) circle (1.5pt);
\Ticks{1,2,3}{-1,1,2,3}
\node[cDef] at (0.6,3.2){$(C_f)$};
\node[cRem] at (3.2,2.6){$(\Delta)$};
\node[cThm] at (3.5,2){$(T)$};
\end{GR}

\QQ{6}
\pa مجموعة التعريف متناظرة بالنسبة لـ $0$ ، ومن أجل كل $x\neq0$: $|-x|=|x|$ و $(-x)^2=x^2$ ، إذن $h(-x)=h(x)$: \bxt{$h$ دالة زوجية}.

\pb من أجل $x>0$: $h(x)=-x+1+\frac{\ln x}{x^2}=-\Big(x-1-\frac{\ln x}{x^2}\Big)=-f(x)$.

إذن على $]0;+\infty[$ ، $(C_h)$ هو نظير $(C_f)$ بالنسبة لمحور الفواصل، ثم نكمل الجزء الخاص بـ $]-\infty;0[$ بالتناظر بالنسبة لمحور التراتيب (لأن $h$ زوجية).
\end{hal}

% ===================== بكالوريا 2021 م1 =====================
\begin{tamrin}{بكالوريا 2021 -- الموضوع الأول}
\textbf{I)} الدالة $g$ معرفة على $\R$ بـ: \ $g(x)=2x^3-2x^2+3x-2$.
\QQ{1} بيّن أن الدالة $g$ متزايدة تماما على $\R$.
\QQ{2}
\begin{sq}
\qa بيّن أن المعادلة $g(x)=0$ تقبل حلا وحيدا $\alpha$ يحقق: $0{,}7<\alpha<0{,}8$.
\qb استنتج حسب قيم العدد الحقيقي $x$ إشارة $g(x)$.
\end{sq}

\textbf{II)} الدالة العددية $f$ معرفة على $]-\infty;0[\,\cup\,]0;+\infty[$ بـ: \ $f(x)=2x-1+\ln\Big(1+\dfrac{1-x}{x^2}\Big)$.

$(C_f)$ تمثيلها البياني في المستوي المنسوب إلى المعلم المتعامد والمتجانس $(O;\vec i,\vec j)$.
\QQ{1}
\begin{sq}
\qa بيّن أن $\lim\limits_{x\to0}f(x)=+\infty$ ، ثم فسّر النتيجة هندسيا.
\qb أحسب $\lim\limits_{x\to-\infty}f(x)$ و $\lim\limits_{x\to+\infty}f(x)$.
\end{sq}
\QQ{2}
\begin{sq}
\qa بيّن أنه من أجل كل عدد حقيقي غير معدوم $x$: \ $f'(x)=\dfrac{g(x)}{x(x^2-x+1)}$.
\qb استنتج أن $f$ متزايدة تماما على كل من $]-\infty;0[$ و $[\alpha;+\infty[$ ومتناقصة تماما على $]0;\alpha]$.
\qc شكّل جدول تغيّرات الدالة $f$.
\end{sq}
\QQ{3} بيّن أن المستقيم $(\Delta)$ ذا المعادلة $y=2x-1$ مقارب مائل لـ $(C_f)$ ، ثم ادرس وضعية $(C_f)$ بالنسبة لـ $(\Delta)$.
\QQ{4} بيّن أن $(C_f)$ يقبل مماسا $(T)$ موازيا لـ $(\Delta)$ في النقطة $A$ ذات الفاصلة $2$ ، ثم أكتب معادلة له.
\QQ{5} بيّن أن $(C_f)$ يقطع حامل محور الفواصل في نقطة وحيدة فاصلتها $\beta$ تحقق: $-0{,}5<\beta<-0{,}4$.
\QQ{6} أرسم $(\Delta)$ ، $(T)$ والمنحنى $(C_f)$. (نأخذ: $f(\alpha)\approx0{,}87$).
\end{tamrin}

\begin{hal}
\textbf{الجزء I}

\QQ{1} $g'(x)=6x^2-4x+3$ ، مميزه $\Delta'=(-2)^2-6\times3=-14<0$ ، إذن $g'(x)$ له إشارة معامل $x^2$ أي $g'(x)>0$ من أجل كل $x\in\R$: \bxt{$g$ متزايدة تماما على $\R$}.

\QQ{2}
\pa $g$ مستمرة ومتزايدة تماما على $[0{,}7;0{,}8]$ و: $g(0{,}7)=-0{,}194<0$ ، $g(0{,}8)=0{,}144>0$.
حسب مبرهنة القيم المتوسطة المعادلة $g(x)=0$ تقبل حلا وحيدا $\alpha$ في $]0{,}7;0{,}8[$ ، وهو وحيد في $\R$ لأن $g$ متزايدة تماما على $\R$: \bx{0{,}7<\alpha<0{,}8}.

\pb
\begin{TAB}
\tkzTabInit[lgt=2,espcl=3]{$x$/1,$g(x)$/1}{$-\infty$,$\alpha$,$+\infty$}
\tkzTabLine{,-,z,+,}
\end{TAB}

\textbf{الجزء II}

\textbf{ملاحظة:} $1+\frac{1-x}{x^2}=\frac{x^2-x+1}{x^2}$ و $x^2-x+1>0$ دائما (مميزه $-3<0$) ، إذن $f(x)=2x-1+\ln(x^2-x+1)-\ln(x^2)$.

\QQ{1}
\pa لما $x\to0$: $1-x\to1$ و $x^2\to0^+$ ، إذن $\frac{1-x}{x^2}\to+\infty$ ومنه $\ln\big(1+\frac{1-x}{x^2}\big)\to+\infty$ ، و $2x-1\to-1$: \bx{\lim_{x\to0}f(x)=+\infty}.

\textbf{التفسير:} المستقيم ذو المعادلة $x=0$ مقارب عمودي لـ $(C_f)$.

\pb $\frac{1-x}{x^2}=\frac{1}{x^2}-\frac1x\to0$ عند $\pm\infty$ ، إذن $\ln\big(1+\frac{1-x}{x^2}\big)\to\ln1=0$ ، ومنه:
\bx{\lim_{x\to-\infty}f(x)=-\infty} \ و \ \bx{\lim_{x\to+\infty}f(x)=+\infty}.

\QQ{2}
\pa باستعمال الكتابة السابقة:
\[f'(x)=2+\frac{2x-1}{x^2-x+1}-\frac{2}{x}=\frac{2x(x^2-x+1)+x(2x-1)-2(x^2-x+1)}{x(x^2-x+1)}.\]
البسط: $2x^3-2x^2+2x+2x^2-x-2x^2+2x-2=2x^3-2x^2+3x-2=g(x)$ ، إذن $f'(x)=\dfrac{g(x)}{x(x^2-x+1)}$.

\pb لأن $x^2-x+1>0$ فإن إشارة $f'(x)$ هي إشارة $\frac{g(x)}{x}$:
\begin{itemize}
\item على $]-\infty;0[$: $g(x)<0$ (لأن $x<\alpha$) و $x<0$ إذن $f'(x)>0$: $f$ متزايدة تماما.
\item على $]0;\alpha[$: $g(x)<0$ و $x>0$ إذن $f'(x)<0$: $f$ متناقصة تماما.
\item على $]\alpha;+\infty[$: $g(x)>0$ و $x>0$ إذن $f'(x)>0$: $f$ متزايدة تماما.
\end{itemize}
\pc
\begin{TAB}
\tkzTabInit[lgt=2,espcl=2.8]{$x$/1,$f'(x)$/1,$f(x)$/2}{$-\infty$,$0$,$\alpha$,$+\infty$}
\tkzTabLine{,+,d,-,z,+,}
\tkzTabVar{-/$-\infty$,+D+/$+\infty$/$+\infty$,-/$f(\alpha)$,+/$+\infty$}
\end{TAB}

\QQ{3} $f(x)-(2x-1)=\ln\Big(1+\dfrac{1-x}{x^2}\Big)$ و $\lim\limits_{x\to\pm\infty}\ln\Big(1+\dfrac{1-x}{x^2}\Big)=0$ ، إذن \bxt{$(\Delta): y=2x-1$ مقارب مائل لـ $(C_f)$ عند $\pm\infty$}.

الوضعية: $\ln\big(1+\frac{1-x}{x^2}\big)>0\iff 1+\frac{1-x}{x^2}>1\iff\frac{1-x}{x^2}>0\iff x<1$.
\begin{itemize}
\item على $]-\infty;0[$ و $]0;1[$: $(C_f)$ فوق $(\Delta)$.
\item على $]1;+\infty[$: $(C_f)$ تحت $(\Delta)$.
\item $(C_f)$ يقطع $(\Delta)$ في النقطة $(1;1)$.
\end{itemize}

\QQ{4} $f'(2)=\dfrac{g(2)}{2(4-2+1)}=\dfrac{16-8+6-2}{6}=\dfrac{12}{6}=2$ وهو ميل $(\Delta)$ ، إذن المماس في $A$ موازٍ لـ $(\Delta)$.

$f(2)=3+\ln\big(1-\frac14\big)=3+\ln\frac34$ ، ومنه $y=2(x-2)+3+\ln\frac34$: \bx{(T):\ y=2x-1+\ln\frac34}.

\QQ{5} على $]-\infty;0[$: $f$ مستمرة ومتزايدة تماما وتأخذ قيمها في $\R$ ، و:
$f(-0{,}5)=-2+\ln7\approx-0{,}05<0$ ، $f(-0{,}4)=-1{,}8+\ln9{,}75\approx0{,}48>0$.
إذن حسب مبرهنة القيم المتوسطة يوجد حل وحيد $\beta\in\,]-0{,}5;-0{,}4[$.

على $]0;+\infty[$: القيمة الحدية الصغرى لـ $f$ هي $f(\alpha)\approx0{,}87>0$ ، إذن $f(x)>0$ ولا يوجد تقاطع.

إذن $(C_f)$ يقطع حامل محور الفواصل في نقطة وحيدة فاصلتها \bx{-0{,}5<\beta<-0{,}4}.

\QQ{6} الرسم:
\begin{GR}{-3}{3}{-7}{7}{0.65}
\draw[cRem,thick] (-3,-7)--(3,5);
\draw[cThm,thick] (-3,-7.288)--(3,4.712);
\draw[cDef,thick] plot coordinates {(-3.0000,-6.6323) (-3.0000,-6.6323) (-3.0000,-6.6323) (-3.0000,-6.6323) (-3.0000,-6.6323) (-3.0000,-6.6323) (-3.0000,-6.6323) (-3.0000,-6.6323) (-3.0000,-6.6323) (-3.0000,-6.6323) (-3.0000,-6.6323) (-3.0000,-6.6323) (-3.0000,-6.6323) (-3.0000,-6.6323) (-3.0000,-6.6323) (-3.0000,-6.6323) (-3.0000,-6.6322) (-3.0000,-6.6322) (-3.0000,-6.6322) (-3.0000,-6.6322) (-3.0000,-6.6322) (-3.0000,-6.6322) (-3.0000,-6.6322) (-3.0000,-6.6322) (-3.0000,-6.6322) (-3.0000,-6.6322) (-3.0000,-6.6322) (-3.0000,-6.6322) (-3.0000,-6.6322) (-3.0000,-6.6322) (-3.0000,-6.6322) (-3.0000,-6.6322) (-3.0000,-6.6322) (-3.0000,-6.6322) (-3.0000,-6.6322) (-3.0000,-6.6322) (-3.0000,-6.6322) (-3.0000,-6.6322) (-3.0000,-6.6322) (-3.0000,-6.6322) (-3.0000,-6.6322) (-3.0000,-6.6322) (-3.0000,-6.6322) (-3.0000,-6.6322) (-3.0000,-6.6322) (-3.0000,-6.6322) (-3.0000,-6.6322) (-2.9999,-6.6322) (-2.9999,-6.6322) (-2.9999,-6.6322) (-2.9999,-6.6322) (-2.9999,-6.6321) (-2.9999,-6.6321) (-2.9999,-6.6321) (-2.9999,-6.6321) (-2.9999,-6.6321) (-2.9999,-6.6321) (-2.9999,-6.6321) (-2.9999,-6.6321) (-2.9999,-6.6321) (-2.9999,-6.6321) (-2.9999,-6.6321) (-2.9999,-6.6321) (-2.9999,-6.6321) (-2.9999,-6.6320) (-2.9999,-6.6320) (-2.9999,-6.6320) (-2.9999,-6.6320) (-2.9999,-6.6320) (-2.9999,-6.6320) (-2.9999,-6.6320) (-2.9998,-6.6319) (-2.9998,-6.6319) (-2.9998,-6.6319) (-2.9998,-6.6319) (-2.9998,-6.6319) (-2.9998,-6.6319) (-2.9998,-6.6318) (-2.9998,-6.6318) (-2.9998,-6.6318) (-2.9998,-6.6318) (-2.9998,-6.6318) (-2.9997,-6.6317) (-2.9997,-6.6317) (-2.9997,-6.6317) (-2.9997,-6.6317) (-2.9997,-6.6316) (-2.9997,-6.6316) (-2.9997,-6.6316) (-2.9996,-6.6315) (-2.9996,-6.6315) (-2.9996,-6.6315) (-2.9996,-6.6314) (-2.9996,-6.6314) (-2.9996,-6.6313) (-2.9995,-6.6313) (-2.9995,-6.6312) (-2.9995,-6.6312) (-2.9995,-6.6311) (-2.9994,-6.6311) (-2.9994,-6.6310) (-2.9994,-6.6310) (-2.9994,-6.6309) (-2.9993,-6.6308) (-2.9993,-6.6308) (-2.9993,-6.6307) (-2.9992,-6.6306) (-2.9992,-6.6306) (-2.9992,-6.6305) (-2.9991,-6.6304) (-2.9991,-6.6303) (-2.9990,-6.6302) (-2.9990,-6.6301) (-2.9989,-6.6300) (-2.9989,-6.6299) (-2.9988,-6.6298) (-2.9988,-6.6297) (-2.9987,-6.6295) (-2.9987,-6.6294) (-2.9986,-6.6293) (-2.9985,-6.6291) (-2.9985,-6.6290) (-2.9984,-6.6288) (-2.9983,-6.6287) (-2.9982,-6.6285) (-2.9981,-6.6283) (-2.9980,-6.6281) (-2.9980,-6.6279) (-2.9979,-6.6277) (-2.9978,-6.6275) (-2.9977,-6.6273) (-2.9975,-6.6270) (-2.9974,-6.6268) (-2.9973,-6.6265) (-2.9972,-6.6263) (-2.9970,-6.6260) (-2.9969,-6.6257) (-2.9968,-6.6254) (-2.9966,-6.6250) (-2.9964,-6.6247) (-2.9963,-6.6243) (-2.9961,-6.6240) (-2.9959,-6.6236) (-2.9957,-6.6231) (-2.9955,-6.6227) (-2.9953,-6.6223) (-2.9951,-6.6218) (-2.9948,-6.6213) (-2.9946,-6.6208) (-2.9943,-6.6202) (-2.9941,-6.6197) (-2.9938,-6.6191) (-2.9935,-6.6184) (-2.9932,-6.6178) (-2.9929,-6.6171) (-2.9925,-6.6164) (-2.9925,-6.6163) (-2.9922,-6.6156) (-2.9918,-6.6148) (-2.9914,-6.6140) (-2.9910,-6.6131) (-2.9906,-6.6122) (-2.9901,-6.6113) (-2.9897,-6.6103) (-2.9892,-6.6092) (-2.9887,-6.6081) (-2.9881,-6.6070) (-2.9876,-6.6058) (-2.9870,-6.6046) (-2.9864,-6.6032) (-2.9857,-6.6019) (-2.9850,-6.6004) (-2.9850,-6.6003) (-2.9843,-6.5989) (-2.9836,-6.5973) (-2.9828,-6.5957) (-2.9820,-6.5939) (-2.9811,-6.5921) (-2.9802,-6.5902) (-2.9793,-6.5882) (-2.9783,-6.5861) (-2.9775,-6.5843) (-2.9773,-6.5840) (-2.9762,-6.5817) (-2.9751,-6.5793) (-2.9739,-6.5768) (-2.9727,-6.5741) (-2.9714,-6.5714) (-2.9700,-6.5685) (-2.9699,-6.5683) (-2.9686,-6.5655) (-2.9671,-6.5623) (-2.9656,-6.5590) (-2.9640,-6.5555) (-2.9624,-6.5522) (-2.9622,-6.5519) (-2.9605,-6.5480) (-2.9586,-6.5441) (-2.9566,-6.5399) (-2.9549,-6.5362) (-2.9546,-6.5355) (-2.9524,-6.5309) (-2.9502,-6.5261) (-2.9478,-6.5211) (-2.9474,-6.5202) (-2.9453,-6.5158) (-2.9427,-6.5103) (-2.9400,-6.5045) (-2.9399,-6.5041) (-2.9324,-6.4881) (-2.9248,-6.4721) (-2.9173,-6.4560) (-2.9098,-6.4400) (-2.9023,-6.4239) (-2.8948,-6.4078) (-2.8873,-6.3918) (-2.8797,-6.3757) (-2.8722,-6.3596) (-2.8647,-6.3435) (-2.8572,-6.3274) (-2.8497,-6.3113) (-2.8422,-6.2952) (-2.8346,-6.2791) (-2.8271,-6.2630) (-2.8196,-6.2469) (-2.8121,-6.2308) (-2.8046,-6.2146) (-2.7971,-6.1985) (-2.7895,-6.1823) (-2.7820,-6.1662) (-2.7745,-6.1500) (-2.7670,-6.1339) (-2.7595,-6.1177) (-2.7520,-6.1015) (-2.7444,-6.0853) (-2.7369,-6.0691) (-2.7294,-6.0530) (-2.7219,-6.0368) (-2.7144,-6.0205) (-2.7069,-6.0043) (-2.6993,-5.9881) (-2.6918,-5.9719) (-2.6843,-5.9557) (-2.6768,-5.9394) (-2.6693,-5.9232) (-2.6618,-5.9069) (-2.6543,-5.8906) (-2.6467,-5.8744) (-2.6392,-5.8581) (-2.6317,-5.8418) (-2.6242,-5.8255) (-2.6167,-5.8092) (-2.6092,-5.7929) (-2.6016,-5.7766) (-2.5941,-5.7603) (-2.5866,-5.7440) (-2.5791,-5.7276) (-2.5716,-5.7113) (-2.5641,-5.6950) (-2.5565,-5.6786) (-2.5490,-5.6622) (-2.5415,-5.6459) (-2.5340,-5.6295) (-2.5265,-5.6131) (-2.5190,-5.5967) (-2.5114,-5.5803) (-2.5039,-5.5639) (-2.4964,-5.5475) (-2.4889,-5.5310) (-2.4814,-5.5146) (-2.4739,-5.4982) (-2.4663,-5.4817) (-2.4588,-5.4652) (-2.4513,-5.4488) (-2.4438,-5.4323) (-2.4363,-5.4158) (-2.4288,-5.3993) (-2.4212,-5.3828) (-2.4137,-5.3663) (-2.4062,-5.3498) (-2.3987,-5.3332) (-2.3912,-5.3167) (-2.3837,-5.3001) (-2.3761,-5.2836) (-2.3686,-5.2670) (-2.3611,-5.2504) (-2.3536,-5.2338) (-2.3461,-5.2172) (-2.3386,-5.2006) (-2.3311,-5.1840) (-2.3235,-5.1674) (-2.3160,-5.1507) (-2.3085,-5.1341) (-2.3010,-5.1174) (-2.2935,-5.1007) (-2.2860,-5.0840) (-2.2784,-5.0674) (-2.2709,-5.0507) (-2.2634,-5.0339) (-2.2559,-5.0172) (-2.2484,-5.0005) (-2.2409,-4.9837) (-2.2333,-4.9670) (-2.2258,-4.9502) (-2.2183,-4.9334) (-2.2108,-4.9166) (-2.2033,-4.8998) (-2.1958,-4.8830) (-2.1882,-4.8662) (-2.1807,-4.8493) (-2.1732,-4.8325) (-2.1657,-4.8156) (-2.1582,-4.7987) (-2.1507,-4.7818) (-2.1431,-4.7649) (-2.1356,-4.7480) (-2.1281,-4.7311) (-2.1206,-4.7141) (-2.1131,-4.6972) (-2.1056,-4.6802) (-2.0980,-4.6632) (-2.0905,-4.6462) (-2.0830,-4.6292) (-2.0755,-4.6122) (-2.0680,-4.5951) (-2.0605,-4.5781) (-2.0529,-4.5610) (-2.0454,-4.5439) (-2.0379,-4.5269) (-2.0304,-4.5097) (-2.0229,-4.4926) (-2.0154,-4.4755) (-2.0078,-4.4583) (-2.0003,-4.4411) (-1.9928,-4.4240) (-1.9853,-4.4068) (-1.9778,-4.3895) (-1.9703,-4.3723) (-1.9628,-4.3551) (-1.9552,-4.3378) (-1.9477,-4.3205) (-1.9402,-4.3032) (-1.9327,-4.2859) (-1.9252,-4.2685) (-1.9177,-4.2512) (-1.9101,-4.2338) (-1.9026,-4.2164) (-1.8951,-4.1990) (-1.8876,-4.1816) (-1.8801,-4.1642) (-1.8726,-4.1467) (-1.8650,-4.1292) (-1.8575,-4.1117) (-1.8500,-4.0942) (-1.8425,-4.0767) (-1.8350,-4.0591) (-1.8275,-4.0415) (-1.8199,-4.0240) (-1.8124,-4.0063) (-1.8049,-3.9887) (-1.7974,-3.9710) (-1.7899,-3.9534) (-1.7824,-3.9357) (-1.7748,-3.9179) (-1.7673,-3.9002) (-1.7598,-3.8824) (-1.7523,-3.8647) (-1.7448,-3.8468) (-1.7373,-3.8290) (-1.7297,-3.8112) (-1.7222,-3.7933) (-1.7147,-3.7754) (-1.7072,-3.7575) (-1.6997,-3.7395) (-1.6922,-3.7215) (-1.6846,-3.7035) (-1.6771,-3.6855) (-1.6696,-3.6675) (-1.6621,-3.6494) (-1.6546,-3.6313) (-1.6471,-3.6132) (-1.6396,-3.5950) (-1.6320,-3.5769) (-1.6245,-3.5586) (-1.6170,-3.5404) (-1.6095,-3.5222) (-1.6020,-3.5039) (-1.5945,-3.4856) (-1.5869,-3.4672) (-1.5794,-3.4488) (-1.5719,-3.4304) (-1.5644,-3.4120) (-1.5569,-3.3935) (-1.5494,-3.3750) (-1.5418,-3.3565) (-1.5343,-3.3379) (-1.5268,-3.3194) (-1.5193,-3.3007) (-1.5118,-3.2821) (-1.5043,-3.2634) (-1.4967,-3.2447) (-1.4892,-3.2259) (-1.4817,-3.2071) (-1.4742,-3.1883) (-1.4667,-3.1694) (-1.4592,-3.1505) (-1.4516,-3.1316) (-1.4441,-3.1126) (-1.4366,-3.0936) (-1.4291,-3.0746) (-1.4216,-3.0555) (-1.4141,-3.0364) (-1.4065,-3.0172) (-1.3990,-2.9980) (-1.3915,-2.9787) (-1.3840,-2.9595) (-1.3765,-2.9401) (-1.3690,-2.9208) (-1.3614,-2.9013) (-1.3539,-2.8819) (-1.3464,-2.8624) (-1.3389,-2.8428) (-1.3314,-2.8233) (-1.3239,-2.8036) (-1.3164,-2.7839) (-1.3088,-2.7642) (-1.3013,-2.7444) (-1.2938,-2.7246) (-1.2863,-2.7047) (-1.2788,-2.6848) (-1.2713,-2.6648) (-1.2637,-2.6448) (-1.2562,-2.6247) (-1.2487,-2.6045) (-1.2412,-2.5843) (-1.2337,-2.5641) (-1.2262,-2.5438) (-1.2186,-2.5234) (-1.2111,-2.5030) (-1.2036,-2.4825) (-1.1961,-2.4620) (-1.1886,-2.4414) (-1.1811,-2.4207) (-1.1735,-2.4000) (-1.1660,-2.3792) (-1.1585,-2.3583) (-1.1510,-2.3374) (-1.1435,-2.3164) (-1.1360,-2.2954) (-1.1284,-2.2743) (-1.1209,-2.2531) (-1.1134,-2.2318) (-1.1059,-2.2105) (-1.0984,-2.1890) (-1.0909,-2.1676) (-1.0833,-2.1460) (-1.0758,-2.1244) (-1.0683,-2.1026) (-1.0608,-2.0808) (-1.0533,-2.0590) (-1.0458,-2.0370) (-1.0382,-2.0149) (-1.0307,-1.9928) (-1.0232,-1.9706) (-1.0157,-1.9483) (-1.0082,-1.9259) (-1.0007,-1.9034) (-0.9932,-1.8808) (-0.9856,-1.8581) (-0.9781,-1.8353) (-0.9706,-1.8125) (-0.9631,-1.7895) (-0.9556,-1.7664) (-0.9481,-1.7432) (-0.9405,-1.7199) (-0.9330,-1.6965) (-0.9255,-1.6730) (-0.9180,-1.6494) (-0.9105,-1.6256) (-0.9030,-1.6017) (-0.8954,-1.5778) (-0.8879,-1.5537) (-0.8804,-1.5294) (-0.8729,-1.5051) (-0.8654,-1.4806) (-0.8579,-1.4560) (-0.8503,-1.4312) (-0.8428,-1.4063) (-0.8353,-1.3813) (-0.8278,-1.3561) (-0.8203,-1.3308) (-0.8128,-1.3053) (-0.8052,-1.2797) (-0.7977,-1.2539) (-0.7902,-1.2280) (-0.7827,-1.2019) (-0.7752,-1.1756) (-0.7677,-1.1491) (-0.7601,-1.1225) (-0.7526,-1.0957) (-0.7451,-1.0687) (-0.7376,-1.0416) (-0.7301,-1.0142) (-0.7226,-0.9867) (-0.7150,-0.9589) (-0.7075,-0.9310) (-0.7000,-0.9028) (-0.6925,-0.8744) (-0.6850,-0.8458) (-0.6775,-0.8170) (-0.6699,-0.7880) (-0.6624,-0.7587) (-0.6549,-0.7291) (-0.6474,-0.6994) (-0.6399,-0.6693) (-0.6324,-0.6390) (-0.6249,-0.6085) (-0.6173,-0.5776) (-0.6098,-0.5465) (-0.6023,-0.5151) (-0.5948,-0.4834) (-0.5873,-0.4514) (-0.5798,-0.4190) (-0.5722,-0.3864) (-0.5647,-0.3534) (-0.5572,-0.3201) (-0.5497,-0.2864) (-0.5422,-0.2523) (-0.5347,-0.2179) (-0.5271,-0.1831) (-0.5196,-0.1479) (-0.5121,-0.1123) (-0.5046,-0.0763) (-0.4971,-0.0398) (-0.4896,-0.0029) (-0.4820,0.0344) (-0.4745,0.0723) (-0.4670,0.1106) (-0.4595,0.1495) (-0.4520,0.1889) (-0.4445,0.2288) (-0.4369,0.2693) (-0.4294,0.3104) (-0.4219,0.3520) (-0.4144,0.3944) (-0.4069,0.4373) (-0.3994,0.4810) (-0.3918,0.5253) (-0.3843,0.5704) (-0.3768,0.6163) (-0.3693,0.6629) (-0.3618,0.7104) (-0.3543,0.7588) (-0.3467,0.8080) (-0.3392,0.8582) (-0.3317,0.9094) (-0.3242,0.9616) (-0.3167,1.0149) (-0.3092,1.0693) (-0.3017,1.1249) (-0.2941,1.1817) (-0.2866,1.2399) (-0.2791,1.2994) (-0.2716,1.3604) (-0.2641,1.4230) (-0.2566,1.4872) (-0.2490,1.5531) (-0.2415,1.6208) (-0.2340,1.6905) (-0.2265,1.7623) (-0.2190,1.8363) (-0.2115,1.9126) (-0.2039,1.9915) (-0.1964,2.0732) (-0.1889,2.1578) (-0.1814,2.2456) (-0.1739,2.3368) (-0.1664,2.4318) (-0.1588,2.5310) (-0.1513,2.6346) (-0.1438,2.7432) (-0.1363,2.8573) (-0.1288,2.9775) (-0.1213,3.1046) (-0.1137,3.2394) (-0.1062,3.3830) (-0.0987,3.5366) (-0.0912,3.7020) (-0.0837,3.8810) (-0.0762,4.0762) (-0.0686,4.2911) (-0.0611,4.5301) (-0.0536,4.7996) (-0.0461,5.1089) (-0.0386,5.4721) (-0.0311,5.9127) (-0.0235,6.4741) (-0.0160,7.2504)};
\draw[cDef,thick] plot coordinates {(0.0146,7.4628) (0.0153,7.3771) (0.0160,7.2912) (0.0160,7.2824) (0.0167,7.2050) (0.0174,7.1186) (0.0182,7.0321) (0.0190,6.9453) (0.0199,6.8584) (0.0208,6.7713) (0.0217,6.6840) (0.0227,6.5966) (0.0235,6.5212) (0.0237,6.5091) (0.0248,6.4215) (0.0259,6.3337) (0.0271,6.2459) (0.0283,6.1579) (0.0296,6.0699) (0.0310,5.9818) (0.0311,5.9749) (0.0324,5.8936) (0.0339,5.8054) (0.0354,5.7171) (0.0370,5.6288) (0.0386,5.5493) (0.0388,5.5404) (0.0405,5.4521) (0.0424,5.3637) (0.0444,5.2753) (0.0461,5.2012) (0.0464,5.1869) (0.0486,5.0986) (0.0508,5.0102) (0.0532,4.9220) (0.0536,4.9070) (0.0557,4.8337) (0.0583,4.7455) (0.0610,4.6574) (0.0611,4.6526) (0.0686,4.4288) (0.0762,4.2291) (0.0837,4.0491) (0.0912,3.8854) (0.0987,3.7353) (0.1062,3.5970) (0.1137,3.4688) (0.1213,3.3495) (0.1288,3.2379) (0.1363,3.1332) (0.1438,3.0348) (0.1513,2.9418) (0.1588,2.8540) (0.1664,2.7706) (0.1739,2.6915) (0.1814,2.6162) (0.1889,2.5445) (0.1964,2.4760) (0.2039,2.4106) (0.2115,2.3480) (0.2190,2.2880) (0.2265,2.2305) (0.2340,2.1753) (0.2415,2.1223) (0.2490,2.0713) (0.2566,2.0223) (0.2641,1.9751) (0.2716,1.9297) (0.2791,1.8859) (0.2866,1.8437) (0.2941,1.8030) (0.3017,1.7637) (0.3092,1.7258) (0.3167,1.6892) (0.3242,1.6539) (0.3317,1.6198) (0.3392,1.5868) (0.3467,1.5550) (0.3543,1.5242) (0.3618,1.4945) (0.3693,1.4657) (0.3768,1.4380) (0.3843,1.4112) (0.3918,1.3853) (0.3994,1.3602) (0.4069,1.3360) (0.4144,1.3127) (0.4219,1.2902) (0.4294,1.2684) (0.4369,1.2474) (0.4445,1.2271) (0.4520,1.2076) (0.4595,1.1888) (0.4670,1.1706) (0.4745,1.1531) (0.4820,1.1363) (0.4896,1.1201) (0.4971,1.1045) (0.5046,1.0895) (0.5121,1.0752) (0.5196,1.0614) (0.5271,1.0482) (0.5347,1.0355) (0.5422,1.0234) (0.5497,1.0118) (0.5572,1.0007) (0.5647,0.9902) (0.5722,0.9801) (0.5798,0.9706) (0.5873,0.9615) (0.5948,0.9529) (0.6023,0.9448) (0.6098,0.9371) (0.6173,0.9299) (0.6249,0.9231) (0.6324,0.9167) (0.6399,0.9108) (0.6474,0.9053) (0.6549,0.9001) (0.6624,0.8954) (0.6699,0.8911) (0.6775,0.8872) (0.6850,0.8836) (0.6925,0.8804) (0.7000,0.8776) (0.7075,0.8752) (0.7150,0.8731) (0.7226,0.8713) (0.7301,0.8699) (0.7376,0.8688) (0.7451,0.8680) (0.7526,0.8676) (0.7601,0.8675) (0.7677,0.8677) (0.7752,0.8682) (0.7827,0.8690) (0.7902,0.8701) (0.7977,0.8715) (0.8052,0.8731) (0.8128,0.8751) (0.8203,0.8773) (0.8278,0.8798) (0.8353,0.8825) (0.8428,0.8855) (0.8503,0.8888) (0.8579,0.8923) (0.8654,0.8961) (0.8729,0.9001) (0.8804,0.9043) (0.8879,0.9088) (0.8954,0.9135) (0.9030,0.9184) (0.9105,0.9235) (0.9180,0.9288) (0.9255,0.9344) (0.9330,0.9402) (0.9405,0.9461) (0.9481,0.9523) (0.9556,0.9587) (0.9631,0.9652) (0.9706,0.9719) (0.9781,0.9789) (0.9856,0.9859) (0.9932,0.9932) (1.0007,1.0007) (1.0082,1.0083) (1.0157,1.0161) (1.0232,1.0240) (1.0307,1.0321) (1.0382,1.0404) (1.0458,1.0488) (1.0533,1.0573) (1.0608,1.0661) (1.0683,1.0749) (1.0758,1.0839) (1.0833,1.0930) (1.0909,1.1023) (1.0984,1.1117) (1.1059,1.1212) (1.1134,1.1309) (1.1209,1.1407) (1.1284,1.1506) (1.1360,1.1606) (1.1435,1.1707) (1.1510,1.1810) (1.1585,1.1913) (1.1660,1.2018) (1.1735,1.2124) (1.1811,1.2231) (1.1886,1.2339) (1.1961,1.2448) (1.2036,1.2558) (1.2111,1.2668) (1.2186,1.2780) (1.2262,1.2893) (1.2337,1.3007) (1.2412,1.3121) (1.2487,1.3236) (1.2562,1.3353) (1.2637,1.3470) (1.2713,1.3588) (1.2788,1.3706) (1.2863,1.3826) (1.2938,1.3946) (1.3013,1.4067) (1.3088,1.4189) (1.3164,1.4311) (1.3239,1.4434) (1.3314,1.4558) (1.3389,1.4683) (1.3464,1.4808) (1.3539,1.4933) (1.3614,1.5060) (1.3690,1.5187) (1.3765,1.5314) (1.3840,1.5443) (1.3915,1.5571) (1.3990,1.5701) (1.4065,1.5831) (1.4141,1.5961) (1.4216,1.6092) (1.4291,1.6223) (1.4366,1.6355) (1.4441,1.6488) (1.4516,1.6621) (1.4592,1.6754) (1.4667,1.6888) (1.4742,1.7022) (1.4817,1.7157) (1.4892,1.7292) (1.4967,1.7428) (1.5043,1.7564) (1.5118,1.7700) (1.5193,1.7837) (1.5268,1.7974) (1.5343,1.8112) (1.5418,1.8250) (1.5494,1.8388) (1.5569,1.8527) (1.5644,1.8666) (1.5719,1.8806) (1.5794,1.8945) (1.5869,1.9085) (1.5945,1.9226) (1.6020,1.9366) (1.6095,1.9507) (1.6170,1.9649) (1.6245,1.9790) (1.6320,1.9932) (1.6396,2.0074) (1.6471,2.0216) (1.6546,2.0359) (1.6621,2.0502) (1.6696,2.0645) (1.6771,2.0789) (1.6846,2.0932) (1.6922,2.1076) (1.6997,2.1220) (1.7072,2.1365) (1.7147,2.1509) (1.7222,2.1654) (1.7297,2.1799) (1.7373,2.1944) (1.7448,2.2090) (1.7523,2.2235) (1.7598,2.2381) (1.7673,2.2527) (1.7748,2.2674) (1.7824,2.2820) (1.7899,2.2967) (1.7974,2.3113) (1.8049,2.3260) (1.8124,2.3407) (1.8199,2.3555) (1.8275,2.3702) (1.8350,2.3850) (1.8425,2.3997) (1.8500,2.4145) (1.8575,2.4293) (1.8650,2.4441) (1.8726,2.4590) (1.8801,2.4738) (1.8876,2.4887) (1.8951,2.5035) (1.9026,2.5184) (1.9101,2.5333) (1.9177,2.5482) (1.9252,2.5632) (1.9327,2.5781) (1.9402,2.5930) (1.9477,2.6080) (1.9552,2.6230) (1.9628,2.6379) (1.9703,2.6529) (1.9778,2.6679) (1.9853,2.6829) (1.9928,2.6980) (2.0003,2.7130) (2.0078,2.7280) (2.0154,2.7431) (2.0229,2.7581) (2.0304,2.7732) (2.0379,2.7883) (2.0454,2.8033) (2.0529,2.8184) (2.0605,2.8335) (2.0680,2.8486) (2.0755,2.8638) (2.0830,2.8789) (2.0905,2.8940) (2.0980,2.9091) (2.1056,2.9243) (2.1131,2.9394) (2.1206,2.9546) (2.1281,2.9697) (2.1356,2.9849) (2.1431,3.0001) (2.1507,3.0153) (2.1582,3.0305) (2.1657,3.0457) (2.1732,3.0608) (2.1807,3.0761) (2.1882,3.0913) (2.1958,3.1065) (2.2033,3.1217) (2.2108,3.1369) (2.2183,3.1522) (2.2258,3.1674) (2.2333,3.1826) (2.2409,3.1979) (2.2484,3.2131) (2.2559,3.2284) (2.2634,3.2436) (2.2709,3.2589) (2.2784,3.2742) (2.2860,3.2894) (2.2935,3.3047) (2.3010,3.3200) (2.3085,3.3353) (2.3160,3.3505) (2.3235,3.3658) (2.3311,3.3811) (2.3386,3.3964) (2.3461,3.4117) (2.3536,3.4270) (2.3611,3.4423) (2.3686,3.4576) (2.3761,3.4729) (2.3837,3.4882) (2.3912,3.5036) (2.3987,3.5189) (2.4062,3.5342) (2.4137,3.5495) (2.4212,3.5648) (2.4288,3.5802) (2.4363,3.5955) (2.4438,3.6108) (2.4513,3.6262) (2.4588,3.6415) (2.4663,3.6569) (2.4739,3.6722) (2.4814,3.6875) (2.4889,3.7029) (2.4964,3.7182) (2.5039,3.7336) (2.5114,3.7489) (2.5190,3.7643) (2.5265,3.7796) (2.5340,3.7950) (2.5415,3.8103) (2.5490,3.8257) (2.5565,3.8411) (2.5641,3.8564) (2.5716,3.8718) (2.5791,3.8872) (2.5866,3.9025) (2.5941,3.9179) (2.6016,3.9333) (2.6092,3.9486) (2.6167,3.9640) (2.6242,3.9794) (2.6317,3.9947) (2.6392,4.0101) (2.6467,4.0255) (2.6543,4.0409) (2.6618,4.0562) (2.6693,4.0716) (2.6768,4.0870) (2.6843,4.1024) (2.6918,4.1178) (2.6993,4.1331) (2.7069,4.1485) (2.7144,4.1639) (2.7219,4.1793) (2.7294,4.1947) (2.7369,4.2101) (2.7444,4.2254) (2.7520,4.2408) (2.7595,4.2562) (2.7670,4.2716) (2.7745,4.2870) (2.7820,4.3024) (2.7895,4.3178) (2.7971,4.3331) (2.8046,4.3485) (2.8121,4.3639) (2.8196,4.3793) (2.8271,4.3947) (2.8346,4.4101) (2.8422,4.4255) (2.8497,4.4409) (2.8572,4.4563) (2.8647,4.4717) (2.8722,4.4870) (2.8797,4.5024) (2.8873,4.5178) (2.8948,4.5332) (2.9023,4.5486) (2.9098,4.5640) (2.9173,4.5794) (2.9248,4.5948) (2.9324,4.6102) (2.9399,4.6256) (2.9474,4.6409) (2.9549,4.6563) (2.9624,4.6717) (2.9699,4.6871) (2.9775,4.7025) (2.9850,4.7179) (2.9925,4.7333) (3.0000,4.7487)};
\fill (1,1) circle (2pt); \fill (2,2.712) circle (2pt) node[left]{$A$};
\fill (0.7592,0.8675) circle (2pt); \fill (-0.489,0) circle (2pt);
\Ticks{-2,-1,1,2}{-6,-4,-2,2,4,6}
\node[cDef] at (-1.2,4){$(C_f)$};
\node[cRem] at (2.3,5.6){$(\Delta)$};
\node[cThm] at (2.7,3.2){$(T)$};
\end{GR}
\end{hal}

% ===================== بكالوريا 2022 م1 =====================
\begin{tamrin}{بكالوريا 2022 -- الموضوع الأول}
\textbf{I)} $g$ الدالة العددية المعرفة على $]0;+\infty[$ بـ: \ $g(x)=\dfrac{2x^2-2x-1}{x^2}+\ln x$.
\QQ{1} بيّن أن الدالة $g$ متزايدة تماما على $]0;+\infty[$.
\QQ{2}
\begin{sq}
\qa بيّن أن المعادلة $g(x)=0$ تقبل حلا وحيدا $\alpha$ حيث: $1{,}2<\alpha<1{,}3$.
\qb استنتج إشارة $g(x)$ على $]0;+\infty[$.
\end{sq}

\textbf{II)} نعتبر الدالة العددية $f$ المعرفة على $]0;+\infty[$ بـ: \ $f(x)=\Big(\dfrac1x-2-\ln x\Big)e^{-x}$.

$(C_f)$ تمثيلها البياني في مستوٍ منسوب إلى المعلم المتعامد والمتجانس $(O;\vec i,\vec j)$.
\QQ{1}
\begin{sq}
\qa أحسب $\lim\limits_{x\to0}f(x)$ ، ثم بيّن أن: $\lim\limits_{x\to+\infty}f(x)=0$.
\qb فسّر النتيجتين السابقتين بيانيا.
\end{sq}
\QQ{2}
\begin{sq}
\qa بيّن أنه من أجل كل عدد حقيقي $x$ موجب تماما: \ $f'(x)=\dfrac{g(x)}{e^x}$.
\qb استنتج اتجاه تغيّر الدالة $f$ وشكّل جدول تغيّراتها.
\end{sq}
\QQ{3} أنشئ المنحنى $(C_f)$. (نأخذ: $f(\alpha)\approx-0{,}4$ و $f(0{,}65)\approx0$).
\end{tamrin}

\begin{hal}
\textbf{الجزء I}

\QQ{1} نكتب $g(x)=2-\frac2x-\frac{1}{x^2}+\ln x$ ، إذن:
\[g'(x)=\frac{2}{x^2}+\frac{2}{x^3}+\frac1x=\frac{x^2+2x+2}{x^3}.\]
$x^2+2x+2=(x+1)^2+1>0$ و $x^3>0$ ، إذن $g'(x)>0$: \bxt{$g$ متزايدة تماما على $]0;+\infty[$}.

\QQ{2}
\pa $g$ مستمرة ومتزايدة تماما على $[1{,}2;1{,}3]$ ، و $g(1{,}2)\approx-0{,}18<0$ ، $g(1{,}3)\approx0{,}13>0$.
حسب مبرهنة القيم المتوسطة: المعادلة $g(x)=0$ تقبل حلا وحيدا \bx{\alpha\in\,]1{,}2;1{,}3[}.

\pb
\begin{TAB}
\tkzTabInit[lgt=2,espcl=3]{$x$/1,$g(x)$/1}{$0$,$\alpha$,$+\infty$}
\tkzTabLine{d,-,z,+,}
\end{TAB}

\textbf{الجزء II}

\QQ{1}
\pa لما $x\to0$: $\frac1x\to+\infty$ و $-\ln x\to+\infty$ ، إذن $\frac1x-2-\ln x\to+\infty$ ، و $e^{-x}\to1$: \bx{\lim_{x\to0}f(x)=+\infty}.

لما $x\to+\infty$: نكتب $f(x)=\Big(\frac1x-2\Big)e^{-x}-\frac{\ln x}{x}\cdot\frac{x}{e^x}$.
لدينا $\big(\frac1x-2\big)\to-2$ و $e^{-x}\to0$ ، $\frac{\ln x}{x}\to0$ و $\frac{x}{e^x}\to0$ ، إذن \bx{\lim_{x\to+\infty}f(x)=0}.

\pb المستقيم $x=0$ مقارب عمودي لـ $(C_f)$ ، والمستقيم $y=0$ مقارب أفقي لـ $(C_f)$ بجوار $+\infty$.

\QQ{2}
\pa $f=u\cdot v$ مع $u=\frac1x-2-\ln x$ ، $u'=-\frac{1}{x^2}-\frac1x$ و $v=e^{-x}$ ، $v'=-e^{-x}$:
\begin{align*}
f'(x)&=\Big(-\frac{1}{x^2}-\frac1x\Big)e^{-x}-\Big(\frac1x-2-\ln x\Big)e^{-x}
=e^{-x}\Big(2-\frac2x-\frac{1}{x^2}+\ln x\Big)=\frac{g(x)}{e^x}.
\end{align*}
\pb $e^x>0$ ، فإشارة $f'(x)$ هي إشارة $g(x)$: \bxt{$f$ متناقصة تماما على $]0;\alpha]$ ومتزايدة تماما على $[\alpha;+\infty[$}.
\begin{TAB}
\tkzTabInit[lgt=2,espcl=3]{$x$/1,$f'(x)$/1,$f(x)$/2}{$0$,$\alpha$,$+\infty$}
\tkzTabLine{d,-,z,+,}
\tkzTabVar{D+/$+\infty$,-/$f(\alpha)$,+/$0$}
\end{TAB}

\QQ{3} الرسم ( $f(\alpha)\approx-0{,}4$ ، $(C_f)$ يقطع محور الفواصل عند $x\approx0{,}65$ ، ويقترب من محور الفواصل من الأسفل عند $+\infty$ ):
\begin{GR}{-0.5}{4}{-1}{3}{2}
\draw[cDef,thick] plot coordinates {(0.2200,3.2579) (0.2300,3.0331) (0.2400,2.8270) (0.2500,2.6372) (0.2600,2.4621) (0.2700,2.3001) (0.2800,2.1497) (0.2900,2.0099) (0.3000,1.8797) (0.3100,1.7581) (0.3200,1.6443) (0.3300,1.5378) (0.3400,1.4378) (0.3500,1.3438) (0.3600,1.2554) (0.3700,1.1721) (0.3800,1.0936) (0.3900,1.0195) (0.4000,0.9494) (0.4100,0.8831) (0.4200,0.8203) (0.4300,0.7608) (0.4400,0.7044) (0.4500,0.6508) (0.4600,0.6000) (0.4700,0.5517) (0.4800,0.5057) (0.4900,0.4620) (0.5000,0.4204) (0.5100,0.3808) (0.5200,0.3430) (0.5300,0.3071) (0.5400,0.2727) (0.5500,0.2400) (0.5600,0.2088) (0.5700,0.1790) (0.5800,0.1505) (0.5900,0.1234) (0.6000,0.0974) (0.6100,0.0726) (0.6200,0.0489) (0.6300,0.0263) (0.6400,0.0046) (0.6500,-0.0161) (0.6600,-0.0358) (0.6700,-0.0547) (0.6800,-0.0728) (0.6900,-0.0901) (0.7000,-0.1066) (0.7100,-0.1224) (0.7200,-0.1376) (0.7300,-0.1520) (0.7400,-0.1658) (0.7500,-0.1790) (0.7600,-0.1916) (0.7700,-0.2037) (0.7800,-0.2152) (0.7900,-0.2262) (0.8000,-0.2367) (0.8100,-0.2468) (0.8200,-0.2563) (0.8300,-0.2655) (0.8400,-0.2742) (0.8500,-0.2825) (0.8600,-0.2905) (0.8700,-0.2980) (0.8800,-0.3052) (0.8900,-0.3120) (0.9000,-0.3186) (0.9100,-0.3248) (0.9200,-0.3306) (0.9300,-0.3362) (0.9400,-0.3415) (0.9500,-0.3465) (0.9600,-0.3513) (0.9700,-0.3558) (0.9800,-0.3601) (0.9900,-0.3641) (1.0000,-0.3679) (1.0100,-0.3714) (1.0200,-0.3748) (1.0300,-0.3780) (1.0400,-0.3809) (1.0500,-0.3837) (1.0600,-0.3863) (1.0700,-0.3887) (1.0800,-0.3909) (1.0900,-0.3930) (1.1000,-0.3949) (1.1100,-0.3966) (1.1200,-0.3982) (1.1300,-0.3997) (1.1400,-0.4010) (1.1500,-0.4022) (1.1600,-0.4033) (1.1700,-0.4042) (1.1800,-0.4050) (1.1900,-0.4057) (1.2000,-0.4063) (1.2100,-0.4068) (1.2200,-0.4072) (1.2300,-0.4075) (1.2400,-0.4076) (1.2500,-0.4077) (1.2600,-0.4077) (1.2700,-0.4077) (1.2800,-0.4075) (1.2900,-0.4072) (1.3000,-0.4069) (1.3100,-0.4065) (1.3200,-0.4061) (1.3300,-0.4055) (1.3400,-0.4049) (1.3500,-0.4042) (1.3600,-0.4035) (1.3700,-0.4027) (1.3800,-0.4019) (1.3900,-0.4010) (1.4000,-0.4000) (1.4100,-0.3990) (1.4200,-0.3980) (1.4300,-0.3969) (1.4400,-0.3957) (1.4500,-0.3945) (1.4600,-0.3933) (1.4700,-0.3920) (1.4800,-0.3907) (1.4900,-0.3894) (1.5000,-0.3880) (1.5100,-0.3866) (1.5200,-0.3851) (1.5300,-0.3836) (1.5400,-0.3821) (1.5500,-0.3806) (1.5600,-0.3790) (1.5700,-0.3774) (1.5800,-0.3758) (1.5900,-0.3742) (1.6000,-0.3725) (1.6100,-0.3708) (1.6200,-0.3691) (1.6300,-0.3674) (1.6400,-0.3656) (1.6500,-0.3639) (1.6600,-0.3621) (1.6700,-0.3603) (1.6800,-0.3585) (1.6900,-0.3567) (1.7000,-0.3548) (1.7100,-0.3530) (1.7200,-0.3511) (1.7300,-0.3493) (1.7400,-0.3474) (1.7500,-0.3455) (1.7600,-0.3436) (1.7700,-0.3417) (1.7800,-0.3398) (1.7900,-0.3379) (1.8000,-0.3359) (1.8100,-0.3340) (1.8200,-0.3321) (1.8300,-0.3301) (1.8400,-0.3282) (1.8500,-0.3262) (1.8600,-0.3243) (1.8700,-0.3223) (1.8800,-0.3203) (1.8900,-0.3184) (1.9000,-0.3164) (1.9100,-0.3145) (1.9200,-0.3125) (1.9300,-0.3105) (1.9400,-0.3086) (1.9500,-0.3066) (1.9600,-0.3046) (1.9700,-0.3027) (1.9800,-0.3007) (1.9900,-0.2988) (2.0000,-0.2968) (2.0100,-0.2949) (2.0200,-0.2929) (2.0300,-0.2910) (2.0400,-0.2890) (2.0500,-0.2871) (2.0600,-0.2851) (2.0700,-0.2832) (2.0800,-0.2813) (2.0900,-0.2794) (2.1000,-0.2775) (2.1100,-0.2755) (2.1200,-0.2736) (2.1300,-0.2717) (2.1400,-0.2698) (2.1500,-0.2680) (2.1600,-0.2661) (2.1700,-0.2642) (2.1800,-0.2623) (2.1900,-0.2605) (2.2000,-0.2586) (2.2100,-0.2568) (2.2200,-0.2549) (2.2300,-0.2531) (2.2400,-0.2512) (2.2500,-0.2494) (2.2600,-0.2476) (2.2700,-0.2458) (2.2800,-0.2440) (2.2900,-0.2422) (2.3000,-0.2404) (2.3100,-0.2387) (2.3200,-0.2369) (2.3300,-0.2351) (2.3400,-0.2334) (2.3500,-0.2316) (2.3600,-0.2299) (2.3700,-0.2282) (2.3800,-0.2265) (2.3900,-0.2248) (2.4000,-0.2231) (2.4100,-0.2214) (2.4200,-0.2197) (2.4300,-0.2180) (2.4400,-0.2163) (2.4500,-0.2147) (2.4600,-0.2130) (2.4700,-0.2114) (2.4800,-0.2098) (2.4900,-0.2082) (2.5000,-0.2065) (2.5100,-0.2049) (2.5200,-0.2034) (2.5300,-0.2018) (2.5400,-0.2002) (2.5500,-0.1986) (2.5600,-0.1971) (2.5700,-0.1955) (2.5800,-0.1940) (2.5900,-0.1925) (2.6000,-0.1909) (2.6100,-0.1894) (2.6200,-0.1879) (2.6300,-0.1864) (2.6400,-0.1850) (2.6500,-0.1835) (2.6600,-0.1820) (2.6700,-0.1806) (2.6800,-0.1791) (2.6900,-0.1777) (2.7000,-0.1763) (2.7100,-0.1749) (2.7200,-0.1734) (2.7300,-0.1720) (2.7400,-0.1707) (2.7500,-0.1693) (2.7600,-0.1679) (2.7700,-0.1665) (2.7800,-0.1652) (2.7900,-0.1638) (2.8000,-0.1625) (2.8100,-0.1612) (2.8200,-0.1599) (2.8300,-0.1586) (2.8400,-0.1573) (2.8500,-0.1560) (2.8600,-0.1547) (2.8700,-0.1534) (2.8800,-0.1522) (2.8900,-0.1509) (2.9000,-0.1497) (2.9100,-0.1484) (2.9200,-0.1472) (2.9300,-0.1460) (2.9400,-0.1448) (2.9500,-0.1436) (2.9600,-0.1424) (2.9700,-0.1412) (2.9800,-0.1400) (2.9900,-0.1388) (3.0000,-0.1377) (3.0100,-0.1365) (3.0200,-0.1354) (3.0300,-0.1342) (3.0400,-0.1331) (3.0500,-0.1320) (3.0600,-0.1309) (3.0700,-0.1298) (3.0800,-0.1287) (3.0900,-0.1276) (3.1000,-0.1265) (3.1100,-0.1255) (3.1200,-0.1244) (3.1300,-0.1234) (3.1400,-0.1223) (3.1500,-0.1213) (3.1600,-0.1202) (3.1700,-0.1192) (3.1800,-0.1182) (3.1900,-0.1172) (3.2000,-0.1162) (3.2100,-0.1152) (3.2200,-0.1142) (3.2300,-0.1132) (3.2400,-0.1123) (3.2500,-0.1113) (3.2600,-0.1104) (3.2700,-0.1094) (3.2800,-0.1085) (3.2900,-0.1075) (3.3000,-0.1066) (3.3100,-0.1057) (3.3200,-0.1048) (3.3300,-0.1039) (3.3400,-0.1030) (3.3500,-0.1021) (3.3600,-0.1012) (3.3700,-0.1004) (3.3800,-0.0995) (3.3900,-0.0986) (3.4000,-0.0978) (3.4100,-0.0969) (3.4200,-0.0961) (3.4300,-0.0953) (3.4400,-0.0944) (3.4500,-0.0936) (3.4600,-0.0928) (3.4700,-0.0920) (3.4800,-0.0912) (3.4900,-0.0904) (3.5000,-0.0896) (3.5100,-0.0888) (3.5200,-0.0880) (3.5300,-0.0873) (3.5400,-0.0865) (3.5500,-0.0858) (3.5600,-0.0850) (3.5700,-0.0843) (3.5800,-0.0835) (3.5900,-0.0828) (3.6000,-0.0821) (3.6100,-0.0813) (3.6200,-0.0806) (3.6300,-0.0799) (3.6400,-0.0792) (3.6500,-0.0785) (3.6600,-0.0778) (3.6700,-0.0771) (3.6800,-0.0765) (3.6900,-0.0758) (3.7000,-0.0751) (3.7100,-0.0744) (3.7200,-0.0738) (3.7300,-0.0731) (3.7400,-0.0725) (3.7500,-0.0718) (3.7600,-0.0712) (3.7700,-0.0706) (3.7800,-0.0700) (3.7900,-0.0693) (3.8000,-0.0687) (3.8100,-0.0681) (3.8200,-0.0675) (3.8300,-0.0669) (3.8400,-0.0663) (3.8500,-0.0657) (3.8600,-0.0651) (3.8700,-0.0646) (3.8800,-0.0640) (3.8900,-0.0634) (3.9000,-0.0628) (3.9100,-0.0623) (3.9200,-0.0617) (3.9300,-0.0612) (3.9400,-0.0606) (3.9500,-0.0601) (3.9600,-0.0595) (3.9700,-0.0590) (3.9800,-0.0585) (3.9900,-0.0580) (4.0000,-0.0574)};
\fill (0.65,0) circle (1.5pt) node[above right]{$0{,}65$};
\fill (1.2555,-0.4078) circle (1.5pt) node[below]{$\alpha$};
\Ticks{1,2,3}{-1,1,2}
\node[cDef] at (0.6,2.4){$(C_f)$};
\end{GR}
\end{hal}

% ===================== بكالوريا 2023 م2 =====================
\begin{tamrin}{بكالوريا 2023 -- الموضوع الثاني}
$f$ الدالة المعرفة على المجال $]0;+\infty[$ بـ: \ $f(x)=\big((\ln x)^2-3\big)\ln x$.

$(C_f)$ تمثيلها البياني في المستوي منسوب إلى المعلم المتعامد والمتجانس $(O;\vec i,\vec j)$.
\QQ{1}
\begin{sq}
\qa أحسب $\lim\limits_{x\to0}f(x)$ ، ثم فسّر النتيجة هندسيا.
\qb أحسب $\lim\limits_{x\to+\infty}f(x)$.
\end{sq}
\QQ{2}
\begin{sq}
\qa بيّن أنه من أجل كل $x$ من المجال $]0;+\infty[$: \ $f'(x)=\dfrac{3(-1+\ln x)(1+\ln x)}{x}$.
\qb حلّ في المجال $]0;+\infty[$ المتراجحة ذات المجهول $x$: \ $(-1+\ln x)(1+\ln x)>0$.
\qc استنتج أن الدالة $f$ متزايدة تماما على كل من المجالين $\left]0;e^{-1}\right]$ و $[e;+\infty[$ ومتناقصة تماما على المجال $\left[e^{-1};e\right]$ ، ثم شكّل جدول تغيّراتها.
\end{sq}
\QQ{3}
\begin{sq}
\qa أكتب معادلة لـ $(T)$ مماس $(C_f)$ في النقطة ذات الفاصلة $1$.
\qb عيّن فواصل نقط تقاطع المنحنى $(C_f)$ مع حامل محور الفواصل.
\qc أرسم $(T)$ و $(C_f)$ على المجال $\left]0;e^2\right]$.
\end{sq}
\QQ{4} $h$ الدالة المعرفة على $]0;+\infty[$ بـ: $h(x)=\big((\ln x)^2-3\big)|\ln x|$ ، و $(C_h)$ تمثيلها البياني في المعلم السابق.

اشرح كيف يمكن رسم $(C_h)$ انطلاقا من $(C_f)$ ، ثم أرسمه على المجال $\left]0;e^2\right]$.
\end{tamrin}

\begin{hal}
\QQ{1}
\pa لما $x\to0$: $\ln x\to-\infty$ ومنه $(\ln x)^2-3\to+\infty$ ، إذن \bx{\lim_{x\to0}f(x)=(+\infty)(-\infty)=-\infty}.

\textbf{التفسير:} المستقيم ذو المعادلة $x=0$ مقارب عمودي لـ $(C_f)$.

\pb لما $x\to+\infty$: $\ln x\to+\infty$ و $(\ln x)^2-3\to+\infty$ ، إذن \bx{\lim_{x\to+\infty}f(x)=+\infty}.

\QQ{2}
\pa $f(x)=(\ln x)^3-3\ln x$ ، إذن:
\[f'(x)=3(\ln x)^2\cdot\frac1x-\frac3x=\frac{3\big((\ln x)^2-1\big)}{x}=\frac{3(-1+\ln x)(1+\ln x)}{x}.\]
\pb نضع $t=\ln x$: $(t-1)(t+1)>0\iff t<-1$ أو $t>1$ ، أي $\ln x<-1$ أو $\ln x>1$ ، أي $x<e^{-1}$ أو $x>e$:
\[\bx{S=\left]0;e^{-1}\right[\,\cup\,]e;+\infty[}\]
\pc لأن $\frac3x>0$ فإن إشارة $f'(x)$ هي إشارة $(-1+\ln x)(1+\ln x)$: موجبة على $\left]0;e^{-1}\right[$ و $]e;+\infty[$ ، وسالبة على $\left]e^{-1};e\right[$ ، ومنه النتيجة المطلوبة.

$f(e^{-1})=\big(1-3\big)(-1)=2$ ، $f(e)=(1-3)\times1=-2$.
\begin{TAB}
\tkzTabInit[lgt=2,espcl=2.5]{$x$/1,$f'(x)$/1,$f(x)$/2}{$0$,$e^{-1}$,$e$,$+\infty$}
\tkzTabLine{d,+,z,-,z,+,}
\tkzTabVar{D-/$-\infty$,+/$2$,-/$-2$,+/$+\infty$}
\end{TAB}

\QQ{3}
\pa $f(1)=(0-3)\times0=0$ و $f'(1)=\frac{3(-1)(1)}{1}=-3$ ، إذن $y=-3(x-1)+0$: \bx{(T):\ y=-3x+3}.

\pb $f(x)=0\iff\ln x=0$ أو $(\ln x)^2=3\iff\ln x=0$ أو $\ln x=\sqrt3$ أو $\ln x=-\sqrt3$:
\[\bx{x=1\quad;\quad x=e^{\sqrt3}\approx5{,}65\quad;\quad x=e^{-\sqrt3}\approx0{,}18}\]
\pc الرسم ( $f(e^2)=(4-3)\times2=2$ ):
\begin{GR}{-0.5}{7.6}{-3.4}{4}{1}
\draw[cThm,thick] (0,3)--(2.13,-3.4);
\draw[cDef,thick] plot coordinates {(0.1182,-3.3277) (0.1238,-2.8510) (0.1296,-2.4007) (0.1306,-2.3271) (0.1357,-1.9760) (0.1421,-1.5765) (0.1488,-1.2016) (0.1491,-1.1820) (0.1676,-0.3382) (0.1862,0.2922) (0.2047,0.7672) (0.2232,1.1261) (0.2417,1.3967) (0.2602,1.5990) (0.2787,1.7477) (0.2973,1.8540) (0.3158,1.9265) (0.3343,1.9716) (0.3528,1.9947) (0.3713,1.9997) (0.3898,1.9901) (0.4084,1.9684) (0.4269,1.9369) (0.4454,1.8973) (0.4639,1.8511) (0.4824,1.7995) (0.5009,1.7435) (0.5195,1.6839) (0.5380,1.6216) (0.5565,1.5570) (0.5750,1.4907) (0.5935,1.4231) (0.6120,1.3545) (0.6306,1.2854) (0.6491,1.2159) (0.6676,1.1463) (0.6861,1.0767) (0.7046,1.0074) (0.7231,0.9384) (0.7417,0.8699) (0.7602,0.8020) (0.7787,0.7348) (0.7972,0.6683) (0.8157,0.6026) (0.8342,0.5378) (0.8528,0.4738) (0.8713,0.4108) (0.8898,0.3487) (0.9083,0.2876) (0.9268,0.2275) (0.9453,0.1685) (0.9639,0.1104) (0.9824,0.0534) (1.0009,-0.0027) (1.0194,-0.0576) (1.0379,-0.1116) (1.0564,-0.1645) (1.0750,-0.2165) (1.0935,-0.2674) (1.1120,-0.3172) (1.1305,-0.3661) (1.1490,-0.4140) (1.1675,-0.4610) (1.1861,-0.5069) (1.2046,-0.5519) (1.2231,-0.5960) (1.2416,-0.6391) (1.2601,-0.6813) (1.2786,-0.7225) (1.2972,-0.7629) (1.3157,-0.8024) (1.3342,-0.8410) (1.3527,-0.8787) (1.3712,-0.9156) (1.3897,-0.9517) (1.4082,-0.9869) (1.4268,-1.0213) (1.4453,-1.0550) (1.4638,-1.0878) (1.4823,-1.1198) (1.5008,-1.1511) (1.5193,-1.1817) (1.5379,-1.2115) (1.5564,-1.2405) (1.5749,-1.2689) (1.5934,-1.2965) (1.6119,-1.3235) (1.6304,-1.3497) (1.6490,-1.3753) (1.6675,-1.4003) (1.6860,-1.4245) (1.7045,-1.4482) (1.7230,-1.4712) (1.7415,-1.4936) (1.7601,-1.5153) (1.7786,-1.5365) (1.7971,-1.5571) (1.8156,-1.5771) (1.8341,-1.5965) (1.8526,-1.6154) (1.8712,-1.6337) (1.8897,-1.6515) (1.9082,-1.6687) (1.9267,-1.6854) (1.9452,-1.7016) (1.9637,-1.7172) (1.9823,-1.7324) (2.0008,-1.7470) (2.0193,-1.7612) (2.0378,-1.7749) (2.0563,-1.7881) (2.0748,-1.8008) (2.0934,-1.8131) (2.1119,-1.8249) (2.1304,-1.8363) (2.1489,-1.8473) (2.1674,-1.8578) (2.1859,-1.8678) (2.2045,-1.8775) (2.2230,-1.8867) (2.2415,-1.8956) (2.2600,-1.9040) (2.2785,-1.9121) (2.2970,-1.9197) (2.3156,-1.9270) (2.3341,-1.9339) (2.3526,-1.9404) (2.3711,-1.9465) (2.3896,-1.9523) (2.4081,-1.9577) (2.4267,-1.9628) (2.4452,-1.9676) (2.4637,-1.9719) (2.4822,-1.9760) (2.5007,-1.9797) (2.5192,-1.9831) (2.5378,-1.9862) (2.5563,-1.9889) (2.5748,-1.9913) (2.5933,-1.9935) (2.6118,-1.9953) (2.6303,-1.9968) (2.6488,-1.9980) (2.6674,-1.9989) (2.6859,-1.9996) (2.7044,-1.9999) (2.7229,-2.0000) (2.7414,-1.9998) (2.7599,-1.9993) (2.7785,-1.9986) (2.7970,-1.9975) (2.8155,-1.9963) (2.8340,-1.9947) (2.8525,-1.9929) (2.8710,-1.9909) (2.8896,-1.9886) (2.9081,-1.9860) (2.9266,-1.9832) (2.9451,-1.9802) (2.9636,-1.9770) (2.9821,-1.9735) (3.0007,-1.9697) (3.0192,-1.9658) (3.0377,-1.9616) (3.0562,-1.9572) (3.0747,-1.9526) (3.0932,-1.9477) (3.1118,-1.9427) (3.1303,-1.9374) (3.1488,-1.9320) (3.1673,-1.9263) (3.1858,-1.9204) (3.2043,-1.9144) (3.2229,-1.9081) (3.2414,-1.9016) (3.2599,-1.8950) (3.2784,-1.8881) (3.2969,-1.8811) (3.3154,-1.8739) (3.3340,-1.8664) (3.3525,-1.8589) (3.3710,-1.8511) (3.3895,-1.8431) (3.4080,-1.8350) (3.4265,-1.8267) (3.4451,-1.8183) (3.4636,-1.8096) (3.4821,-1.8008) (3.5006,-1.7919) (3.5191,-1.7828) (3.5376,-1.7735) (3.5562,-1.7640) (3.5747,-1.7544) (3.5932,-1.7447) (3.6117,-1.7348) (3.6302,-1.7247) (3.6487,-1.7145) (3.6673,-1.7042) (3.6858,-1.6936) (3.7043,-1.6830) (3.7228,-1.6722) (3.7413,-1.6613) (3.7598,-1.6502) (3.7784,-1.6390) (3.7969,-1.6277) (3.8154,-1.6162) (3.8339,-1.6046) (3.8524,-1.5928) (3.8709,-1.5809) (3.8895,-1.5689) (3.9080,-1.5568) (3.9265,-1.5446) (3.9450,-1.5322) (3.9635,-1.5197) (3.9820,-1.5071) (4.0005,-1.4943) (4.0191,-1.4814) (4.0376,-1.4685) (4.0561,-1.4554) (4.0746,-1.4421) (4.0931,-1.4288) (4.1116,-1.4154) (4.1302,-1.4018) (4.1487,-1.3882) (4.1672,-1.3744) (4.1857,-1.3605) (4.2042,-1.3465) (4.2227,-1.3325) (4.2413,-1.3183) (4.2598,-1.3040) (4.2783,-1.2896) (4.2968,-1.2751) (4.3153,-1.2605) (4.3338,-1.2458) (4.3524,-1.2310) (4.3709,-1.2161) (4.3894,-1.2011) (4.4079,-1.1860) (4.4264,-1.1708) (4.4449,-1.1556) (4.4635,-1.1402) (4.4820,-1.1248) (4.5005,-1.1092) (4.5190,-1.0936) (4.5375,-1.0779) (4.5560,-1.0621) (4.5746,-1.0462) (4.5931,-1.0302) (4.6116,-1.0142) (4.6301,-0.9980) (4.6486,-0.9818) (4.6671,-0.9655) (4.6857,-0.9491) (4.7042,-0.9326) (4.7227,-0.9161) (4.7412,-0.8995) (4.7597,-0.8828) (4.7782,-0.8660) (4.7968,-0.8491) (4.8153,-0.8322) (4.8338,-0.8152) (4.8523,-0.7981) (4.8708,-0.7810) (4.8893,-0.7638) (4.9079,-0.7465) (4.9264,-0.7291) (4.9449,-0.7117) (4.9634,-0.6942) (4.9819,-0.6766) (5.0004,-0.6590) (5.0190,-0.6413) (5.0375,-0.6235) (5.0560,-0.6057) (5.0745,-0.5878) (5.0930,-0.5698) (5.1115,-0.5518) (5.1301,-0.5337) (5.1486,-0.5155) (5.1671,-0.4973) (5.1856,-0.4791) (5.2041,-0.4607) (5.2226,-0.4423) (5.2412,-0.4239) (5.2597,-0.4053) (5.2782,-0.3868) (5.2967,-0.3681) (5.3152,-0.3495) (5.3337,-0.3307) (5.3522,-0.3119) (5.3708,-0.2931) (5.3893,-0.2741) (5.4078,-0.2552) (5.4263,-0.2362) (5.4448,-0.2171) (5.4633,-0.1980) (5.4819,-0.1788) (5.5004,-0.1596) (5.5189,-0.1403) (5.5374,-0.1210) (5.5559,-0.1016) (5.5744,-0.0822) (5.5930,-0.0627) (5.6115,-0.0431) (5.6300,-0.0236) (5.6485,-0.0040) (5.6670,0.0157) (5.6855,0.0354) (5.7041,0.0552) (5.7226,0.0750) (5.7411,0.0949) (5.7596,0.1148) (5.7781,0.1347) (5.7966,0.1547) (5.8152,0.1747) (5.8337,0.1948) (5.8522,0.2149) (5.8707,0.2351) (5.8892,0.2553) (5.9077,0.2755) (5.9263,0.2958) (5.9448,0.3161) (5.9633,0.3365) (5.9818,0.3569) (6.0003,0.3773) (6.0188,0.3978) (6.0374,0.4184) (6.0559,0.4389) (6.0744,0.4595) (6.0929,0.4802) (6.1114,0.5008) (6.1299,0.5215) (6.1485,0.5423) (6.1670,0.5631) (6.1855,0.5839) (6.2040,0.6048) (6.2225,0.6256) (6.2410,0.6466) (6.2596,0.6675) (6.2781,0.6885) (6.2966,0.7096) (6.3151,0.7306) (6.3336,0.7517) (6.3521,0.7728) (6.3707,0.7940) (6.3892,0.8152) (6.4077,0.8364) (6.4262,0.8577) (6.4447,0.8790) (6.4632,0.9003) (6.4818,0.9216) (6.5003,0.9430) (6.5188,0.9644) (6.5373,0.9859) (6.5558,1.0074) (6.5743,1.0289) (6.5928,1.0504) (6.6114,1.0719) (6.6299,1.0935) (6.6484,1.1151) (6.6669,1.1368) (6.6854,1.1585) (6.7039,1.1802) (6.7225,1.2019) (6.7410,1.2236) (6.7595,1.2454) (6.7780,1.2672) (6.7965,1.2890) (6.8150,1.3109) (6.8336,1.3328) (6.8521,1.3547) (6.8706,1.3766) (6.8891,1.3986) (6.9076,1.4206) (6.9261,1.4426) (6.9447,1.4646) (6.9632,1.4867) (6.9817,1.5087) (7.0002,1.5308) (7.0187,1.5530) (7.0372,1.5751) (7.0558,1.5973) (7.0743,1.6195) (7.0928,1.6417) (7.1113,1.6639) (7.1298,1.6862) (7.1483,1.7085) (7.1669,1.7308) (7.1854,1.7531) (7.2039,1.7754) (7.2224,1.7978) (7.2409,1.8202) (7.2594,1.8426) (7.2780,1.8650) (7.2965,1.8875) (7.3150,1.9099) (7.3335,1.9324) (7.3520,1.9549) (7.3705,1.9775) (7.3891,2.0000)};
\fill (0.3679,2) circle (2pt); \fill (2.7183,-2) circle (2pt); \fill (7.389,2) circle (2pt);
\fill (1,0) circle (2pt); \fill (5.652,0) circle (2pt); \fill (0.1769,0) circle (2pt);
\draw[dashed] (7.389,0)--(7.389,2) node[above]{};
\node[below,font=\scriptsize] at (7.389,0){$e^2$};
\node[below,font=\scriptsize] at (2.7183,0){$e$};
\Ticks{1,2,3,4,5,6}{-3,-2,-1,1,2,3}
\node[cDef] at (4.3,-1.3){$(C_f)$};
\node[cThm] at (1.8,-1.8){$(T)$};
\end{GR}

\QQ{4} \textbf{شرح:}
\begin{itemize}
\item إذا كان $x\in[1;+\infty[$: $\ln x\geqslant0$ إذن $|\ln x|=\ln x$ و $h(x)=f(x)$: $(C_h)$ ينطبق على $(C_f)$.
\item إذا كان $x\in\,]0;1[$: $\ln x<0$ إذن $|\ln x|=-\ln x$ و $h(x)=-f(x)$: $(C_h)$ نظير $(C_f)$ بالنسبة لمحور الفواصل.
\end{itemize}
\begin{GR}{-0.5}{7.6}{-3.4}{4}{1}
\draw[cDef!40,dashed] plot coordinates {(0.1182,-3.3277) (0.1238,-2.8510) (0.1296,-2.4007) (0.1306,-2.3271) (0.1357,-1.9760) (0.1421,-1.5765) (0.1488,-1.2016) (0.1491,-1.1820) (0.1676,-0.3382) (0.1862,0.2922) (0.2047,0.7672) (0.2232,1.1261) (0.2417,1.3967) (0.2602,1.5990) (0.2787,1.7477) (0.2973,1.8540) (0.3158,1.9265) (0.3343,1.9716) (0.3528,1.9947) (0.3713,1.9997) (0.3898,1.9901) (0.4084,1.9684) (0.4269,1.9369) (0.4454,1.8973) (0.4639,1.8511) (0.4824,1.7995) (0.5009,1.7435) (0.5195,1.6839) (0.5380,1.6216) (0.5565,1.5570) (0.5750,1.4907) (0.5935,1.4231) (0.6120,1.3545) (0.6306,1.2854) (0.6491,1.2159) (0.6676,1.1463) (0.6861,1.0767) (0.7046,1.0074) (0.7231,0.9384) (0.7417,0.8699) (0.7602,0.8020) (0.7787,0.7348) (0.7972,0.6683) (0.8157,0.6026) (0.8342,0.5378) (0.8528,0.4738) (0.8713,0.4108) (0.8898,0.3487) (0.9083,0.2876) (0.9268,0.2275) (0.9453,0.1685) (0.9639,0.1104) (0.9824,0.0534) (1.0009,-0.0027) (1.0194,-0.0576) (1.0379,-0.1116) (1.0564,-0.1645) (1.0750,-0.2165) (1.0935,-0.2674) (1.1120,-0.3172) (1.1305,-0.3661) (1.1490,-0.4140) (1.1675,-0.4610) (1.1861,-0.5069) (1.2046,-0.5519) (1.2231,-0.5960) (1.2416,-0.6391) (1.2601,-0.6813) (1.2786,-0.7225) (1.2972,-0.7629) (1.3157,-0.8024) (1.3342,-0.8410) (1.3527,-0.8787) (1.3712,-0.9156) (1.3897,-0.9517) (1.4082,-0.9869) (1.4268,-1.0213) (1.4453,-1.0550) (1.4638,-1.0878) (1.4823,-1.1198) (1.5008,-1.1511) (1.5193,-1.1817) (1.5379,-1.2115) (1.5564,-1.2405) (1.5749,-1.2689) (1.5934,-1.2965) (1.6119,-1.3235) (1.6304,-1.3497) (1.6490,-1.3753) (1.6675,-1.4003) (1.6860,-1.4245) (1.7045,-1.4482) (1.7230,-1.4712) (1.7415,-1.4936) (1.7601,-1.5153) (1.7786,-1.5365) (1.7971,-1.5571) (1.8156,-1.5771) (1.8341,-1.5965) (1.8526,-1.6154) (1.8712,-1.6337) (1.8897,-1.6515) (1.9082,-1.6687) (1.9267,-1.6854) (1.9452,-1.7016) (1.9637,-1.7172) (1.9823,-1.7324) (2.0008,-1.7470) (2.0193,-1.7612) (2.0378,-1.7749) (2.0563,-1.7881) (2.0748,-1.8008) (2.0934,-1.8131) (2.1119,-1.8249) (2.1304,-1.8363) (2.1489,-1.8473) (2.1674,-1.8578) (2.1859,-1.8678) (2.2045,-1.8775) (2.2230,-1.8867) (2.2415,-1.8956) (2.2600,-1.9040) (2.2785,-1.9121) (2.2970,-1.9197) (2.3156,-1.9270) (2.3341,-1.9339) (2.3526,-1.9404) (2.3711,-1.9465) (2.3896,-1.9523) (2.4081,-1.9577) (2.4267,-1.9628) (2.4452,-1.9676) (2.4637,-1.9719) (2.4822,-1.9760) (2.5007,-1.9797) (2.5192,-1.9831) (2.5378,-1.9862) (2.5563,-1.9889) (2.5748,-1.9913) (2.5933,-1.9935) (2.6118,-1.9953) (2.6303,-1.9968) (2.6488,-1.9980) (2.6674,-1.9989) (2.6859,-1.9996) (2.7044,-1.9999) (2.7229,-2.0000) (2.7414,-1.9998) (2.7599,-1.9993) (2.7785,-1.9986) (2.7970,-1.9975) (2.8155,-1.9963) (2.8340,-1.9947) (2.8525,-1.9929) (2.8710,-1.9909) (2.8896,-1.9886) (2.9081,-1.9860) (2.9266,-1.9832) (2.9451,-1.9802) (2.9636,-1.9770) (2.9821,-1.9735) (3.0007,-1.9697) (3.0192,-1.9658) (3.0377,-1.9616) (3.0562,-1.9572) (3.0747,-1.9526) (3.0932,-1.9477) (3.1118,-1.9427) (3.1303,-1.9374) (3.1488,-1.9320) (3.1673,-1.9263) (3.1858,-1.9204) (3.2043,-1.9144) (3.2229,-1.9081) (3.2414,-1.9016) (3.2599,-1.8950) (3.2784,-1.8881) (3.2969,-1.8811) (3.3154,-1.8739) (3.3340,-1.8664) (3.3525,-1.8589) (3.3710,-1.8511) (3.3895,-1.8431) (3.4080,-1.8350) (3.4265,-1.8267) (3.4451,-1.8183) (3.4636,-1.8096) (3.4821,-1.8008) (3.5006,-1.7919) (3.5191,-1.7828) (3.5376,-1.7735) (3.5562,-1.7640) (3.5747,-1.7544) (3.5932,-1.7447) (3.6117,-1.7348) (3.6302,-1.7247) (3.6487,-1.7145) (3.6673,-1.7042) (3.6858,-1.6936) (3.7043,-1.6830) (3.7228,-1.6722) (3.7413,-1.6613) (3.7598,-1.6502) (3.7784,-1.6390) (3.7969,-1.6277) (3.8154,-1.6162) (3.8339,-1.6046) (3.8524,-1.5928) (3.8709,-1.5809) (3.8895,-1.5689) (3.9080,-1.5568) (3.9265,-1.5446) (3.9450,-1.5322) (3.9635,-1.5197) (3.9820,-1.5071) (4.0005,-1.4943) (4.0191,-1.4814) (4.0376,-1.4685) (4.0561,-1.4554) (4.0746,-1.4421) (4.0931,-1.4288) (4.1116,-1.4154) (4.1302,-1.4018) (4.1487,-1.3882) (4.1672,-1.3744) (4.1857,-1.3605) (4.2042,-1.3465) (4.2227,-1.3325) (4.2413,-1.3183) (4.2598,-1.3040) (4.2783,-1.2896) (4.2968,-1.2751) (4.3153,-1.2605) (4.3338,-1.2458) (4.3524,-1.2310) (4.3709,-1.2161) (4.3894,-1.2011) (4.4079,-1.1860) (4.4264,-1.1708) (4.4449,-1.1556) (4.4635,-1.1402) (4.4820,-1.1248) (4.5005,-1.1092) (4.5190,-1.0936) (4.5375,-1.0779) (4.5560,-1.0621) (4.5746,-1.0462) (4.5931,-1.0302) (4.6116,-1.0142) (4.6301,-0.9980) (4.6486,-0.9818) (4.6671,-0.9655) (4.6857,-0.9491) (4.7042,-0.9326) (4.7227,-0.9161) (4.7412,-0.8995) (4.7597,-0.8828) (4.7782,-0.8660) (4.7968,-0.8491) (4.8153,-0.8322) (4.8338,-0.8152) (4.8523,-0.7981) (4.8708,-0.7810) (4.8893,-0.7638) (4.9079,-0.7465) (4.9264,-0.7291) (4.9449,-0.7117) (4.9634,-0.6942) (4.9819,-0.6766) (5.0004,-0.6590) (5.0190,-0.6413) (5.0375,-0.6235) (5.0560,-0.6057) (5.0745,-0.5878) (5.0930,-0.5698) (5.1115,-0.5518) (5.1301,-0.5337) (5.1486,-0.5155) (5.1671,-0.4973) (5.1856,-0.4791) (5.2041,-0.4607) (5.2226,-0.4423) (5.2412,-0.4239) (5.2597,-0.4053) (5.2782,-0.3868) (5.2967,-0.3681) (5.3152,-0.3495) (5.3337,-0.3307) (5.3522,-0.3119) (5.3708,-0.2931) (5.3893,-0.2741) (5.4078,-0.2552) (5.4263,-0.2362) (5.4448,-0.2171) (5.4633,-0.1980) (5.4819,-0.1788) (5.5004,-0.1596) (5.5189,-0.1403) (5.5374,-0.1210) (5.5559,-0.1016) (5.5744,-0.0822) (5.5930,-0.0627) (5.6115,-0.0431) (5.6300,-0.0236) (5.6485,-0.0040) (5.6670,0.0157) (5.6855,0.0354) (5.7041,0.0552) (5.7226,0.0750) (5.7411,0.0949) (5.7596,0.1148) (5.7781,0.1347) (5.7966,0.1547) (5.8152,0.1747) (5.8337,0.1948) (5.8522,0.2149) (5.8707,0.2351) (5.8892,0.2553) (5.9077,0.2755) (5.9263,0.2958) (5.9448,0.3161) (5.9633,0.3365) (5.9818,0.3569) (6.0003,0.3773) (6.0188,0.3978) (6.0374,0.4184) (6.0559,0.4389) (6.0744,0.4595) (6.0929,0.4802) (6.1114,0.5008) (6.1299,0.5215) (6.1485,0.5423) (6.1670,0.5631) (6.1855,0.5839) (6.2040,0.6048) (6.2225,0.6256) (6.2410,0.6466) (6.2596,0.6675) (6.2781,0.6885) (6.2966,0.7096) (6.3151,0.7306) (6.3336,0.7517) (6.3521,0.7728) (6.3707,0.7940) (6.3892,0.8152) (6.4077,0.8364) (6.4262,0.8577) (6.4447,0.8790) (6.4632,0.9003) (6.4818,0.9216) (6.5003,0.9430) (6.5188,0.9644) (6.5373,0.9859) (6.5558,1.0074) (6.5743,1.0289) (6.5928,1.0504) (6.6114,1.0719) (6.6299,1.0935) (6.6484,1.1151) (6.6669,1.1368) (6.6854,1.1585) (6.7039,1.1802) (6.7225,1.2019) (6.7410,1.2236) (6.7595,1.2454) (6.7780,1.2672) (6.7965,1.2890) (6.8150,1.3109) (6.8336,1.3328) (6.8521,1.3547) (6.8706,1.3766) (6.8891,1.3986) (6.9076,1.4206) (6.9261,1.4426) (6.9447,1.4646) (6.9632,1.4867) (6.9817,1.5087) (7.0002,1.5308) (7.0187,1.5530) (7.0372,1.5751) (7.0558,1.5973) (7.0743,1.6195) (7.0928,1.6417) (7.1113,1.6639) (7.1298,1.6862) (7.1483,1.7085) (7.1669,1.7308) (7.1854,1.7531) (7.2039,1.7754) (7.2224,1.7978) (7.2409,1.8202) (7.2594,1.8426) (7.2780,1.8650) (7.2965,1.8875) (7.3150,1.9099) (7.3335,1.9324) (7.3520,1.9549) (7.3705,1.9775) (7.3891,2.0000)};
\draw[cThm,very thick] plot coordinates {(0.1079,4.3618) (0.1121,3.9150) (0.1129,3.8311) (0.1182,3.3277) (0.1238,2.8510) (0.1296,2.4007) (0.1306,2.3271) (0.1357,1.9760) (0.1421,1.5765) (0.1488,1.2016) (0.1491,1.1820) (0.1676,0.3382) (0.1862,-0.2922) (0.2047,-0.7672) (0.2232,-1.1261) (0.2417,-1.3967) (0.2602,-1.5990) (0.2787,-1.7477) (0.2973,-1.8540) (0.3158,-1.9265) (0.3343,-1.9716) (0.3528,-1.9947) (0.3713,-1.9997) (0.3898,-1.9901) (0.4084,-1.9684) (0.4269,-1.9369) (0.4454,-1.8973) (0.4639,-1.8511) (0.4824,-1.7995) (0.5009,-1.7435) (0.5195,-1.6839) (0.5380,-1.6216) (0.5565,-1.5570) (0.5750,-1.4907) (0.5935,-1.4231) (0.6120,-1.3545) (0.6306,-1.2854) (0.6491,-1.2159) (0.6676,-1.1463) (0.6861,-1.0767) (0.7046,-1.0074) (0.7231,-0.9384) (0.7417,-0.8699) (0.7602,-0.8020) (0.7787,-0.7348) (0.7972,-0.6683) (0.8157,-0.6026) (0.8342,-0.5378) (0.8528,-0.4738) (0.8713,-0.4108) (0.8898,-0.3487) (0.9083,-0.2876) (0.9268,-0.2275) (0.9453,-0.1685) (0.9639,-0.1104) (0.9824,-0.0534) (1.0009,-0.0027) (1.0194,-0.0576) (1.0379,-0.1116) (1.0564,-0.1645) (1.0750,-0.2165) (1.0935,-0.2674) (1.1120,-0.3172) (1.1305,-0.3661) (1.1490,-0.4140) (1.1675,-0.4610) (1.1861,-0.5069) (1.2046,-0.5519) (1.2231,-0.5960) (1.2416,-0.6391) (1.2601,-0.6813) (1.2786,-0.7225) (1.2972,-0.7629) (1.3157,-0.8024) (1.3342,-0.8410) (1.3527,-0.8787) (1.3712,-0.9156) (1.3897,-0.9517) (1.4082,-0.9869) (1.4268,-1.0213) (1.4453,-1.0550) (1.4638,-1.0878) (1.4823,-1.1198) (1.5008,-1.1511) (1.5193,-1.1817) (1.5379,-1.2115) (1.5564,-1.2405) (1.5749,-1.2689) (1.5934,-1.2965) (1.6119,-1.3235) (1.6304,-1.3497) (1.6490,-1.3753) (1.6675,-1.4003) (1.6860,-1.4245) (1.7045,-1.4482) (1.7230,-1.4712) (1.7415,-1.4936) (1.7601,-1.5153) (1.7786,-1.5365) (1.7971,-1.5571) (1.8156,-1.5771) (1.8341,-1.5965) (1.8526,-1.6154) (1.8712,-1.6337) (1.8897,-1.6515) (1.9082,-1.6687) (1.9267,-1.6854) (1.9452,-1.7016) (1.9637,-1.7172) (1.9823,-1.7324) (2.0008,-1.7470) (2.0193,-1.7612) (2.0378,-1.7749) (2.0563,-1.7881) (2.0748,-1.8008) (2.0934,-1.8131) (2.1119,-1.8249) (2.1304,-1.8363) (2.1489,-1.8473) (2.1674,-1.8578) (2.1859,-1.8678) (2.2045,-1.8775) (2.2230,-1.8867) (2.2415,-1.8956) (2.2600,-1.9040) (2.2785,-1.9121) (2.2970,-1.9197) (2.3156,-1.9270) (2.3341,-1.9339) (2.3526,-1.9404) (2.3711,-1.9465) (2.3896,-1.9523) (2.4081,-1.9577) (2.4267,-1.9628) (2.4452,-1.9676) (2.4637,-1.9719) (2.4822,-1.9760) (2.5007,-1.9797) (2.5192,-1.9831) (2.5378,-1.9862) (2.5563,-1.9889) (2.5748,-1.9913) (2.5933,-1.9935) (2.6118,-1.9953) (2.6303,-1.9968) (2.6488,-1.9980) (2.6674,-1.9989) (2.6859,-1.9996) (2.7044,-1.9999) (2.7229,-2.0000) (2.7414,-1.9998) (2.7599,-1.9993) (2.7785,-1.9986) (2.7970,-1.9975) (2.8155,-1.9963) (2.8340,-1.9947) (2.8525,-1.9929) (2.8710,-1.9909) (2.8896,-1.9886) (2.9081,-1.9860) (2.9266,-1.9832) (2.9451,-1.9802) (2.9636,-1.9770) (2.9821,-1.9735) (3.0007,-1.9697) (3.0192,-1.9658) (3.0377,-1.9616) (3.0562,-1.9572) (3.0747,-1.9526) (3.0932,-1.9477) (3.1118,-1.9427) (3.1303,-1.9374) (3.1488,-1.9320) (3.1673,-1.9263) (3.1858,-1.9204) (3.2043,-1.9144) (3.2229,-1.9081) (3.2414,-1.9016) (3.2599,-1.8950) (3.2784,-1.8881) (3.2969,-1.8811) (3.3154,-1.8739) (3.3340,-1.8664) (3.3525,-1.8589) (3.3710,-1.8511) (3.3895,-1.8431) (3.4080,-1.8350) (3.4265,-1.8267) (3.4451,-1.8183) (3.4636,-1.8096) (3.4821,-1.8008) (3.5006,-1.7919) (3.5191,-1.7828) (3.5376,-1.7735) (3.5562,-1.7640) (3.5747,-1.7544) (3.5932,-1.7447) (3.6117,-1.7348) (3.6302,-1.7247) (3.6487,-1.7145) (3.6673,-1.7042) (3.6858,-1.6936) (3.7043,-1.6830) (3.7228,-1.6722) (3.7413,-1.6613) (3.7598,-1.6502) (3.7784,-1.6390) (3.7969,-1.6277) (3.8154,-1.6162) (3.8339,-1.6046) (3.8524,-1.5928) (3.8709,-1.5809) (3.8895,-1.5689) (3.9080,-1.5568) (3.9265,-1.5446) (3.9450,-1.5322) (3.9635,-1.5197) (3.9820,-1.5071) (4.0005,-1.4943) (4.0191,-1.4814) (4.0376,-1.4685) (4.0561,-1.4554) (4.0746,-1.4421) (4.0931,-1.4288) (4.1116,-1.4154) (4.1302,-1.4018) (4.1487,-1.3882) (4.1672,-1.3744) (4.1857,-1.3605) (4.2042,-1.3465) (4.2227,-1.3325) (4.2413,-1.3183) (4.2598,-1.3040) (4.2783,-1.2896) (4.2968,-1.2751) (4.3153,-1.2605) (4.3338,-1.2458) (4.3524,-1.2310) (4.3709,-1.2161) (4.3894,-1.2011) (4.4079,-1.1860) (4.4264,-1.1708) (4.4449,-1.1556) (4.4635,-1.1402) (4.4820,-1.1248) (4.5005,-1.1092) (4.5190,-1.0936) (4.5375,-1.0779) (4.5560,-1.0621) (4.5746,-1.0462) (4.5931,-1.0302) (4.6116,-1.0142) (4.6301,-0.9980) (4.6486,-0.9818) (4.6671,-0.9655) (4.6857,-0.9491) (4.7042,-0.9326) (4.7227,-0.9161) (4.7412,-0.8995) (4.7597,-0.8828) (4.7782,-0.8660) (4.7968,-0.8491) (4.8153,-0.8322) (4.8338,-0.8152) (4.8523,-0.7981) (4.8708,-0.7810) (4.8893,-0.7638) (4.9079,-0.7465) (4.9264,-0.7291) (4.9449,-0.7117) (4.9634,-0.6942) (4.9819,-0.6766) (5.0004,-0.6590) (5.0190,-0.6413) (5.0375,-0.6235) (5.0560,-0.6057) (5.0745,-0.5878) (5.0930,-0.5698) (5.1115,-0.5518) (5.1301,-0.5337) (5.1486,-0.5155) (5.1671,-0.4973) (5.1856,-0.4791) (5.2041,-0.4607) (5.2226,-0.4423) (5.2412,-0.4239) (5.2597,-0.4053) (5.2782,-0.3868) (5.2967,-0.3681) (5.3152,-0.3495) (5.3337,-0.3307) (5.3522,-0.3119) (5.3708,-0.2931) (5.3893,-0.2741) (5.4078,-0.2552) (5.4263,-0.2362) (5.4448,-0.2171) (5.4633,-0.1980) (5.4819,-0.1788) (5.5004,-0.1596) (5.5189,-0.1403) (5.5374,-0.1210) (5.5559,-0.1016) (5.5744,-0.0822) (5.5930,-0.0627) (5.6115,-0.0431) (5.6300,-0.0236) (5.6485,-0.0040) (5.6670,0.0157) (5.6855,0.0354) (5.7041,0.0552) (5.7226,0.0750) (5.7411,0.0949) (5.7596,0.1148) (5.7781,0.1347) (5.7966,0.1547) (5.8152,0.1747) (5.8337,0.1948) (5.8522,0.2149) (5.8707,0.2351) (5.8892,0.2553) (5.9077,0.2755) (5.9263,0.2958) (5.9448,0.3161) (5.9633,0.3365) (5.9818,0.3569) (6.0003,0.3773) (6.0188,0.3978) (6.0374,0.4184) (6.0559,0.4389) (6.0744,0.4595) (6.0929,0.4802) (6.1114,0.5008) (6.1299,0.5215) (6.1485,0.5423) (6.1670,0.5631) (6.1855,0.5839) (6.2040,0.6048) (6.2225,0.6256) (6.2410,0.6466) (6.2596,0.6675) (6.2781,0.6885) (6.2966,0.7096) (6.3151,0.7306) (6.3336,0.7517) (6.3521,0.7728) (6.3707,0.7940) (6.3892,0.8152) (6.4077,0.8364) (6.4262,0.8577) (6.4447,0.8790) (6.4632,0.9003) (6.4818,0.9216) (6.5003,0.9430) (6.5188,0.9644) (6.5373,0.9859) (6.5558,1.0074) (6.5743,1.0289) (6.5928,1.0504) (6.6114,1.0719) (6.6299,1.0935) (6.6484,1.1151) (6.6669,1.1368) (6.6854,1.1585) (6.7039,1.1802) (6.7225,1.2019) (6.7410,1.2236) (6.7595,1.2454) (6.7780,1.2672) (6.7965,1.2890) (6.8150,1.3109) (6.8336,1.3328) (6.8521,1.3547) (6.8706,1.3766) (6.8891,1.3986) (6.9076,1.4206) (6.9261,1.4426) (6.9447,1.4646) (6.9632,1.4867) (6.9817,1.5087) (7.0002,1.5308) (7.0187,1.5530) (7.0372,1.5751) (7.0558,1.5973) (7.0743,1.6195) (7.0928,1.6417) (7.1113,1.6639) (7.1298,1.6862) (7.1483,1.7085) (7.1669,1.7308) (7.1854,1.7531) (7.2039,1.7754) (7.2224,1.7978) (7.2409,1.8202) (7.2594,1.8426) (7.2780,1.8650) (7.2965,1.8875) (7.3150,1.9099) (7.3335,1.9324) (7.3520,1.9549) (7.3705,1.9775) (7.3891,2.0000)};
\Ticks{1,2,3,4,5,6}{-3,-2,-1,1,2,3}
\node[cThm] at (5.2,1.6){$(C_h)$};
\node[cDef!60] at (0.8,3.6){$(C_f)$};
\end{GR}
\end{hal}

% ===================== بكالوريا 2024 م2 =====================
\begin{tamrin}{بكالوريا 2024 -- الموضوع الثاني}
\textbf{I)} $g$ الدالة المعرّفة على $]0;+\infty[$ كما يلي: \ $g(x)=\dfrac12x^3+\dfrac12-\ln x$ ، و $(C_g)$ تمثيلها البياني كما في الشكل.
\begin{GR}{-0.5}{2}{-0.4}{3.2}{1.4}
\draw[cDef,thick] plot coordinates {(0.0498,3.4998) (0.0499,3.4978) (0.0549,3.4026) (0.0599,3.3156) (0.0649,3.2356) (0.0698,3.1616) (0.0748,3.0927) (0.0798,3.0282) (0.0848,2.9676) (0.0898,2.9105) (0.0948,2.8565) (0.0998,2.8053) (0.1048,2.7566) (0.1097,2.7102) (0.1147,2.6659) (0.1197,2.6234) (0.1247,2.5827) (0.1297,2.5436) (0.1347,2.5060) (0.1397,2.4698) (0.1447,2.4349) (0.1496,2.4011) (0.1546,2.3685) (0.1596,2.3370) (0.1646,2.3064) (0.1696,2.2768) (0.1746,2.2480) (0.1796,2.2201) (0.1846,2.1929) (0.1895,2.1665) (0.1945,2.1408) (0.1995,2.1158) (0.2045,2.0914) (0.2095,2.0676) (0.2145,2.0444) (0.2195,2.0218) (0.2245,1.9997) (0.2294,1.9781) (0.2344,1.9570) (0.2394,1.9364) (0.2444,1.9162) (0.2494,1.8965) (0.2544,1.8771) (0.2594,1.8582) (0.2644,1.8397) (0.2693,1.8215) (0.2743,1.8037) (0.2793,1.7863) (0.2843,1.7692) (0.2893,1.7524) (0.2943,1.7359) (0.2993,1.7198) (0.3043,1.7040) (0.3092,1.6884) (0.3142,1.6731) (0.3192,1.6581) (0.3242,1.6434) (0.3292,1.6289) (0.3342,1.6147) (0.3392,1.6008) (0.3442,1.5870) (0.3491,1.5735) (0.3541,1.5603) (0.3591,1.5472) (0.3641,1.5344) (0.3691,1.5218) (0.3741,1.5094) (0.3791,1.4973) (0.3841,1.4853) (0.3890,1.4735) (0.3940,1.4619) (0.3990,1.4505) (0.4040,1.4393) (0.4090,1.4283) (0.4140,1.4174) (0.4190,1.4067) (0.4240,1.3962) (0.4289,1.3859) (0.4339,1.3757) (0.4389,1.3657) (0.4439,1.3559) (0.4489,1.3462) (0.4539,1.3367) (0.4589,1.3273) (0.4639,1.3181) (0.4688,1.3090) (0.4738,1.3001) (0.4788,1.2913) (0.4838,1.2827) (0.4888,1.2742) (0.4938,1.2659) (0.4988,1.2576) (0.5038,1.2496) (0.5087,1.2416) (0.5137,1.2338) (0.5187,1.2262) (0.5237,1.2186) (0.5287,1.2112) (0.5337,1.2040) (0.5387,1.1968) (0.5437,1.1898) (0.5486,1.1829) (0.5536,1.1761) (0.5586,1.1694) (0.5636,1.1629) (0.5686,1.1565) (0.5736,1.1502) (0.5786,1.1440) (0.5836,1.1380) (0.5885,1.1320) (0.5935,1.1262) (0.5985,1.1205) (0.6035,1.1149) (0.6085,1.1094) (0.6135,1.1040) (0.6185,1.0988) (0.6235,1.0936) (0.6284,1.0886) (0.6334,1.0837) (0.6384,1.0789) (0.6434,1.0742) (0.6484,1.0696) (0.6534,1.0651) (0.6584,1.0607) (0.6634,1.0564) (0.6683,1.0522) (0.6733,1.0482) (0.6783,1.0442) (0.6833,1.0403) (0.6883,1.0366) (0.6933,1.0329) (0.6983,1.0294) (0.7033,1.0259) (0.7082,1.0226) (0.7132,1.0194) (0.7182,1.0162) (0.7232,1.0132) (0.7282,1.0103) (0.7332,1.0074) (0.7382,1.0047) (0.7432,1.0021) (0.7481,0.9995) (0.7531,0.9971) (0.7581,0.9948) (0.7631,0.9925) (0.7681,0.9904) (0.7731,0.9884) (0.7781,0.9865) (0.7831,0.9846) (0.7880,0.9829) (0.7930,0.9813) (0.7980,0.9797) (0.8030,0.9783) (0.8080,0.9770) (0.8130,0.9757) (0.8180,0.9746) (0.8230,0.9735) (0.8279,0.9726) (0.8329,0.9717) (0.8379,0.9710) (0.8429,0.9703) (0.8479,0.9698) (0.8529,0.9693) (0.8579,0.9690) (0.8629,0.9687) (0.8678,0.9686) (0.8728,0.9685) (0.8778,0.9685) (0.8828,0.9687) (0.8878,0.9689) (0.8928,0.9692) (0.8978,0.9696) (0.9028,0.9702) (0.9077,0.9708) (0.9127,0.9715) (0.9177,0.9723) (0.9227,0.9732) (0.9277,0.9742) (0.9327,0.9754) (0.9377,0.9766) (0.9427,0.9779) (0.9476,0.9793) (0.9526,0.9808) (0.9576,0.9824) (0.9626,0.9841) (0.9676,0.9859) (0.9726,0.9878) (0.9776,0.9898) (0.9826,0.9919) (0.9875,0.9941) (0.9925,0.9964) (0.9975,0.9988) (1.0025,1.0013) (1.0075,1.0039) (1.0125,1.0066) (1.0175,1.0093) (1.0225,1.0122) (1.0274,1.0152) (1.0324,1.0183) (1.0374,1.0215) (1.0424,1.0248) (1.0474,1.0282) (1.0524,1.0317) (1.0574,1.0353) (1.0624,1.0390) (1.0673,1.0428) (1.0723,1.0467) (1.0773,1.0507) (1.0823,1.0548) (1.0873,1.0590) (1.0923,1.0633) (1.0973,1.0677) (1.1023,1.0722) (1.1072,1.0769) (1.1122,1.0816) (1.1172,1.0864) (1.1222,1.0913) (1.1272,1.0964) (1.1322,1.1015) (1.1372,1.1067) (1.1422,1.1121) (1.1471,1.1175) (1.1521,1.1231) (1.1571,1.1287) (1.1621,1.1345) (1.1671,1.1403) (1.1721,1.1463) (1.1771,1.1524) (1.1821,1.1586) (1.1870,1.1648) (1.1920,1.1712) (1.1970,1.1777) (1.2020,1.1843) (1.2070,1.1911) (1.2120,1.1979) (1.2170,1.2048) (1.2220,1.2118) (1.2269,1.2190) (1.2319,1.2262) (1.2369,1.2336) (1.2419,1.2411) (1.2469,1.2486) (1.2519,1.2563) (1.2569,1.2641) (1.2619,1.2720) (1.2668,1.2800) (1.2718,1.2882) (1.2768,1.2964) (1.2818,1.3047) (1.2868,1.3132) (1.2918,1.3218) (1.2968,1.3305) (1.3018,1.3392) (1.3067,1.3481) (1.3117,1.3572) (1.3167,1.3663) (1.3217,1.3755) (1.3267,1.3849) (1.3317,1.3943) (1.3367,1.4039) (1.3417,1.4136) (1.3466,1.4234) (1.3516,1.4333) (1.3566,1.4434) (1.3616,1.4535) (1.3666,1.4638) (1.3716,1.4742) (1.3766,1.4847) (1.3816,1.4953) (1.3865,1.5060) (1.3915,1.5168) (1.3965,1.5278) (1.4015,1.5389) (1.4065,1.5501) (1.4115,1.5614) (1.4165,1.5728) (1.4215,1.5844) (1.4264,1.5960) (1.4314,1.6078) (1.4364,1.6197) (1.4414,1.6317) (1.4464,1.6439) (1.4514,1.6562) (1.4564,1.6685) (1.4614,1.6810) (1.4663,1.6937) (1.4713,1.7064) (1.4763,1.7193) (1.4813,1.7323) (1.4863,1.7454) (1.4913,1.7586) (1.4963,1.7720) (1.5013,1.7854) (1.5062,1.7990) (1.5112,1.8128) (1.5162,1.8266) (1.5212,1.8406) (1.5262,1.8547) (1.5312,1.8689) (1.5362,1.8832) (1.5412,1.8977) (1.5461,1.9123) (1.5511,1.9270) (1.5561,1.9419) (1.5611,1.9568) (1.5661,1.9719) (1.5711,1.9872) (1.5761,2.0025) (1.5811,2.0180) (1.5860,2.0336) (1.5910,2.0494) (1.5960,2.0652) (1.6010,2.0812) (1.6060,2.0973) (1.6110,2.1136) (1.6160,2.1300) (1.6210,2.1465) (1.6259,2.1631) (1.6309,2.1799) (1.6359,2.1968) (1.6409,2.2139) (1.6459,2.2310) (1.6509,2.2483) (1.6559,2.2658) (1.6609,2.2833) (1.6658,2.3010) (1.6708,2.3189) (1.6758,2.3368) (1.6808,2.3549) (1.6858,2.3732) (1.6908,2.3915) (1.6958,2.4101) (1.7008,2.4287) (1.7057,2.4475) (1.7107,2.4664) (1.7157,2.4854) (1.7207,2.5046) (1.7257,2.5239) (1.7307,2.5434) (1.7357,2.5630) (1.7407,2.5827) (1.7456,2.6026) (1.7506,2.6226) (1.7556,2.6427) (1.7606,2.6630) (1.7656,2.6835) (1.7706,2.7040) (1.7756,2.7247) (1.7806,2.7456) (1.7855,2.7666) (1.7905,2.7877) (1.7955,2.8090) (1.8005,2.8304) (1.8055,2.8519) (1.8105,2.8736) (1.8155,2.8955) (1.8205,2.9174) (1.8254,2.9396) (1.8304,2.9618) (1.8354,2.9842) (1.8404,3.0068) (1.8454,3.0295) (1.8504,3.0524) (1.8554,3.0753) (1.8604,3.0985) (1.8653,3.1218) (1.8703,3.1452) (1.8753,3.1688) (1.8803,3.1925) (1.8853,3.2164) (1.8903,3.2404) (1.8953,3.2646) (1.9003,3.2889) (1.9052,3.3133) (1.9102,3.3379) (1.9152,3.3627) (1.9202,3.3876) (1.9252,3.4127) (1.9302,3.4379) (1.9352,3.4633) (1.9402,3.4888)};
\Ticks{0.5,1,1.5}{1,2,3}
\node[cDef] at (1.75,2.6){$(C_g)$};
\end{GR}
\par بقراءة بيانية، عيّن إشارة $g(x)$.

\textbf{II)} $f$ الدالة المعرّفة على $]0;+\infty[$ بـ: \ $f(x)=-x-\dfrac{\ln x}{x^2}$.

$(C_f)$ تمثيلها البياني في المستوي المنسوب إلى المعلم المتعامد والمتجانس $(O;\vec i,\vec j)$. (وحدة الطول $2\text{cm}$).
\QQ{1} احسب $\lim\limits_{x\to0}f(x)$ و $\lim\limits_{x\to+\infty}f(x)$.
\QQ{2}
\begin{sq}
\qa بيّن أنّه من أجل كل $x$ من $]0;+\infty[$: \ $f'(x)=\dfrac{-2g(x)}{x^3}$.
\qb استنتج اتجاه تغيّر الدالة $f$ ثم شكّل جدول تغيّراتها.
\qc أثبت أن المعادلة $f(x)=0$ تقبل حلا وحيدا $\alpha$ حيث $0{,}70<\alpha<0{,}71$.
\end{sq}
\QQ{3}
\begin{sq}
\qa بيّن أن المنحنى $(C_f)$ يقبل مستقيما مقاربا مائلا $(\Delta)$ ، يُطلب تعيين معادلة له.
\qb أدرس الوضع النسبي للمنحنى $(C_f)$ والمستقيم $(\Delta)$.
\end{sq}
\QQ{4} بيّن أن المنحنى $(C_f)$ يقبل مماسا $(T)$ معامل توجيهه $(-1)$ ، يُطلب تعيين معادلة له.
\QQ{5}
\begin{sq}
\qa أرسم كلا من $(\Delta)$ ، $(T)$ و $(C_f)$.
\qb $m$ وسيط حقيقي، عيّن بيانيا قيم $m$ التي من أجلها تقبل المعادلة $\dfrac{\ln x}{x^2}=m$ حلّين مختلفين.
\end{sq}
\end{tamrin}

\begin{molahada}
في نص المجلة كُتبت الدالة $f(x)=x-\dfrac{\ln x}{x^2}$ ، وهذا لا يتوافق مع المعطى $f'(x)=\dfrac{-2g(x)}{x^3}$ ولا مع وجود مماس معامل توجيهه $(-1)$ ولا مع المجال $]0{,}70;0{,}71[$ للحل $\alpha$. العبارة الصحيحة المتوافقة مع كل الأسئلة هي $f(x)=-x-\dfrac{\ln x}{x^2}$ وقد اعتمدناها. كما أضفنا $\lim\limits_{x\to0}f(x)$ في السؤال 1 (سقطت من النص).
\end{molahada}

\begin{hal}
\textbf{الجزء I:} المنحنى $(C_g)$ يقع كليا فوق محور الفواصل، إذن: \bxt{$g(x)>0$ من أجل كل $x\in ]0;+\infty[$}.

\textbf{الجزء II}

\QQ{1} لما $x\to0$: $\frac{\ln x}{x^2}\to-\infty$ (لأن $\ln x\to-\infty$ و $\frac1{x^2}\to+\infty$) ، إذن $-\frac{\ln x}{x^2}\to+\infty$: \bx{\lim_{x\to0}f(x)=+\infty}.

لما $x\to+\infty$: $\frac{\ln x}{x^2}=\frac{\ln x}{x}\cdot\frac1x\to0$ و $-x\to-\infty$: \bx{\lim_{x\to+\infty}f(x)=-\infty}.

\QQ{2}
\pa $\Big(\dfrac{\ln x}{x^2}\Big)'=\dfrac{1-2\ln x}{x^3}$ ، إذن:
\[f'(x)=-1-\frac{1-2\ln x}{x^3}=\frac{-x^3-1+2\ln x}{x^3}=\frac{-2\big(\frac12x^3+\frac12-\ln x\big)}{x^3}=\frac{-2g(x)}{x^3}.\]
\pb $g(x)>0$ و $x^3>0$ ، إذن $f'(x)<0$: \bxt{$f$ متناقصة تماما على $]0;+\infty[$}.
\begin{TAB}
\tkzTabInit[lgt=2,espcl=4]{$x$/1,$f'(x)$/1,$f(x)$/2}{$0$,$+\infty$}
\tkzTabLine{d,-,}
\tkzTabVar{D+/$+\infty$,-/$-\infty$}
\end{TAB}
\pc $f$ مستمرة ومتناقصة تماما على $[0{,}70;0{,}71]$ ، و $f(0{,}70)\approx0{,}028>0$ ، $f(0{,}71)\approx-0{,}031<0$.
حسب مبرهنة القيم المتوسطة المعادلة $f(x)=0$ تقبل حلا وحيدا $\alpha$ (وحيد على $]0;+\infty[$ لأن $f$ رتيبة تماما): \bx{0{,}70<\alpha<0{,}71}.

\QQ{3}
\pa $f(x)-(-x)=-\dfrac{\ln x}{x^2}$ و $\lim\limits_{x\to+\infty}\Big(-\dfrac{\ln x}{x^2}\Big)=0$ ، إذن \bx{(\Delta):\ y=-x} مستقيم مقارب مائل لـ $(C_f)$ بجوار $+\infty$.

\pb إشارة $-\frac{\ln x}{x^2}$ هي إشارة $-\ln x$:
\begin{itemize}
\item على $]0;1[$: $(C_f)$ فوق $(\Delta)$ ؛ \quad على $]1;+\infty[$: $(C_f)$ تحت $(\Delta)$ ؛
\item $(C_f)$ يقطع $(\Delta)$ في النقطة $(1;-1)$.
\end{itemize}

\QQ{4} $f'(x)=-1\iff\dfrac{-x^3-1+2\ln x}{x^3}=-1\iff-x^3-1+2\ln x=-x^3\iff\ln x=\dfrac12\iff x=\sqrt e$.

$f(\sqrt e)=-\sqrt e-\dfrac{1}{2e}$ ، ومنه $y=-(x-\sqrt e)-\sqrt e-\dfrac1{2e}$: \bx{(T):\ y=-x-\frac{1}{2e}}.

\QQ{5}
\pa الرسم ( $\frac{1}{2e}\approx0{,}18$ ):
\begin{GR}{-0.5}{3.2}{-3.6}{2.6}{1.6}
\draw[cRem,thick] (0,0)--(3.2,-3.2);
\draw[cThm,thick] (0,-0.184)--(3.2,-3.384);
\draw[cDef,thick] plot coordinates {(0.4737,2.8559) (0.4817,2.6661) (0.4897,2.4876) (0.4977,2.3193) (0.5057,2.1606) (0.5137,2.0108) (0.5217,1.8693) (0.5297,1.7355) (0.5377,1.6088) (0.5457,1.4888) (0.5537,1.3750) (0.5617,1.2670) (0.5696,1.1645) (0.5776,1.0671) (0.5856,0.9744) (0.5936,0.8862) (0.6016,0.8022) (0.6096,0.7221) (0.6176,0.6457) (0.6256,0.5727) (0.6336,0.5030) (0.6416,0.4364) (0.6496,0.3727) (0.6576,0.3117) (0.6656,0.2533) (0.6736,0.1973) (0.6816,0.1436) (0.6896,0.0921) (0.6976,0.0426) (0.7056,-0.0050) (0.7136,-0.0507) (0.7216,-0.0947) (0.7295,-0.1371) (0.7375,-0.1779) (0.7455,-0.2172) (0.7535,-0.2552) (0.7615,-0.2918) (0.7695,-0.3271) (0.7775,-0.3613) (0.7855,-0.3943) (0.7935,-0.4262) (0.8015,-0.4571) (0.8095,-0.4870) (0.8175,-0.5160) (0.8255,-0.5441) (0.8335,-0.5713) (0.8415,-0.5977) (0.8495,-0.6234) (0.8575,-0.6483) (0.8655,-0.6726) (0.8735,-0.6961) (0.8815,-0.7190) (0.8894,-0.7414) (0.8974,-0.7631) (0.9054,-0.7843) (0.9134,-0.8049) (0.9214,-0.8250) (0.9294,-0.8447) (0.9374,-0.8639) (0.9454,-0.8826) (0.9534,-0.9009) (0.9614,-0.9188) (0.9694,-0.9363) (0.9774,-0.9535) (0.9854,-0.9702) (0.9934,-0.9867) (1.0014,-1.0028) (1.0094,-1.0185) (1.0174,-1.0340) (1.0254,-1.0492) (1.0334,-1.0641) (1.0414,-1.0787) (1.0493,-1.0931) (1.0573,-1.1072) (1.0653,-1.1211) (1.0733,-1.1348) (1.0813,-1.1482) (1.0893,-1.1614) (1.0973,-1.1744) (1.1053,-1.1873) (1.1133,-1.1999) (1.1213,-1.2124) (1.1293,-1.2246) (1.1373,-1.2368) (1.1453,-1.2487) (1.1533,-1.2605) (1.1613,-1.2722) (1.1693,-1.2837) (1.1773,-1.2950) (1.1853,-1.3062) (1.1933,-1.3173) (1.2013,-1.3283) (1.2092,-1.3392) (1.2172,-1.3499) (1.2252,-1.3606) (1.2332,-1.3711) (1.2412,-1.3815) (1.2492,-1.3918) (1.2572,-1.4020) (1.2652,-1.4122) (1.2732,-1.4222) (1.2812,-1.4322) (1.2892,-1.4420) (1.2972,-1.4518) (1.3052,-1.4615) (1.3132,-1.4712) (1.3212,-1.4807) (1.3292,-1.4902) (1.3372,-1.4997) (1.3452,-1.5090) (1.3532,-1.5183) (1.3612,-1.5276) (1.3691,-1.5368) (1.3771,-1.5459) (1.3851,-1.5549) (1.3931,-1.5640) (1.4011,-1.5729) (1.4091,-1.5818) (1.4171,-1.5907) (1.4251,-1.5995) (1.4331,-1.6083) (1.4411,-1.6171) (1.4491,-1.6257) (1.4571,-1.6344) (1.4651,-1.6430) (1.4731,-1.6516) (1.4811,-1.6601) (1.4891,-1.6686) (1.4971,-1.6771) (1.5051,-1.6855) (1.5131,-1.6940) (1.5211,-1.7023) (1.5290,-1.7107) (1.5370,-1.7190) (1.5450,-1.7273) (1.5530,-1.7355) (1.5610,-1.7438) (1.5690,-1.7520) (1.5770,-1.7602) (1.5850,-1.7684) (1.5930,-1.7765) (1.6010,-1.7846) (1.6090,-1.7927) (1.6170,-1.8008) (1.6250,-1.8088) (1.6330,-1.8169) (1.6410,-1.8249) (1.6490,-1.8329) (1.6570,-1.8409) (1.6650,-1.8489) (1.6730,-1.8568) (1.6810,-1.8648) (1.6889,-1.8727) (1.6969,-1.8806) (1.7049,-1.8885) (1.7129,-1.8964) (1.7209,-1.9042) (1.7289,-1.9121) (1.7369,-1.9199) (1.7449,-1.9278) (1.7529,-1.9356) (1.7609,-1.9434) (1.7689,-1.9512) (1.7769,-1.9590) (1.7849,-1.9667) (1.7929,-1.9745) (1.8009,-1.9823) (1.8089,-1.9900) (1.8169,-1.9978) (1.8249,-2.0055) (1.8329,-2.0132) (1.8409,-2.0209) (1.8488,-2.0286) (1.8568,-2.0363) (1.8648,-2.0440) (1.8728,-2.0517) (1.8808,-2.0594) (1.8888,-2.0671) (1.8968,-2.0747) (1.9048,-2.0824) (1.9128,-2.0901) (1.9208,-2.0977) (1.9288,-2.1054) (1.9368,-2.1130) (1.9448,-2.1207) (1.9528,-2.1283) (1.9608,-2.1359) (1.9688,-2.1435) (1.9768,-2.1512) (1.9848,-2.1588) (1.9928,-2.1664) (2.0008,-2.1740) (2.0087,-2.1816) (2.0167,-2.1892) (2.0247,-2.1968) (2.0327,-2.2044) (2.0407,-2.2120) (2.0487,-2.2196) (2.0567,-2.2272) (2.0647,-2.2348) (2.0727,-2.2424) (2.0807,-2.2499) (2.0887,-2.2575) (2.0967,-2.2651) (2.1047,-2.2727) (2.1127,-2.2803) (2.1207,-2.2878) (2.1287,-2.2954) (2.1367,-2.3030) (2.1447,-2.3105) (2.1527,-2.3181) (2.1607,-2.3257) (2.1686,-2.3332) (2.1766,-2.3408) (2.1846,-2.3484) (2.1926,-2.3559) (2.2006,-2.3635) (2.2086,-2.3711) (2.2166,-2.3786) (2.2246,-2.3862) (2.2326,-2.3937) (2.2406,-2.4013) (2.2486,-2.4089) (2.2566,-2.4164) (2.2646,-2.4240) (2.2726,-2.4315) (2.2806,-2.4391) (2.2886,-2.4466) (2.2966,-2.4542) (2.3046,-2.4618) (2.3126,-2.4693) (2.3206,-2.4769) (2.3285,-2.4844) (2.3365,-2.4920) (2.3445,-2.4996) (2.3525,-2.5071) (2.3605,-2.5147) (2.3685,-2.5222) (2.3765,-2.5298) (2.3845,-2.5373) (2.3925,-2.5449) (2.4005,-2.5525) (2.4085,-2.5600) (2.4165,-2.5676) (2.4245,-2.5751) (2.4325,-2.5827) (2.4405,-2.5903) (2.4485,-2.5978) (2.4565,-2.6054) (2.4645,-2.6130) (2.4725,-2.6205) (2.4805,-2.6281) (2.4884,-2.6357) (2.4964,-2.6432) (2.5044,-2.6508) (2.5124,-2.6584) (2.5204,-2.6659) (2.5284,-2.6735) (2.5364,-2.6811) (2.5444,-2.6887) (2.5524,-2.6962) (2.5604,-2.7038) (2.5684,-2.7114) (2.5764,-2.7190) (2.5844,-2.7265) (2.5924,-2.7341) (2.6004,-2.7417) (2.6084,-2.7493) (2.6164,-2.7569) (2.6244,-2.7645) (2.6324,-2.7720) (2.6404,-2.7796) (2.6483,-2.7872) (2.6563,-2.7948) (2.6643,-2.8024) (2.6723,-2.8100) (2.6803,-2.8176) (2.6883,-2.8252) (2.6963,-2.8327) (2.7043,-2.8403) (2.7123,-2.8479) (2.7203,-2.8555) (2.7283,-2.8631) (2.7363,-2.8707) (2.7443,-2.8783) (2.7523,-2.8859) (2.7603,-2.8935) (2.7683,-2.9011) (2.7763,-2.9087) (2.7843,-2.9164) (2.7923,-2.9240) (2.8003,-2.9316) (2.8082,-2.9392) (2.8162,-2.9468) (2.8242,-2.9544) (2.8322,-2.9620) (2.8402,-2.9696) (2.8482,-2.9772) (2.8562,-2.9849) (2.8642,-2.9925) (2.8722,-3.0001) (2.8802,-3.0077) (2.8882,-3.0153) (2.8962,-3.0230) (2.9042,-3.0306) (2.9122,-3.0382) (2.9202,-3.0458) (2.9282,-3.0535) (2.9362,-3.0611) (2.9442,-3.0687) (2.9522,-3.0764) (2.9602,-3.0840) (2.9681,-3.0916) (2.9761,-3.0993) (2.9841,-3.1069) (2.9921,-3.1145) (3.0001,-3.1222) (3.0081,-3.1298) (3.0161,-3.1375) (3.0241,-3.1451) (3.0321,-3.1528) (3.0401,-3.1604) (3.0481,-3.1681) (3.0561,-3.1757) (3.0641,-3.1834) (3.0721,-3.1910) (3.0801,-3.1987) (3.0881,-3.2063) (3.0961,-3.2140) (3.1041,-3.2216) (3.1121,-3.2293) (3.1201,-3.2369) (3.1280,-3.2446) (3.1360,-3.2523) (3.1440,-3.2599) (3.1520,-3.2676) (3.1600,-3.2752) (3.1680,-3.2829) (3.1760,-3.2906) (3.1840,-3.2982) (3.1920,-3.3059) (3.2000,-3.3136)};
\fill (1,-1) circle (1.5pt); \fill (0.7035,0) circle (1.5pt) node[above right]{$\alpha$};
\fill (1.6487,-1.8327) circle (1.5pt);
\Ticks{1,2,3}{-3,-2,-1,1,2}
\node[cDef] at (0.75,2){$(C_f)$};
\node[cRem] at (2.6,-2.2){$(\Delta)$};
\node[cThm] at (2.1,-2.6){$(T)$};
\end{GR}

\pb $\dfrac{\ln x}{x^2}=m\iff-x-\dfrac{\ln x}{x^2}=-x-m\iff f(x)=-x-m$.

الحلول هي فواصل نقط تقاطع $(C_f)$ مع المستقيم $(D_m):\ y=-x-m$ الموازي لـ $(\Delta)$. لدينا $(D_0)=(\Delta)$ و $(D_{\frac{1}{2e}})=(T)$ ، وبيانيا $(C_f)$ يقع فوق $(T)$ ويمسه عند $x=\sqrt e$:
\begin{itemize}
\item إذا كان $m\leqslant0$: $(D_m)$ يقطع $(C_f)$ في نقطة واحدة: حلّ وحيد.
\item إذا كان $0<m<\frac{1}{2e}$: $(D_m)$ محصور بين $(\Delta)$ و $(T)$ فيقطع $(C_f)$ في نقطتين: حلّان مختلفان.
\item إذا كان $m=\frac{1}{2e}$: حلّ وحيد $\sqrt e$ ، وإذا كان $m>\frac1{2e}$: لا توجد حلول.
\end{itemize}
إذن: \bx{m\in\left]0;\frac{1}{2e}\right[}.
\end{hal}

% ===================== بكالوريا 2025 م2 =====================
\begin{tamrin}{بكالوريا 2025 -- الموضوع الثاني}
\textbf{I)} $g$ الدالة العددية المعرّفة على $]-4;+\infty[$ بـ: \ $g(x)=\dfrac{x^2+(x^2+8x)\ln(x+4)}{(x+4)^2}$.

تمثيلها البياني $(C_g)$ يقطع حامل محور الفواصل في نقطتين فاصلتاهما $\alpha$ و $0$ كما في الشكل المقابل.
\begin{GR}{-4}{4}{-1.5}{2.4}{0.8}
\draw[dashed,gray] (-4,-1.5)--(-4,2.4);
\draw[cDef,thick] plot coordinates {(-2.7160,2.2977) (-2.6959,2.0411) (-2.6759,1.8024) (-2.6558,1.5802) (-2.6358,1.3732) (-2.6157,1.1805) (-2.5957,1.0010) (-2.5756,0.8338) (-2.5556,0.6781) (-2.5355,0.5330) (-2.5155,0.3979) (-2.4954,0.2720) (-2.4754,0.1548) (-2.4553,0.0458) (-2.4353,-0.0558) (-2.4152,-0.1502) (-2.3952,-0.2380) (-2.3751,-0.3196) (-2.3551,-0.3954) (-2.3350,-0.4658) (-2.3150,-0.5310) (-2.2950,-0.5914) (-2.2749,-0.6474) (-2.2549,-0.6991) (-2.2348,-0.7469) (-2.2148,-0.7909) (-2.1947,-0.8314) (-2.1747,-0.8686) (-2.1546,-0.9027) (-2.1346,-0.9339) (-2.1145,-0.9623) (-2.0945,-0.9882) (-2.0744,-1.0116) (-2.0544,-1.0327) (-2.0343,-1.0517) (-2.0143,-1.0686) (-1.9942,-1.0836) (-1.9742,-1.0967) (-1.9541,-1.1082) (-1.9341,-1.1180) (-1.9140,-1.1263) (-1.8940,-1.1332) (-1.8740,-1.1387) (-1.8539,-1.1430) (-1.8339,-1.1460) (-1.8138,-1.1479) (-1.7938,-1.1487) (-1.7737,-1.1485) (-1.7537,-1.1474) (-1.7336,-1.1453) (-1.7136,-1.1424) (-1.6935,-1.1387) (-1.6735,-1.1342) (-1.6534,-1.1290) (-1.6334,-1.1231) (-1.6133,-1.1166) (-1.5933,-1.1095) (-1.5732,-1.1018) (-1.5532,-1.0936) (-1.5331,-1.0849) (-1.5131,-1.0757) (-1.4930,-1.0660) (-1.4730,-1.0560) (-1.4530,-1.0455) (-1.4329,-1.0346) (-1.4129,-1.0235) (-1.3928,-1.0119) (-1.3728,-1.0001) (-1.3527,-0.9880) (-1.3327,-0.9756) (-1.3126,-0.9630) (-1.2926,-0.9501) (-1.2725,-0.9370) (-1.2525,-0.9237) (-1.2324,-0.9102) (-1.2124,-0.8965) (-1.1923,-0.8827) (-1.1723,-0.8686) (-1.1522,-0.8545) (-1.1322,-0.8402) (-1.1121,-0.8258) (-1.0921,-0.8113) (-1.0720,-0.7967) (-1.0520,-0.7819) (-1.0320,-0.7671) (-1.0119,-0.7522) (-0.9919,-0.7373) (-0.9718,-0.7223) (-0.9518,-0.7072) (-0.9317,-0.6920) (-0.9117,-0.6769) (-0.8916,-0.6617) (-0.8716,-0.6464) (-0.8515,-0.6311) (-0.8315,-0.6158) (-0.8114,-0.6005) (-0.7914,-0.5852) (-0.7713,-0.5698) (-0.7513,-0.5545) (-0.7312,-0.5391) (-0.7112,-0.5238) (-0.6911,-0.5085) (-0.6711,-0.4931) (-0.6510,-0.4778) (-0.6310,-0.4625) (-0.6110,-0.4472) (-0.5909,-0.4320) (-0.5709,-0.4167) (-0.5508,-0.4015) (-0.5308,-0.3863) (-0.5107,-0.3712) (-0.4907,-0.3561) (-0.4706,-0.3410) (-0.4506,-0.3259) (-0.4305,-0.3109) (-0.4105,-0.2959) (-0.3904,-0.2810) (-0.3704,-0.2661) (-0.3503,-0.2513) (-0.3303,-0.2365) (-0.3102,-0.2217) (-0.2902,-0.2070) (-0.2701,-0.1923) (-0.2501,-0.1777) (-0.2300,-0.1632) (-0.2100,-0.1487) (-0.1900,-0.1342) (-0.1699,-0.1198) (-0.1499,-0.1055) (-0.1298,-0.0912) (-0.1098,-0.0769) (-0.0897,-0.0628) (-0.0697,-0.0486) (-0.0496,-0.0346) (-0.0296,-0.0206) (-0.0095,-0.0066) (0.0105,0.0073) (0.0306,0.0211) (0.0506,0.0349) (0.0707,0.0486) (0.0907,0.0623) (0.1108,0.0759) (0.1308,0.0894) (0.1509,0.1029) (0.1709,0.1163) (0.1910,0.1297) (0.2110,0.1430) (0.2310,0.1562) (0.2511,0.1694) (0.2711,0.1825) (0.2912,0.1956) (0.3112,0.2086) (0.3313,0.2215) (0.3513,0.2344) (0.3714,0.2472) (0.3914,0.2600) (0.4115,0.2727) (0.4315,0.2853) (0.4516,0.2979) (0.4716,0.3104) (0.4917,0.3229) (0.5117,0.3353) (0.5318,0.3476) (0.5518,0.3599) (0.5719,0.3721) (0.5919,0.3843) (0.6120,0.3964) (0.6320,0.4084) (0.6520,0.4204) (0.6721,0.4323) (0.6921,0.4442) (0.7122,0.4560) (0.7322,0.4678) (0.7523,0.4795) (0.7723,0.4911) (0.7924,0.5027) (0.8124,0.5142) (0.8325,0.5257) (0.8525,0.5371) (0.8726,0.5485) (0.8926,0.5598) (0.9127,0.5710) (0.9327,0.5822) (0.9528,0.5934) (0.9728,0.6044) (0.9929,0.6155) (1.0129,0.6265) (1.0330,0.6374) (1.0530,0.6483) (1.0730,0.6591) (1.0931,0.6698) (1.1131,0.6806) (1.1332,0.6912) (1.1532,0.7018) (1.1733,0.7124) (1.1933,0.7229) (1.2134,0.7333) (1.2334,0.7438) (1.2535,0.7541) (1.2735,0.7644) (1.2936,0.7747) (1.3136,0.7849) (1.3337,0.7950) (1.3537,0.8051) (1.3738,0.8152) (1.3938,0.8252) (1.4139,0.8352) (1.4339,0.8451) (1.4540,0.8550) (1.4740,0.8648) (1.4940,0.8746) (1.5141,0.8843) (1.5341,0.8940) (1.5542,0.9036) (1.5742,0.9132) (1.5943,0.9227) (1.6143,0.9322) (1.6344,0.9417) (1.6544,0.9511) (1.6745,0.9604) (1.6945,0.9698) (1.7146,0.9790) (1.7346,0.9883) (1.7547,0.9975) (1.7747,1.0066) (1.7948,1.0157) (1.8148,1.0248) (1.8349,1.0338) (1.8549,1.0428) (1.8750,1.0517) (1.8950,1.0606) (1.9150,1.0695) (1.9351,1.0783) (1.9551,1.0870) (1.9752,1.0958) (1.9952,1.1045) (2.0153,1.1131) (2.0353,1.1217) (2.0554,1.1303) (2.0754,1.1388) (2.0955,1.1473) (2.1155,1.1558) (2.1356,1.1642) (2.1556,1.1726) (2.1757,1.1809) (2.1957,1.1893) (2.2158,1.1975) (2.2358,1.2058) (2.2559,1.2139) (2.2759,1.2221) (2.2960,1.2302) (2.3160,1.2383) (2.3360,1.2464) (2.3561,1.2544) (2.3761,1.2624) (2.3962,1.2703) (2.4162,1.2782) (2.4363,1.2861) (2.4563,1.2939) (2.4764,1.3017) (2.4964,1.3095) (2.5165,1.3172) (2.5365,1.3249) (2.5566,1.3326) (2.5766,1.3403) (2.5967,1.3479) (2.6167,1.3554) (2.6368,1.3630) (2.6568,1.3705) (2.6769,1.3780) (2.6969,1.3854) (2.7170,1.3928) (2.7370,1.4002) (2.7570,1.4075) (2.7771,1.4149) (2.7971,1.4221) (2.8172,1.4294) (2.8372,1.4366) (2.8573,1.4438) (2.8773,1.4510) (2.8974,1.4581) (2.9174,1.4652) (2.9375,1.4723) (2.9575,1.4793) (2.9776,1.4864) (2.9976,1.4934) (3.0177,1.5003) (3.0377,1.5072) (3.0578,1.5142) (3.0778,1.5210) (3.0979,1.5279) (3.1179,1.5347) (3.1380,1.5415) (3.1580,1.5482) (3.1780,1.5550) (3.1981,1.5617) (3.2181,1.5684) (3.2382,1.5750) (3.2582,1.5817) (3.2783,1.5883) (3.2983,1.5948) (3.3184,1.6014) (3.3384,1.6079) (3.3585,1.6144) (3.3785,1.6209) (3.3986,1.6273) (3.4186,1.6337) (3.4387,1.6401) (3.4587,1.6465) (3.4788,1.6529) (3.4988,1.6592) (3.5189,1.6655) (3.5389,1.6717) (3.5590,1.6780) (3.5790,1.6842) (3.5990,1.6904) (3.6191,1.6966) (3.6391,1.7027) (3.6592,1.7089) (3.6792,1.7150) (3.6993,1.7211) (3.7193,1.7271) (3.7394,1.7332) (3.7594,1.7392) (3.7795,1.7452) (3.7995,1.7511) (3.8196,1.7571) (3.8396,1.7630) (3.8597,1.7689) (3.8797,1.7748) (3.8998,1.7806) (3.9198,1.7865) (3.9399,1.7923) (3.9599,1.7981) (3.9800,1.8038) (4.0000,1.8096)};
\fill (-2.4465,0) circle (2pt) node[above right]{$\alpha$}; \fill (0,0) circle (2pt);
\Ticks{-3,-2,-1,1,2,3}{-1,1,2}
\node[cDef] at (2.2,1.5){$(C_g)$};
\end{GR}
\QQ{1} بقراءة بيانية حدّد إشارة $g(x)$ على $]-4;+\infty[$.
\QQ{2} تحقّق أنّ $-2{,}5<\alpha<-2{,}4$.

\textbf{II)} $f$ الدالة المعرّفة على $]-4;+\infty[$ بـ: \ $f(x)=\dfrac{x^2\ln(x+4)}{x+4}$ ، و $(C_f)$ تمثيلها البياني في المستوي المنسوب إلى المعلم المتعامد والمتجانس $(O;\vec i,\vec j)$.
\QQ{1} أحسب $\lim\limits_{x\to-4}f(x)$ وبيّن أنّ $\lim\limits_{x\to+\infty}f(x)=+\infty$.
\QQ{2}
\begin{sq}
\qa بيّن أنّه من أجل كلّ $x$ من $]-4;+\infty[$: \ $f'(x)=g(x)$.
\qb استنتج اتّجاه تغيّر الدالّة $f$ ثمّ شكّل جدول تغيّراتها.
\end{sq}
\QQ{3}
\begin{sq}
\qa عيّن فاصلتي نقطتي تقاطع $(C_f)$ مع حامل محور الفواصل.
\qb أحسب $f(2)$ ، $f(4)$ ثمّ ارسم $(C_f)$. (نأخذ $f(\alpha)\approx1{,}7$).
\end{sq}
\QQ{4} عيّن قيم العدد الحقيقي الموجب تماما $m$ حتّى يكون للمعادلة $f(x)=\ln(m)$ ثلاثة حلول مختلفة.
\QQ{5} $h$ الدالّة المعرّفة على $]-4;+\infty[$ بـ: $h(x)=\dfrac{(x^2+1)\ln(x+4)}{x+4}$ ، و $(C_h)$ تمثيلها البياني.
\begin{sq}
\qa بيّن أنّه من أجل كلّ $x$ من $[-3;0]$: \ $h(x)-f(x)\geqslant0$.
\end{sq}
\end{tamrin}

\begin{hal}
\textbf{الجزء I}

\QQ{1} من الشكل:
\begin{TAB}
\tkzTabInit[lgt=2,espcl=2.5]{$x$/1,$g(x)$/1}{$-4$,$\alpha$,$0$,$+\infty$}
\tkzTabLine{d,+,z,-,z,+,}
\end{TAB}
\QQ{2} $g$ مستمرة على $[-2{,}5;-2{,}4]$ ، و $g(-2{,}5)\approx0{,}30>0$ ، $g(-2{,}4)\approx-0{,}22<0$ ، إذن $g(-2{,}5)\times g(-2{,}4)<0$ ، وبما أن $(C_g)$ يقطع محور الفواصل مرة واحدة على $]-4;0[$ فإن \bx{-2{,}5<\alpha<-2{,}4}.

\textbf{الجزء II}

\QQ{1} لما $x\to-4$ ( $x>-4$ ): $x^2\to16$ ، $\ln(x+4)\to-\infty$ و $x+4\to0^+$ ، إذن \bx{\lim_{x\to-4}f(x)=-\infty} ، والمستقيم $x=-4$ مقارب عمودي لـ $(C_f)$.

لما $x\to+\infty$: $f(x)=\dfrac{x^2}{x+4}\cdot\ln(x+4)$ مع $\dfrac{x^2}{x+4}=\dfrac{x}{1+\frac4x}\to+\infty$ و $\ln(x+4)\to+\infty$ ، إذن \bx{\lim_{x\to+\infty}f(x)=+\infty}.

\QQ{2}
\pa $f=\frac uv$ مع $u=x^2\ln(x+4)$ ، $u'=2x\ln(x+4)+\frac{x^2}{x+4}$ ، $v=x+4$ ، $v'=1$:
\begin{align*}
f'(x)&=\frac{\Big[2x\ln(x+4)+\frac{x^2}{x+4}\Big](x+4)-x^2\ln(x+4)}{(x+4)^2}
=\frac{(2x^2+8x)\ln(x+4)+x^2-x^2\ln(x+4)}{(x+4)^2}\\
&=\frac{x^2+(x^2+8x)\ln(x+4)}{(x+4)^2}=g(x).
\end{align*}
\pb إشارة $f'(x)$ هي إشارة $g(x)$: \bxt{$f$ متزايدة تماما على $]-4;\alpha]$ و $[0;+\infty[$ ، ومتناقصة تماما على $[\alpha;0]$}. و $f(0)=0$.
\begin{TAB}
\tkzTabInit[lgt=2,espcl=2.5]{$x$/1,$f'(x)$/1,$f(x)$/2}{$-4$,$\alpha$,$0$,$+\infty$}
\tkzTabLine{d,+,z,-,z,+,}
\tkzTabVar{D-/$-\infty$,+/$f(\alpha)$,-/$0$,+/$+\infty$}
\end{TAB}

\QQ{3}
\pa $f(x)=0\iff x^2\ln(x+4)=0\iff x=0$ أو $\ln(x+4)=0\iff x=0$ أو $x=-3$. \ \bx{x=-3\ ;\ x=0}.

\pb $f(2)=\dfrac{4\ln6}{6}=\dfrac23\ln6\approx1{,}19$ \quad ، \quad $f(4)=\dfrac{16\ln8}{8}=2\ln8=6\ln2\approx4{,}16$.
\begin{GR}{-4}{4.5}{-3}{5}{0.85}
\draw[dashed,cThm] (-4,-3)--(-4,5);
\draw[cDef,thick] plot coordinates {(-3.2109,-3.0941) (-3.1896,-2.6386) (-3.1683,-2.2236) (-3.1470,-1.8455) (-3.1257,-1.5006) (-3.1044,-1.1860) (-3.0831,-0.8990) (-3.0618,-0.6370) (-3.0405,-0.3980) (-3.0192,-0.1798) (-2.9979,0.0191) (-2.9766,0.2005) (-2.9553,0.3659) (-2.9340,0.5164) (-2.9127,0.6533) (-2.8914,0.7777) (-2.8701,0.8906) (-2.8488,0.9928) (-2.8275,1.0853) (-2.8062,1.1686) (-2.7849,1.2437) (-2.7636,1.3110) (-2.7423,1.3711) (-2.7210,1.4246) (-2.6997,1.4720) (-2.6784,1.5137) (-2.6571,1.5501) (-2.6358,1.5817) (-2.6145,1.6087) (-2.5932,1.6316) (-2.5718,1.6506) (-2.5505,1.6659) (-2.5292,1.6779) (-2.5079,1.6869) (-2.4866,1.6929) (-2.4653,1.6963) (-2.4440,1.6972) (-2.4227,1.6958) (-2.4014,1.6923) (-2.3801,1.6869) (-2.3588,1.6796) (-2.3375,1.6707) (-2.3162,1.6602) (-2.2949,1.6482) (-2.2736,1.6350) (-2.2523,1.6205) (-2.2310,1.6050) (-2.2097,1.5884) (-2.1884,1.5709) (-2.1671,1.5525) (-2.1458,1.5333) (-2.1245,1.5135) (-2.1032,1.4929) (-2.0819,1.4718) (-2.0606,1.4502) (-2.0393,1.4281) (-2.0180,1.4056) (-1.9967,1.3828) (-1.9754,1.3596) (-1.9541,1.3361) (-1.9328,1.3124) (-1.9115,1.2884) (-1.8902,1.2644) (-1.8689,1.2401) (-1.8476,1.2158) (-1.8263,1.1914) (-1.8050,1.1670) (-1.7837,1.1425) (-1.7624,1.1180) (-1.7411,1.0936) (-1.7198,1.0692) (-1.6985,1.0449) (-1.6772,1.0207) (-1.6559,0.9965) (-1.6346,0.9725) (-1.6133,0.9487) (-1.5920,0.9250) (-1.5707,0.9014) (-1.5494,0.8781) (-1.5281,0.8549) (-1.5068,0.8320) (-1.4855,0.8092) (-1.4642,0.7867) (-1.4429,0.7644) (-1.4216,0.7424) (-1.4003,0.7206) (-1.3790,0.6991) (-1.3577,0.6779) (-1.3364,0.6569) (-1.3151,0.6362) (-1.2938,0.6158) (-1.2725,0.5957) (-1.2512,0.5759) (-1.2299,0.5564) (-1.2086,0.5372) (-1.1873,0.5183) (-1.1660,0.4997) (-1.1447,0.4815) (-1.1234,0.4636) (-1.1021,0.4460) (-1.0808,0.4287) (-1.0595,0.4117) (-1.0382,0.3951) (-1.0169,0.3789) (-0.9956,0.3629) (-0.9743,0.3473) (-0.9530,0.3321) (-0.9317,0.3172) (-0.9104,0.3026) (-0.8891,0.2884) (-0.8678,0.2745) (-0.8465,0.2610) (-0.8252,0.2478) (-0.8039,0.2349) (-0.7826,0.2224) (-0.7613,0.2103) (-0.7400,0.1985) (-0.7187,0.1870) (-0.6974,0.1759) (-0.6761,0.1652) (-0.6548,0.1548) (-0.6335,0.1447) (-0.6122,0.1350) (-0.5909,0.1256) (-0.5696,0.1166) (-0.5483,0.1079) (-0.5270,0.0996) (-0.5057,0.0916) (-0.4844,0.0839) (-0.4631,0.0766) (-0.4418,0.0696) (-0.4205,0.0630) (-0.3992,0.0567) (-0.3779,0.0507) (-0.3566,0.0451) (-0.3353,0.0398) (-0.3140,0.0349) (-0.2927,0.0303) (-0.2714,0.0260) (-0.2501,0.0220) (-0.2288,0.0184) (-0.2075,0.0151) (-0.1862,0.0122) (-0.1649,0.0095) (-0.1436,0.0072) (-0.1223,0.0052) (-0.1010,0.0036) (-0.0797,0.0022) (-0.0584,0.0012) (-0.0371,0.0005) (-0.0158,0.0001) (0.0055,0.0000) (0.0268,0.0002) (0.0481,0.0008) (0.0694,0.0017) (0.0907,0.0028) (0.1120,0.0043) (0.1333,0.0061) (0.1546,0.0082) (0.1759,0.0106) (0.1972,0.0133) (0.2185,0.0163) (0.2398,0.0196) (0.2612,0.0232) (0.2825,0.0271) (0.3038,0.0313) (0.3251,0.0358) (0.3464,0.0406) (0.3677,0.0456) (0.3890,0.0510) (0.4103,0.0566) (0.4316,0.0626) (0.4529,0.0688) (0.4742,0.0753) (0.4955,0.0821) (0.5168,0.0891) (0.5381,0.0965) (0.5594,0.1041) (0.5807,0.1120) (0.6020,0.1202) (0.6233,0.1286) (0.6446,0.1374) (0.6659,0.1464) (0.6872,0.1556) (0.7085,0.1652) (0.7298,0.1750) (0.7511,0.1850) (0.7724,0.1954) (0.7937,0.2060) (0.8150,0.2168) (0.8363,0.2279) (0.8576,0.2393) (0.8789,0.2509) (0.9002,0.2628) (0.9215,0.2749) (0.9428,0.2873) (0.9641,0.3000) (0.9854,0.3129) (1.0067,0.3260) (1.0280,0.3394) (1.0493,0.3531) (1.0706,0.3670) (1.0919,0.3811) (1.1132,0.3955) (1.1345,0.4101) (1.1558,0.4249) (1.1771,0.4400) (1.1984,0.4554) (1.2197,0.4710) (1.2410,0.4868) (1.2623,0.5028) (1.2836,0.5191) (1.3049,0.5356) (1.3262,0.5523) (1.3475,0.5693) (1.3688,0.5865) (1.3901,0.6039) (1.4114,0.6216) (1.4327,0.6394) (1.4540,0.6575) (1.4753,0.6759) (1.4966,0.6944) (1.5179,0.7132) (1.5392,0.7322) (1.5605,0.7514) (1.5818,0.7708) (1.6031,0.7904) (1.6244,0.8103) (1.6457,0.8303) (1.6670,0.8506) (1.6883,0.8711) (1.7096,0.8918) (1.7309,0.9127) (1.7522,0.9338) (1.7735,0.9552) (1.7948,0.9767) (1.8161,0.9984) (1.8374,1.0204) (1.8587,1.0425) (1.8800,1.0649) (1.9013,1.0874) (1.9226,1.1102) (1.9439,1.1331) (1.9652,1.1563) (1.9865,1.1796) (2.0078,1.2032) (2.0291,1.2269) (2.0504,1.2508) (2.0717,1.2750) (2.0930,1.2993) (2.1143,1.3238) (2.1356,1.3485) (2.1569,1.3734) (2.1782,1.3985) (2.1995,1.4238) (2.2208,1.4492) (2.2421,1.4749) (2.2634,1.5007) (2.2847,1.5267) (2.3060,1.5529) (2.3273,1.5793) (2.3486,1.6059) (2.3699,1.6326) (2.3912,1.6595) (2.4125,1.6866) (2.4338,1.7139) (2.4551,1.7414) (2.4764,1.7690) (2.4977,1.7968) (2.5190,1.8248) (2.5403,1.8530) (2.5616,1.8813) (2.5829,1.9098) (2.6042,1.9385) (2.6255,1.9674) (2.6468,1.9964) (2.6681,2.0256) (2.6894,2.0550) (2.7107,2.0845) (2.7320,2.1142) (2.7533,2.1441) (2.7746,2.1741) (2.7959,2.2043) (2.8172,2.2347) (2.8385,2.2652) (2.8598,2.2959) (2.8811,2.3268) (2.9024,2.3578) (2.9237,2.3890) (2.9450,2.4203) (2.9663,2.4518) (2.9876,2.4835) (3.0089,2.5153) (3.0302,2.5472) (3.0515,2.5794) (3.0728,2.6117) (3.0942,2.6441) (3.1155,2.6767) (3.1368,2.7094) (3.1581,2.7423) (3.1794,2.7754) (3.2007,2.8086) (3.2220,2.8420) (3.2433,2.8755) (3.2646,2.9091) (3.2859,2.9429) (3.3072,2.9769) (3.3285,3.0110) (3.3498,3.0453) (3.3711,3.0797) (3.3924,3.1142) (3.4137,3.1489) (3.4350,3.1837) (3.4563,3.2187) (3.4776,3.2539) (3.4989,3.2891) (3.5202,3.3245) (3.5415,3.3601) (3.5628,3.3958) (3.5841,3.4316) (3.6054,3.4676) (3.6267,3.5037) (3.6480,3.5400) (3.6693,3.5764) (3.6906,3.6129) (3.7119,3.6496) (3.7332,3.6864) (3.7545,3.7233) (3.7758,3.7604) (3.7971,3.7976) (3.8184,3.8350) (3.8397,3.8725) (3.8610,3.9101) (3.8823,3.9479) (3.9036,3.9857) (3.9249,4.0238) (3.9462,4.0619) (3.9675,4.1002) (3.9888,4.1386) (4.0101,4.1771) (4.0314,4.2158) (4.0527,4.2546) (4.0740,4.2935) (4.0953,4.3326) (4.1166,4.3718) (4.1379,4.4111) (4.1592,4.4505) (4.1805,4.4901) (4.2018,4.5298) (4.2231,4.5696) (4.2444,4.6095) (4.2657,4.6496) (4.2870,4.6898) (4.3083,4.7301) (4.3296,4.7705) (4.3509,4.8111) (4.3722,4.8518) (4.3935,4.8926) (4.4148,4.9335) (4.4361,4.9745) (4.4574,5.0157) (4.4787,5.0570) (4.5000,5.0984)};
\fill (-3,0) circle (2pt); \fill (0,0) circle (2pt);
\fill (-2.4465,1.6972) circle (2pt);
\fill (2,1.1945) circle (2pt); \fill (4,4.1589) circle (2pt);
\draw[dashed,gray] (-2.4465,0)--(-2.4465,1.6972);
\node[below,font=\scriptsize] at (-2.4465,0){$\alpha$};
\Ticks{-3,-2,-1,1,2,3,4}{-2,-1,1,2,3,4}
\node[cDef] at (2.6,3.4){$(C_f)$};
\end{GR}

\QQ{4} حلول $f(x)=\ln m$ هي فواصل نقط تقاطع $(C_f)$ مع المستقيم الأفقي $y=\ln m$.
من جدول التغيرات يكون ثلاثة حلول مختلفة (واحد في كل من $]-4;\alpha[$ ، $]\alpha;0[$ ، $]0;+\infty[$) إذا وفقط إذا كان:
\[0<\ln m<f(\alpha)\iff e^0<m<e^{f(\alpha)}\iff\bx{1<m<e^{f(\alpha)}}\qquad\big(e^{f(\alpha)}\approx5{,}5\big).\]

\QQ{5}
\pa $h(x)-f(x)=\dfrac{(x^2+1-x^2)\ln(x+4)}{x+4}=\dfrac{\ln(x+4)}{x+4}$.

من أجل $x\in[-3;0]$: $1\leqslant x+4\leqslant4$ ومنه $\ln(x+4)\geqslant\ln1=0$ و $x+4>0$ ، إذن \bx{h(x)-f(x)\geqslant0}.
\end{hal}

% ===================== بكالوريا 2026 م1 =====================
\begin{tamrin}{بكالوريا 2026 -- الموضوع الأول}
\textbf{I)} الدالة $g$ المعرفة على $]0;+\infty[$ بـ: $g(x)=2x-2-4\ln x$.
\QQ{1} ادرس اتجاه تغيّر الدالة $g$ وشكّل جدول تغيّراتها.
\QQ{2}
\begin{sq}
\qa احسب $g(1)$ ثم بيّن أن المعادلة $g(x)=0$ تقبل حلا وحيدا $\alpha$ في المجال $[3{,}5;3{,}6]$.
\qb استنتج حسب قيم $x$ إشارة $g(x)$.
\end{sq}

\textbf{II)} الدالة $f$ المعرفة على $[0;+\infty[$ بـ: $f(x)=x^2+2x-4x\ln x$ من أجل $x>0$ و $f(0)=0$.

$(C_f)$ تمثيلها البياني في المعلم المتعامد والمتجانس $(O;\vec i,\vec j)$.
\QQ{1} احسب $\lim\limits_{x\to0}\dfrac{f(x)}{x}$ ( $x>0$ ) ، فسّر النتيجة هندسيا واحسب $\lim\limits_{x\to+\infty}f(x)$.
\QQ{2} بيّن أنه من أجل كل $x>0$: $f'(x)=g(x)$ ، ثم استنتج اتجاه تغيّر الدالة $f$ وشكّل جدول تغيّراتها.
\QQ{3}
\begin{sq}
\qa بيّن أن $(C_f)$ يقبل نقطة انعطاف $A$ ، يُطلب تعيين إحداثييها.
\qb عيّن معادلة لـ $(T)$ مماس المنحنى $(C_f)$ عند النقطة ذات الفاصلة $2$.
\qc عيّن بيانيا قيم الوسيط الحقيقي $m$ حتى تقبل المعادلة $f(x)=m$ ثلاثة حلول.
\end{sq}
\QQ{4} ارسم $(T)$ و $(C_f)$. (نأخذ $f(\alpha)\simeq1{,}7$).
\end{tamrin}

\begin{hal}
\textbf{الجزء I}

\QQ{1} $g'(x)=2-\frac4x=\frac{2(x-2)}{x}$ ، إشارتها إشارة $x-2$ (لأن $x>0$).

النهايات: $\lim\limits_{x\to0}g(x)=+\infty$ (لأن $-4\ln x\to+\infty$) ، و $g(x)=x\Big(2-\frac2x-4\frac{\ln x}{x}\Big)$ مع $\frac{\ln x}{x}\to0$ ، إذن $\lim\limits_{x\to+\infty}g(x)=+\infty$.

$g(2)=2-4\ln2\approx-0{,}77$. \bxt{$g$ متناقصة تماما على $]0;2]$ ومتزايدة تماما على $[2;+\infty[$}.
\begin{TAB}
\tkzTabInit[lgt=2,espcl=3]{$x$/1,$g'(x)$/1,$g(x)$/2}{$0$,$2$,$+\infty$}
\tkzTabLine{d,-,z,+,}
\tkzTabVar{D+/$+\infty$,-/$2-4\ln2$,+/$+\infty$}
\end{TAB}

\QQ{2}
\pa $g(1)=2-2-4\ln1=0$ ، وهو الحل الوحيد على $]0;2]$ لأن $g$ متناقصة تماما عليه.

على $[2;+\infty[$: $g$ مستمرة ومتزايدة تماما، و $g(3{,}5)\approx-0{,}011<0$ ، $g(3{,}6)\approx0{,}076>0$.
حسب مبرهنة القيم المتوسطة: \bxt{$g(x)=0$ تقبل حلا وحيدا $\alpha\in[3{,}5;3{,}6]$} (إضافة إلى الحل $1$).

\pb
\begin{TAB}
\tkzTabInit[lgt=2,espcl=2.5]{$x$/1,$g(x)$/1}{$0$,$1$,$\alpha$,$+\infty$}
\tkzTabLine{d,+,z,-,z,+,}
\end{TAB}

\textbf{الجزء II}

\QQ{1} من أجل $x>0$: $\dfrac{f(x)-f(0)}{x-0}=\dfrac{f(x)}{x}=x+2-4\ln x$ ، و $-4\ln x\to+\infty$ لما $x\to0$ ، إذن:
\[\bx{\lim_{x\to0}\frac{f(x)}{x}=+\infty}\]
\textbf{التفسير:} $f$ غير قابلة للاشتقاق عند $0$ ، و $(C_f)$ يقبل عند المبدأ $O$ نصف مماس عمودي (موازٍ لمحور التراتيب).

عند $+\infty$: $f(x)=x^2\Big(1+\frac2x-4\frac{\ln x}{x}\Big)$ ، والقوس يؤول إلى $1$ ، إذن \bx{\lim_{x\to+\infty}f(x)=+\infty}.

\QQ{2} $f'(x)=2x+2-4\Big(\ln x+x\cdot\frac1x\Big)=2x+2-4\ln x-4=2x-2-4\ln x=g(x)$.

إشارة $f'$ هي إشارة $g$: \bxt{$f$ متزايدة تماما على $[0;1]$ و $[\alpha;+\infty[$ ، ومتناقصة تماما على $[1;\alpha]$}.

$f(1)=1+2-0=3$. ولأن $g(\alpha)=0$ أي $\ln\alpha=\frac{\alpha-1}{2}$ فإن $f(\alpha)=\alpha^2+2\alpha-2\alpha(\alpha-1)=-\alpha^2+4\alpha\approx1{,}7$.
\begin{TAB}
\tkzTabInit[lgt=2,espcl=2.5]{$x$/1,$f'(x)$/1,$f(x)$/2}{$0$,$1$,$\alpha$,$+\infty$}
\tkzTabLine{,+,z,-,z,+,}
\tkzTabVar{-/$0$,+/$3$,-/$f(\alpha)$,+/$+\infty$}
\end{TAB}

\QQ{3}
\pa $f''(x)=g'(x)=\dfrac{2(x-2)}{x}$ تنعدم عند $x=2$ مغيّرة إشارتها (سالبة قبلها وموجبة بعدها) ، إذن $(C_f)$ يقبل نقطة انعطاف فاصلتها $2$:
$f(2)=4+4-8\ln2=8-8\ln2$ ، ومنه \bx{A(2;\ 8-8\ln2)\approx(2;\ 2{,}45)}.

\pb $f'(2)=g(2)=2-4\ln2$ ، إذن $y=(2-4\ln2)(x-2)+8-8\ln2=(2-4\ln2)x-4+8\ln2+8-8\ln2$:
\[\bx{(T):\ y=(2-4\ln2)x+4}\]
\pc حلول $f(x)=m$ هي فواصل نقط تقاطع $(C_f)$ مع المستقيم الأفقي $y=m$. من جدول التغيرات ( $f(0)=0<f(\alpha)$ ):
تقبل المعادلة ثلاثة حلول (في $]0;1[$ ، $]1;\alpha[$ و $]\alpha;+\infty[$) إذا وفقط إذا كان \bx{f(\alpha)<m<3}.

\QQ{4} الرسم:
\begin{GR}{-0.5}{6}{-1}{5.5}{1.1}
\draw[cThm,thick] (0,4)--(6,-0.636);
\draw[cDef,thick] plot coordinates {(0.0000,0.0000) (0.0000,0.0006) (0.0000,0.0006) (0.0000,0.0006) (0.0000,0.0006) (0.0000,0.0007) (0.0000,0.0007) (0.0000,0.0007) (0.0000,0.0008) (0.0000,0.0008) (0.0000,0.0008) (0.0000,0.0009) (0.0000,0.0009) (0.0000,0.0009) (0.0000,0.0010) (0.0000,0.0010) (0.0000,0.0011) (0.0000,0.0011) (0.0000,0.0012) (0.0000,0.0012) (0.0000,0.0013) (0.0000,0.0013) (0.0000,0.0014) (0.0000,0.0014) (0.0000,0.0015) (0.0000,0.0016) (0.0000,0.0016) (0.0000,0.0017) (0.0000,0.0018) (0.0000,0.0018) (0.0000,0.0019) (0.0000,0.0020) (0.0001,0.0021) (0.0001,0.0022) (0.0001,0.0023) (0.0001,0.0024) (0.0001,0.0025) (0.0001,0.0026) (0.0001,0.0027) (0.0001,0.0028) (0.0001,0.0029) (0.0001,0.0031) (0.0001,0.0032) (0.0001,0.0033) (0.0001,0.0035) (0.0001,0.0036) (0.0001,0.0038) (0.0001,0.0039) (0.0001,0.0041) (0.0001,0.0043) (0.0001,0.0044) (0.0001,0.0046) (0.0001,0.0048) (0.0001,0.0050) (0.0001,0.0052) (0.0001,0.0055) (0.0002,0.0057) (0.0002,0.0059) (0.0002,0.0062) (0.0002,0.0064) (0.0002,0.0067) (0.0002,0.0070) (0.0002,0.0073) (0.0002,0.0076) (0.0002,0.0079) (0.0002,0.0082) (0.0002,0.0086) (0.0003,0.0089) (0.0003,0.0093) (0.0003,0.0097) (0.0003,0.0101) (0.0003,0.0105) (0.0003,0.0110) (0.0003,0.0114) (0.0004,0.0119) (0.0004,0.0124) (0.0004,0.0129) (0.0004,0.0134) (0.0004,0.0140) (0.0004,0.0146) (0.0005,0.0152) (0.0005,0.0158) (0.0005,0.0165) (0.0005,0.0172) (0.0006,0.0179) (0.0006,0.0186) (0.0006,0.0194) (0.0006,0.0202) (0.0007,0.0210) (0.0007,0.0219) (0.0007,0.0228) (0.0008,0.0237) (0.0008,0.0247) (0.0008,0.0257) (0.0009,0.0267) (0.0009,0.0278) (0.0010,0.0290) (0.0010,0.0302) (0.0011,0.0314) (0.0011,0.0327) (0.0012,0.0340) (0.0012,0.0354) (0.0013,0.0368) (0.0013,0.0383) (0.0014,0.0399) (0.0015,0.0415) (0.0015,0.0432) (0.0016,0.0449) (0.0017,0.0467) (0.0018,0.0486) (0.0019,0.0506) (0.0020,0.0526) (0.0020,0.0547) (0.0021,0.0569) (0.0022,0.0592) (0.0023,0.0616) (0.0025,0.0640) (0.0026,0.0666) (0.0027,0.0692) (0.0028,0.0720) (0.0030,0.0748) (0.0031,0.0778) (0.0032,0.0809) (0.0034,0.0841) (0.0036,0.0874) (0.0037,0.0909) (0.0039,0.0945) (0.0041,0.0982) (0.0043,0.1020) (0.0045,0.1060) (0.0047,0.1102) (0.0049,0.1145) (0.0052,0.1190) (0.0054,0.1236) (0.0057,0.1284) (0.0059,0.1334) (0.0062,0.1386) (0.0065,0.1440) (0.0068,0.1495) (0.0071,0.1553) (0.0075,0.1613) (0.0078,0.1675) (0.0082,0.1739) (0.0086,0.1805) (0.0090,0.1874) (0.0094,0.1946) (0.0099,0.2020) (0.0103,0.2096) (0.0108,0.2175) (0.0113,0.2258) (0.0119,0.2343) (0.0124,0.2431) (0.0130,0.2522) (0.0136,0.2616) (0.0143,0.2714) (0.0150,0.2815) (0.0150,0.2828) (0.0157,0.2919) (0.0164,0.3027) (0.0172,0.3139) (0.0180,0.3255) (0.0188,0.3374) (0.0197,0.3497) (0.0207,0.3625) (0.0217,0.3757) (0.0227,0.3893) (0.0238,0.4034) (0.0249,0.4179) (0.0261,0.4329) (0.0273,0.4484) (0.0286,0.4644) (0.0299,0.4809) (0.0301,0.4826) (0.0314,0.4979) (0.0328,0.5155) (0.0344,0.5336) (0.0360,0.5522) (0.0377,0.5715) (0.0395,0.5913) (0.0414,0.6117) (0.0433,0.6328) (0.0451,0.6514) (0.0454,0.6544) (0.0476,0.6767) (0.0498,0.6997) (0.0522,0.7233) (0.0546,0.7476) (0.0572,0.7725) (0.0599,0.7982) (0.0602,0.8002) (0.0628,0.8246) (0.0657,0.8517) (0.0689,0.8795) (0.0721,0.9080) (0.0752,0.9343) (0.0755,0.9373) (0.0791,0.9673) (0.0829,0.9981) (0.0868,1.0297) (0.0902,1.0567) (0.0909,1.0620) (0.0952,1.0951) (0.0997,1.1290) (0.1044,1.1636) (0.1053,1.1695) (0.1094,1.1990) (0.1146,1.2352) (0.1200,1.2721) (0.1203,1.2741) (0.1353,1.3717) (0.1504,1.4630) (0.1654,1.5487) (0.1805,1.6294) (0.1955,1.7055) (0.2105,1.7775) (0.2256,1.8456) (0.2406,1.9101) (0.2556,1.9714) (0.2707,2.0295) (0.2857,2.0848) (0.3008,2.1373) (0.3158,2.1873) (0.3308,2.2349) (0.3459,2.2802) (0.3609,2.3233) (0.3759,2.3644) (0.3910,2.4035) (0.4060,2.4407) (0.4211,2.4762) (0.4361,2.5100) (0.4511,2.5422) (0.4662,2.5728) (0.4812,2.6019) (0.4962,2.6296) (0.5113,2.6559) (0.5263,2.6809) (0.5414,2.7046) (0.5564,2.7272) (0.5714,2.7485) (0.5865,2.7687) (0.6015,2.7878) (0.6165,2.8059) (0.6316,2.8230) (0.6466,2.8391) (0.6617,2.8542) (0.6767,2.8684) (0.6917,2.8817) (0.7068,2.8942) (0.7218,2.9058) (0.7368,2.9167) (0.7519,2.9268) (0.7669,2.9361) (0.7820,2.9447) (0.7970,2.9526) (0.8120,2.9598) (0.8271,2.9663) (0.8421,2.9722) (0.8571,2.9775) (0.8722,2.9822) (0.8872,2.9863) (0.9023,2.9898) (0.9173,2.9928) (0.9323,2.9952) (0.9474,2.9971) (0.9624,2.9986) (0.9774,2.9995) (0.9925,2.9999) (1.0075,2.9999) (1.0226,2.9995) (1.0376,2.9986) (1.0526,2.9973) (1.0677,2.9956) (1.0827,2.9935) (1.0977,2.9910) (1.1128,2.9882) (1.1278,2.9850) (1.1429,2.9814) (1.1579,2.9775) (1.1729,2.9733) (1.1880,2.9687) (1.2030,2.9639) (1.2180,2.9587) (1.2331,2.9532) (1.2481,2.9475) (1.2632,2.9415) (1.2782,2.9352) (1.2932,2.9287) (1.3083,2.9220) (1.3233,2.9149) (1.3383,2.9077) (1.3534,2.9002) (1.3684,2.8926) (1.3835,2.8847) (1.3985,2.8766) (1.4135,2.8683) (1.4286,2.8598) (1.4436,2.8512) (1.4586,2.8423) (1.4737,2.8333) (1.4887,2.8242) (1.5038,2.8149) (1.5188,2.8054) (1.5338,2.7958) (1.5489,2.7861) (1.5639,2.7762) (1.5789,2.7662) (1.5940,2.7561) (1.6090,2.7458) (1.6241,2.7355) (1.6391,2.7250) (1.6541,2.7145) (1.6692,2.7038) (1.6842,2.6931) (1.6992,2.6823) (1.7143,2.6714) (1.7293,2.6604) (1.7444,2.6493) (1.7594,2.6382) (1.7744,2.6271) (1.7895,2.6158) (1.8045,2.6045) (1.8195,2.5932) (1.8346,2.5818) (1.8496,2.5704) (1.8647,2.5590) (1.8797,2.5475) (1.8947,2.5359) (1.9098,2.5244) (1.9248,2.5128) (1.9398,2.5013) (1.9549,2.4897) (1.9699,2.4781) (1.9850,2.4664) (2.0000,2.4548) (2.0150,2.4432) (2.0301,2.4316) (2.0451,2.4200) (2.0602,2.4084) (2.0752,2.3968) (2.0902,2.3852) (2.1053,2.3737) (2.1203,2.3622) (2.1353,2.3507) (2.1504,2.3392) (2.1654,2.3278) (2.1805,2.3163) (2.1955,2.3050) (2.2105,2.2937) (2.2256,2.2824) (2.2406,2.2711) (2.2556,2.2599) (2.2707,2.2488) (2.2857,2.2377) (2.3008,2.2267) (2.3158,2.2157) (2.3308,2.2048) (2.3459,2.1940) (2.3609,2.1832) (2.3759,2.1725) (2.3910,2.1618) (2.4060,2.1513) (2.4211,2.1408) (2.4361,2.1304) (2.4511,2.1201) (2.4662,2.1098) (2.4812,2.0997) (2.4962,2.0896) (2.5113,2.0796) (2.5263,2.0697) (2.5414,2.0599) (2.5564,2.0502) (2.5714,2.0406) (2.5865,2.0311) (2.6015,2.0217) (2.6165,2.0125) (2.6316,2.0033) (2.6466,1.9942) (2.6617,1.9852) (2.6767,1.9764) (2.6917,1.9676) (2.7068,1.9590) (2.7218,1.9505) (2.7368,1.9421) (2.7519,1.9339) (2.7669,1.9257) (2.7820,1.9177) (2.7970,1.9098) (2.8120,1.9021) (2.8271,1.8944) (2.8421,1.8869) (2.8571,1.8796) (2.8722,1.8724) (2.8872,1.8653) (2.9023,1.8583) (2.9173,1.8515) (2.9323,1.8448) (2.9474,1.8383) (2.9624,1.8319) (2.9774,1.8257) (2.9925,1.8196) (3.0075,1.8137) (3.0226,1.8079) (3.0376,1.8023) (3.0526,1.7968) (3.0677,1.7915) (3.0827,1.7864) (3.0977,1.7814) (3.1128,1.7765) (3.1278,1.7718) (3.1429,1.7673) (3.1579,1.7630) (3.1729,1.7588) (3.1880,1.7548) (3.2030,1.7509) (3.2180,1.7472) (3.2331,1.7437) (3.2481,1.7404) (3.2632,1.7372) (3.2782,1.7342) (3.2932,1.7314) (3.3083,1.7288) (3.3233,1.7263) (3.3383,1.7241) (3.3534,1.7220) (3.3684,1.7201) (3.3835,1.7183) (3.3985,1.7168) (3.4135,1.7154) (3.4286,1.7143) (3.4436,1.7133) (3.4586,1.7125) (3.4737,1.7119) (3.4887,1.7115) (3.5038,1.7113) (3.5188,1.7113) (3.5338,1.7114) (3.5489,1.7118) (3.5639,1.7124) (3.5789,1.7131) (3.5940,1.7141) (3.6090,1.7153) (3.6241,1.7166) (3.6391,1.7182) (3.6541,1.7200) (3.6692,1.7220) (3.6842,1.7242) (3.6992,1.7265) (3.7143,1.7291) (3.7293,1.7320) (3.7444,1.7350) (3.7594,1.7382) (3.7744,1.7416) (3.7895,1.7453) (3.8045,1.7492) (3.8195,1.7532) (3.8346,1.7575) (3.8496,1.7621) (3.8647,1.7668) (3.8797,1.7717) (3.8947,1.7769) (3.9098,1.7823) (3.9248,1.7879) (3.9398,1.7937) (3.9549,1.7998) (3.9699,1.8061) (3.9850,1.8126) (4.0000,1.8193) (4.0150,1.8262) (4.0301,1.8334) (4.0451,1.8408) (4.0602,1.8485) (4.0752,1.8563) (4.0902,1.8644) (4.1053,1.8728) (4.1203,1.8813) (4.1353,1.8901) (4.1504,1.8991) (4.1654,1.9084) (4.1805,1.9179) (4.1955,1.9276) (4.2105,1.9376) (4.2256,1.9478) (4.2406,1.9582) (4.2556,1.9689) (4.2707,1.9798) (4.2857,1.9910) (4.3008,2.0024) (4.3158,2.0140) (4.3308,2.0259) (4.3459,2.0381) (4.3609,2.0504) (4.3759,2.0631) (4.3910,2.0759) (4.4060,2.0890) (4.4211,2.1024) (4.4361,2.1160) (4.4511,2.1299) (4.4662,2.1440) (4.4812,2.1583) (4.4962,2.1729) (4.5113,2.1878) (4.5263,2.2029) (4.5414,2.2182) (4.5564,2.2339) (4.5714,2.2497) (4.5865,2.2658) (4.6015,2.2822) (4.6165,2.2988) (4.6316,2.3157) (4.6466,2.3329) (4.6617,2.3503) (4.6767,2.3679) (4.6917,2.3859) (4.7068,2.4040) (4.7218,2.4225) (4.7368,2.4412) (4.7519,2.4601) (4.7669,2.4794) (4.7820,2.4988) (4.7970,2.5186) (4.8120,2.5386) (4.8271,2.5589) (4.8421,2.5794) (4.8571,2.6002) (4.8722,2.6213) (4.8872,2.6426) (4.9023,2.6642) (4.9173,2.6861) (4.9323,2.7083) (4.9474,2.7307) (4.9624,2.7534) (4.9774,2.7763) (4.9925,2.7995) (5.0075,2.8230) (5.0226,2.8468) (5.0376,2.8708) (5.0526,2.8951) (5.0677,2.9197) (5.0827,2.9446) (5.0977,2.9697) (5.1128,2.9951) (5.1278,3.0208) (5.1429,3.0467) (5.1579,3.0730) (5.1729,3.0995) (5.1880,3.1263) (5.2030,3.1533) (5.2180,3.1807) (5.2331,3.2083) (5.2481,3.2362) (5.2632,3.2644) (5.2782,3.2928) (5.2932,3.3216) (5.3083,3.3506) (5.3233,3.3799) (5.3383,3.4095) (5.3534,3.4394) (5.3684,3.4695) (5.3835,3.5000) (5.3985,3.5307) (5.4135,3.5617) (5.4286,3.5930) (5.4436,3.6246) (5.4586,3.6564) (5.4737,3.6886) (5.4887,3.7210) (5.5038,3.7538) (5.5188,3.7868) (5.5338,3.8201) (5.5489,3.8537) (5.5639,3.8875) (5.5789,3.9217) (5.5940,3.9562) (5.6090,3.9909) (5.6241,4.0260) (5.6391,4.0613) (5.6541,4.0969) (5.6692,4.1328) (5.6842,4.1690) (5.6992,4.2055) (5.7143,4.2423) (5.7293,4.2794) (5.7444,4.3168) (5.7594,4.3545) (5.7744,4.3925) (5.7895,4.4307) (5.8045,4.4693) (5.8195,4.5082) (5.8346,4.5473) (5.8496,4.5868) (5.8647,4.6265) (5.8797,4.6666) (5.8947,4.7069) (5.9098,4.7476) (5.9248,4.7885) (5.9398,4.8298) (5.9549,4.8713) (5.9699,4.9132) (5.9850,4.9553) (6.0000,4.9978)};
\fill (1,3) circle (2pt); \fill (2,2.4548) circle (2pt) node[above right]{$A$};
\fill (3.5129,1.7112) circle (2pt);
\draw[dashed,gray] (3.5129,0)--(3.5129,1.7112);
\node[below,font=\scriptsize] at (3.5129,0){$\alpha$};
\Ticks{1,2,3,4,5}{1,2,3,4,5}
\node[cDef] at (5.3,4.3){$(C_f)$};
\node[cThm] at (5.3,0.6){$(T)$};
\end{GR}
\end{hal}

\section{تمارين البكالوريا -- شعبة تقني رياضي}
% graph frame with different units on the axes: {xmin}{xmax}{ymin}{ymax}{x unit}{y unit}{y grid step}
\newenvironment{GRtm}[7]{\par\begin{center}\begin{LTR}\begin{tikzpicture}[x=#5cm,y=#6cm,font=\footnotesize]
 \draw[gray!25,very thin,xstep=1,ystep=#7] (#1,#3) grid (#2,#4);
 \draw[-{Stealth[length=2mm]}] (#1,0)--(#2+0.2,0) node[below left]{$x$};
 \draw[-{Stealth[length=2mm]}] (0,#3)--(0,#4+0.1) node[below left]{$y$};
 \node[below left] at (0,0){$O$};
 \begin{scope}\clip (#1,#3) rectangle (#2,#4);}
{\end{scope}\end{tikzpicture}\end{LTR}\end{center}}

%==================== التمرين 1 ====================
\begin{tamrin}{بكالوريا 2009 -- الموضوع الثاني -- تقني رياضي}
\textbf{I)} $g$ دالة معرفة على المجال $[1;+\infty[$ بـ: $g(x)=2x+\ln x$.
\QQ{1} أحسب نهاية الدالة $g$ عندما يؤول $x$ إلى $+\infty$.
\QQ{2} أدرس اتجاه تغير الدالة $g$.
\QQ{3} بيّن أنه من أجل كل عدد حقيقي $x$ من المجال $[1;+\infty[$ فإن: $g(x)\neq 0$.

\textbf{II)} نعتبر الدالة $f$ المعرفة على المجال $[1;+\infty[$ كما يلي: $f(x)=\dfrac{6\ln x}{2x+\ln x}$.

$(C_f)$ تمثيلها البياني في المستوي المنسوب إلى المعلم المتعامد والمتجانس $(O;\vec{i},\vec{j})$.
\QQ{1}
\begin{sq}
\qa بيّن أنه يمكن كتابة $f(x)$ على الشكل $f(x)=\dfrac{6\,\frac{\ln x}{x}}{2+\frac{\ln x}{x}}$ من أجل كل $x\in[1;+\infty[$.
\qb أحسب $\displaystyle\lim_{x\to+\infty}f(x)$. ماذا تستنتج؟
\end{sq}
\QQ{2} أدرس اتجاه تغير الدالة $f$، ثم شكّل جدول تغيراتها.
\QQ{3} ما هي قيم العدد الحقيقي $k$ بحيث تقبل المعادلة $f(x)=k$ حلين متمايزين؟
\QQ{4} أكتب معادلة لـ $(\Delta_1)$ مماس المنحنى $(C_f)$ في النقطة التي فاصلتها $1$.
\QQ{5} لتكن الدالة $h$ المعرفة على المجال $[1;+\infty[$ بـ: $h(x)=f\left(e^x\right)$، و$(C_h)$ تمثيلها البياني في المعلم السابق.
\begin{sq}
\qa شكّل جدول تغيرات الدالة $h$.
\qb أكتب معادلة لـ $(\Delta_2)$ مماس المنحنى $(C_h)$ في النقطة التي فاصلتها $1$.
\end{sq}
\QQ{6} أرسم كلا من: $(\Delta_1)$، $(\Delta_2)$، $(C_f)$ و$(C_h)$.
\end{tamrin}

\begin{hal}
\textbf{الجزء الأول:}
\QQ{1} لدينا $\displaystyle\lim_{x\to+\infty}2x=+\infty$ و$\displaystyle\lim_{x\to+\infty}\ln x=+\infty$، إذن بالجمع: \bx{\lim_{x\to+\infty}g(x)=+\infty}.
\QQ{2} الدالة $g$ قابلة للاشتقاق على $[1;+\infty[$ و: $g'(x)=2+\dfrac{1}{x}$.

بما أن $x>0$ فإن $\dfrac1x>0$، إذن $g'(x)>0$، ومنه \bxt{$g$ متزايدة تماما على $[1;+\infty[$}.
\begin{TAB}
\tkzTabInit[lgt=2,espcl=4]{$x$ /0.8, $g'(x)$ /0.8, $g(x)$ /1.6}{$1$, $+\infty$}
\tkzTabLine{, +, }
\tkzTabVar{-/$2$, +/$+\infty$}
\end{TAB}
\QQ{3} $g$ متزايدة تماما على $[1;+\infty[$، إذن من أجل كل $x\geqslant 1$: $g(x)\geqslant g(1)=2$، أي $g(x)>0$، ومنه \bx{g(x)\neq 0}.

\textbf{الجزء الثاني:}
\QQ{1}\pa من أجل $x\in[1;+\infty[$ لدينا $x>0$، نقسم البسط والمقام على $x$:
\[f(x)=\frac{6\ln x}{2x+\ln x}=\frac{\frac{6\ln x}{x}}{\frac{2x}{x}+\frac{\ln x}{x}}=\frac{6\,\frac{\ln x}{x}}{2+\frac{\ln x}{x}}.\]
\pb نعلم أن $\displaystyle\lim_{x\to+\infty}\frac{\ln x}{x}=0$، إذن البسط يؤول إلى $0$ والمقام إلى $2$، ومنه \bx{\lim_{x\to+\infty}f(x)=0}.

نستنتج أن $(C_f)$ يقبل مستقيما مقاربا أفقيا معادلته $y=0$ (حامل محور الفواصل) بجوار $+\infty$.
\QQ{2} $f$ قابلة للاشتقاق على $[1;+\infty[$ (لأن $2x+\ln x=g(x)>0$)، ومن أجل كل $x\geqslant1$:
\begin{align*}
f'(x)&=6\cdot\frac{\frac1x(2x+\ln x)-\ln x\left(2+\frac1x\right)}{(2x+\ln x)^2}
=6\cdot\frac{2+\frac{\ln x}{x}-2\ln x-\frac{\ln x}{x}}{(2x+\ln x)^2}=\frac{12(1-\ln x)}{(2x+\ln x)^2}.
\end{align*}
المقام موجب تماما، إذن إشارة $f'(x)$ من إشارة $1-\ln x$:
\begin{itemize}
\item $1-\ln x>0 \iff x<e$، إذن $f$ متزايدة تماما على $[1;e]$.
\item $1-\ln x<0 \iff x>e$، إذن $f$ متناقصة تماما على $[e;+\infty[$.
\end{itemize}
$f(1)=\dfrac{0}{2}=0$ و$f(e)=\dfrac{6}{2e+1}\approx 0{,}93$.
\begin{TAB}
\tkzTabInit[lgt=2,espcl=3.5]{$x$ /0.8, $f'(x)$ /0.8, $f(x)$ /1.8}{$1$, $e$, $+\infty$}
\tkzTabLine{, +, z, -, }
\tkzTabVar{-/$0$, +/$\frac{6}{2e+1}$, -/$0$}
\end{TAB}
\QQ{3} حلول المعادلة $f(x)=k$ هي فواصل نقط تقاطع $(C_f)$ مع المستقيم ذي المعادلة $y=k$.
\begin{itemize}
\item على $[1;e]$: $f$ مستمرة ومتزايدة تماما من $0$ إلى $\frac{6}{2e+1}$، إذن كل $k\in\left]0;\frac{6}{2e+1}\right[$ له سابقة وحيدة في $]1;e[$.
\item على $[e;+\infty[$: $f$ مستمرة ومتناقصة تماما من $\frac{6}{2e+1}$ نحو $0$ (دون أن تبلغه)، إذن كل $k\in\left]0;\frac{6}{2e+1}\right[$ له سابقة وحيدة في $]e;+\infty[$.
\item من أجل $k=0$ أو $k=\frac{6}{2e+1}$ يوجد حل وحيد، وفي غير ذلك لا يوجد حل.
\end{itemize}
إذن تقبل المعادلة حلين متمايزين إذا وفقط إذا: \bx{k\in\left]0;\frac{6}{2e+1}\right[}.
\QQ{4} $f(1)=0$ و$f'(1)=\dfrac{12(1-0)}{(2+0)^2}=3$، إذن $(\Delta_1): y=f'(1)(x-1)+f(1)$، أي \bx{(\Delta_1):\ y=3x-3}.
\QQ{5}\pa $h(x)=f(e^x)$؛ من أجل $x\geqslant1$ لدينا $e^x\geqslant e>1$، إذن $h$ معرفة وقابلة للاشتقاق على $[1;+\infty[$ و:
\[h'(x)=\left(e^x\right)'\cdot f'\left(e^x\right)=e^x f'\left(e^x\right).\]
بما أن $e^x>0$، فإشارة $h'(x)$ من إشارة $f'(e^x)$. ولما كان $e^x\geqslant e$ و$f'\leqslant0$ على $[e;+\infty[$ (تنعدم فقط عند $e$) فإن $h'(x)\leqslant 0$ وتنعدم فقط عند $x=1$؛ إذن $h$ متناقصة تماما على $[1;+\infty[$.

$h(1)=f(e)=\dfrac{6}{2e+1}$، و$\displaystyle\lim_{x\to+\infty}h(x)=\lim_{t\to+\infty}f(t)=0$ (بوضع $t=e^x$، $t\to+\infty$).
\begin{TAB}
\tkzTabInit[lgt=2,espcl=4]{$x$ /0.8, $h'(x)$ /0.8, $h(x)$ /1.8}{$1$, $+\infty$}
\tkzTabLine{z, -, }
\tkzTabVar{+/$\frac{6}{2e+1}$, -/$0$}
\end{TAB}
\pb $h'(1)=e\,f'(e)=e\times0=0$ و$h(1)=\dfrac{6}{2e+1}$، إذن \bx{(\Delta_2):\ y=\frac{6}{2e+1}} (مماس أفقي، وهو نفسه مماس $(C_f)$ عند النقطة ذات الفاصلة $e$).

ملاحظة: $h(x)=\dfrac{6x}{2e^x+x}$.
\QQ{6} الرسم: ($(C_f)$ بالأزرق، $(C_h)$ بالأحمر، $(\Delta_1)$ بالأخضر، $(\Delta_2)$ بالبرتقالي؛ سلم المحور العمودي مكبر).
\begin{GRtm}{-0.5}{9}{-0.6}{1.3}{1.5}{5}{0.2}
 \draw[cEx,thick] (-0.5,0.932)--(9,0.932) node[pos=0.9,above]{$(\Delta_2)$};
 \draw[cRem,thick,domain=0.85:1.45] plot(\x,{3*\x-3}) node[right]{$(\Delta_1)$};
 \draw[cDef,very thick,domain=1:9,samples=150] plot(\x,{6*ln(\x)/(2*\x+ln(\x))});
 \draw[cThm,very thick,domain=1:9,samples=150] plot(\x,{6*\x/(2*exp(\x)+\x)});
 \node[cDef] at (6.5,0.62){$(C_f)$};
 \node[cThm] at (2.6,0.4){$(C_h)$};
 \draw[dashed] (2.718,0)--(2.718,0.932);
 \node[below,font=\scriptsize] at (2.718,0){$e$};
 \foreach \t in {1,2,3,4,5,6,7,8} \node[below,font=\scriptsize] at (\t,0){$\t$};
 \foreach \t in {0.2,0.4,0.6,0.8,1} \node[left,font=\scriptsize] at (0,\t){$\t$};
\end{GRtm}
\end{hal}

%==================== التمرين 2 ====================
\begin{tamrin}{بكالوريا 2011 -- الموضوع الأول -- تقني رياضي}
$f$ دالة عددية معرفة على المجال $]0;+\infty[$ كما يلي: $f(x)=\dfrac{a+b\ln 2x}{4x^2}$، حيث $a$ و$b$ عددان حقيقيان.

$(C_f)$ تمثيلها البياني في المستوي المنسوب إلى المعلم المتعامد والمتجانس $(O;\vec{i},\vec{j})$.
\QQ{1} عيّن $a$ و$b$ بحيث يكون مماس المنحنى $(C_f)$ في النقطة $A\left(\frac12;1\right)$ موازيا لحامل محور الفواصل.
\QQ{2} $g$ الدالة العددية المعرفة على المجال $]0;+\infty[$ كما يلي: $g(x)=\dfrac{1+2\ln 2x}{4x^2}$، وليكن $(C_g)$ تمثيلها البياني في المستوي المنسوب إلى المعلم السابق.
\begin{sq}
\qa أحسب $\displaystyle\lim_{\substack{x\to0\\x>0}}g(x)$ و$\displaystyle\lim_{x\to+\infty}g(x)$. فسّر النتيجتين هندسيا.
\qb أدرس اتجاه تغير الدالة $g$، ثم شكّل جدول تغيراتها.
\qc حلّ في المجال $]0;+\infty[$ المعادلة $g(x)=0$.
\qd أنشئ $(C_g)$.
\end{sq}
\end{tamrin}

\begin{hal}
\QQ{1} الشرطان: $A\in(C_f)$ أي $f\left(\frac12\right)=1$، والمماس عند $A$ أفقي أي $f'\left(\frac12\right)=0$.
\begin{itemize}
\item $f\left(\frac12\right)=\dfrac{a+b\ln 1}{4\times\frac14}=a$، إذن $a=1$.
\item $f'(x)=\dfrac{\frac{b}{x}\cdot4x^2-8x(a+b\ln2x)}{16x^4}=\dfrac{4bx-8ax-8bx\ln 2x}{16x^4}=\dfrac{b-2a-2b\ln2x}{4x^3}$.

$f'\left(\frac12\right)=\dfrac{b-2a-0}{4\times\frac18}=2(b-2a)=0$، إذن $b=2a=2$.
\end{itemize}
ومنه \bx{a=1\ ;\ b=2}، أي $f(x)=\dfrac{1+2\ln2x}{4x^2}=g(x)$.
\QQ{2}\pa \textbf{عند $0$:} $\displaystyle\lim_{x\to0^+}\ln 2x=-\infty$ إذن البسط $1+2\ln2x\to-\infty$، و$4x^2\to0^+$، ومنه \bx{\lim_{x\to0^+}g(x)=-\infty}.

التفسير: المستقيم ذو المعادلة $x=0$ (حامل محور التراتيب) مقارب عمودي لـ $(C_g)$.

\textbf{عند $+\infty$:} نكتب $g(x)=\dfrac{1}{4x^2}+\dfrac{\ln 2x}{2x^2}=\dfrac{1}{4x^2}+\dfrac{\ln 2x}{2x}\cdot\dfrac{1}{x}$.

بما أن $\displaystyle\lim_{x\to+\infty}\frac{\ln 2x}{2x}=0$ و$\displaystyle\lim_{x\to+\infty}\frac1x=0$، فإن \bx{\lim_{x\to+\infty}g(x)=0}.

التفسير: المستقيم ذو المعادلة $y=0$ (حامل محور الفواصل) مقارب أفقي لـ $(C_g)$ بجوار $+\infty$.
\pb من السؤال 1 (مع $a=1$ و$b=2$): $g'(x)=\dfrac{2-2-4\ln 2x}{4x^3}=\dfrac{-\ln 2x}{x^3}$.

بما أن $x^3>0$، فإشارة $g'(x)$ من إشارة $-\ln 2x$:
\begin{itemize}
\item $\ln 2x<0 \iff 0<2x<1 \iff 0<x<\frac12$: $g'(x)>0$، إذن $g$ متزايدة تماما على $\left]0;\frac12\right]$.
\item $\ln 2x>0 \iff x>\frac12$: $g'(x)<0$، إذن $g$ متناقصة تماما على $\left[\frac12;+\infty\right[$.
\end{itemize}
$g\left(\frac12\right)=1$.
\begin{TAB}
\tkzTabInit[lgt=2,espcl=3.5]{$x$ /0.8, $g'(x)$ /0.8, $g(x)$ /1.6}{$0$, $\frac12$, $+\infty$}
\tkzTabLine{d, +, z, -, }
\tkzTabVar{D-/$-\infty$, +/$1$, -/$0$}
\end{TAB}
\pc $g(x)=0 \iff 1+2\ln 2x=0$ (لأن $4x^2\neq0$) $\iff \ln 2x=-\dfrac12 \iff 2x=e^{-\frac12} \iff x=\dfrac{1}{2\sqrt e}$.

ومنه مجموعة الحلول: \bx{S=\left\{\frac{1}{2\sqrt e}\right\}}، مع $\dfrac{1}{2\sqrt e}\approx0{,}30$.
\pd الإنشاء:
\begin{GR}{-0.3}{4}{-2.2}{1.5}{2.8}
 \draw[cDef,very thick,domain=0.16:4,samples=200] plot(\x,{(1+2*ln(2*\x))/(4*\x*\x)});
 \draw[dashed] (0.5,0)--(0.5,1)--(0,1);
 \fill[cThm] (0.5,1) circle (1.5pt) node[above right]{$A$};
 \fill[cThm] (0.303,0) circle (1.5pt);
 \node[cDef] at (2.5,0.35){$(C_g)$};
 \Ticks{1,2,3}{-2,-1,1}
 \node[below,font=\scriptsize] at (0.5,0){$\frac12$};
\end{GR}
\end{hal}

%==================== التمرين 3 ====================
\begin{tamrin}{بكالوريا 2012 -- الموضوع الثاني -- تقني رياضي}
\textbf{I)} $g$ دالة معرفة على المجال $]0;+\infty[$ كما يلي: $g(x)=x^2+a+b\ln x$، حيث $a$ و$b$ عددان حقيقيان.
\QQ{1} عيّن $a$ و$b$ علما أن التمثيل البياني للدالة $g$ يقبل في النقطة $A(1;-1)$ مماسا معامل توجيهه $4$.
\QQ{2} نضع: $a=-2$ و$b=2$.
\begin{sq}
\qa أدرس تغيرات الدالة $g$، ثم شكّل جدول تغيراتها.
\qb بيّن أن المعادلة $g(x)=0$ تقبل حلا وحيدا $\alpha$ على $]0;+\infty[$، ثم استنتج إشارة $g(x)$ على $]0;+\infty[$.
\end{sq}

\textbf{II)} $f$ هي الدالة المعرفة على المجال $]0;+\infty[$ كما يلي: $f(x)=x-2-\dfrac{2\ln x}{x}$.

$(C_f)$ تمثيلها البياني في المستوي المنسوب إلى المعلم المتعامد والمتجانس $(O;\vec{i},\vec{j})$ (وحدة الطول $2cm$).
\QQ{1}
\begin{sq}
\qa أحسب $\displaystyle\lim_{x\to+\infty}f(x)$ و$\displaystyle\lim_{\substack{x\to0\\x>0}}f(x)$.
\qb أحسب $f'(x)$، ثم تحقق أن: $f'(x)=\dfrac{g(x)}{x^2}$.
\qc استنتج إشارة $f'(x)$، ثم شكّل جدول تغيرات الدالة $f$.
\end{sq}
\QQ{2}
\begin{sq}
\qa بيّن أن المستقيم $(\Delta)$ ذا المعادلة $y=x-2$ مقارب لـ $(C_f)$، ثم أدرس وضعية $(C_f)$ بالنسبة لـ $(\Delta)$.
\qb بيّن أن $(C_f)$ يقبل مماسا $(T)$ يوازي $(\Delta)$، ثم جد معادلة له.
\qc نأخذ $\alpha=1{,}25$. بيّن أن المعادلة $f(x)=0$ تقبل حلين $x_1$ و$x_2$ حيث: $0{,}6<x_1<0{,}7$ و$2{,}7<x_2<2{,}8$، ثم أرسم كلا من $(\Delta)$، $(T)$ و$(C_f)$.
\end{sq}
\QQ{3} ناقش بيانيا، حسب قيم الوسيط الحقيقي $m$، عدد حلول المعادلة: $(m+2)x+2\ln x=0$.
\end{tamrin}

\begin{hal}
\textbf{الجزء الأول:}
\QQ{1} $g'(x)=2x+\dfrac{b}{x}$. الشرطان: $g(1)=-1$ و$g'(1)=4$.
\begin{itemize}
\item $g(1)=1+a+b\ln1=1+a=-1$، إذن $a=-2$.
\item $g'(1)=2+b=4$، إذن $b=2$.
\end{itemize}
ومنه \bx{a=-2\ ;\ b=2} و$g(x)=x^2-2+2\ln x$.
\QQ{2}\pa \textbf{النهايات:} $\displaystyle\lim_{x\to0^+}\ln x=-\infty$ و$x^2-2\to-2$، إذن $\displaystyle\lim_{x\to0^+}g(x)=-\infty$؛ و$\displaystyle\lim_{x\to+\infty}g(x)=+\infty$ (مجموع نهايتين $+\infty$).

\textbf{المشتقة:} $g'(x)=2x+\dfrac2x>0$ لأن $x>0$، إذن $g$ متزايدة تماما على $]0;+\infty[$.
\begin{TAB}
\tkzTabInit[lgt=2,espcl=3.5]{$x$ /0.8, $g'(x)$ /0.8, $g(x)$ /1.6}{$0$, $\alpha$, $+\infty$}
\tkzTabLine{d, +, t, +, }
\tkzTabVar{D-/$-\infty$, R/, +/$+\infty$}
\tkzTabVal{1}{3}{0.5}{}{$0$}
\end{TAB}
\pb $g$ مستمرة ومتزايدة تماما على $]0;+\infty[$ وتأخذ قيمها في $]-\infty;+\infty[$ الذي يشمل $0$، إذن حسب مبرهنة القيم المتوسطة (حالة الرتابة) المعادلة $g(x)=0$ تقبل حلا وحيدا $\alpha$.

تحقق: $g(1{,}24)\approx-0{,}03<0$ و$g(1{,}25)\approx0{,}01>0$، إذن $1{,}24<\alpha<1{,}25$ (أي $\alpha\approx1{,}25$).

\textbf{الإشارة:} \bxt{$g(x)<0$ على $]0;\alpha[ ; g(\alpha)=0 ; g(x)>0$ على $]\alpha;+\infty[$}.

\textbf{الجزء الثاني:}
\QQ{1}\pa \textbf{عند $+\infty$:} $\displaystyle\lim_{x\to+\infty}\frac{\ln x}{x}=0$ و$\displaystyle\lim_{x\to+\infty}(x-2)=+\infty$، إذن \bx{\lim_{x\to+\infty}f(x)=+\infty}.

\textbf{عند $0$:} $-\dfrac{2\ln x}{x}=-2\ln x\times\dfrac1x$، مع $-2\ln x\to+\infty$ و$\dfrac1x\to+\infty$، إذن \bx{\lim_{x\to0^+}f(x)=+\infty}، والمستقيم $x=0$ مقارب عمودي.
\pb $f'(x)=1-2\cdot\dfrac{\frac1x\cdot x-\ln x}{x^2}=1-\dfrac{2-2\ln x}{x^2}=\dfrac{x^2-2+2\ln x}{x^2}$، أي \bx{f'(x)=\frac{g(x)}{x^2}}.
\pc $x^2>0$، إذن إشارة $f'(x)$ هي إشارة $g(x)$: $f$ متناقصة تماما على $]0;\alpha]$ ومتزايدة تماما على $[\alpha;+\infty[$.

حساب $f(\alpha)$: من $g(\alpha)=0$ لدينا $2\ln\alpha=2-\alpha^2$، إذن $f(\alpha)=\alpha-2-\dfrac{2-\alpha^2}{\alpha}=\dfrac{2\alpha^2-2\alpha-2}{\alpha}$، ومن أجل $\alpha=1{,}25$: $f(\alpha)\approx-1{,}11$.
\begin{TAB}
\tkzTabInit[lgt=2,espcl=3.5]{$x$ /0.8, $f'(x)$ /0.8, $f(x)$ /1.8}{$0$, $\alpha$, $+\infty$}
\tkzTabLine{d, -, z, +, }
\tkzTabVar{D+/$+\infty$, -/$f(\alpha)$, +/$+\infty$}
\end{TAB}
\QQ{2}\pa لدينا:
\[f(x)-(x-2)=-\frac{2\ln x}{x}\qquad\text{و}\qquad\lim_{x\to+\infty}\left(-\frac{2\ln x}{x}\right)=0.\]
إذن المستقيم \bx{(\Delta):y=x-2} مقارب مائل لـ $(C_f)$ بجوار $+\infty$.

\textbf{الوضعية:} إشارة $f(x)-(x-2)$ من إشارة $-\ln x$ (لأن $x>0$):
\begin{itemize}
\item على $]0;1[$: $\ln x<0$ إذن $(C_f)$ فوق $(\Delta)$.
\item على $]1;+\infty[$: $\ln x>0$ إذن $(C_f)$ تحت $(\Delta)$.
\item عند $x=1$: $(C_f)$ يقطع $(\Delta)$ في النقطة $(1;-1)$.
\end{itemize}
\pb $(T)\parallel(\Delta)$ معناه $f'(x)=1$، أي $\dfrac{x^2-2+2\ln x}{x^2}=1 \iff 2\ln x-2=0 \iff x=e$.

$f(e)=e-2-\dfrac2e$، إذن $(T): y=1\cdot(x-e)+e-2-\dfrac2e$، أي \bx{(T):\ y=x-2-\frac{2}{e}}.
\pc $f$ مستمرة ومتناقصة تماما على $]0;\alpha]$ و$[0{,}6;0{,}7]\subset]0;\alpha]$، مع $f(0{,}6)\approx0{,}30>0$ و$f(0{,}7)\approx-0{,}28<0$، إذن يوجد حل وحيد $x_1$ في $]0;\alpha]$ و$0{,}6<x_1<0{,}7$.

$f$ مستمرة ومتزايدة تماما على $[\alpha;+\infty[$، مع $f(2{,}7)\approx-0{,}04<0$ و$f(2{,}8)\approx0{,}06>0$، إذن يوجد حل وحيد $x_2$ و$2{,}7<x_2<2{,}8$.

الرسم (وحدة الطول $2cm$):
\begin{GR}{-0.5}{6}{-2.5}{4.5}{1.75}
 \draw[cRem,thick,domain=-0.5:6] plot(\x,{\x-2});
 \node[cRem,above left] at (5.3,3.3){$(\Delta)$};
 \draw[cEx,thick,domain=-0.5:6] plot(\x,{\x-2-2/2.718});
 \node[cEx,below right] at (5,2.2){$(T)$};
 \draw[cDef,very thick,domain=0.33:6,samples=200] plot(\x,{\x-2-2*ln(\x)/\x});
 \fill (2.718,{2.718-2-2/2.718}) circle (1.5pt);
 \fill (1,-1) circle (1.5pt);
 \node[cDef] at (0.9,2.6){$(C_f)$};
 \Ticks{1,2,3,4,5}{-2,-1,1,2,3,4}
\end{GR}
\QQ{3} من أجل $x>0$:
\[(m+2)x+2\ln x=0 \iff m+2+\frac{2\ln x}{x}=0 \iff x-2-\frac{2\ln x}{x}=x+m \iff f(x)=x+m.\]
حلول المعادلة هي فواصل نقط تقاطع $(C_f)$ مع المستقيم $(D_m): y=x+m$ الموازي لـ $(\Delta)$ و$(T)$.
\begin{itemize}
\item $m<-2-\frac2e$: $(D_m)$ تحت $(T)$، لا توجد حلول.
\item $m=-2-\frac2e$: $(D_m)=(T)$، حل وحيد ($x=e$).
\item $-2-\frac2e<m<-2$: $(D_m)$ بين $(T)$ و$(\Delta)$، حلان متمايزان.
\item $m\geqslant-2$: $(D_m)$ على $(\Delta)$ أو فوقه، حل وحيد.
\end{itemize}
\end{hal}

%==================== التمرين 4 ====================
\begin{tamrin}{بكالوريا 2013 -- الموضوع الأول -- تقني رياضي}
\textbf{I)} الدالة $g$ معرفة على المجال $]-1;+\infty[$ كما يلي: $g(x)=(x+1)^2-2+\ln(x+1)$.
\QQ{1} أدرس اتجاه تغير الدالة $g$ على المجال $]-1;+\infty[$.
\QQ{2} بيّن أن المعادلة $g(x)=0$ تقبل حلا وحيدا $\alpha$ حيث $0{,}31<\alpha<0{,}32$، وأن: $\ln(\alpha+1)=2-(\alpha+1)^2$.
\QQ{3} استنتج حسب قيم $x$ إشارة $g(x)$.

\textbf{II)} الدالة $f$ معرفة على المجال $]-1;+\infty[$ كما يلي: $f(x)=(x+1)^2+\left(2-\ln(x+1)\right)^2$.

$(C_f)$ تمثيلها البياني في المستوي المنسوب إلى المعلم المتعامد والمتجانس $(O;\vec{i},\vec{j})$.
\QQ{1} أحسب $\displaystyle\lim_{x\to+\infty}f(x)$ و$\displaystyle\lim_{\substack{x\to-1\\x>-1}}f(x)$.
\QQ{2} أثبت أنه من أجل كل $x$ من $]-1;+\infty[$: $f'(x)=\dfrac{2g(x)}{x+1}$.
\QQ{3} أدرس اتجاه تغير الدالة $f$، ثم شكّل جدول تغيراتها.
\QQ{4} بيّن أن: $f(\alpha)=(\alpha+1)^2\left(1+(\alpha+1)^2\right)$، ثم استنتج حصرا للعدد $f(\alpha)$.
\QQ{5} مثّل المنحنى على المجال $]-1;2]$.

\textbf{III)} $(\Gamma)$ المنحنى الممثل للدالة $h$ المعرفة على المجال $]-1;+\infty[$ بـ: $h(x)=\ln(x+1)$.

$A$ النقطة ذات الإحداثيين $(-1;2)$ و$M$ نقطة من $(\Gamma)$ فاصلتها $x$.
\QQ{1} أثبت أن المسافة $AM$ تعطى بالعبارة: $AM=\sqrt{f(x)}$.
\QQ{2} الدالة $k$ معرفة على المجال $]-1;+\infty[$ بالعبارة: $k(x)=\sqrt{f(x)}$.
\begin{sq}
\qa بيّن أن للدالتين $k$ و$f$ نفس اتجاه التغير على المجال $]-1;+\infty[$.
\qb عيّن إحداثيي النقطة $B$ من $(\Gamma)$، بحيث تكون المسافة $AM$ أصغر ما يمكن.
\qc بيّن أن: $AB=(\alpha+1)\sqrt{1+(\alpha+1)^2}$.
\end{sq}
\end{tamrin}

\begin{hal}
\textbf{الجزء الأول:}
\QQ{1} $g$ قابلة للاشتقاق على $]-1;+\infty[$ و: $g'(x)=2(x+1)+\dfrac{1}{x+1}$.

بما أن $x+1>0$ فإن الحدين موجبان تماما، إذن $g'(x)>0$ و\bxt{$g$ متزايدة تماما على $]-1;+\infty[$}.

النهايتان: $\displaystyle\lim_{x\to-1^+}g(x)=-\infty$ (لأن $(x+1)^2\to0$ و$\ln(x+1)\to-\infty$)، و$\displaystyle\lim_{x\to+\infty}g(x)=+\infty$.
\begin{TAB}
\tkzTabInit[lgt=2,espcl=3.5]{$x$ /0.8, $g'(x)$ /0.8, $g(x)$ /1.6}{$-1$, $\alpha$, $+\infty$}
\tkzTabLine{d, +, t, +, }
\tkzTabVar{D-/$-\infty$, R/, +/$+\infty$}
\tkzTabVal{1}{3}{0.5}{}{$0$}
\end{TAB}
\QQ{2} $g$ مستمرة ومتزايدة تماما على $]-1;+\infty[$ وتأخذ قيمها في $\R$، إذن المعادلة $g(x)=0$ تقبل حلا وحيدا $\alpha$.

$g(0{,}31)\approx-0{,}014<0$ و$g(0{,}32)\approx0{,}020>0$، إذن \bx{0{,}31<\alpha<0{,}32}.

ولدينا $g(\alpha)=0$ أي $(\alpha+1)^2-2+\ln(\alpha+1)=0$، ومنه \bx{\ln(\alpha+1)=2-(\alpha+1)^2}.
\QQ{3} $g$ متزايدة تماما و$g(\alpha)=0$، إذن: $g(x)<0$ على $]-1;\alpha[$، $g(x)>0$ على $]\alpha;+\infty[$.

\textbf{الجزء الثاني:}
\QQ{1} \textbf{عند $+\infty$:} $(x+1)^2\to+\infty$ و$\left(2-\ln(x+1)\right)^2\geqslant0$، إذن \bx{\lim_{x\to+\infty}f(x)=+\infty}.

\textbf{عند $-1$:} $\ln(x+1)\to-\infty$ إذن $2-\ln(x+1)\to+\infty$ ومربعه $\to+\infty$، و$(x+1)^2\to0$، إذن \bx{\lim_{x\to-1^+}f(x)=+\infty}؛ المستقيم $x=-1$ مقارب عمودي.
\QQ{2} $f$ قابلة للاشتقاق على $]-1;+\infty[$ و:
\begin{align*}
f'(x)&=2(x+1)+2\left(2-\ln(x+1)\right)\cdot\left(-\frac{1}{x+1}\right)\\
&=\frac{2(x+1)^2-4+2\ln(x+1)}{x+1}=\frac{2\left[(x+1)^2-2+\ln(x+1)\right]}{x+1}.
\end{align*}
إذن \bx{f'(x)=\frac{2g(x)}{x+1}}.
\QQ{3} $x+1>0$، إذن إشارة $f'(x)$ هي إشارة $g(x)$: $f$ متناقصة تماما على $]-1;\alpha]$ ومتزايدة تماما على $[\alpha;+\infty[$.
\begin{TAB}
\tkzTabInit[lgt=2,espcl=3.5]{$x$ /0.8, $f'(x)$ /0.8, $f(x)$ /1.8}{$-1$, $\alpha$, $+\infty$}
\tkzTabLine{d, -, z, +, }
\tkzTabVar{D+/$+\infty$, -/$f(\alpha)$, +/$+\infty$}
\end{TAB}
\QQ{4} من السؤال I-2: $2-\ln(\alpha+1)=(\alpha+1)^2$، إذن:
\[f(\alpha)=(\alpha+1)^2+\left((\alpha+1)^2\right)^2=(\alpha+1)^2+(\alpha+1)^4=(\alpha+1)^2\left(1+(\alpha+1)^2\right).\]
\textbf{الحصر:} الدالة $u:t\mapsto(t+1)^2\left(1+(t+1)^2\right)$ متزايدة تماما على $[0{,}31;0{,}32]$ (جداء دالتين موجبتين متزايدتين)، و$0{,}31<\alpha<0{,}32$، إذن $u(0{,}31)<f(\alpha)<u(0{,}32)$:
\[1{,}31^2\left(1+1{,}31^2\right)<f(\alpha)<1{,}32^2\left(1+1{,}32^2\right),\quad\text{أي}\quad \bx{4{,}66<f(\alpha)<4{,}78}.\]
\QQ{5} التمثيل على $]-1;2]$ (سلم المحور العمودي مصغر). من القيم: $f(-0{,}5)\approx7{,}50$، $f(0)=5$، $f(1)\approx5{,}71$، $f(2)\approx9{,}81$.
\begin{GRtm}{-1.2}{2.3}{-0.5}{13}{3}{0.5}{1}
 \draw[dashed,cThm] (-1,-0.5)--(-1,13);
 \draw[cDef,very thick,domain=-0.93:2,samples=200] plot(\x,{(\x+1)^2+(2-ln(\x+1))^2});
 \fill[cDef] (2,9.8125) circle (1.8pt);
 \fill (0.314,4.709) circle (1.5pt);
 \draw[dashed] (0.314,0)--(0.314,4.709);
 \node[below,font=\scriptsize] at (0.314,0){$\alpha$};
 \node[cDef] at (1.3,8.5){$(C_f)$};
 \foreach \t in {-1,1,2} \node[below,font=\scriptsize] at (\t,0){$\t$};
 \foreach \t in {2,4,6,8,10,12} \node[left,font=\scriptsize] at (0,\t){$\t$};
\end{GRtm}

\textbf{الجزء الثالث:}
\QQ{1} $A(-1;2)$ و$M\left(x;\ln(x+1)\right)$، إذن:
\[AM^2=(x+1)^2+\left(\ln(x+1)-2\right)^2=(x+1)^2+\left(2-\ln(x+1)\right)^2=f(x).\]
وبما أن $AM\geqslant0$: \bx{AM=\sqrt{f(x)}}.
\QQ{2}\pa من جدول التغيرات $f(x)\geqslant f(\alpha)>0$، إذن $k$ قابلة للاشتقاق على $]-1;+\infty[$ و $k'(x)=\dfrac{f'(x)}{2\sqrt{f(x)}}$.

بما أن $2\sqrt{f(x)}>0$، فإن $k'(x)$ و$f'(x)$ لهما نفس الإشارة، إذن \bxt{$k$ و $f$ لهما نفس اتجاه التغير}.
\pb $AM=k(x)$ أصغر ما يمكن عندما تبلغ $k$ قيمتها الحدية الصغرى، أي عند $x=\alpha$ (نفس تغيرات $f$). إذن $B\left(\alpha;\ln(\alpha+1)\right)$، أي \bx{B\left(\alpha\ ;\ 2-(\alpha+1)^2\right)}.
\pc $AB=k(\alpha)=\sqrt{f(\alpha)}=\sqrt{(\alpha+1)^2\left(1+(\alpha+1)^2\right)}=|\alpha+1|\sqrt{1+(\alpha+1)^2}$، وبما أن $\alpha+1>0$: \bx{AB=(\alpha+1)\sqrt{1+(\alpha+1)^2}}.
\end{hal}

%==================== التمرين 5 ====================
\begin{tamrin}{بكالوريا 2014 -- الموضوع الثاني -- تقني رياضي}
المستوي منسوب إلى المعلم المتعامد والمتجانس $(O;\vec{i},\vec{j})$.

\textbf{I)} $g$ الدالة المعرفة على المجال $]0;3]$ بـ: $g(x)=x\ln x+x$.
\QQ{1} أدرس تغيرات الدالة $g$.
\QQ{2}
\begin{sq}
\qa بيّن أن المعادلة $g(x)=2$ تقبل حلا وحيدا $\alpha$ في $]0;3]$، ثم تحقق أن: $1{,}45<\alpha<1{,}46$.
\qb استنتج إشارة $g(x)-2$.
\end{sq}

\textbf{II)} التمثيل المقابل $(C_f)$ هو للدالة $f$ المعرفة على المجال $]0;3]$ بـ: $f(x)=|x-2|\ln x$.
\begin{GR}{-0.2}{3.3}{-2}{1.4}{2}
 \draw[cDef,very thick,domain=0.2:2,samples=150] plot(\x,{(2-\x)*ln(\x)});
 \draw[cDef,very thick,domain=2:3,samples=80] plot(\x,{(\x-2)*ln(\x)});
 \draw[thick,-{Stealth[length=2mm]}] (2,0)--(3,0.693);
 \draw[thick,-{Stealth[length=2mm]}] (2,0)--(1,0.693);
 \node[cDef] at (0.65,-1.5){$(C_f)$};
 \Ticks{1,2,3}{-1,1}
\end{GR}
\QQ{1} باستعمال $(C_f)$ ضع تخمينا حول قابلية اشتقاق الدالة $f$ عند $2$.
\QQ{2} أثبت صحة تخمينك.
\QQ{3} أدرس تغيرات الدالة $f$.

\textbf{III)} $h$ الدالة المعرفة على المجال $\left[0;\frac{\pi}{2}\right[$ كما يلي: $h(x)=(2-\cos x)\ln(\cos x)$.
\QQ{1} بيّن أن المستقيم $(\Delta)$ ذا المعادلة $x=\frac{\pi}{2}$ مقارب للمنحنى $(C_h)$، حيث $(C_h)$ هو المنحنى الممثل للدالة $h$.
\QQ{2} أدرس اتجاه تغير الدالة $h$، ثم شكّل جدول تغيراتها وأرسم $(\Delta)$ و$(C_h)$.
\end{tamrin}

\begin{hal}
\textbf{الجزء الأول:}
\QQ{1} $g$ قابلة للاشتقاق على $]0;3]$ و: $g'(x)=\ln x+x\cdot\dfrac1x+1=\ln x+2$.

$g'(x)=0\iff \ln x=-2\iff x=e^{-2}$؛ $g'(x)<0$ على $]0;e^{-2}[$ و$g'(x)>0$ على $]e^{-2};3]$.

إذن $g$ متناقصة تماما على $]0;e^{-2}]$ ومتزايدة تماما على $[e^{-2};3]$.

النهاية: $\displaystyle\lim_{x\to0^+}x\ln x=0$ إذن $\displaystyle\lim_{x\to0^+}g(x)=0$. القيم: $g(e^{-2})=-2e^{-2}+e^{-2}=-e^{-2}$، $g(3)=3\ln3+3\approx6{,}30$.
\begin{TAB}
\tkzTabInit[lgt=2,espcl=3.5]{$x$ /0.8, $g'(x)$ /0.8, $g(x)$ /1.6}{$0$, $e^{-2}$, $3$}
\tkzTabLine{d, -, z, +, }
\tkzTabVar{D+/$0$, -/$-e^{-2}$, +/$3+3\ln3$}
\end{TAB}
\QQ{2}\pa \textbf{على $]0;e^{-2}]$:} $g$ متناقصة تماما و$\displaystyle\lim_{x\to0^+}g(x)=0$، إذن $g(x)<0<2$، فالمعادلة $g(x)=2$ لا تقبل حلولا في هذا المجال.

\textbf{على $[e^{-2};3]$:} $g$ مستمرة ومتزايدة تماما، و$2\in\left[-e^{-2};3+3\ln3\right]$، إذن (مبرهنة القيم المتوسطة) المعادلة $g(x)=2$ تقبل حلا وحيدا $\alpha$.

إذن المعادلة تقبل حلا وحيدا $\alpha$ في $]0;3]$. ولدينا $g(1{,}45)\approx1{,}989<2$ و$g(1{,}46)\approx2{,}013>2$، إذن \bx{1{,}45<\alpha<1{,}46}.
\pb من الدراسة السابقة: \bxt{$g(x)-2<0$ على $]0;\alpha[ ; g(x)-2>0$ على $]\alpha;3]$} و$g(\alpha)-2=0$.

\textbf{الجزء الثاني:}
\QQ{1} \textbf{التخمين:} يظهر أن $(C_f)$ يقبل عند النقطة $(2;0)$ نصفي مماسين مختلفين (نقطة زاوية)، إذن نخمن أن $f$ غير قابلة للاشتقاق عند $2$.
\QQ{2} $f(2)=0$، و$\dfrac{f(x)-f(2)}{x-2}=\dfrac{|x-2|\ln x}{x-2}$:
\begin{itemize}
\item من أجل $x>2$: $|x-2|=x-2$، إذن $\displaystyle\lim_{x\to2^+}\frac{f(x)-f(2)}{x-2}=\lim_{x\to2^+}\ln x=\ln2$.
\item من أجل $x<2$: $|x-2|=-(x-2)$، إذن $\displaystyle\lim_{x\to2^-}\frac{f(x)-f(2)}{x-2}=\lim_{x\to2^-}(-\ln x)=-\ln2$.
\end{itemize}
بما أن $\ln 2\neq-\ln2$، فإن \bxt{$f$ غير قابلة للاشتقاق عند $2$}، والتخمين صحيح. ($(C_f)$ يقبل عند $(2;0)$ نصفي مماسين معاملا توجيههما $\ln2$ و$-\ln2$.)
\QQ{3} \textbf{على $]0;2]$:} $f(x)=(2-x)\ln x$، و$f'(x)=-\ln x+\dfrac{2-x}{x}=\dfrac{2-x-x\ln x}{x}=\dfrac{2-g(x)}{x}$.

\textbf{على $[2;3]$:} $f(x)=(x-2)\ln x$، و$f'(x)=\ln x+\dfrac{x-2}{x}=\dfrac{x\ln x+x-2}{x}=\dfrac{g(x)-2}{x}$.

بما أن $x>0$ ومن إشارة $g(x)-2$ (مع $\alpha<2$):
\begin{itemize}
\item على $]0;\alpha[$: $g(x)-2<0$ إذن $f'(x)>0$: $f$ متزايدة تماما.
\item على $]\alpha;2[$: $g(x)-2>0$ إذن $f'(x)<0$: $f$ متناقصة تماما.
\item على $]2;3]$: $g(x)-2>0$ إذن $f'(x)>0$: $f$ متزايدة تماما.
\end{itemize}
النهاية: $\displaystyle\lim_{x\to0^+}(2-x)\ln x=-\infty$. ومن $g(\alpha)=2$: $\ln\alpha=\dfrac{2}{\alpha}-1=\dfrac{2-\alpha}{\alpha}$، إذن $f(\alpha)=\dfrac{(2-\alpha)^2}{\alpha}\approx0{,}20$؛ $f(3)=\ln3$.
\begin{TAB}
\tkzTabInit[lgt=2,espcl=2.8]{$x$ /0.8, $f'(x)$ /0.8, $f(x)$ /1.8}{$0$, $\alpha$, $2$, $3$}
\tkzTabLine{d, +, z, -, d, +, }
\tkzTabVar{D-/$-\infty$, +/$f(\alpha)$, -/$0$, +/$\ln3$}
\end{TAB}
(الخطان المزدوجان عند $2$ في سطر $f'(x)$ يعنيان أن $f'(2)$ غير موجود، والدالة $f$ مستمرة عند $2$.)

\textbf{الجزء الثالث:}
\QQ{1} لما $x\to\frac{\pi}{2}^-$: $\cos x\to0^+$، إذن $\ln(\cos x)\to-\infty$ و$2-\cos x\to2$، ومنه \bx{\lim_{x\to\frac{\pi}{2}^-}h(x)=-\infty}.

إذن المستقيم $(\Delta): x=\frac{\pi}{2}$ مقارب عمودي لـ $(C_h)$.
\QQ{2} $h$ قابلة للاشتقاق على $\left[0;\frac{\pi}{2}\right[$ (لأن $\cos x>0$) و:
\[h'(x)=\sin x\ln(\cos x)+(2-\cos x)\cdot\frac{-\sin x}{\cos x}=\frac{\sin x}{\cos x}\left[\cos x\ln(\cos x)+\cos x-2\right].\]
نلاحظ أن $\cos x\ln(\cos x)+\cos x=g(\cos x)$، إذن \bx{h'(x)=\tan x\left[g(\cos x)-2\right]}.

من أجل $x\in\left]0;\frac{\pi}{2}\right[$: $\tan x>0$، و$\cos x\in]0;1[\subset]0;\alpha[$ (لأن $\alpha>1$)، إذن $g(\cos x)-2<0$، ومنه $h'(x)<0$؛ و$h'(0)=0$.

إذن $h$ متناقصة تماما على $\left[0;\frac{\pi}{2}\right[$، مع $h(0)=(2-1)\ln1=0$.
\begin{TAB}
\tkzTabInit[lgt=2,espcl=4]{$x$ /0.8, $h'(x)$ /0.8, $h(x)$ /1.6}{$0$, $\frac{\pi}{2}$}
\tkzTabLine{z, -, d}
\tkzTabVar{+/$0$, -D/$-\infty$}
\end{TAB}
الرسم (قيم مساعدة: $h(0{,}5)\approx-0{,}15$، $h(1)\approx-0{,}90$، $h(1{,}3)\approx-2{,}28$، $h(1{,}5)\approx-5{,}11$):
\begin{GR}{-0.3}{2}{-5}{0.6}{1.5}
 \draw[cRem,thick,dashed] (1.5708,-5)--(1.5708,0.6);
 \draw[cDef,very thick,domain=0:1.5,samples=150] plot(\x,{(2-cos(deg(\x)))*ln(cos(deg(\x)))});
 \node[cRem,right] at (1.5708,-1){$(\Delta)$};
 \node[cDef] at (0.7,-1.5){$(C_h)$};
 \node[below right,font=\scriptsize] at (1.5708,0){$\frac{\pi}{2}$};
 \Ticks{1}{-4,-3,-2,-1}
\end{GR}
\end{hal}

%==================== التمرين 6 ====================
\begin{tamrin}{بكالوريا 2015 -- الموضوع الأول -- تقني رياضي}
\textbf{I)} $h$ الدالة المعرفة على المجال $]-2;+\infty[$ كما يلي: $h(x)=(x+2)^2+2-2\ln(x+2)$.
\QQ{1} أحسب $\displaystyle\lim_{x\to+\infty}h(x)$ و$\displaystyle\lim_{\substack{x\to-2\\x>-2}}h(x)$.
\QQ{2} أدرس اتجاه تغير الدالة $h$، ثم شكّل جدول تغيراتها.
\QQ{3} استنتج أنه من أجل كل $x$ من $]-2;+\infty[$: $h(x)>0$.

\textbf{II)} $f$ الدالة المعرفة على المجال $]-2;+\infty[$ كما يلي: $f(x)=x+1+\dfrac{2}{x+2}\ln(x+2)$.

$(C_f)$ تمثيلها البياني في المستوي المنسوب إلى المعلم المتعامد والمتجانس $(O;\vec{i},\vec{j})$ (وحدة الطول $1cm$).
\QQ{1} أحسب $\displaystyle\lim_{\substack{x\to-2\\x>-2}}f(x)$ وفسّر النتيجة هندسيا، ثم أحسب $\displaystyle\lim_{x\to+\infty}f(x)$.
\QQ{2}
\begin{sq}
\qa بيّن أنه من أجل كل $x$ من المجال $]-2;+\infty[$: $f'(x)=\dfrac{h(x)}{(x+2)^2}$.
\qb أدرس اتجاه تغير الدالة $f$ على المجال $]-2;+\infty[$، ثم شكّل جدول تغيراتها.
\end{sq}
\QQ{3}
\begin{sq}
\qa بيّن أن المستقيم $(\Delta)$ ذا المعادلة $y=x+1$ مقارب مائل للمنحنى $(C_f)$ بجوار $+\infty$.
\qb أدرس وضعية المنحنى $(C_f)$ بالنسبة لـ $(\Delta)$.
\end{sq}
\QQ{4}
\begin{sq}
\qa أثبت أن المنحنى $(C_f)$ يقبل نقطة انعطاف $A$ يطلب تعيين إحداثييها.
\qb أرسم المستقيمين المقاربين والمنحنى $(C_f)$.
\end{sq}

\textbf{III)} $g$ الدالة المعرفة على المجال $]-2;+\infty[$ بـ: $g(x)=|x+1|+\dfrac{2}{x+2}\left|\ln(x+2)\right|$.
\QQ{1} أحسب $\displaystyle\lim_{\substack{x\to-1\\x<-1}}\frac{g(x)-g(-1)}{x+1}$ و$\displaystyle\lim_{\substack{x\to-1\\x>-1}}\frac{g(x)-g(-1)}{x+1}$. ماذا تستنتج بالنسبة إلى $g$؟
\QQ{2} أعط تفسيرا هندسيا لهذه النتيجة.
\QQ{3} انطلاقا من المنحنى $(C_f)$ أرسم المنحنى $(C_g)$ الممثل للدالة $g$ في نفس المعلم السابق.
\end{tamrin}
\begin{molahada}
في نص المصدر كُتبت عبارة الجزء III على الشكل «$f(x)=|x+1|+\dots$»، والصواب $g(x)=|x+1|+\dfrac{2}{x+2}\left|\ln(x+2)\right|$ لأن الدالة المقصودة هي $g$.
\end{molahada}

\begin{hal}
\textbf{الجزء الأول:}
\QQ{1} \textbf{عند $-2$:} $(x+2)^2\to0$ و$\ln(x+2)\to-\infty$ إذن $-2\ln(x+2)\to+\infty$، ومنه \bx{\lim_{x\to-2^+}h(x)=+\infty}.

\textbf{عند $+\infty$:} حالة عدم تعيين ($+\infty-\infty$)؛ نكتب:
\[h(x)=(x+2)^2\left[1+\frac{2}{(x+2)^2}-2\cdot\frac{\ln(x+2)}{x+2}\cdot\frac{1}{x+2}\right],\]
وبما أن $\displaystyle\lim_{t\to+\infty}\frac{\ln t}{t}=0$ فإن ما بين القوسين يؤول إلى $1$، ومنه \bx{\lim_{x\to+\infty}h(x)=+\infty}.
\QQ{2} $h'(x)=2(x+2)-\dfrac{2}{x+2}=\dfrac{2\left[(x+2)^2-1\right]}{x+2}=\dfrac{2(x+1)(x+3)}{x+2}$.

على $]-2;+\infty[$: $x+2>0$ و$x+3>0$، إذن إشارة $h'(x)$ هي إشارة $x+1$: $h$ متناقصة تماما على $]-2;-1]$ ومتزايدة تماما على $[-1;+\infty[$؛ $h(-1)=1+2-2\ln1=3$.
\begin{TAB}
\tkzTabInit[lgt=2,espcl=3.5]{$x$ /0.8, $h'(x)$ /0.8, $h(x)$ /1.6}{$-2$, $-1$, $+\infty$}
\tkzTabLine{d, -, z, +, }
\tkzTabVar{D+/$+\infty$, -/$3$, +/$+\infty$}
\end{TAB}
\QQ{3} القيمة الحدية الصغرى للدالة $h$ هي $h(-1)=3$، إذن من أجل كل $x>-2$: $h(x)\geqslant3>0$، أي \bx{h(x)>0}.

\textbf{الجزء الثاني:}
\QQ{1} \textbf{عند $-2$:} $\dfrac{2}{x+2}\to+\infty$ و$\ln(x+2)\to-\infty$، إذن جداؤهما $\to-\infty$، و$x+1\to-1$، ومنه \bx{\lim_{x\to-2^+}f(x)=-\infty}.

التفسير: المستقيم ذو المعادلة $x=-2$ مقارب عمودي لـ $(C_f)$.

\textbf{عند $+\infty$:} $\displaystyle\lim_{x\to+\infty}\frac{\ln(x+2)}{x+2}=0$ و$x+1\to+\infty$، إذن \bx{\lim_{x\to+\infty}f(x)=+\infty}.
\QQ{2}\pa $f'(x)=1+2\cdot\dfrac{\frac{1}{x+2}(x+2)-\ln(x+2)}{(x+2)^2}=1+\dfrac{2-2\ln(x+2)}{(x+2)^2}=\dfrac{(x+2)^2+2-2\ln(x+2)}{(x+2)^2}$،

أي \bx{f'(x)=\frac{h(x)}{(x+2)^2}}.
\pb $h(x)>0$ و$(x+2)^2>0$، إذن $f'(x)>0$ و$f$ متزايدة تماما على $]-2;+\infty[$.
\begin{TAB}
\tkzTabInit[lgt=2,espcl=5]{$x$ /0.8, $f'(x)$ /0.8, $f(x)$ /1.6}{$-2$, $+\infty$}
\tkzTabLine{d, +, }
\tkzTabVar{D-/$-\infty$, +/$+\infty$}
\end{TAB}
\QQ{3}\pa لدينا:
\[f(x)-(x+1)=\frac{2\ln(x+2)}{x+2}\qquad\text{و}\qquad\lim_{x\to+\infty}\frac{2\ln(x+2)}{x+2}=0.\]
إذن المستقيم \bx{(\Delta): y=x+1} مقارب مائل لـ $(C_f)$ بجوار $+\infty$.
\pb إشارة $f(x)-(x+1)$ من إشارة $\ln(x+2)$ (لأن $x+2>0$):
\begin{itemize}
\item $-2<x<-1$: $0<x+2<1$ إذن $\ln(x+2)<0$: $(C_f)$ تحت $(\Delta)$.
\item $x>-1$: $\ln(x+2)>0$: $(C_f)$ فوق $(\Delta)$.
\item $x=-1$: $(C_f)$ يقطع $(\Delta)$ في النقطة $(-1;0)$.
\end{itemize}
\QQ{4}\pa $f'(x)=1+\dfrac{2\left(1-\ln(x+2)\right)}{(x+2)^2}$. بوضع $u=x+2$:
\[\left(\frac{1-\ln u}{u^2}\right)'=\frac{-\frac1u\cdot u^2-2u(1-\ln u)}{u^4}=\frac{2\ln u-3}{u^3},\quad\text{إذن}\quad f''(x)=\frac{4\ln(x+2)-6}{(x+2)^3}.\]
$f''(x)=0\iff \ln(x+2)=\dfrac32\iff x=e^{\frac32}-2\approx2{,}48$، و$f''(x)$ تنعدم عند هذه القيمة مغيرة إشارتها (سالبة قبلها وموجبة بعدها، لأن $(x+2)^3>0$).

إذن $(C_f)$ يقبل نقطة انعطاف: $f\left(e^{\frac32}-2\right)=e^{\frac32}-1+\dfrac{2\cdot\frac32}{e^{\frac32}}$، أي
\[\bx{A\left(e^{\frac32}-2\ ;\ e^{\frac32}-1+3e^{-\frac32}\right)}\approx(2{,}48\ ;\ 4{,}15).\]
\pb الرسم (وحدة الطول $1cm$):
\begin{GR}{-3}{7}{-5}{8}{1}
 \draw[cThm,thick,dashed] (-2,-5)--(-2,8);
 \draw[cRem,thick,domain=-3:7] plot(\x,{\x+1});
 \node[cRem,above left] at (5.5,6.5){$(\Delta)$};
 \draw[cDef,very thick,domain=-1.97:7,samples=250] plot(\x,{\x+1+2*ln(\x+2)/(\x+2)});
 \fill (2.4817,4.1511) circle (2pt) node[below right]{$A$};
 \fill (-1,0) circle (1.5pt);
 \node[cDef] at (-1.2,-3){$(C_f)$};
 \Ticks{-1,1,2,3,4,5,6}{-4,-3,-2,-1,1,2,3,4,5,6,7}
\end{GR}

\textbf{الجزء الثالث:}
لدينا $f(-1)=0$ و$f$ متزايدة تماما، إذن $f(x)<0$ على $]-2;-1[$ و$f(x)\geqslant0$ على $[-1;+\infty[$.
\begin{itemize}
\item من أجل $x\geqslant-1$: $x+1\geqslant0$ و$\ln(x+2)\geqslant0$، إذن $g(x)=x+1+\dfrac{2\ln(x+2)}{x+2}=f(x)$.
\item من أجل $-2<x<-1$: $x+1<0$ و$\ln(x+2)<0$، إذن $g(x)=-(x+1)-\dfrac{2\ln(x+2)}{x+2}=-f(x)$.
\end{itemize}
\QQ{1} $g(-1)=f(-1)=0$.
\begin{itemize}
\item $\displaystyle\lim_{x\to-1^+}\frac{g(x)-g(-1)}{x+1}=\lim_{x\to-1^+}\frac{f(x)-f(-1)}{x-(-1)}=f'(-1)=\frac{h(-1)}{1}=3$.
\item $\displaystyle\lim_{x\to-1^-}\frac{g(x)-g(-1)}{x+1}=\lim_{x\to-1^-}\frac{-\left(f(x)-f(-1)\right)}{x-(-1)}=-f'(-1)=-3$.
\end{itemize}
النهايتان مختلفتان، إذن \bxt{$g$ غير قابلة للاشتقاق عند $-1$}.
\QQ{2} $(C_g)$ يقبل عند النقطة $(-1;0)$ نصفي مماسين معاملا توجيههما $3$ (على اليمين) و$-3$ (على اليسار)؛ هذه النقطة نقطة زاوية.
\QQ{3} على $[-1;+\infty[$: $(C_g)$ ينطبق على $(C_f)$؛ وعلى $]-2;-1[$: $(C_g)$ هو نظير $(C_f)$ بالنسبة إلى حامل محور الفواصل.
\begin{GR}{-3}{7}{-2}{8}{1}
 \draw[cThm,thick,dashed] (-2,-2)--(-2,8);
 \draw[cDef,thick,dashed,domain=-1.9:-1,samples=80] plot(\x,{\x+1+2*ln(\x+2)/(\x+2)});
 \draw[cThm,very thick,domain=-1.97:-1,samples=150] plot(\x,{-(\x+1+2*ln(\x+2)/(\x+2))});
 \draw[cThm,very thick,domain=-1:7,samples=200] plot(\x,{\x+1+2*ln(\x+2)/(\x+2)});
 \draw[gray,thick,-{Stealth[length=2mm]}] (-1,0)--(-0.4,1.8);
 \draw[gray,thick,-{Stealth[length=2mm]}] (-1,0)--(-1.6,1.8);
 \node[cThm] at (4.5,7){$(C_g)$};
 \Ticks{-1,1,2,3,4,5,6}{-1,1,2,3,4,5,6,7}
\end{GR}
(بالأحمر $(C_g)$، وبالمتقطع الأزرق جزء $(C_f)$ على $]-2;-1[$، وبالرمادي نصفا المماسين.)
\end{hal}

%==================== التمرين 7 ====================
\begin{tamrin}{بكالوريا 2016 -- الموضوع الأول -- تقني رياضي}
\textbf{I)} $g$ الدالة المعرفة على المجال $]-1;+\infty[$ بـ: $g(x)=\dfrac{x-1}{x+1}+\ln(x+1)$.
\QQ{1}
\begin{sq}
\qa أحسب $\displaystyle\lim_{\substack{x\to-1\\x>-1}}g(x)$ و$\displaystyle\lim_{x\to+\infty}g(x)$.
\qb أدرس اتجاه تغير الدالة $g$ على المجال $]-1;+\infty[$، ثم شكّل جدول تغيراتها.
\end{sq}
\QQ{2}
\begin{sq}
\qa بيّن أن المعادلة $g(x)=0$ تقبل حلا وحيدا $\alpha$ حيث: $0{,}4<\alpha<0{,}5$.
\qb استنتج إشارة $g(x)$ على المجال $]-1;+\infty[$.
\end{sq}

\textbf{II)} $f$ الدالة المعرفة على المجال $]-1;+\infty[$ كما يلي: $f(x)=1+(x-1)\ln(x+1)$.

$(C_f)$ تمثيلها البياني في المستوي المنسوب إلى المعلم المتعامد والمتجانس $(O;\vec{i},\vec{j})$.
\QQ{1} أحسب $\displaystyle\lim_{\substack{x\to-1\\x>-1}}f(x)$ وفسّر النتيجة هندسيا، ثم أحسب $\displaystyle\lim_{x\to+\infty}f(x)$.
\QQ{2}
\begin{sq}
\qa أدرس اتجاه تغير الدالة $f$ على المجال $]-1;+\infty[$، ثم شكّل جدول تغيراتها.
\qb بيّن أن $f(\alpha)=-\alpha+4-\dfrac{4}{\alpha+1}$، ثم أعط حصرا لـ $f(\alpha)$ (تدور النتائج إلى $10^{-2}$).
\end{sq}
\QQ{3} ليكن $a$ عددا حقيقيا من المجال $]-1;+\infty[$. نسمي $(T_a)$ مماس المنحنى $(C_f)$ عند النقطة ذات الفاصلة $a$.

نضع من أجل كل عدد حقيقي $x$ من المجال $]-1;+\infty[$: $h(x)=f(x)-\left[f'(a)(x-a)+f(a)\right]$.
\begin{sq}
\qa تحقق أنه من أجل كل $x$ من المجال $]-1;+\infty[$: $h'(x)=f'(x)-f'(a)$.
\qb باستعمال اتجاه تغير الدالة $g$، عيّن إشارة $h'(x)$ حسب قيم $x$ واستنتج اتجاه تغير الدالة $h$ على $]-1;+\infty[$.
\qc حدد الوضع النسبي للمنحنى $(C_f)$ والمستقيم $(T_a)$.
\end{sq}
\QQ{4}
\begin{sq}
\qa بيّن أنه يوجد مماسان $(T_a)$ يشملان النقطة $A(1;0)$ يطلب تعيين معادلتيهما.
\qb أرسم المماسين والمنحنى $(C_f)$.
\end{sq}
\end{tamrin}

\begin{hal}
\textbf{الجزء الأول:}
\QQ{1}\pa \textbf{عند $-1$:} $x-1\to-2$ و$x+1\to0^+$ إذن $\dfrac{x-1}{x+1}\to-\infty$، و$\ln(x+1)\to-\infty$، ومنه \bx{\lim_{x\to-1^+}g(x)=-\infty}.

\textbf{عند $+\infty$:} $\dfrac{x-1}{x+1}\to1$ و$\ln(x+1)\to+\infty$، ومنه \bx{\lim_{x\to+\infty}g(x)=+\infty}.
\pb $g'(x)=\dfrac{(x+1)-(x-1)}{(x+1)^2}+\dfrac{1}{x+1}=\dfrac{2}{(x+1)^2}+\dfrac{x+1}{(x+1)^2}=\dfrac{x+3}{(x+1)^2}$.

على $]-1;+\infty[$: $x+3>0$، إذن $g'(x)>0$ و$g$ متزايدة تماما.
\begin{TAB}
\tkzTabInit[lgt=2,espcl=3.5]{$x$ /0.8, $g'(x)$ /0.8, $g(x)$ /1.6}{$-1$, $\alpha$, $+\infty$}
\tkzTabLine{d, +, t, +, }
\tkzTabVar{D-/$-\infty$, R/, +/$+\infty$}
\tkzTabVal{1}{3}{0.5}{}{$0$}
\end{TAB}
\QQ{2}\pa $g$ مستمرة ومتزايدة تماما على $]-1;+\infty[$ وتأخذ قيمها في $\R$، إذن المعادلة $g(x)=0$ تقبل حلا وحيدا $\alpha$.

$g(0{,}4)\approx-0{,}09<0$ و$g(0{,}5)\approx0{,}07>0$، إذن \bx{0{,}4<\alpha<0{,}5}.
\pb $g(x)<0$ على $]-1;\alpha[$، $g(\alpha)=0$، $g(x)>0$ على $]\alpha;+\infty[$.

\textbf{الجزء الثاني:}
\QQ{1} \textbf{عند $-1$:} $x-1\to-2$ و$\ln(x+1)\to-\infty$، إذن $(x-1)\ln(x+1)\to+\infty$، ومنه \bx{\lim_{x\to-1^+}f(x)=+\infty}.

التفسير: المستقيم ذو المعادلة $x=-1$ مقارب عمودي لـ $(C_f)$.

\textbf{عند $+\infty$:} $x-1\to+\infty$ و$\ln(x+1)\to+\infty$، ومنه \bx{\lim_{x\to+\infty}f(x)=+\infty}.
\QQ{2}\pa $f'(x)=\ln(x+1)+(x-1)\cdot\dfrac{1}{x+1}=g(x)$.

إذن $f$ متناقصة تماما على $]-1;\alpha]$ ومتزايدة تماما على $[\alpha;+\infty[$.
\begin{TAB}
\tkzTabInit[lgt=2,espcl=3.5]{$x$ /0.8, $f'(x)$ /0.8, $f(x)$ /1.8}{$-1$, $\alpha$, $+\infty$}
\tkzTabLine{d, -, z, +, }
\tkzTabVar{D+/$+\infty$, -/$f(\alpha)$, +/$+\infty$}
\end{TAB}
\pb من $g(\alpha)=0$: $\ln(\alpha+1)=-\dfrac{\alpha-1}{\alpha+1}$، إذن $f(\alpha)=1-\dfrac{(\alpha-1)^2}{\alpha+1}$.

ولدينا $(\alpha-1)^2=\left[(\alpha+1)-2\right]^2=(\alpha+1)^2-4(\alpha+1)+4$، إذن:
\[f(\alpha)=1-\left[(\alpha+1)-4+\frac{4}{\alpha+1}\right]=-\alpha+4-\frac{4}{\alpha+1}.\]
\textbf{الحصر:} نعتبر $\varphi(t)=-t+4-\dfrac{4}{t+1}$ على $[0{,}4;0{,}5]$: $\varphi'(t)=-1+\dfrac{4}{(t+1)^2}=\dfrac{4-(t+1)^2}{(t+1)^2}>0$ (لأن $t+1<2$)، إذن $\varphi$ متزايدة تماما، ومنه $\varphi(0{,}4)<f(\alpha)<\varphi(0{,}5)$:
\[\varphi(0{,}4)\approx0{,}743\quad;\quad\varphi(0{,}5)\approx0{,}833\qquad\text{إذن}\qquad\bx{0{,}74<f(\alpha)<0{,}84}.\]
\QQ{3}\pa $f'(a)$ و$f(a)$ و$a$ ثوابت، إذن $\left[f'(a)(x-a)+f(a)\right]'=f'(a)$.

ومنه \bx{h'(x)=f'(x)-f'(a)=g(x)-g(a)}.
\pb $g$ متزايدة تماما على $]-1;+\infty[$:
\begin{itemize}
\item إذا كان $x<a$ فإن $g(x)<g(a)$ أي $h'(x)<0$: $h$ متناقصة تماما على $]-1;a]$.
\item إذا كان $x>a$ فإن $g(x)>g(a)$ أي $h'(x)>0$: $h$ متزايدة تماما على $[a;+\infty[$.
\end{itemize}
\pc $h$ تقبل قيمة حدية صغرى عند $a$ هي $h(a)=f(a)-\left[0+f(a)\right]=0$، إذن $h(x)\geqslant0$ من أجل كل $x>-1$.

ومنه \bxt{$(C_f)$ يقع فوق $(T_a)$} ويمسه فقط في النقطة ذات الفاصلة $a$.
\QQ{4}\pa $(T_a): y=f'(a)(x-a)+f(a)$. $A(1;0)\in(T_a)\iff f'(a)(1-a)+f(a)=0$:
\begin{align*}
&(1-a)\left[\frac{a-1}{a+1}+\ln(a+1)\right]+1+(a-1)\ln(a+1)=0\\
\iff&-\frac{(a-1)^2}{a+1}+(1-a)\ln(a+1)+1+(a-1)\ln(a+1)=0
\iff 1-\frac{(a-1)^2}{a+1}=0\\
\iff&(a-1)^2=a+1\iff a^2-3a=0\iff a(a-3)=0\iff a=0\ \text{أو}\ a=3.
\end{align*}
كلا القيمتين في $]-1;+\infty[$، إذن يوجد مماسان:
\begin{itemize}
\item $a=0$: $f(0)=1$، $f'(0)=g(0)=-1$: \bx{(T_0):\ y=-x+1}.
\item $a=3$: $f(3)=1+2\ln4=1+4\ln2$، $f'(3)=g(3)=\frac12+2\ln2$:
\[y=\left(\tfrac12+2\ln2\right)(x-3)+1+4\ln2,\quad\text{أي}\quad \bx{(T_3):\ y=\left(\tfrac12+2\ln2\right)x-\tfrac12-2\ln2}.\]
\end{itemize}
\pb الرسم ($f(\alpha)\approx0{,}80$، $f(1)=1$، $f(3)\approx3{,}77$):
\begin{GR}{-1.5}{4.5}{-1.5}{5.5}{1.5}
 \draw[cThm,thick,dashed] (-1,-1.5)--(-1,5.5);
 \draw[cRem,thick,domain=-1.5:3] plot(\x,{-\x+1});
 \node[cRem,right] at (-1.3,2.3){$(T_0)$};
 \draw[cEx,thick,domain=0:4.5] plot(\x,{1.8863*(\x-1)});
 \node[cEx,right] at (3.6,4.9){$(T_3)$};
 \draw[cDef,very thick,domain=-0.975:4.5,samples=250] plot(\x,{1+(\x-1)*ln(\x+1)});
 \fill (1,0) circle (2pt) node[below right]{$A$};
 \fill (0,1) circle (1.5pt);
 \fill (3,3.7726) circle (1.5pt);
 \node[cDef] at (-0.5,4){$(C_f)$};
 \Ticks{-1,1,2,3,4}{-1,1,2,3,4,5}
\end{GR}
\end{hal}

%==================== التمرين 8 ====================
\begin{tamrin}{بكالوريا 2016 -- الموضوع الثاني -- تقني رياضي}
\textbf{I)} نعتبر الدالة $g$ المعرفة على المجال $]0;+\infty[$ بـ: $g(x)=x-x\ln x$.
\QQ{1}
\begin{sq}
\qa أحسب $\displaystyle\lim_{\substack{x\to0\\x>0}}g(x)$ و$\displaystyle\lim_{x\to+\infty}g(x)$.
\qb أدرس اتجاه تغير الدالة $g$ على المجال $]0;+\infty[$، ثم شكّل جدول تغيراتها.
\end{sq}
\QQ{2} بيّن أن المعادلة $g(x)=-1$ تقبل حلا وحيدا $\alpha$ حيث: $3{,}5<\alpha<3{,}6$.
\QQ{3} استنتج إشارة العبارة $g(x)+1$ على المجال $]0;+\infty[$.

\textbf{II)} $f$ الدالة المعرفة على المجال $]0;+\infty[$ كما يلي: $f(x)=\dfrac{\ln x}{x+1}$.

$(C_f)$ تمثيلها البياني في المستوي المنسوب إلى المعلم المتعامد $(O;\vec{i},\vec{j})$، حيث: $\|\vec{i}\|=2cm$ و$\|\vec{j}\|=4cm$.
\QQ{1} بيّن أن $(C_f)$ يقبل مستقيمين مقاربين معادلتاهما $x=0$ و$y=0$.
\QQ{2}
\begin{sq}
\qa برهن أنه من أجل كل عدد حقيقي $x$ من المجال $]0;+\infty[$: $f'(x)=\dfrac{g(x)+1}{x(x+1)^2}$.
\qb بيّن أن الدالة $f$ متزايدة تماما على المجال $]0;\alpha]$ ومتناقصة تماما على المجال $[\alpha;+\infty[$، ثم شكّل جدول تغيراتها.
\qc أكتب معادلة المماس $(T)$ للمنحنى $(C_f)$ عند النقطة ذات الفاصلة $1$.
\qd أحسب $\displaystyle\lim_{x\to\alpha}\frac{f(x)-f(\alpha)}{x-\alpha}$، ثم فسّر النتيجة هندسيا.
\end{sq}
\QQ{3}
\begin{sq}
\qa بيّن أن: $f(\alpha)=\dfrac{1}{\alpha}$.
\qb استنتج حصرا لـ $f(\alpha)$ (تدور النتائج إلى $10^{-2}$).
\qc أرسم $(C_f)$.
\end{sq}
\QQ{4} نعتبر المعادلة ذات المجهول الحقيقي الموجب تماما $x$ و$m$ وسيط حقيقي: $x^2+x-2m(x+1)=\ln\left(x^2\right)\quad\dots(E)$.
\begin{sq}
\qa تحقق أن المعادلة $(E)$ يؤول حلها إلى حل المعادلة: $f(x)=\dfrac12x-m$.
\qb عيّن بيانيا قيم $m$ التي من أجلها تقبل المعادلة $(E)$ حلين متمايزين.
\end{sq}
\QQ{5} $h$ هي الدالة المعرفة على $\R^*$ كما يلي: $h(x)=\dfrac{\ln|x|}{-|x|-1}$، و$(C_h)$ تمثيلها البياني.
\begin{sq}
\qa بيّن أن الدالة $h$ زوجية.
\qb أرسم في نفس المعلم المنحنى $(C_h)$ مستعينا بالمنحنى $(C_f)$.
\end{sq}
\end{tamrin}

\begin{hal}
\textbf{الجزء الأول:}
\QQ{1}\pa $\displaystyle\lim_{x\to0^+}x\ln x=0$ إذن \bx{\lim_{x\to0^+}g(x)=0}.

$g(x)=x(1-\ln x)$ مع $x\to+\infty$ و$1-\ln x\to-\infty$، إذن \bx{\lim_{x\to+\infty}g(x)=-\infty}.
\pb $g'(x)=1-\left(\ln x+x\cdot\dfrac1x\right)=-\ln x$.

$g'(x)>0$ على $]0;1[$ و$g'(x)<0$ على $]1;+\infty[$: $g$ متزايدة تماما على $]0;1]$ ومتناقصة تماما على $[1;+\infty[$؛ $g(1)=1$.
\begin{TAB}
\tkzTabInit[lgt=2,espcl=3.5]{$x$ /0.8, $g'(x)$ /0.8, $g(x)$ /1.6}{$0$, $1$, $+\infty$}
\tkzTabLine{d, +, z, -, }
\tkzTabVar{D-/$0$, +/$1$, -/$-\infty$}
\end{TAB}
\QQ{2} \textbf{على $]0;1]$:} $g(x)>0>-1$، فلا حلول.

\textbf{على $[1;+\infty[$:} $g$ مستمرة ومتناقصة تماما من $1$ إلى $-\infty$، و$-1\in]-\infty;1]$، إذن المعادلة $g(x)=-1$ تقبل حلا وحيدا $\alpha$.

$g(3{,}5)\approx-0{,}88>-1$ و$g(3{,}6)\approx-1{,}01<-1$، إذن \bx{3{,}5<\alpha<3{,}6}.
\QQ{3} على $]0;1]$: $g(x)>0$ إذن $g(x)+1>0$؛ على $[1;\alpha[$: $g(x)>g(\alpha)=-1$ إذن $g(x)+1>0$؛ على $]\alpha;+\infty[$: $g(x)<-1$ إذن $g(x)+1<0$.

\bxt{$g(x)+1>0$ على $]0;\alpha[ ; g(x)+1<0$ على $]\alpha;+\infty[$}، و$g(\alpha)+1=0$.

\textbf{الجزء الثاني:}
\QQ{1} $\displaystyle\lim_{x\to0^+}\ln x=-\infty$ و$x+1\to1$ إذن $\displaystyle\lim_{x\to0^+}f(x)=-\infty$: المستقيم \bx{x=0} مقارب عمودي.

$f(x)=\dfrac{\ln x}{x}\cdot\dfrac{x}{x+1}$ مع $\dfrac{\ln x}{x}\to0$ و$\dfrac{x}{x+1}\to1$، إذن $\displaystyle\lim_{x\to+\infty}f(x)=0$: المستقيم \bx{y=0} مقارب أفقي بجوار $+\infty$.
\QQ{2}\pa $f'(x)=\dfrac{\frac1x(x+1)-\ln x}{(x+1)^2}=\dfrac{x+1-x\ln x}{x(x+1)^2}=\dfrac{(x-x\ln x)+1}{x(x+1)^2}$، أي \bx{f'(x)=\frac{g(x)+1}{x(x+1)^2}}.
\pb $x(x+1)^2>0$، إذن إشارة $f'(x)$ من إشارة $g(x)+1$: $f$ متزايدة تماما على $]0;\alpha]$ ومتناقصة تماما على $[\alpha;+\infty[$.
\begin{TAB}
\tkzTabInit[lgt=2,espcl=3.5]{$x$ /0.8, $f'(x)$ /0.8, $f(x)$ /1.8}{$0$, $\alpha$, $+\infty$}
\tkzTabLine{d, +, z, -, }
\tkzTabVar{D-/$-\infty$, +/$f(\alpha)$, -/$0$}
\end{TAB}
\pc $f(1)=0$ و$f'(1)=\dfrac{g(1)+1}{1\times4}=\dfrac{2}{4}=\dfrac12$، إذن \bx{(T):\ y=\frac12x-\frac12}.
\pd $f$ قابلة للاشتقاق عند $\alpha$، إذن $\displaystyle\lim_{x\to\alpha}\frac{f(x)-f(\alpha)}{x-\alpha}=f'(\alpha)=\frac{g(\alpha)+1}{\alpha(\alpha+1)^2}=\bx{0}$.

التفسير: $(C_f)$ يقبل عند النقطة ذات الفاصلة $\alpha$ مماسا أفقيا معادلته $y=f(\alpha)$.
\QQ{3}\pa $g(\alpha)=-1\iff \alpha-\alpha\ln\alpha=-1\iff \ln\alpha=\dfrac{\alpha+1}{\alpha}$، إذن $f(\alpha)=\dfrac{\alpha+1}{\alpha}\times\dfrac{1}{\alpha+1}=\bx{\frac1\alpha}$.
\pb $3{,}5<\alpha<3{,}6\Rightarrow \dfrac{1}{3{,}6}<\dfrac1\alpha<\dfrac{1}{3{,}5}$، أي $0{,}277\ldots<f(\alpha)<0{,}285\ldots$، إذن \bx{0{,}27<f(\alpha)<0{,}29}.
\pc الرسم مع $(C_h)$ (السؤال 5) و$(T)$ في نفس المعلم (الرسم بنصف السلم المطلوب: $1cm$ أفقيا و$2cm$ عموديا):
\begin{GRtm}{-8}{8}{-2}{2}{1}{2}{0.5}
 \draw[cEx,thick,domain=-1:6] plot(\x,{0.5*\x-0.5});
 \node[cEx,right] at (4.8,1.9){$(T)$};
 \draw[cDef,very thick,domain=0.11:8,samples=200] plot(\x,{ln(\x)/(\x+1)});
 \draw[cThm,very thick,domain=0.11:8,samples=200] plot(\x,{-ln(\x)/(\x+1)});
 \draw[cThm,very thick,domain=-8:-0.11,samples=200] plot(\x,{-ln(-\x)/(-\x+1)});
 \fill[cDef] (3.591,0.2785) circle (1.5pt);
 \node[cDef] at (6,0.5){$(C_f)$};
 \node[cThm] at (6,-0.5){$(C_h)$};
 \node[cThm] at (-6,-0.5){$(C_h)$};
 \foreach \t in {-6,-4,-2,2,4,6} \node[below,font=\scriptsize] at (\t,0){$\t$};
 \node[below,font=\scriptsize] at (1,0){$1$};
 \foreach \t in {-1.5,-1,-0.5,0.5,1,1.5} \node[left,font=\scriptsize] at (0,\t){$\t$};
\end{GRtm}
\QQ{4}\pa من أجل $x>0$: $\ln\left(x^2\right)=2\ln x$، إذن:
\[(E)\iff x(x+1)-2m(x+1)=2\ln x\iff \frac{x}{2}-m=\frac{\ln x}{x+1}\iff f(x)=\frac12x-m,\]
وذلك بالقسمة على $2(x+1)>0$.
\pb حلول $(E)$ هي فواصل نقط تقاطع $(C_f)$ مع المستقيم $(D_m): y=\frac12x-m$ الموازي لـ $(T)$ (لـ $(T)$: $m=\frac12$).

من الرسم، $(C_f)$ يقع تحت $(T)$ ويمسه في النقطة $(1;0)$:
\begin{itemize}
\item $m<\frac12$: $(D_m)$ فوق $(T)$، لا توجد حلول.
\item $m=\frac12$: حل وحيد $x=1$.
\item $m>\frac12$: $(D_m)$ تحت $(T)$ ويقطع $(C_f)$ في نقطتين.
\end{itemize}
إذن تقبل $(E)$ حلين متمايزين من أجل \bx{m>\frac12}.
\QQ{5}\pa $\R^*$ متناظر بالنسبة إلى $0$، ومن أجل كل $x\in\R^*$: $h(-x)=\dfrac{\ln|-x|}{-|-x|-1}=\dfrac{\ln|x|}{-|x|-1}=h(x)$، إذن \bxt{$h$ زوجية}.
\pb من أجل $x>0$: $h(x)=\dfrac{\ln x}{-(x+1)}=-f(x)$، إذن على $]0;+\infty[$ $(C_h)$ هو نظير $(C_f)$ بالنسبة إلى حامل محور الفواصل، وبما أن $h$ زوجية نكمل الرسم على $]-\infty;0[$ بالتناظر بالنسبة إلى حامل محور التراتيب (انظر الرسم أعلاه، $(C_h)$ بالأحمر).
\end{hal}

%==================== التمرين 9 ====================
\begin{tamrin}{بكالوريا 2017 -- الموضوع الأول -- تقني رياضي}
$f$ الدالة العددية المعرفة على $D_f$ حيث $D_f=]-\infty;1[\cup]2;+\infty[$ كما يلي: $f(x)=-2x+3+2\ln\left(\dfrac{x-1}{x-2}\right)$.

$(C_f)$ تمثيلها البياني في المستوي المنسوب إلى المعلم المتعامد والمتجانس $(O;\vec{i},\vec{j})$.
\QQ{1}
\begin{sq}
\qa أحسب النهايتين: $\displaystyle\lim_{\substack{x\to1\\x<1}}f(x)$، $\displaystyle\lim_{\substack{x\to2\\x>2}}f(x)$، ثم فسّر النتيجتين بيانيا.
\qb أحسب: $\displaystyle\lim_{x\to-\infty}f(x)$ و$\displaystyle\lim_{x\to+\infty}f(x)$.
\end{sq}
\QQ{2} بيّن أنه من أجل كل $x$ من $D_f$: $f'(x)=-2-\dfrac{2}{(x-1)(x-2)}$، ثم شكّل جدول تغيرات الدالة $f$.
\QQ{3}
\begin{sq}
\qa تحقق أن: من أجل كل عدد حقيقي $x$ من $D_f$: $(3-x)\in D_f$ و$f(3-x)+f(x)=0$.
\qb استنتج أن $(C_f)$ يقبل مركز تناظر يطلب تعيين إحداثييه.
\end{sq}
\QQ{4} أثبت أن المعادلة $f(x)=0$ تقبل حلا وحيدا $\alpha$ على المجال $]0{,}45;0{,}46[$، ثم استنتج أنها تقبل حلا آخر $\beta$ يطلب تعيين حصر له.
\QQ{5} بيّن أن المستقيم $(\Delta)$ ذا المعادلة $y=-2x+3$ مقارب مائل لـ $(C_f)$، ثم أدرس وضعية $(C_f)$ بالنسبة لـ $(\Delta)$.
\QQ{6} أرسم $(\Delta)$ و$(C_f)$.
\end{tamrin}

\begin{hal}
\QQ{1}\pa \textbf{عند $1^-$:} $x-1\to0^-$ و$x-2\to-1$، إذن $\dfrac{x-1}{x-2}\to0^+$ و$\ln\left(\dfrac{x-1}{x-2}\right)\to-\infty$؛ و$-2x+3\to1$. ومنه \bx{\lim_{x\to1^-}f(x)=-\infty}.

\textbf{عند $2^+$:} $x-1\to1$ و$x-2\to0^+$، إذن $\dfrac{x-1}{x-2}\to+\infty$ و$\ln\left(\dfrac{x-1}{x-2}\right)\to+\infty$؛ و$-2x+3\to-1$. ومنه \bx{\lim_{x\to2^+}f(x)=+\infty}.

التفسير: المستقيمان ذوا المعادلتين $x=1$ و$x=2$ مقاربان عموديان لـ $(C_f)$.
\pb $\displaystyle\lim_{x\to\pm\infty}\frac{x-1}{x-2}=1$ إذن $\ln\left(\dfrac{x-1}{x-2}\right)\to\ln1=0$، ومنه:
\[\bx{\lim_{x\to-\infty}f(x)=+\infty}\qquad;\qquad\bx{\lim_{x\to+\infty}f(x)=-\infty}.\]
\QQ{2} نضع $u(x)=\dfrac{x-1}{x-2}$ (موجبة تماما على $D_f$)، $u'(x)=\dfrac{(x-2)-(x-1)}{(x-2)^2}=\dfrac{-1}{(x-2)^2}$، إذن:
\[\frac{u'(x)}{u(x)}=\frac{-1}{(x-2)^2}\cdot\frac{x-2}{x-1}=\frac{-1}{(x-1)(x-2)},\qquad\text{ومنه}\qquad \bx{f'(x)=-2-\frac{2}{(x-1)(x-2)}}.\]
على $D_f$: $x-1$ و$x-2$ لهما نفس الإشارة، إذن $(x-1)(x-2)>0$ و$f'(x)<0$: $f$ متناقصة تماما على $]-\infty;1[$ وعلى $]2;+\infty[$.
\begin{TAB}
\tkzTabInit[lgt=2,espcl=2.8]{$x$ /0.8, $f'(x)$ /0.8, $f(x)$ /1.8}{$-\infty$, $1$, $2$, $+\infty$}
\tkzTabLine{, -, d, h, d, -, }
\tkzTabVar{+/$+\infty$, -DH/$-\infty$, D+/$+\infty$, -/$-\infty$}
\end{TAB}
\QQ{3}\pa إذا كان $x<1$ فإن $3-x>2$، وإذا كان $x>2$ فإن $3-x<1$؛ إذن $(3-x)\in D_f$.
\begin{align*}
f(3-x)&=-2(3-x)+3+2\ln\left(\frac{2-x}{1-x}\right)=2x-3+2\ln\left(\frac{x-2}{x-1}\right)\\
&=2x-3-2\ln\left(\frac{x-1}{x-2}\right)=-\left[-2x+3+2\ln\left(\frac{x-1}{x-2}\right)\right]=-f(x).
\end{align*}
ومنه \bx{f(3-x)+f(x)=0}.
\pb العلاقة من الشكل $f(2a-x)+f(x)=2b$ مع $a=\frac32$ و$b=0$، إذن النقطة \bx{\Omega\left(\tfrac32;0\right)} مركز تناظر لـ $(C_f)$.
\QQ{4} $f$ مستمرة ومتناقصة تماما على $]-\infty;1[$ وتأخذ قيمها في $\R$، إذن المعادلة $f(x)=0$ تقبل حلا وحيدا $\alpha$ في $]-\infty;1[$.

$f(0{,}45)\approx0{,}028>0$ و$f(0{,}46)\approx-0{,}016<0$، إذن $\alpha\in]0{,}45;0{,}46[$.

من 3-أ: $f(3-\alpha)=-f(\alpha)=0$ و$3-\alpha\in]2;+\infty[$، وبما أن $f$ متناقصة تماما على $]2;+\infty[$ فهذا الحل وحيد فيه: $\beta=3-\alpha$.

من $0{,}45<\alpha<0{,}46$: \bx{2{,}54<\beta<2{,}55}.
\QQ{5} لدينا:
\[f(x)-(-2x+3)=2\ln\left(\frac{x-1}{x-2}\right)\qquad\text{و}\qquad\lim_{x\to\pm\infty}2\ln\left(\frac{x-1}{x-2}\right)=2\ln1=0.\]
إذن المستقيم \bx{(\Delta):y=-2x+3} مقارب مائل لـ $(C_f)$ بجوار $-\infty$ وبجوار $+\infty$.

\textbf{الوضعية:} $\dfrac{x-1}{x-2}=1+\dfrac{1}{x-2}$.
\begin{itemize}
\item على $]-\infty;1[$: $x-2<0$ إذن $\dfrac{x-1}{x-2}<1$ و$\ln\left(\dfrac{x-1}{x-2}\right)<0$: $(C_f)$ تحت $(\Delta)$.
\item على $]2;+\infty[$: $x-2>0$ إذن $\dfrac{x-1}{x-2}>1$ و$\ln\left(\dfrac{x-1}{x-2}\right)>0$: $(C_f)$ فوق $(\Delta)$.
\end{itemize}
\QQ{6} الرسم:
\begin{GR}{-3}{6}{-6}{8}{0.9}
 \draw[cThm,thick,dashed] (1,-6)--(1,8);
 \draw[cThm,thick,dashed] (2,-6)--(2,8);
 \draw[cRem,thick,domain=-2.5:5.5] plot(\x,{-2*\x+3});
 \node[cRem,right] at (3.6,-4.2){$(\Delta)$};
 \draw[cDef,very thick,domain=-3:0.9,samples=200] plot(\x,{-2*\x+3+2*ln((\x-1)/(\x-2))});
 \draw[cDef,very thick,domain=0.9:0.996,samples=100] plot(\x,{-2*\x+3+2*ln((\x-1)/(\x-2))});
 \draw[cDef,very thick,domain=2.004:2.05,samples=100] plot(\x,{-2*\x+3+2*ln((\x-1)/(\x-2))});
 \draw[cDef,very thick,domain=2.05:6,samples=200] plot(\x,{-2*\x+3+2*ln((\x-1)/(\x-2))});
 \fill (1.5,0) circle (2pt) node[below]{$\Omega$};
 \fill (0.456,0) circle (1.5pt) node[above right]{$\alpha$};
 \fill (2.544,0) circle (1.5pt) node[below left]{$\beta$};
 \node[cDef] at (-2,6.2){$(C_f)$};
 \Ticks{-2,-1,3,4,5}{-5,-4,-3,-2,-1,1,2,3,4,5,6,7}
\end{GR}
\end{hal}

%==================== التمرين 10 ====================
\begin{tamrin}{بكالوريا 2017 -- الموضوع الأول (الدورة الاستثنائية) -- تقني رياضي}
\textbf{I)} $g$ الدالة العددية المعرفة على المجال $]0;+\infty[$ كما يلي: $g(x)=-\dfrac12+\dfrac{2-\ln x}{x^2}$.
\QQ{1} أحسب $\displaystyle\lim_{\substack{x\to0\\x>0}}g(x)$ و$\displaystyle\lim_{x\to+\infty}g(x)$.
\QQ{2} أدرس اتجاه تغير الدالة $g$، ثم شكّل جدول تغيراتها.
\QQ{3} بيّن أن المعادلة $g(x)=0$ تقبل حلا وحيدا $\alpha$ حيث: $1{,}71<\alpha<1{,}72$، ثم استنتج إشارة $g(x)$ حسب قيم $x$.

\textbf{II)} نعتبر الدالة العددية $f$ المعرفة على المجال $]0;+\infty[$ كما يلي: $f(x)=-\dfrac12x+2+\dfrac{-1+\ln x}{x}$.

$(C_f)$ تمثيلها البياني في المستوي المنسوب إلى المعلم المتعامد والمتجانس $(O;\vec{i},\vec{j})$.
\QQ{1}
\begin{sq}
\qa أحسب $\displaystyle\lim_{\substack{x\to0\\x>0}}f(x)$ و$\displaystyle\lim_{x\to+\infty}f(x)$.
\qb أدرس اتجاه تغير الدالة $f$، ثم شكّل جدول تغيراتها.
\end{sq}
\QQ{2}
\begin{sq}
\qa بيّن أن المستقيم $(\Delta)$ ذا المعادلة $y=-\frac12x+2$ مقارب مائل للمنحنى $(C_f)$.
\qb أدرس وضعية المنحنى $(C_f)$ بالنسبة للمستقيم $(\Delta)$.
\end{sq}
\QQ{3} نقبل أن $f(\alpha)\approx0{,}87$ و$f(\gamma)=f(\beta)=0$ حيث: $0{,}77<\beta<0{,}78$ و$4{,}19<\gamma<4{,}22$.

أنشئ في المعلم السابق المستقيم $(\Delta)$ والمنحنى $(C_f)$.
\end{tamrin}
\begin{molahada}
في نص المصدر ورد الحصر $1{,}76<\beta<1{,}78$، وهو خطأ: $f(1{,}76)\approx0{,}87\neq0$ (هذه القيمة قريبة من القيمة العظمى $f(\alpha)$). الصواب: $f(0{,}77)\approx-0{,}023<0$ و$f(0{,}78)\approx0{,}009>0$، إذن $0{,}77<\beta<0{,}78$.
\end{molahada}

\begin{hal}
\textbf{الجزء الأول:}
\QQ{1} \textbf{عند $0$:} $2-\ln x\to+\infty$ و$\dfrac{1}{x^2}\to+\infty$، إذن \bx{\lim_{x\to0^+}g(x)=+\infty}.

\textbf{عند $+\infty$:} $\dfrac{2-\ln x}{x^2}=\dfrac{2}{x^2}-\dfrac{\ln x}{x}\cdot\dfrac1x\to0$، إذن \bx{\lim_{x\to+\infty}g(x)=-\frac12}.
\QQ{2} $g'(x)=\dfrac{-\frac1x\cdot x^2-2x(2-\ln x)}{x^4}=\dfrac{-x-4x+2x\ln x}{x^4}=\dfrac{2\ln x-5}{x^3}$.

$x^3>0$، إذن إشارة $g'(x)$ من إشارة $2\ln x-5$: $2\ln x-5>0\iff x>e^{\frac52}$.

$g$ متناقصة تماما على $\left]0;e^{\frac52}\right]$ ومتزايدة تماما على $\left[e^{\frac52};+\infty\right[$، و$g\left(e^{\frac52}\right)=-\dfrac12+\dfrac{2-\frac52}{e^5}=-\dfrac12-\dfrac{1}{2e^5}$.
\begin{TAB}
\tkzTabInit[lgt=2,espcl=3.5]{$x$ /0.8, $g'(x)$ /0.8, $g(x)$ /1.8}{$0$, $e^{\frac52}$, $+\infty$}
\tkzTabLine{d, -, z, +, }
\tkzTabVar{D+/$+\infty$, -/$-\frac12-\frac{1}{2e^5}$, +/$-\frac12$}
\end{TAB}
\QQ{3} \textbf{على $\left[e^{\frac52};+\infty\right[$:} $g$ متزايدة تماما نحو $-\frac12$، إذن $g(x)<-\frac12<0$: لا حلول.

\textbf{على $\left]0;e^{\frac52}\right]$:} $g$ مستمرة ومتناقصة تماما من $+\infty$ إلى $-\frac12-\frac{1}{2e^5}<0$، إذن المعادلة $g(x)=0$ تقبل حلا وحيدا $\alpha$.

$g(1{,}71)\approx0{,}0005>0$ و$g(1{,}72)\approx-0{,}007<0$، إذن \bx{1{,}71<\alpha<1{,}72}.

\textbf{الإشارة:} \bxt{$g(x)>0$ على $]0;\alpha[ ; g(x)<0$ على $]\alpha;+\infty[$} و$g(\alpha)=0$.

\textbf{الجزء الثاني:}
\QQ{1}\pa \textbf{عند $0$:} $-1+\ln x\to-\infty$ و$\dfrac1x\to+\infty$، إذن $\dfrac{-1+\ln x}{x}\to-\infty$، ومنه \bx{\lim_{x\to0^+}f(x)=-\infty} (المستقيم $x=0$ مقارب عمودي).

\textbf{عند $+\infty$:} $\dfrac{-1+\ln x}{x}=-\dfrac1x+\dfrac{\ln x}{x}\to0$ و$-\frac12x+2\to-\infty$، ومنه \bx{\lim_{x\to+\infty}f(x)=-\infty}.
\pb $f'(x)=-\dfrac12+\dfrac{\frac1x\cdot x-(-1+\ln x)}{x^2}=-\dfrac12+\dfrac{2-\ln x}{x^2}=g(x)$.

إذن $f$ متزايدة تماما على $]0;\alpha]$ ومتناقصة تماما على $[\alpha;+\infty[$، مع $f(\alpha)\approx0{,}87$.
\begin{TAB}
\tkzTabInit[lgt=2,espcl=3.5]{$x$ /0.8, $f'(x)$ /0.8, $f(x)$ /1.8}{$0$, $\alpha$, $+\infty$}
\tkzTabLine{d, +, z, -, }
\tkzTabVar{D-/$-\infty$, +/$f(\alpha)$, -/$-\infty$}
\end{TAB}
\QQ{2}\pa لدينا:
\[f(x)-\left(-\frac12x+2\right)=\frac{\ln x-1}{x}\qquad\text{و}\qquad\lim_{x\to+\infty}\frac{\ln x-1}{x}=0.\]
إذن المستقيم \bx{(\Delta): y=-\tfrac12x+2} مقارب مائل لـ $(C_f)$ بجوار $+\infty$.
\pb إشارة $f(x)-\left(-\frac12x+2\right)$ من إشارة $\ln x-1$ (لأن $x>0$):
\begin{itemize}
\item على $]0;e[$: $\ln x<1$: $(C_f)$ تحت $(\Delta)$.
\item على $]e;+\infty[$: $\ln x>1$: $(C_f)$ فوق $(\Delta)$.
\item عند $x=e$: $(C_f)$ يقطع $(\Delta)$ في النقطة $\left(e;2-\frac e2\right)\approx(2{,}72;0{,}64)$.
\end{itemize}
\QQ{3} الرسم:
\begin{GR}{-0.5}{7}{-3}{2.5}{1.5}
 \draw[cRem,thick,domain=-0.5:7] plot(\x,{-0.5*\x+2});
 \node[cRem,above right] at (0.3,1.85){$(\Delta)$};
 \draw[cDef,very thick,domain=0.33:7,samples=250] plot(\x,{-0.5*\x+2+(ln(\x)-1)/\x});
 \fill (1.7106,0.8739) circle (1.5pt);
 \draw[dashed] (1.7106,0)--(1.7106,0.8739);
 \node[below,font=\scriptsize] at (1.7106,0){$\alpha$};
 \fill (0.7771,0) circle (1.5pt) node[above left,font=\scriptsize]{$\beta$};
 \fill (4.2077,0) circle (1.5pt) node[above right,font=\scriptsize]{$\gamma$};
 \fill (2.718,0.641) circle (1.5pt);
 \node[cDef] at (5.5,-1.6){$(C_f)$};
 \Ticks{1,2,3,4,5,6}{-2,-1,1,2}
\end{GR}
\end{hal}

% ======================= part4 : pages 89-93 (تقني رياضي) =======================
\begin{tamrin}{بكالوريا 2018 -- الموضوع الثاني -- تقني رياضي}
\textbf{I)} $g$ الدالة العددية المعرفة على المجال $]0;1[$ كما يلي: $g(x)=2-x+\ln x$.

\textbf{1)}
\begin{sq}
\qa ادرس اتجاه تغير الدالة $g$ على المجال $]0;1[$.
\qb بيّن أن المعادلة $g(x)=0$ تقبل حلا وحيدا $\alpha$ حيث: $0{,}15<\alpha<0{,}16$.
\end{sq}
\textbf{2)} استنتج حسب قيم $x$ إشارة $g(x)$ على المجال $]0;1[$.

\textbf{II)} نعتبر الدالة العددية $f$ المعرفة على المجال $]1;+\infty[$ بـ: $f(x)=\dfrac{1-2x+\ln x}{x-1}$.

$(C_f)$ تمثيلها البياني في المستوي المنسوب إلى المعلم المتعامد والمتجانس $(O;\vec i,\vec j)$.

\textbf{1)} احسب $\displaystyle\lim_{x\to1^+}f(x)$ و $\displaystyle\lim_{x\to+\infty}f(x)$ (يمكن كتابة $f(x)$ على الشكل $f(x)=\dfrac{1-2x}{x-1}+\dfrac{\ln x}{x-1}$)، ثم فسّر النتيجتين بيانيا.

\textbf{2)}
\begin{sq}
\qa بيّن أنه من أجل كل عدد حقيقي $x$ من المجال $]1;+\infty[$: $f'(x)=\dfrac{g\left(\frac1x\right)}{(x-1)^2}$.
\qb بيّن أن الدالة $f$ متزايدة تماما على المجال $\left]1;\frac1\alpha\right]$ ومتناقصة تماما على المجال $\left[\frac1\alpha;+\infty\right[$، ثم شكّل جدول تغيراتها.
\end{sq}
\textbf{3)} ادرس الوضع النسبي للمنحنى $(C_f)$ والمستقيم $(\Delta)$ ذي المعادلة $y=-2$.

\textbf{4)} ارسم المستقيمين المقاربين والمنحنى $(C_f)$. (يعطى $f\left(\frac1\alpha\right)\approx-1{,}8$).

\textbf{5)} عيّن بيانيا قيم الوسيط الحقيقي $m$ حتى تقبل المعادلة $|f(x)|=m$ حلين متمايزين.
\end{tamrin}

\begin{molahada}
في نص المجلة كُتب في السؤال II-2-أ: $f(x)=\dfrac{g\left(\frac1x\right)}{(x-1)^2}$، والصواب $f'(x)=\dfrac{g\left(\frac1x\right)}{(x-1)^2}$ (المطلوب هو المشتقة)، وقد صُحّح في النص أعلاه.
\end{molahada}

\begin{hal}
\textbf{الجزء I}

\QQ{1}
\pa $g$ قابلة للاشتقاق على $]0;1[$ ولدينا: $g'(x)=-1+\dfrac1x=\dfrac{1-x}{x}$.

من أجل $0<x<1$: $1-x>0$ و $x>0$، إذن $g'(x)>0$، ومنه \bxt{$g$ متزايدة تماما على $]0;1[$}.

\pb النهايتان عند طرفي المجال: $\displaystyle\lim_{x\to0^+}g(x)=2-0+(-\infty)=-\infty$ و $\displaystyle\lim_{x\to1^-}g(x)=2-1+\ln1=1$.

الدالة $g$ مستمرة ومتزايدة تماما على $]0;1[$، وحسابيا:
\[g(0{,}15)=1{,}85+\ln0{,}15\approx-0{,}047<0,\qquad g(0{,}16)=1{,}84+\ln0{,}16\approx0{,}007>0.\]
إذن $g(0{,}15)\times g(0{,}16)<0$، وحسب مبرهنة القيم المتوسطة (مع الرتابة التامة) المعادلة $g(x)=0$ تقبل حلا وحيدا $\alpha$ حيث \bx{0{,}15<\alpha<0{,}16}.

\QQ{2} بما أن $g$ متزايدة تماما و $g(\alpha)=0$:
\begin{itemize}
\item إذا كان $0<x<\alpha$ فإن $g(x)<g(\alpha)$ أي $g(x)<0$؛
\item إذا كان $\alpha<x<1$ فإن $g(x)>0$؛ و $g(\alpha)=0$.
\end{itemize}
\begin{TAB}
\tkzTabInit[lgt=2,espcl=3]{$x$/1,$g(x)$/1}{$0$,$\alpha$,$1$}
\tkzTabLine{d,-,z,+,}
\end{TAB}

\textbf{الجزء II}

\QQ{1} \textbf{عند $1^+$:} $\displaystyle\lim_{x\to1^+}(1-2x+\ln x)=1-2+0=-1$ و $\displaystyle\lim_{x\to1^+}(x-1)=0^+$، إذن \bx{\lim_{x\to1^+}f(x)=-\infty}.

\textbf{عند $+\infty$:} $\displaystyle\lim_{x\to+\infty}\frac{1-2x}{x-1}=\lim_{x\to+\infty}\frac{-2x}{x}=-2$، و $\dfrac{\ln x}{x-1}=\dfrac{\ln x}{x}\times\dfrac{x}{x-1}$ مع $\displaystyle\lim_{x\to+\infty}\frac{\ln x}{x}=0$ و $\displaystyle\lim_{x\to+\infty}\frac{x}{x-1}=1$، إذن \bx{\lim_{x\to+\infty}f(x)=-2}.

\textbf{التفسير البياني:} المستقيم ذو المعادلة $x=1$ مقارب عمودي لـ $(C_f)$، والمستقيم $(\Delta):y=-2$ مقارب أفقي لـ $(C_f)$ بجوار $+\infty$.

\QQ{2}
\pa $f$ قابلة للاشتقاق على $]1;+\infty[$ (حاصل قسمة دالتين قابلتين للاشتقاق، والمقام لا ينعدم):
\begin{align*}
f'(x)&=\frac{\left(-2+\frac1x\right)(x-1)-(1-2x+\ln x)\times1}{(x-1)^2}
=\frac{-2x+2+1-\frac1x-1+2x-\ln x}{(x-1)^2}\\
&=\frac{2-\frac1x-\ln x}{(x-1)^2}.
\end{align*}
ومن جهة أخرى: $g\left(\frac1x\right)=2-\frac1x+\ln\frac1x=2-\frac1x-\ln x$. إذن \bx{f'(x)=\frac{g\left(\frac1x\right)}{(x-1)^2}}.

\pb لأن $(x-1)^2>0$ فإن إشارة $f'(x)$ هي إشارة $g\left(\frac1x\right)$. لما $x>1$ يكون $0<\frac1x<1$، وحسب الجزء I:
\begin{itemize}
\item $g\left(\frac1x\right)>0\iff\frac1x>\alpha\iff x<\frac1\alpha$، إذن $f'(x)>0$ على $\left]1;\frac1\alpha\right[$؛
\item $g\left(\frac1x\right)<0\iff\frac1x<\alpha\iff x>\frac1\alpha$، إذن $f'(x)<0$ على $\left]\frac1\alpha;+\infty\right[$؛ و $f'\left(\frac1\alpha\right)=0$.
\end{itemize}
ومنه $f$ متزايدة تماما على $\left]1;\frac1\alpha\right]$ ومتناقصة تماما على $\left[\frac1\alpha;+\infty\right[$.

\textbf{حساب $f\left(\frac1\alpha\right)$:} من $g(\alpha)=0$ نجد $\ln\alpha=\alpha-2$، أي $\ln\frac1\alpha=2-\alpha$، ومنه
\[f\left(\tfrac1\alpha\right)=\frac{1-\frac2\alpha+2-\alpha}{\frac1\alpha-1}=\frac{-\alpha^2+3\alpha-2}{1-\alpha}=\frac{-(\alpha-1)(\alpha-2)}{1-\alpha}=\alpha-2\approx-1{,}84.\]
\begin{TAB}
\tkzTabInit[lgt=2,espcl=3.5]{$x$/1,$f'(x)$/1,$f(x)$/2}{$1$,$\frac1\alpha$,$+\infty$}
\tkzTabLine{d,+,z,-,}
\tkzTabVar{D-/$-\infty$,+/$f\left(\frac1\alpha\right)$,-/$-2$}
\end{TAB}

\QQ{3} ندرس إشارة الفرق:
\[f(x)-(-2)=\frac{1-2x+\ln x+2(x-1)}{x-1}=\frac{\ln x-1}{x-1}.\]
بما أن $x-1>0$ فإن الإشارة هي إشارة $\ln x-1$، و $\ln x-1>0\iff x>e$.
\begin{itemize}
\item على $]1;e[$: $f(x)+2<0$، إذن $(C_f)$ تحت $(\Delta)$؛
\item على $]e;+\infty[$: $f(x)+2>0$، إذن $(C_f)$ فوق $(\Delta)$؛
\item $(C_f)$ يقطع $(\Delta)$ في النقطة \bx{(e;-2)}.
\end{itemize}

\QQ{4} الرسم:
\begin{GR}{-0.5}{10}{-6}{1}{1.2}
\Ticks{1,2,3,4,5,6,7,8,9}{-5,-4,-3,-1}
\node[left,font=\scriptsize] at (0,-2){$-2$};
\draw[cThm,dashed,thick] (1,-6)--(1,1);
\draw[cRem!80!black,dashed,thick] (0,-2)--(10,-2) node[pos=0.9,below]{$(\Delta)$};
\draw[cDef,very thick,domain=1.12:10,samples=200,smooth,variable=\x] plot (\x,{(1-2*\x+ln(\x))/(\x-1)});
\fill (2.718,-2) circle (1.5pt) node[below]{$e$};
\draw[densely dotted] (6.305,0)--(6.305,-1.841);
\fill (6.305,-1.841) circle (1.5pt);
\node[below,font=\scriptsize] at (6.305,0){$\frac1\alpha$};
\node[cDef] at (2.2,-4.5){$(C_f)$};
\end{GR}

\QQ{5} من الجدول: $f(x)<0$ على $]1;+\infty[$ (أكبر قيمة لها $\alpha-2<0$)، إذن $|f(x)|=-f(x)$، والمنحنى $(C_{|f|})$ هو نظير $(C_f)$ بالنسبة إلى محور الفواصل.

الدالة $|f|$ متناقصة تماما من $+\infty$ إلى $2-\alpha\approx1{,}84$ على $\left]1;\frac1\alpha\right]$، ثم متزايدة تماما من $2-\alpha$ نحو $2$ (دون أن تبلغها) على $\left[\frac1\alpha;+\infty\right[$. حلول المعادلة $|f(x)|=m$ هي فواصل نقط تقاطع $(C_{|f|})$ مع المستقيم $y=m$:
\begin{itemize}
\item $m<2-\alpha$: لا حلول؛ $m=2-\alpha$: حل وحيد $\frac1\alpha$؛
\item $2-\alpha<m<2$: حلان متمايزان؛
\item $m\geqslant2$: حل وحيد.
\end{itemize}
إذن المعادلة تقبل حلين متمايزين من أجل \bx{m\in\left]2-\alpha;2\right[} أي تقريبا $m\in\left]1{,}8;2\right[$.
\end{hal}

% ------------------------------------------------------------------
\begin{tamrin}{بكالوريا 2019 -- الموضوع الثاني -- تقني رياضي}
\textbf{I)} $g$ الدالة العددية المعرفة والمتزايدة تماما على $]0;+\infty[$ بـ: $g(x)=(x+1)(x+e)-e(x\ln x)$.

احسب $\displaystyle\lim_{x\to0^+}g(x)$ ثم استنتج إشارة $g(x)$ على المجال $]0;+\infty[$.

\textbf{II)} $f$ الدالة العددية المعرفة على $]0;+\infty[$ كما يلي: $f(x)=\ln(x+1)+\dfrac{e\ln x}{x+1}$.

$(C_f)$ تمثيلها البياني في المستوي المنسوب إلى المعلم المتعامد والمتجانس $(O;\vec i,\vec j)$.

\textbf{1)}
\begin{sq}
\qa احسب $\displaystyle\lim_{x\to0^+}f(x)$ ثم بيّن أن: $\displaystyle\lim_{x\to+\infty}f(x)=+\infty$.
\qb بيّن أنه من أجل كل $x$ من $]0;+\infty[$: $f'(x)=\dfrac{g(x)}{x(x+1)^2}$.
\qc استنتج اتجاه تغير الدالة $f$، ثم شكّل جدول تغيراتها.
\end{sq}
\textbf{2)} اكتب معادلة لـ $(T)$ مماس المنحنى $(C_f)$ في النقطة ذات الفاصلة $1$.

\textbf{3)}
\begin{sq}
\qa بيّن أن المنحنى $(C_f)$ يقطع حامل محور الفواصل في نقطة وحيدة $A$ فاصلتها $\alpha$.
\qb تحقق أن $0{,}7<\alpha<0{,}8$.
\end{sq}
\textbf{4)} $(\Gamma)$ التمثيل البياني للدالة $x\mapsto\ln(x+1)$ على المجال $[0;+\infty[$.
\begin{sq}
\qa احسب $\displaystyle\lim_{x\to+\infty}\left[f(x)-\ln(x+1)\right]$ ثم فسّر النتيجة بيانيا.
\qb ادرس الوضع النسبي للمنحنيين $(C_f)$ و $(\Gamma)$.
\qc ارسم المماس $(T)$ و $(\Gamma)$ و $(C_f)$.
\end{sq}
\textbf{5)} $m$ وسيط حقيقي، عيّن قيم $m$ بحيث تقبل المعادلة $f(x)=\dfrac{1+e}{2}x-m$ حلين متمايزين.

\textbf{6)} نقبل أنه من أجل كل $x$ من المجال $]1;+\infty[$: $\ln x<x+1$.

بيّن أنه من أجل كل $x$ من المجال $]1;+\infty[$: $\ln2<f(x)<e+\ln(x+1)$.
\end{tamrin}

\begin{hal}
\textbf{الجزء I}

لدينا $\displaystyle\lim_{x\to0^+}x\ln x=0$ و $\displaystyle\lim_{x\to0^+}(x+1)(x+e)=1\times e=e$، إذن \bx{\lim_{x\to0^+}g(x)=e}.

الدالة $g$ متزايدة تماما على $]0;+\infty[$، إذن من أجل كل $x>0$: $g(x)>\displaystyle\lim_{t\to0^+}g(t)=e>0$. ومنه \bx{g(x)>0} على $]0;+\infty[$.

\textbf{الجزء II}

\QQ{1}
\pa $\displaystyle\lim_{x\to0^+}\ln(x+1)=0$ و $\displaystyle\lim_{x\to0^+}\frac{e\ln x}{x+1}=\frac{e\times(-\infty)}{1}=-\infty$، إذن \bx{\lim_{x\to0^+}f(x)=-\infty} (المستقيم $x=0$ مقارب عمودي).

$\displaystyle\lim_{x\to+\infty}\ln(x+1)=+\infty$، و $\dfrac{e\ln x}{x+1}=e\times\dfrac{\ln x}{x}\times\dfrac{x}{x+1}$ تؤول إلى $e\times0\times1=0$، إذن \bx{\lim_{x\to+\infty}f(x)=+\infty}.

\pb $f$ قابلة للاشتقاق على $]0;+\infty[$:
\begin{align*}
f'(x)&=\frac{1}{x+1}+e\cdot\frac{\frac1x(x+1)-\ln x}{(x+1)^2}
=\frac{x(x+1)+e(x+1)-ex\ln x}{x(x+1)^2}\\
&=\frac{(x+1)(x+e)-ex\ln x}{x(x+1)^2}=\frac{g(x)}{x(x+1)^2}.
\end{align*}
(ضربنا البسط والمقام في $x(x+1)^2$ بعد توحيد المقامات.)

\pc لأن $x(x+1)^2>0$ و $g(x)>0$ على $]0;+\infty[$، فإن $f'(x)>0$، إذن \bxt{$f$ متزايدة تماما على $]0;+\infty[$}.
\begin{TAB}
\tkzTabInit[lgt=2,espcl=5]{$x$/1,$f'(x)$/1,$f(x)$/2}{$0$,$+\infty$}
\tkzTabLine{d,+,}
\tkzTabVar{D-/$-\infty$,+/$+\infty$}
\end{TAB}

\QQ{2} $f(1)=\ln2+\dfrac{e\ln1}{2}=\ln2$ و $f'(1)=\dfrac{g(1)}{1\times4}=\dfrac{2(1+e)}{4}=\dfrac{1+e}{2}$.

$(T):y=f'(1)(x-1)+f(1)$ أي \bx{(T):\ y=\frac{1+e}{2}x-\frac{1+e}{2}+\ln2}.

\QQ{3}
\pa $f$ مستمرة ومتزايدة تماما على $]0;+\infty[$ وتأخذ قيمها في $]-\infty;+\infty[$ الذي يشمل $0$، إذن (مبرهنة القيم المتوسطة مع الرتابة) المعادلة $f(x)=0$ تقبل حلا وحيدا $\alpha$، أي أن $(C_f)$ يقطع حامل محور الفواصل في نقطة وحيدة $A(\alpha;0)$.

\pb $f(0{,}7)\approx-0{,}04<0$ و $f(0{,}8)\approx0{,}25>0$، إذن \bx{0{,}7<\alpha<0{,}8}.

\QQ{4}
\pa $f(x)-\ln(x+1)=\dfrac{e\ln x}{x+1}$، وقد رأينا أن نهايتها عند $+\infty$ هي $0$، إذن \bx{\lim_{x\to+\infty}\left[f(x)-\ln(x+1)\right]=0}.

التفسير: المنحنيان $(C_f)$ و $(\Gamma)$ متقاربان بجوار $+\infty$.

\pb إشارة $f(x)-\ln(x+1)=\dfrac{e\ln x}{x+1}$ هي إشارة $\ln x$ (لأن $e>0$ و $x+1>0$):
\begin{itemize}
\item على $]0;1[$: $\ln x<0$، إذن $(C_f)$ تحت $(\Gamma)$؛
\item على $]1;+\infty[$: $\ln x>0$، إذن $(C_f)$ فوق $(\Gamma)$؛
\item يتقاطعان في النقطة $(1;\ln2)$.
\end{itemize}

\pc الرسم ($(T)$ معامل توجيهه $\frac{1+e}{2}\approx1{,}86$ ويمر بـ $(1;\ln2)$):
\begin{GR}{-0.5}{6}{-4}{4}{1.2}
\Ticks{1,2,3,4,5}{-3,-2,-1,1,2,3}
\draw[cRem!80!black,thick,domain=0:6,samples=100,smooth,variable=\x] plot (\x,{ln(\x+1)}) node[pos=0.95,below]{$(\Gamma)$};
\draw[cThm,thick,domain=-0.5:4.5,variable=\x] plot (\x,{1.8591*(\x-1)+0.6931}) node[pos=0.83,left]{$(T)$};
\draw[cDef,very thick,domain=0.07:6,samples=200,smooth,variable=\x] plot (\x,{ln(\x+1)+2.71828*ln(\x)/(\x+1)});
\node[cDef] at (5.4,3.4){$(C_f)$};
\fill (1,0.6931) circle (1.5pt);
\fill (0.7125,0) circle (1.5pt) node[below right]{$\alpha$};
\end{GR}

\QQ{5} حلول المعادلة $f(x)=\dfrac{1+e}{2}x-m$ هي فواصل نقط تقاطع $(C_f)$ مع المستقيم $(D_m):y=\dfrac{1+e}{2}x-m$ الموازي لـ $(T)$ (نفس معامل التوجيه).

المستقيم $(D_m)$ ينطبق على $(T)$ لما $-m=-\dfrac{1+e}{2}+\ln2$ أي $m=\dfrac{1+e}{2}-\ln2$. ومن الرسم المنحنى $(C_f)$ يقع تحت مماسه $(T)$ (يمسّه في النقطة ذات الفاصلة $1$ فقط)، إذن:
\begin{itemize}
\item إذا كان $m<\frac{1+e}{2}-\ln2$: $(D_m)$ فوق $(T)$، لا توجد حلول؛
\item إذا كان $m=\frac{1+e}{2}-\ln2$: حل وحيد هو $1$؛
\item إذا كان $m>\frac{1+e}{2}-\ln2$: $(D_m)$ تحت $(T)$ ويقطع $(C_f)$ في نقطتين، إذن حلان متمايزان.
\end{itemize}
النتيجة: \bx{m>\frac{1+e}{2}-\ln2} (أي $m>1{,}17$ تقريبا).

\QQ{6} ليكن $x>1$.
\begin{itemize}
\item $f$ متزايدة تماما، إذن $f(x)>f(1)=\ln2$.
\item لدينا $\ln x<x+1$ و $\dfrac{e}{x+1}>0$، إذن $\dfrac{e\ln x}{x+1}<e$، وبإضافة $\ln(x+1)$ للطرفين: $f(x)<e+\ln(x+1)$.
\end{itemize}
إذن من أجل كل $x\in]1;+\infty[$: \bx{\ln2<f(x)<e+\ln(x+1)}.
\end{hal}

% ------------------------------------------------------------------
\begin{tamrin}{بكالوريا 2020 -- الموضوع الأول -- تقني رياضي}
\textbf{I)} $g$ الدالة العددية المعرفة على $]0;+\infty[$ بـ: $g(x)=-1+x+2\ln x$.

\textbf{1)} ادرس اتجاه تغيرات الدالة $g$.

\textbf{2)} احسب $g(1)$ ثم استنتج إشارة $g(x)$ على المجال $]0;+\infty[$.

\textbf{II)} $f$ الدالة العددية المعرفة على المجال $]0;+\infty[$ كما يلي: $f(x)=\dfrac{-1+(x-2)\ln x}{x}$.

$(C_f)$ تمثيلها البياني في المستوي المنسوب إلى المعلم المتعامد والمتجانس $(O;\vec i,\vec j)$.

\textbf{1)}
\begin{sq}
\qa احسب $\displaystyle\lim_{x\to0^+}f(x)$ ثم فسّر النتيجة بيانيا.
\qb احسب $\displaystyle\lim_{x\to+\infty}f(x)$.
\end{sq}
\textbf{2)}
\begin{sq}
\qa بيّن أنه من أجل كل $x$ من $]0;+\infty[$: $f'(x)=\dfrac{g(x)}{x^2}$.
\qb عيّن اتجاه تغير الدالة $f$، ثم شكّل جدول تغيراتها.
\end{sq}
\textbf{3)} ليكن $(\Gamma)$ المنحنى البياني للدالة $x\mapsto\ln x$ على المجال $]0;+\infty[$.
\begin{sq}
\qa احسب $\displaystyle\lim_{x\to+\infty}\left[f(x)-\ln x\right]$ ثم فسّر النتيجة بيانيا.
\qb ادرس وضعية المنحنى $(C_f)$ بالنسبة للمنحنى $(\Gamma)$.
\end{sq}
\textbf{4)} بيّن أن المنحنى $(C_f)$ يقطع حامل محور الفواصل في نقطتين فاصلتاهما $\alpha$ و $\beta$، ثم تحقق أن: $0{,}5<\alpha<0{,}6$ و $2{,}9<\beta<3$.

\textbf{5)} ارسم $(\Gamma)$ ثم $(C_f)$.
\end{tamrin}

\begin{hal}
\textbf{الجزء I}

\QQ{1} $g'(x)=1+\dfrac2x>0$ من أجل كل $x>0$، إذن \bxt{$g$ متزايدة تماما على $]0;+\infty[$}.

للتوضيح: $\displaystyle\lim_{x\to0^+}g(x)=-\infty$ و $\displaystyle\lim_{x\to+\infty}g(x)=+\infty$.

\QQ{2} $g(1)=-1+1+2\ln1=0$. وبما أن $g$ متزايدة تماما:
\begin{TAB}
\tkzTabInit[lgt=2,espcl=3]{$x$/1,$g(x)$/1}{$0$,$1$,$+\infty$}
\tkzTabLine{d,-,z,+,}
\end{TAB}
$g(x)<0$ على $]0;1[$، $g(x)>0$ على $]1;+\infty[$، و $g(1)=0$.

\textbf{الجزء II}

\QQ{1}
\pa $\displaystyle\lim_{x\to0^+}(x-2)\ln x=(-2)\times(-\infty)=+\infty$، إذن البسط يؤول إلى $+\infty$، والمقام $x\to0^+$، ومنه \bx{\lim_{x\to0^+}f(x)=+\infty}.

التفسير: المستقيم ذو المعادلة $x=0$ (محور التراتيب) مقارب عمودي لـ $(C_f)$.

\pb نكتب $f(x)=-\dfrac1x+\left(1-\dfrac2x\right)\ln x$. لدينا $\displaystyle\lim_{x\to+\infty}\left(-\frac1x\right)=0$ و $\displaystyle\lim_{x\to+\infty}\left(1-\frac2x\right)=1$ و $\displaystyle\lim_{x\to+\infty}\ln x=+\infty$، إذن \bx{\lim_{x\to+\infty}f(x)=+\infty}.

\QQ{2}
\pa $f$ قابلة للاشتقاق على $]0;+\infty[$:
\begin{align*}f'(x)&=\frac{\left(\ln x+\frac{x-2}{x}\right)x-\left(-1+(x-2)\ln x\right)}{x^2}
=\frac{x\ln x+x-2+1-x\ln x+2\ln x}{x^2}\\&=\frac{-1+x+2\ln x}{x^2}=\frac{g(x)}{x^2}.\end{align*}
\pb $x^2>0$ إذن إشارة $f'(x)$ هي إشارة $g(x)$: $f$ \bxt{متناقصة تماما على $]0;1]$} و \bxt{متزايدة تماما على $[1;+\infty[$}.

$f(1)=\dfrac{-1+(-1)\times0}{1}=-1$.
\begin{TAB}
\tkzTabInit[lgt=2,espcl=3.5]{$x$/1,$f'(x)$/1,$f(x)$/2}{$0$,$1$,$+\infty$}
\tkzTabLine{d,-,z,+,}
\tkzTabVar{D+/$+\infty$,-/$-1$,+/$+\infty$}
\end{TAB}

\QQ{3}
\pa $f(x)-\ln x=\dfrac{-1+(x-2)\ln x-x\ln x}{x}=\dfrac{-1-2\ln x}{x}=-\dfrac1x-2\,\dfrac{\ln x}{x}$.

بما أن $\displaystyle\lim_{x\to+\infty}\frac1x=0$ و $\displaystyle\lim_{x\to+\infty}\frac{\ln x}{x}=0$، فإن \bx{\lim_{x\to+\infty}\left[f(x)-\ln x\right]=0}.

التفسير: المنحنيان $(C_f)$ و $(\Gamma)$ متقاربان بجوار $+\infty$.

\pb إشارة $f(x)-\ln x=\dfrac{-1-2\ln x}{x}$ هي إشارة $-1-2\ln x$ (لأن $x>0$):

$-1-2\ln x>0\iff\ln x<-\dfrac12\iff x<e^{-\frac12}=\dfrac{1}{\sqrt e}$.
\begin{itemize}
\item على $\left]0;\frac1{\sqrt e}\right[$: $(C_f)$ فوق $(\Gamma)$؛
\item على $\left]\frac1{\sqrt e};+\infty\right[$: $(C_f)$ تحت $(\Gamma)$؛
\item يتقاطعان في النقطة \bx{\left(\tfrac{1}{\sqrt e};-\tfrac12\right)} (لأن $\ln e^{-\frac12}=-\frac12$).
\end{itemize}

\QQ{4} \textbf{على $]0;1]$:} $f$ مستمرة ومتناقصة تماما من $+\infty$ إلى $-1$، و $0\in]-1;+\infty[$، إذن المعادلة $f(x)=0$ تقبل حلا وحيدا $\alpha$ في $]0;1[$.

\textbf{على $[1;+\infty[$:} $f$ مستمرة ومتزايدة تماما من $-1$ إلى $+\infty$، إذن تقبل حلا وحيدا $\beta$ في $]1;+\infty[$.

إذن $(C_f)$ يقطع حامل محور الفواصل في نقطتين بالضبط فاصلتاهما $\alpha$ و $\beta$.

\textbf{التحقق:} $f(0{,}5)\approx0{,}08>0$ و $f(0{,}6)\approx-0{,}47<0$، إذن \bx{0{,}5<\alpha<0{,}6}.

$f(2{,}9)\approx-0{,}014<0$ و $f(3)=\dfrac{-1+\ln3}{3}\approx0{,}03>0$، إذن \bx{2{,}9<\beta<3}.

\QQ{5} الرسم:
\begin{GR}{-0.5}{7}{-3}{4}{1.2}
\Ticks{1,2,3,4,5,6}{-2,-1,1,2,3}
\draw[cRem!80!black,thick,domain=0.05:7,samples=150,smooth,variable=\x] plot (\x,{ln(\x)}) node[pos=0.93,above]{$(\Gamma)$};
\draw[cDef,very thick,domain=0.12:7,samples=200,smooth,variable=\x] plot (\x,{(-1+(\x-2)*ln(\x))/\x});
\node[cDef] at (6.3,0.9){$(C_f)$};
\fill (1,-1) circle (1.5pt);
\fill (0.6065,-0.5) circle (1.5pt);
\fill (0.511,0) circle (1.5pt) node[above right]{$\alpha$};
\fill (2.93,0) circle (1.5pt) node[above left]{$\beta$};
\end{GR}
\end{hal}

% ------------------------------------------------------------------
\begin{tamrin}{بكالوريا 2021 -- الموضوع الثاني -- تقني رياضي}
\textbf{I)} الدالة العددية $g$ معرفة على المجال $]0;+\infty[$ بـ: $g(x)=2\ln x-1-\dfrac{1}{x^2}$.

\textbf{1)} بيّن أن الدالة $g$ متزايدة تماما على المجال $]0;+\infty[$.

\textbf{2)}
\begin{sq}
\qa بيّن أن المعادلة $g(x)=0$ تقبل حلا وحيدا $\alpha$ حيث: $1{,}89<\alpha<1{,}90$.
\qb استنتج حسب قيم العدد الحقيقي الموجب تماما $x$ إشارة $g(x)$.
\end{sq}
\textbf{II)} الدالة العددية $f$ معرفة على $]0;+\infty[$ بـ: $f(x)=-x-2+\dfrac{3+2\ln x}{x}$.

$(C_f)$ تمثيلها البياني في المستوي المنسوب إلى المعلم المتعامد والمتجانس $(O;\vec i,\vec j)$. الوحدة $2\,\text{cm}$.

\textbf{1)}
\begin{sq}
\qa احسب $\displaystyle\lim_{x\to0^+}f(x)$ ثم فسّر النتيجة هندسيا.
\qb احسب $\displaystyle\lim_{x\to+\infty}f(x)$.
\end{sq}
\textbf{2)}
\begin{sq}
\qa بيّن أنه من أجل كل $x$ من $]0;+\infty[$: $f'(x)=\dfrac{1}{x^2}\,g\left(\dfrac1x\right)$.
\qb بيّن أن الدالة $f$ متزايدة تماما على المجال $\left]0;\frac1\alpha\right]$ ومتناقصة تماما على المجال $\left[\frac1\alpha;+\infty\right[$.
\qc شكّل جدول تغيرات الدالة $f$.
\end{sq}
\textbf{3)}
\begin{sq}
\qa احسب $\displaystyle\lim_{x\to+\infty}\left[f(x)-(-x-2)\right]$ ثم استنتج أن $(C_f)$ يقبل مستقيما مقاربا $(\Delta)$ يطلب كتابة معادلة له.
\qb ادرس وضعية المنحنى $(C_f)$ بالنسبة لـ $(\Delta)$.
\end{sq}
\textbf{4)} بيّن أن المنحنى $(C_f)$ يقبل نقطة انعطاف $A$ فاصلتها $1$ ثم اكتب معادلة لـ $(T)$ مماس $(C_f)$ عند $A$.

\textbf{5)} ارسم $(T)$ و $(\Delta)$ و $(C_f)$. (نأخذ: $\frac1\alpha\approx0{,}53$ و $f\left(\frac1\alpha\right)\approx0{,}73$).

\textbf{6)} الدالة $h$ معرفة على $\R^*$ بـ: $h(x)=|x|+2-\dfrac{3+\ln(x^2)}{|x|}$، و $(C_h)$ تمثيلها البياني في المعلم السابق.
\begin{sq}
\qa بيّن أن الدالة $h$ زوجية.
\qb تحقق أنه من أجل كل عدد حقيقي $x$ من المجال $]0;+\infty[$: $h(x)=-f(x)$.
\qc اشرح كيفية رسم $(C_h)$ انطلاقا من $(C_f)$ ثم ارسمه.
\end{sq}
\end{tamrin}

\begin{hal}
\textbf{الجزء I}

\QQ{1} $g'(x)=\dfrac2x+\dfrac{2}{x^3}=\dfrac{2(x^2+1)}{x^3}>0$ من أجل كل $x>0$، إذن \bxt{$g$ متزايدة تماما على $]0;+\infty[$}.

\QQ{2}
\pa $g$ مستمرة ومتزايدة تماما على $]0;+\infty[$، و
\[g(1{,}89)\approx-0{,}0068<0,\qquad g(1{,}90)\approx0{,}0067>0.\]
حسب مبرهنة القيم المتوسطة، المعادلة $g(x)=0$ تقبل حلا وحيدا $\alpha$ حيث \bx{1{,}89<\alpha<1{,}90}.

\pb لأن $g$ متزايدة تماما و $g(\alpha)=0$:
\begin{TAB}
\tkzTabInit[lgt=2,espcl=3]{$x$/1,$g(x)$/1}{$0$,$\alpha$,$+\infty$}
\tkzTabLine{d,-,z,+,}
\end{TAB}

\textbf{الجزء II}

\QQ{1}
\pa $\displaystyle\lim_{x\to0^+}(3+2\ln x)=-\infty$ و $\displaystyle\lim_{x\to0^+}x=0^+$، إذن $\displaystyle\lim_{x\to0^+}\frac{3+2\ln x}{x}=-\infty$، ومنه \bx{\lim_{x\to0^+}f(x)=-\infty}.

التفسير الهندسي: المستقيم ذو المعادلة $x=0$ مقارب عمودي لـ $(C_f)$.

\pb $\dfrac{3+2\ln x}{x}=\dfrac3x+2\,\dfrac{\ln x}{x}\to0$ لما $x\to+\infty$، و $\displaystyle\lim_{x\to+\infty}(-x-2)=-\infty$، إذن \bx{\lim_{x\to+\infty}f(x)=-\infty}.

\QQ{2}
\pa $f'(x)=-1+\dfrac{\frac2x\cdot x-(3+2\ln x)}{x^2}=-1+\dfrac{-1-2\ln x}{x^2}=\dfrac{-x^2-1-2\ln x}{x^2}$.

ومن جهة أخرى: $g\left(\frac1x\right)=2\ln\frac1x-1-x^2=-2\ln x-1-x^2$. إذن \bx{f'(x)=\frac{1}{x^2}\,g\left(\frac1x\right)}.

\pb إشارة $f'(x)$ هي إشارة $g\left(\frac1x\right)$ لأن $\frac1{x^2}>0$:
\begin{itemize}
\item $0<x<\frac1\alpha\iff\frac1x>\alpha\iff g\left(\frac1x\right)>0$: $f$ متزايدة تماما على $\left]0;\frac1\alpha\right]$؛
\item $x>\frac1\alpha\iff\frac1x<\alpha\iff g\left(\frac1x\right)<0$: $f$ متناقصة تماما على $\left[\frac1\alpha;+\infty\right[$.
\end{itemize}

\pc
\begin{TAB}
\tkzTabInit[lgt=2,espcl=3.5]{$x$/1,$f'(x)$/1,$f(x)$/2}{$0$,$\frac1\alpha$,$+\infty$}
\tkzTabLine{d,+,z,-,}
\tkzTabVar{D-/$-\infty$,+/$f\left(\frac1\alpha\right)$,-/$-\infty$}
\end{TAB}
(بالاستعانة بـ $2\ln\alpha=1+\frac1{\alpha^2}$ نجد $f\left(\frac1\alpha\right)=2\alpha-2-\frac2\alpha\approx0{,}73$.)

\QQ{3}
\pa $f(x)-(-x-2)=\dfrac{3+2\ln x}{x}=\dfrac3x+2\,\dfrac{\ln x}{x}$، إذن \bx{\lim_{x\to+\infty}\left[f(x)-(-x-2)\right]=0}.

ومنه المستقيم \bx{(\Delta):\ y=-x-2} مقارب مائل لـ $(C_f)$ بجوار $+\infty$.

\pb إشارة الفرق $\dfrac{3+2\ln x}{x}$ هي إشارة $3+2\ln x$: $3+2\ln x>0\iff x>e^{-\frac32}$.
\begin{itemize}
\item على $\left]0;e^{-\frac32}\right[$: $(C_f)$ تحت $(\Delta)$؛
\item على $\left]e^{-\frac32};+\infty\right[$: $(C_f)$ فوق $(\Delta)$؛
\item يتقاطعان في النقطة ذات الفاصلة $e^{-\frac32}\approx0{,}22$.
\end{itemize}

\QQ{4} $f'(x)=-1-\dfrac{1}{x^2}-\dfrac{2\ln x}{x^2}$، ومنه
\[f''(x)=\frac{2}{x^3}-2\cdot\frac{\frac1x\cdot x^2-2x\ln x}{x^4}=\frac{2}{x^3}-\frac{2-4\ln x}{x^3}=\frac{4\ln x}{x^3}.\]
إشارة $f''(x)$ هي إشارة $\ln x$: سالبة على $]0;1[$ وموجبة على $]1;+\infty[$؛ $f''$ تنعدم عند $1$ مغيّرة إشارتها، إذن النقطة $A(1;f(1))$ نقطة انعطاف، مع $f(1)=-1-2+3=0$، أي \bx{A(1;0)}.

$f'(1)=-1-1-0=-2$، إذن $(T):y=-2(x-1)+0$، أي \bx{(T):\ y=-2x+2}.

\QQ{5} الرسم:
\begin{GR}{-0.5}{5}{-6}{2}{1.2}
\Ticks{1,2,3,4}{-5,-4,-3,-2,-1,1}
\draw[cRem!80!black,thick,domain=0:4.5,variable=\x] plot (\x,{-\x-2}); \node[cRem!80!black] at (1.2,-3.8){$(\Delta)$};
\draw[cThm,thick,domain=0:4,variable=\x] plot (\x,{-2*\x+2}); \node[cThm] at (2.3,-3){$(T)$};
\draw[cDef,very thick,domain=0.1:5,samples=200,smooth,variable=\x] plot (\x,{-\x-2+(3+2*ln(\x))/\x});
\node[cDef] at (4.3,-3.9){$(C_f)$};
\fill (1,0) circle (1.5pt) node[above right]{$A$};
\fill (0.5277,0.7347) circle (1.5pt);
\draw[densely dotted] (0.5277,0)--(0.5277,0.7347);
\end{GR}

\QQ{6}
\pa $\R^*$ متناظر بالنسبة إلى $0$، ومن أجل كل $x\in\R^*$: $|-x|=|x|$ و $(-x)^2=x^2$، إذن
\[h(-x)=|-x|+2-\frac{3+\ln((-x)^2)}{|-x|}=|x|+2-\frac{3+\ln(x^2)}{|x|}=h(x).\]
إذن \bxt{$h$ زوجية}.

\pb من أجل $x>0$: $|x|=x$ و $\ln(x^2)=2\ln x$، إذن
\[h(x)=x+2-\frac{3+2\ln x}{x}=-\left(-x-2+\frac{3+2\ln x}{x}\right)=-f(x).\]

\pc \textbf{على $]0;+\infty[$:} $(C_h)$ هو نظير $(C_f)$ بالنسبة إلى محور الفواصل (لأن $h(x)=-f(x)$).

\textbf{على $]-\infty;0[$:} $(C_h)$ هو نظير الجزء السابق بالنسبة إلى محور التراتيب (لأن $h$ زوجية).
\begin{GR}{-5}{5}{-2}{7}{0.9}
\Ticks{-4,-3,-2,-1,1,2,3,4}{-1,1,2,3,4,5,6}
\draw[cThm,very thick,domain=0.1:5,samples=200,smooth,variable=\x] plot (\x,{\x+2-(3+2*ln(\x))/\x});
\draw[cThm,very thick,domain=-5:-0.1,samples=200,smooth,variable=\x] plot (\x,{-\x+2-(3+2*ln(-\x))/(-\x)});
\node[cThm] at (3.8,3.2){$(C_h)$};
\end{GR}
\end{hal}

% ------------------------------------------------------------------
\begin{tamrin}{بكالوريا 2022 -- الموضوع الأول -- تقني رياضي}
$f$ الدالة العددية المعرفة على المجال $]0;+\infty[$ بـ: $f(x)=1+(x-1)\ln x$.

$(C_f)$ تمثيلها البياني في المستوي المنسوب إلى المعلم المتعامد والمتجانس $(O;\vec i,\vec j)$. الوحدة $2\,\text{cm}$.

\textbf{1)}
\begin{sq}
\qa ادرس حسب قيم العدد الحقيقي $x$ الموجب تماما إشارة كلّ من $\dfrac{x-1}{x}$ و $\ln x$.
\qb استنتج حسب قيم العدد الحقيقي $x$ الموجب تماما إشارة $\ln x+\dfrac{x-1}{x}$.
\end{sq}
\textbf{2)}
\begin{sq}
\qa احسب $\displaystyle\lim_{x\to0^+}f(x)$ و $\displaystyle\lim_{x\to+\infty}f(x)$.
\qb ادرس اتجاه تغير الدالة $f$، ثم شكّل جدول تغيراتها.
\end{sq}
\textbf{3)} $h$ الدالة العددية المعرفة على المجال $]0;+\infty[$ كما يلي: $h(x)=x-2+\ln x$.
\begin{sq}
\qa بيّن أن الدالة $h$ متزايدة تماما على المجال $]0;+\infty[$.
\qb برهن أن المعادلة $h(x)=0$ تقبل حلا وحيدا $\alpha$ حيث $1{,}5<\alpha<1{,}6$، ثم بيّن أن: $\ln(\alpha)=2-\alpha$.
\qc بيّن أن $y=\dfrac{-\alpha^2+3\alpha-1}{\alpha}\,x$ هي معادلة لـ $(T)$ مماس المنحنى $(C_f)$ في النقطة ذات الفاصلة $\alpha$.
\end{sq}
\textbf{4)} أنشئ $(T)$ و $(C_f)$ على المجال $]0;4]$. (نأخذ $\dfrac{-\alpha^2+3\alpha-1}{\alpha}\approx0{,}8$).

\textbf{5)}
\begin{sq}
\qa بيّن أنه من أجل كل عدد حقيقي $x$ موجب تماما: $f(x)-x=(x-1)(-1+\ln x)$.
\qb ادرس حسب قيم العدد الحقيقي $x$ الموجب تماما إشارة $f(x)-x$.
\end{sq}
\textbf{6)} $K$ الدالة المعرفة على المجال $]0;+\infty[$ بـ: $K(x)=-\dfrac34x^2+2x+\left(\dfrac12x^2-x\right)\ln x$.

بيّن أنه من أجل كل عدد حقيقي $x$ موجب تماما: $K'(x)=f(x)-x$.

\textbf{7)} $g$ الدالة المعرفة على $]-2;+\infty[$ بـ: $g(x)=(x+1)\ln(x+2)$، و $(C_g)$ تمثيلها البياني في المعلم السابق.
\begin{sq}
\qa بيّن أنه من أجل كل عدد حقيقي $x$ من $]-2;+\infty[$: $g(x)=f(x+2)-1$.
\qb استنتج أن $(C_g)$ صورة $(C_f)$ بانسحاب يطلب تعيين شعاعه. (لا يطلب إنشاء $(C_g)$).
\end{sq}
\end{tamrin}

\begin{hal}
\QQ{1}
\pa بما أن $x>0$، إشارة $\dfrac{x-1}{x}$ هي إشارة $x-1$؛ و $\ln x>0\iff x>1$:
\begin{TAB}
\tkzTabInit[lgt=2.6,espcl=3]{$x$/1,$\frac{x-1}{x}$/1,$\ln x$/1,$\ln x+\frac{x-1}{x}$/1}{$0$,$1$,$+\infty$}
\tkzTabLine{d,-,z,+,}
\tkzTabLine{d,-,z,+,}
\tkzTabLine{d,-,z,+,}
\end{TAB}
\pb العددان $\frac{x-1}{x}$ و $\ln x$ لهما نفس الإشارة دائما، إذن مجموعهما:
\begin{itemize}
\item سالب تماما على $]0;1[$ (مجموع عددين سالبين تماما)؛
\item موجب تماما على $]1;+\infty[$ (مجموع عددين موجبين تماما)؛
\item معدوم عند $x=1$.
\end{itemize}

\QQ{2}
\pa $\displaystyle\lim_{x\to0^+}(x-1)\ln x=(-1)\times(-\infty)=+\infty$، إذن \bx{\lim_{x\to0^+}f(x)=+\infty}.

$\displaystyle\lim_{x\to+\infty}(x-1)=+\infty$ و $\displaystyle\lim_{x\to+\infty}\ln x=+\infty$، إذن \bx{\lim_{x\to+\infty}f(x)=+\infty}.

\pb $f'(x)=1\times\ln x+(x-1)\times\dfrac1x=\ln x+\dfrac{x-1}{x}$، وحسب السؤال 1-ب:

$f$ \bxt{متناقصة تماما على $]0;1]$} و \bxt{متزايدة تماما على $[1;+\infty[$}، مع $f(1)=1$.
\begin{TAB}
\tkzTabInit[lgt=2,espcl=3.5]{$x$/1,$f'(x)$/1,$f(x)$/2}{$0$,$1$,$+\infty$}
\tkzTabLine{d,-,z,+,}
\tkzTabVar{D+/$+\infty$,-/$1$,+/$+\infty$}
\end{TAB}

\QQ{3}
\pa $h'(x)=1+\dfrac1x>0$ من أجل كل $x>0$، إذن $h$ متزايدة تماما على $]0;+\infty[$.

\pb $h$ مستمرة ومتزايدة تماما، و $h(1{,}5)=-0{,}5+\ln1{,}5\approx-0{,}09<0$، $h(1{,}6)=-0{,}4+\ln1{,}6\approx0{,}07>0$. حسب مبرهنة القيم المتوسطة المعادلة $h(x)=0$ تقبل حلا وحيدا $\alpha$ حيث \bx{1{,}5<\alpha<1{,}6}.

$h(\alpha)=0\iff\alpha-2+\ln\alpha=0\iff$ \bx{\ln\alpha=2-\alpha}.

\pc معادلة المماس: $y=f'(\alpha)(x-\alpha)+f(\alpha)$. باستعمال $\ln\alpha=2-\alpha$:
\[f'(\alpha)=\ln\alpha+1-\frac1\alpha=3-\alpha-\frac1\alpha=\frac{-\alpha^2+3\alpha-1}{\alpha},\]
\[f(\alpha)=1+(\alpha-1)(2-\alpha)=-\alpha^2+3\alpha-1=\alpha f'(\alpha).\]
إذن $y=f'(\alpha)x-\alpha f'(\alpha)+\alpha f'(\alpha)=f'(\alpha)\,x$، أي \bx{(T):\ y=\frac{-\alpha^2+3\alpha-1}{\alpha}\,x} (المماس يمر بالمبدأ $O$).

\QQ{4} الرسم ($(T):y\approx0{,}8x$، $\alpha\approx1{,}56$):
\begin{GR}{-0.5}{4.2}{-0.5}{5.5}{1.4}
\Ticks{1,2,3,4}{1,2,3,4,5}
\draw[cThm,thick,domain=0:4.2,variable=\x] plot (\x,{0.8007*\x}) node[pos=0.9,below right]{$(T)$};
\draw[cDef,very thick,domain=0.02:4,samples=200,smooth,variable=\x] plot (\x,{1+(\x-1)*ln(\x)});
\node[cDef] at (3.3,4.4){$(C_f)$};
\fill (1,1) circle (1.5pt);
\fill (1.557,1.2467) circle (1.5pt);
\draw[densely dotted] (1.557,0)--(1.557,1.2467);
\node[below,font=\scriptsize] at (1.557,0){$\alpha$};
\end{GR}

\QQ{5}
\pa $f(x)-x=1-x+(x-1)\ln x=-(x-1)+(x-1)\ln x=$ \bx{(x-1)(-1+\ln x)}.

\pb $x-1=0\iff x=1$ و $-1+\ln x=0\iff x=e$:
\begin{TAB}
\tkzTabInit[lgt=2.6,espcl=2.6]{$x$/1,$x-1$/1,$-1+\ln x$/1,$f(x)-x$/1}{$0$,$1$,$e$,$+\infty$}
\tkzTabLine{d,-,z,+,t,+,}
\tkzTabLine{d,-,t,-,z,+,}
\tkzTabLine{d,+,z,-,z,+,}
\end{TAB}
$f(x)-x>0$ على $]0;1[\cup]e;+\infty[$، $f(x)-x<0$ على $]1;e[$، وينعدم عند $1$ و $e$.

\QQ{6} $K$ قابلة للاشتقاق على $]0;+\infty[$:
\begin{align*}
K'(x)&=-\frac32x+2+(x-1)\ln x+\left(\frac12x^2-x\right)\frac1x
=-\frac32x+2+(x-1)\ln x+\frac12x-1\\
&=-x+1+(x-1)\ln x=f(x)-x.
\end{align*}

\QQ{7}
\pa من أجل $x>-2$ لدينا $x+2>0$، إذن $f(x+2)$ معرفة و
\[f(x+2)-1=1+(x+2-1)\ln(x+2)-1=(x+1)\ln(x+2)=g(x).\]
\pb $g(x)=f(x+2)-1$، إذن $(C_g)$ صورة $(C_f)$ بالانسحاب الذي شعاعه \bx{\vec u=-2\vec i-\vec j} أي $\vec u(-2;-1)$.
\end{hal}

% ------------------------------------------------------------------
\begin{tamrin}{بكالوريا 2023 -- الموضوع الأول -- تقني رياضي}
\textbf{I)} $g$ الدالة المعرفة على $]-1;+\infty[$ بـ: $g(x)=2\ln(x+1)-\dfrac{x}{x+1}$.

$(C_g)$ تمثيلها البياني، يقطع حامل محور الفواصل في النقطتين اللتين فاصلتاهما $\alpha$ و $0$ (لاحظ الشكل).
\begin{GR}{-1.2}{3}{-0.6}{2.6}{1.2}
\Ticks{-1,1,2}{1,2}
\draw[cDef,very thick,domain=-0.93:3,samples=200,smooth,variable=\x] plot (\x,{2*ln(\x+1)-\x/(\x+1)});
\fill (-0.7153,0) circle (1.5pt) node[above left]{$\alpha$};
\node[cDef] at (2.4,1.4){$(C_g)$};
\end{GR}
\textbf{1)} بقراءة بيانية، حدّد حسب قيم $x$ إشارة $g(x)$.

\textbf{2)} تحقق أن: $-0{,}72<\alpha<-0{,}71$.

\textbf{II)} $f$ الدالة المعرفة على $]-1;+\infty[$ بـ: $f(x)=(2x+3)\ln(x+1)-3x$.

$(C_f)$ تمثيلها البياني في المستوي المنسوب إلى المعلم المتعامد والمتجانس $(O;\vec i,\vec j)$ (الوحدة $2\,\text{cm}$).

\textbf{1)}
\begin{sq}
\qa احسب $\displaystyle\lim_{x\to-1^+}f(x)$ ثم فسّر النتيجة هندسيا.
\qb تحقق أنه من أجل كل عدد حقيقي $x$ غير معدوم من المجال $]-1;+\infty[$: $f(x)=x\left[\left(2+\dfrac3x\right)\ln(x+1)-3\right]$، ثم استنتج أن: $\displaystyle\lim_{x\to+\infty}f(x)=+\infty$.
\end{sq}
\textbf{2)}
\begin{sq}
\qa بيّن أنه من أجل كل عدد حقيقي $x$ من المجال $]-1;+\infty[$: $f'(x)=g(x)$.
\qb استنتج أن $f$ متناقصة تماما على $[\alpha;0]$ ومتزايدة تماما على كل من المجالين $]-1;\alpha]$ و $[0;+\infty[$.
\qc شكّل جدول تغيرات الدالة $f$.
\end{sq}
\textbf{3)}
\begin{sq}
\qa ارسم $(C_f)$ في المجال $]-1;4]$. (نأخذ: $f(3)\approx3{,}5$، $f(4)\approx5{,}7$ و $f(\alpha)\approx0{,}2$).
\qb عيّن بيانيا قيم الوسيط الحقيقي $m$ التي من أجلها تقبل المعادلة $f(x)=m$ ثلاثة حلول بالضبط.
\end{sq}
\end{tamrin}

\begin{hal}
\textbf{الجزء I}

\QQ{1} من الشكل: $(C_g)$ فوق محور الفواصل على $]-1;\alpha[$ و $]0;+\infty[$، وتحته على $]\alpha;0[$:
\begin{TAB}
\tkzTabInit[lgt=2,espcl=2.6]{$x$/1,$g(x)$/1}{$-1$,$\alpha$,$0$,$+\infty$}
\tkzTabLine{d,+,z,-,z,+,}
\end{TAB}

\QQ{2} $g(-0{,}72)=2\ln0{,}28+\dfrac{0{,}72}{0{,}28}\approx0{,}025>0$ و $g(-0{,}71)=2\ln0{,}29+\dfrac{0{,}71}{0{,}29}\approx-0{,}027<0$.

$g$ مستمرة على $[-0{,}72;-0{,}71]$ و $g(-0{,}72)\times g(-0{,}71)<0$، وحسب الشكل $g$ متناقصة على هذا المجال، إذن \bx{-0{,}72<\alpha<-0{,}71}.

\textbf{الجزء II}

\QQ{1}
\pa $\displaystyle\lim_{x\to-1^+}(2x+3)=1$ و $\displaystyle\lim_{x\to-1^+}\ln(x+1)=-\infty$ و $\displaystyle\lim_{x\to-1^+}(-3x)=3$، إذن \bx{\lim_{x\to-1^+}f(x)=-\infty}.

التفسير: المستقيم ذو المعادلة $x=-1$ مقارب عمودي لـ $(C_f)$.

\pb من أجل $x\neq0$: $x\left[\left(2+\frac3x\right)\ln(x+1)-3\right]=(2x+3)\ln(x+1)-3x=f(x)$. \ok

لما $x\to+\infty$: $2+\frac3x\to2$ و $\ln(x+1)\to+\infty$، إذن القوس يؤول إلى $+\infty$، وبالضرب في $x\to+\infty$: \bx{\lim_{x\to+\infty}f(x)=+\infty}.

\QQ{2}
\pa $f'(x)=2\ln(x+1)+(2x+3)\cdot\dfrac{1}{x+1}-3=2\ln(x+1)+\dfrac{2x+3-3x-3}{x+1}=2\ln(x+1)-\dfrac{x}{x+1}=g(x)$.

\pb إشارة $f'(x)$ هي إشارة $g(x)$ (الجزء I): $f'(x)>0$ على $]-1;\alpha[$ و $]0;+\infty[$، و $f'(x)<0$ على $]\alpha;0[$. إذن $f$ متزايدة تماما على $]-1;\alpha]$ وعلى $[0;+\infty[$، ومتناقصة تماما على $[\alpha;0]$.

\pc $f(0)=3\ln1-0=0$. باستعمال $g(\alpha)=0$ أي $\ln(\alpha+1)=\dfrac{\alpha}{2(\alpha+1)}$ نجد $f(\alpha)=\dfrac{-\alpha(4\alpha+3)}{2(\alpha+1)}\approx0{,}17$.
\begin{TAB}
\tkzTabInit[lgt=2,espcl=2.8]{$x$/1,$f'(x)$/1,$f(x)$/2}{$-1$,$\alpha$,$0$,$+\infty$}
\tkzTabLine{d,+,z,-,z,+,}
\tkzTabVar{D-/$-\infty$,+/$f(\alpha)$,-/$0$,+/$+\infty$}
\end{TAB}

\QQ{3}
\pa الرسم ($(C_f)$ يمسّ محور الفواصل في $O$ لأن $f'(0)=g(0)=0$):
\begin{GR}{-1}{4.2}{-3}{6}{1}
\Ticks{1,2,3,4}{-2,-1,1,2,3,4,5}
\draw[cThm,dashed,thick] (-1,-3)--(-1,6);
\draw[cDef,very thick,domain=-0.995:4,samples=300,smooth,variable=\x] plot (\x,{(2*\x+3)*ln(\x+1)-3*\x});
\fill (-0.7153,0.1742) circle (1.5pt);
\fill (3,3.4766) circle (1.5pt); \fill (4,5.7038) circle (1.5pt);
\node[cDef] at (3.2,5.2){$(C_f)$};
\end{GR}

\pb حلول المعادلة $f(x)=m$ هي فواصل نقط تقاطع $(C_f)$ مع المستقيم الأفقي $y=m$. من جدول التغيرات والرسم، المستقيم $y=m$ يقطع $(C_f)$ في ثلاث نقاط بالضبط عندما يكون محصورا تماما بين القيمة الحدية الصغرى $0$ والقيمة الحدية الكبرى $f(\alpha)$:
\[\bx{m\in\left]0;f(\alpha)\right[}\quad\text{أي تقريبا } m\in]0;0{,}2[.\]
(من أجل $m=0$ أو $m=f(\alpha)$ يوجد حلان فقط.)
\end{hal}

% ------------------------------------------------------------------
\begin{tamrin}{بكالوريا 2024 -- الموضوع الثاني -- تقني رياضي}
\textbf{I)} يمثل الجدول المقابل تغيرات الدالة $g$ المعرفة على المجال $]-2;+\infty[$ بـ: $g(x)=x+\ln(x+2)$.
\begin{TAB}
\tkzTabInit[lgt=2,espcl=5]{$x$/1,$g'(x)$/1,$g(x)$/2}{$-2$,$+\infty$}
\tkzTabLine{d,+,}
\tkzTabVar{D-/$-\infty$,+/$+\infty$}
\end{TAB}
\textbf{1)} أثبت أن المعادلة $g(x)=0$ تقبل حلا وحيدا $\alpha$ حيث $-0{,}45<\alpha<-0{,}44$.

\textbf{2)} استنتج حسب قيم $x$ إشارة $g(x)$.

\textbf{II)} $f$ الدالة المعرفة على $]-2;+\infty[$ بـ: $f(x)=\left(\dfrac{x+1}{x+2}\right)\left(-1+\ln(x+2)\right)$.

$(C_f)$ تمثيلها البياني في المستوي المنسوب إلى المعلم المتعامد والمتجانس $(O;\vec i,\vec j)$ (وحدة الطول $2\,\text{cm}$).

\textbf{1)} احسب $\displaystyle\lim_{x\to-2^+}f(x)$ و $\displaystyle\lim_{x\to+\infty}f(x)$.

\textbf{2)}
\begin{sq}
\qa بيّن أنه من أجل كل $x$ من $]-2;+\infty[$: $f'(x)=\dfrac{g(x)}{(x+2)^2}$.
\qb استنتج اتجاه تغير الدالة $f$ ثم شكّل جدول تغيراتها.
\end{sq}
\textbf{3)} في الشكل المقابل، $(C_h)$ منحنى الدالة $h$ المعرفة على $]-2;+\infty[$ بـ: $h(x)=-1+\ln(x+2)$.
\begin{GR}{-2.5}{7.5}{-3.5}{1.8}{0.8}
\Ticks{-2,-1,1,2,3,4,5,6,7}{-3,-2,-1,1}
\draw[cDef,very thick,domain=-1.97:7.5,samples=200,smooth,variable=\x] plot (\x,{-1+ln(\x+2)});
\node[cDef] at (6.5,0.8){$(C_h)$};
\end{GR}
\begin{sq}
\qa بيّن أن $(C_h)$ منحن مقارب لـ $(C_f)$ عند $+\infty$.
\qb ادرس الوضع النسبي للمنحنيين $(C_h)$ و $(C_f)$.
\end{sq}
\textbf{4)}
\begin{sq}
\qa أعد رسم $(C_h)$ على ورقة الإجابة ثم ارسم $(C_f)$. (نأخذ: $f(\alpha)\approx-0{,}2$).
\qb ناقش بيانيا حسب قيم الوسيط الحقيقي $m$ عدد حلول المعادلة $f(x)=m$.
\end{sq}
\end{tamrin}

\begin{hal}
\textbf{الجزء I}

\QQ{1} من الجدول $g$ مستمرة ومتزايدة تماما على $]-2;+\infty[$، و
\[g(-0{,}45)=-0{,}45+\ln1{,}55\approx-0{,}012<0,\qquad g(-0{,}44)=-0{,}44+\ln1{,}56\approx0{,}005>0.\]
حسب مبرهنة القيم المتوسطة، المعادلة $g(x)=0$ تقبل حلا وحيدا $\alpha$ حيث \bx{-0{,}45<\alpha<-0{,}44}.

\QQ{2} لأن $g$ متزايدة تماما و $g(\alpha)=0$:
\begin{TAB}
\tkzTabInit[lgt=2,espcl=3]{$x$/1,$g(x)$/1}{$-2$,$\alpha$,$+\infty$}
\tkzTabLine{d,-,z,+,}
\end{TAB}

\textbf{الجزء II}

\QQ{1} \textbf{عند $-2^+$:} $\displaystyle\lim_{x\to-2^+}\frac{x+1}{x+2}=\frac{-1}{0^+}=-\infty$ و $\displaystyle\lim_{x\to-2^+}\left(-1+\ln(x+2)\right)=-\infty$، إذن \bx{\lim_{x\to-2^+}f(x)=+\infty} (المستقيم $x=-2$ مقارب عمودي).

\textbf{عند $+\infty$:} $\displaystyle\lim_{x\to+\infty}\frac{x+1}{x+2}=1$ و $\displaystyle\lim_{x\to+\infty}\left(-1+\ln(x+2)\right)=+\infty$، إذن \bx{\lim_{x\to+\infty}f(x)=+\infty}.

\QQ{2}
\pa $\left(\dfrac{x+1}{x+2}\right)'=\dfrac{(x+2)-(x+1)}{(x+2)^2}=\dfrac{1}{(x+2)^2}$، إذن
\[f'(x)=\frac{-1+\ln(x+2)}{(x+2)^2}+\frac{x+1}{x+2}\cdot\frac{1}{x+2}=\frac{-1+\ln(x+2)+x+1}{(x+2)^2}=\frac{x+\ln(x+2)}{(x+2)^2}=\frac{g(x)}{(x+2)^2}.\]
\pb إشارة $f'(x)$ هي إشارة $g(x)$: $f$ \bxt{متناقصة تماما على $]-2;\alpha]$} و \bxt{متزايدة تماما على $[\alpha;+\infty[$}.

($f(\alpha)$: من $\ln(\alpha+2)=-\alpha$ نجد $f(\alpha)=-\dfrac{(\alpha+1)^2}{\alpha+2}\approx-0{,}2$.)
\begin{TAB}
\tkzTabInit[lgt=2,espcl=3.5]{$x$/1,$f'(x)$/1,$f(x)$/2}{$-2$,$\alpha$,$+\infty$}
\tkzTabLine{d,-,z,+,}
\tkzTabVar{D+/$+\infty$,-/$f(\alpha)$,+/$+\infty$}
\end{TAB}

\QQ{3}
\pa $f(x)-h(x)=\left(-1+\ln(x+2)\right)\left(\dfrac{x+1}{x+2}-1\right)=\dfrac{1-\ln(x+2)}{x+2}=\dfrac{1}{x+2}-\dfrac{\ln(x+2)}{x+2}$.

بوضع $t=x+2\to+\infty$: $\dfrac1t\to0$ و $\dfrac{\ln t}{t}\to0$، إذن $\displaystyle\lim_{x\to+\infty}\left[f(x)-h(x)\right]=0$، ومنه \bxt{$(C_h)$ منحن مقارب لـ $(C_f)$ عند $+\infty$}.

\pb إشارة $f(x)-h(x)=\dfrac{1-\ln(x+2)}{x+2}$ هي إشارة $1-\ln(x+2)$ (لأن $x+2>0$):

$1-\ln(x+2)>0\iff x+2<e\iff x<e-2$.
\begin{itemize}
\item على $]-2;e-2[$: $(C_f)$ فوق $(C_h)$؛
\item على $]e-2;+\infty[$: $(C_f)$ تحت $(C_h)$؛
\item يتقاطعان في النقطة \bx{(e-2;0)} (لأن $h(e-2)=-1+1=0$).
\end{itemize}

\QQ{4}
\pa الرسم ($(C_f)$ يقطع محور الفواصل في $(-1;0)$ و $(e-2;0)$):
\begin{GR}{-2.2}{7.5}{-3}{3.5}{1}
\Ticks{-1,1,2,3,4,5,6,7}{-2,-1,1,2,3}
\draw[cThm,dashed,thick] (-2,-3)--(-2,3.5);
\draw[cRem!80!black,thick,domain=-1.95:7.5,samples=200,smooth,variable=\x] plot (\x,{-1+ln(\x+2)}) node[pos=0.97,below]{$(C_h)$};
\draw[cDef,very thick,domain=-1.85:7.5,samples=300,smooth,variable=\x] plot (\x,{(\x+1)/(\x+2)*(-1+ln(\x+2))});
\node[cDef] at (-1.2,2.6){$(C_f)$};
\fill (-0.4429,-0.1993) circle (1.5pt);
\fill (0.7183,0) circle (1.5pt); \fill (-1,0) circle (1.5pt);
\end{GR}

\pb حلول $f(x)=m$ هي فواصل نقط تقاطع $(C_f)$ مع المستقيم $y=m$:
\begin{itemize}
\item $m<f(\alpha)$: المعادلة لا تقبل حلولا؛
\item $m=f(\alpha)$: حل وحيد (مضاعف) هو $\alpha$؛
\item $m>f(\alpha)$: حلان متمايزان.
\end{itemize}
\end{hal}

% ------------------------------------------------------------------
\begin{tamrin}{بكالوريا 2025 -- الموضوع الثاني -- تقني رياضي}
\textbf{I)} $g$ الدالة المعرفة على $]-1;+\infty[$ بـ: $g(x)=x^2+2x+2-2\ln(x+1)$.

\textbf{1)}
\begin{sq}
\qa بيّن أنه من أجل كل $x$ من $]-1;+\infty[$: $g'(x)=\dfrac{2x(x+2)}{x+1}$.
\qb ادرس اتجاه تغير الدالة $g$ ثم شكّل جدول تغيراتها (لا يطلب حساب النهايات).
\end{sq}
\textbf{2)} استنتج إشارة $g(x)$.

\textbf{II)} $f$ الدالة المعرفة على $]-1;+\infty[$ بـ: $f(x)=x+\dfrac{1+2\ln(x+1)}{x+1}$.

$(C_f)$ تمثيلها البياني في المستوي المنسوب إلى المعلم المتعامد والمتجانس $(O;\vec i,\vec j)$.

\textbf{1)}
\begin{sq}
\qa احسب $\displaystyle\lim_{x\to-1^+}f(x)$ و $\displaystyle\lim_{x\to+\infty}f(x)$.
\qb بيّن أنه من أجل كل $x$ من $]-1;+\infty[$: $f'(x)=\dfrac{g(x)}{(x+1)^2}$.
\qc استنتج اتجاه تغير الدالة $f$ ثم شكّل جدول تغيراتها.
\end{sq}
\textbf{2)}
\begin{sq}
\qa بيّن أن المستقيم $(\Delta)$ ذا المعادلة $y=x$ مقارب مائل للمنحنى $(C_f)$.
\qb بيّن أن المنحنى $(C_f)$ يقبل مماسا $(T)$ موازيا لـ $(\Delta)$؛ يطلب تعيين معادلة له.
\end{sq}
\textbf{3)}
\begin{sq}
\qa بيّن أنه من أجل كل $x$ من $]-1;+\infty[$: $f''(x)=\dfrac{4\ln(x+1)-4}{(x+1)^3}$.
\qb استنتج أن للمنحنى $(C_f)$ نقطة انعطاف، يطلب تعيين إحداثييها.
\end{sq}
\textbf{4)}
\begin{sq}
\qa احسب $f(0)$ ثم ارسم كلا من $(\Delta)$ و $(T)$ و $(C_f)$.
\qb ناقش بيانيا وحسب قيم العدد الحقيقي $m$ عدد حلول المعادلة: $1+2\ln(x+1)=m(x+1)$.
\end{sq}
\end{tamrin}

\begin{hal}
\textbf{الجزء I}

\QQ{1}
\pa $g'(x)=2x+2-\dfrac{2}{x+1}=\dfrac{2(x+1)^2-2}{x+1}=\dfrac{2(x^2+2x)}{x+1}=\dfrac{2x(x+2)}{x+1}$.

\pb على $]-1;+\infty[$: $x+1>0$ و $x+2>0$، إذن إشارة $g'(x)$ هي إشارة $x$: $g$ متناقصة تماما على $]-1;0]$ ومتزايدة تماما على $[0;+\infty[$، و $g(0)=2-2\ln1=2$.
\begin{TAB}
\tkzTabInit[lgt=2,espcl=3.5]{$x$/1,$g'(x)$/1,$g(x)$/2}{$-1$,$0$,$+\infty$}
\tkzTabLine{d,-,z,+,}
\tkzTabVar{D+/ ,-/$2$,+/ }
\end{TAB}

\QQ{2} القيمة الحدية الصغرى لـ $g$ على $]-1;+\infty[$ هي $g(0)=2>0$، إذن من أجل كل $x>-1$: $g(x)\geqslant2>0$، أي \bx{g(x)>0}.

\textbf{الجزء II}

\QQ{1}
\pa \textbf{عند $-1^+$:} $1+2\ln(x+1)\to-\infty$ و $x+1\to0^+$، إذن $\dfrac{1+2\ln(x+1)}{x+1}\to-\infty$، ومنه \bx{\lim_{x\to-1^+}f(x)=-\infty} (المستقيم $x=-1$ مقارب عمودي).

\textbf{عند $+\infty$:} $\dfrac{1+2\ln(x+1)}{x+1}=\dfrac{1}{x+1}+2\,\dfrac{\ln(x+1)}{x+1}\to0$ (بوضع $t=x+1$، $\frac{\ln t}{t}\to0$)، إذن \bx{\lim_{x\to+\infty}f(x)=+\infty}.

\pb
\begin{align*}f'(x)&=1+\frac{\frac{2}{x+1}(x+1)-\left(1+2\ln(x+1)\right)}{(x+1)^2}=1+\frac{1-2\ln(x+1)}{(x+1)^2}\\&=\frac{x^2+2x+2-2\ln(x+1)}{(x+1)^2}=\frac{g(x)}{(x+1)^2}.\end{align*}
\pc $g(x)>0$ و $(x+1)^2>0$، إذن $f'(x)>0$: \bxt{$f$ متزايدة تماما على $]-1;+\infty[$}.
\begin{TAB}
\tkzTabInit[lgt=2,espcl=5]{$x$/1,$f'(x)$/1,$f(x)$/2}{$-1$,$+\infty$}
\tkzTabLine{d,+,}
\tkzTabVar{D-/$-\infty$,+/$+\infty$}
\end{TAB}

\QQ{2}
\pa $f(x)-x=\dfrac{1+2\ln(x+1)}{x+1}$ وقد بيّنا أن نهايتها عند $+\infty$ هي $0$، إذن \bxt{$(\Delta):y=x$ مقارب مائل لـ $(C_f)$} بجوار $+\infty$.

\pb $(T)\parallel(\Delta)$ معناه $f'(x_0)=1$:
\[f'(x_0)=1\iff\frac{1-2\ln(x_0+1)}{(x_0+1)^2}=0\iff\ln(x_0+1)=\frac12\iff x_0=\sqrt e-1.\]
$f(\sqrt e-1)=\sqrt e-1+\dfrac{1+1}{\sqrt e}=\sqrt e-1+\dfrac{2}{\sqrt e}$.

$(T):y=1\cdot\left(x-\sqrt e+1\right)+\sqrt e-1+\dfrac2{\sqrt e}$، أي \bx{(T):\ y=x+\frac{2}{\sqrt e}} ($\frac{2}{\sqrt e}\approx1{,}21$).

\QQ{3}
\pa نشتق $f'(x)=1+\dfrac{1-2\ln(x+1)}{(x+1)^2}$:
\[f''(x)=\frac{-\frac{2}{x+1}(x+1)^2-\left(1-2\ln(x+1)\right)\cdot2(x+1)}{(x+1)^4}=\frac{-2-2+4\ln(x+1)}{(x+1)^3}=\frac{4\ln(x+1)-4}{(x+1)^3}.\]
(بعد الاختزال على $x+1\neq0$.)

\pb $(x+1)^3>0$، إذن إشارة $f''(x)$ هي إشارة $4\left(\ln(x+1)-1\right)$: سالبة لما $x<e-1$ وموجبة لما $x>e-1$. $f''$ تنعدم عند $e-1$ مغيّرة إشارتها، إذن النقطة ذات الفاصلة $e-1$ نقطة انعطاف:

$f(e-1)=e-1+\dfrac{1+2}{e}$، ومنه نقطة الانعطاف \bx{I\left(e-1;\ e-1+\tfrac3e\right)}، أي تقريبا $I(1{,}72;2{,}82)$.

\QQ{4}
\pa $f(0)=0+\dfrac{1+0}{1}=$ \bx{1}.

للرسم نلاحظ أن $(C_f)$ يقطع $(\Delta)$ عند $1+2\ln(x+1)=0$ أي $x=e^{-\frac12}-1\approx-0{,}39$؛ $(C_f)$ تحت $(\Delta)$ قبلها وفوقه بعدها.
\begin{GR}{-1}{5}{-3}{6}{1}
\Ticks{1,2,3,4}{-2,-1,1,2,3,4,5}
\draw[cThm,dashed,thick] (-1,-3)--(-1,6);
\draw[cRem!80!black,thick,domain=-1:5,variable=\x] plot (\x,{\x}) node[pos=0.95,below right]{$(\Delta)$};
\draw[cThm,thick,domain=-1:4.8,variable=\x] plot (\x,{\x+1.2131}) node[pos=0.93,left]{$(T)$};
\draw[cDef,very thick,domain=-0.8:5,samples=300,smooth,variable=\x] plot (\x,{\x+(1+2*ln(\x+1))/(\x+1)});
\node[cDef] at (-0.35,-2.2){$(C_f)$};
\fill (0.6487,1.8618) circle (1.5pt);
\fill (1.7183,2.8219) circle (1.5pt) node[below right]{$I$};
\fill (0,1) circle (1.5pt);
\end{GR}

\pb لأن $x+1\neq0$:
\[1+2\ln(x+1)=m(x+1)\iff\frac{1+2\ln(x+1)}{x+1}=m\iff f(x)=x+m.\]
الحلول هي فواصل نقط تقاطع $(C_f)$ مع المستقيم $(D_m):y=x+m$ الموازي لـ $(\Delta)$ (ميله $1$). المستقيم $(D_m)$ يحصره المستقيمان $(\Delta)$ ($m=0$) و $(T)$ ($m=\frac2{\sqrt e}$):
\begin{itemize}
\item $m\leqslant0$: حل وحيد؛
\item $0<m<\frac{2}{\sqrt e}$: حلان متمايزان؛
\item $m=\frac{2}{\sqrt e}$: حل وحيد هو $\sqrt e-1$ (نقطة التماس)؛
\item $m>\frac{2}{\sqrt e}$: لا حلول.
\end{itemize}
\end{hal}

% ------------------------------------------------------------------
\begin{tamrin}{بكالوريا 2026 -- الموضوع الثاني -- تقني رياضي}
الدالة $f$ المعرفة على $]0;e[\cup]e;+\infty[$ بـ: $f(x)=x+\dfrac{1}{x-x\ln x}$، و $(C_f)$ تمثيلها البياني في المستوي المنسوب إلى المعلم المتعامد والمتجانس $(O;\vec i,\vec j)$.

\textbf{1)} احسب النهايات: $\displaystyle\lim_{x\to0^+}f(x)$، $\displaystyle\lim_{x\to e^-}f(x)$، $\displaystyle\lim_{x\to e^+}f(x)$ و $\displaystyle\lim_{x\to+\infty}f(x)$، ثم فسّر النتائج هندسيا.

\textbf{2)}
\begin{sq}
\qa بيّن أن المستقيم $(\Delta)$ ذا المعادلة $y=x$ مقارب مائل لـ $(C_f)$.
\qb ادرس الوضع النسبي لـ $(C_f)$ و $(\Delta)$.
\end{sq}
\textbf{3)}
\begin{sq}
\qa بيّن أنه من أجل كل $x$ من $]0;e[\cup]e;+\infty[$: $f'(x)=1+\dfrac{\ln x}{(x-x\ln x)^2}$.
\qb نقبل أن إشارة $f'(x)$ موضحة في الجدول المقابل، تحقق أن $0{,}4<\alpha<0{,}5$.
\end{sq}
\begin{TAB}
\tkzTabInit[lgt=2,espcl=3]{$x$/1,$f'(x)$/1}{$0$,$\alpha$,$e$,$+\infty$}
\tkzTabLine{d,-,z,+,d,+,}
\end{TAB}
\begin{sq}
\qc شكّل جدول تغيرات الدالة $f$.
\end{sq}
\textbf{4)}
\begin{sq}
\qa بيّن أن $(C_f)$ يقبل مماسا $(T)$ يوازي $(\Delta)$، يطلب تعيين معادلة له.
\qb أثبت أن المعادلة $f(x)=0$ تقبل حلا وحيدا في المجال $]3;3{,}1[$.
\end{sq}
\textbf{5)}
\begin{sq}
\qa ارسم $(\Delta)$ و $(T)$ و $(C_f)$. (نأخذ $f(\alpha)\simeq1{,}7$).
\qb ناقش بيانيا وحسب قيم العدد الحقيقي $m$ عدد حلول المعادلة: $m(x-x\ln x)=1$.
\end{sq}
\end{tamrin}

\begin{hal}
نلاحظ أن $x-x\ln x=x(1-\ln x)$، وهو موجب تماما على $]0;e[$ (لأن $\ln x<1$) وسالب تماما على $]e;+\infty[$.

\QQ{1}
\begin{itemize}
\item $\displaystyle\lim_{x\to0^+}x\ln x=0$ إذن $x-x\ln x\to0^+$، ومنه $\dfrac{1}{x-x\ln x}\to+\infty$ و \bx{\lim_{x\to0^+}f(x)=+\infty}.
\item $x\to e^-$: $x(1-\ln x)\to e\times0^+=0^+$، إذن \bx{\lim_{x\to e^-}f(x)=+\infty}.
\item $x\to e^+$: $x(1-\ln x)\to0^-$، إذن \bx{\lim_{x\to e^+}f(x)=-\infty}.
\item $x\to+\infty$: $x(1-\ln x)\to+\infty\times(-\infty)=-\infty$، إذن $\dfrac{1}{x-x\ln x}\to0$ و \bx{\lim_{x\to+\infty}f(x)=+\infty}.
\end{itemize}
\textbf{التفسير الهندسي:} المستقيمان $x=0$ و $x=e$ مقاربان عموديان لـ $(C_f)$.

\QQ{2}
\pa $f(x)-x=\dfrac{1}{x-x\ln x}$ ونهايتها عند $+\infty$ هي $0$، إذن \bxt{$(\Delta):y=x$ مقارب مائل لـ $(C_f)$} بجوار $+\infty$.

\pb إشارة $f(x)-x=\dfrac{1}{x(1-\ln x)}$ هي إشارة $1-\ln x$:
\begin{itemize}
\item على $]0;e[$: $f(x)-x>0$، $(C_f)$ فوق $(\Delta)$؛
\item على $]e;+\infty[$: $f(x)-x<0$، $(C_f)$ تحت $(\Delta)$.
\end{itemize}

\QQ{3}
\pa $(x-x\ln x)'=1-\left(\ln x+x\cdot\frac1x\right)=-\ln x$، إذن
\[f'(x)=1-\frac{(x-x\ln x)'}{(x-x\ln x)^2}=1-\frac{-\ln x}{(x-x\ln x)^2}=1+\frac{\ln x}{(x-x\ln x)^2}.\]
\pb $f'(0{,}4)\approx-0{,}56<0$ و $f'(0{,}5)\approx0{,}03>0$، والإشارة تتغير مرة واحدة على $]0;e[$ عند $\alpha$ (حسب الجدول)، إذن \bx{0{,}4<\alpha<0{,}5}.

\pc
\begin{TAB}
\tkzTabInit[lgt=2,espcl=3]{$x$/1,$f'(x)$/1,$f(x)$/2}{$0$,$\alpha$,$e$,$+\infty$}
\tkzTabLine{d,-,z,+,d,+,}
\tkzTabVar{D+/$+\infty$,-/$f(\alpha)$,+D-/$+\infty$/$-\infty$,+/$+\infty$}
\end{TAB}

\QQ{4}
\pa $(T)\parallel(\Delta)$ معناه $f'(x_0)=1$، أي $\dfrac{\ln x_0}{(x_0-x_0\ln x_0)^2}=0\iff\ln x_0=0\iff x_0=1$.

$f(1)=1+\dfrac{1}{1-0}=2$، إذن $(T):y=1\cdot(x-1)+2$، أي \bx{(T):\ y=x+1}.

\pb على $]0;e[$: $f(x)=x+\dfrac{1}{x(1-\ln x)}>0$ (مجموع عددين موجبين تماما)، فلا حلول فيه.

على $]e;+\infty[$: $f$ مستمرة ومتزايدة تماما، و
\[f(3)=3+\frac{1}{3-3\ln3}\approx-0{,}38<0,\qquad f(3{,}1)=3{,}1+\frac{1}{3{,}1-3{,}1\ln3{,}1}\approx0{,}65>0.\]
حسب مبرهنة القيم المتوسطة، المعادلة $f(x)=0$ تقبل حلا وحيدا في المجال \bx{]3;3{,}1[}.

\QQ{5}
\pa الرسم:
\begin{GR}{-0.5}{6.5}{-3}{7}{0.9}
\Ticks{1,2,3,4,5,6}{-2,-1,1,2,3,4,5,6}
\draw[cThm,dashed,thick] (2.718,-3)--(2.718,7);
\draw[cRem!80!black,thick,domain=0:6.5,variable=\x] plot (\x,{\x}) node[pos=0.95,below right]{$(\Delta)$};
\draw[cThm,thick,domain=0:6,variable=\x] plot (\x,{\x+1}) node[pos=0.95,left]{$(T)$};
\draw[cDef,very thick,domain=0.1:2.62,samples=200,smooth,variable=\x] plot (\x,{\x+1/(\x-\x*ln(\x))});
\draw[cDef,very thick,domain=2.79:6.5,samples=200,smooth,variable=\x] plot (\x,{\x+1/(\x-\x*ln(\x))});
\node[cDef] at (4.6,2){$(C_f)$};
\fill (1,2) circle (1.5pt);
\fill (0.4926,1.681) circle (1.5pt);
\fill (3.031,0) circle (1.5pt);
\node[below right,font=\scriptsize] at (2.718,0){$e$};
\end{GR}

\pb إذا كان $x=e$ فإن $x-x\ln x=0$ و $m\times0=0\neq1$، فـ $e$ ليس حلا (وهو خارج مجموعة التعريف أصلا). من أجل $x\in]0;e[\cup]e;+\infty[$:
\[m(x-x\ln x)=1\iff m=\frac{1}{x-x\ln x}\iff f(x)=x+m.\]
الحلول هي فواصل نقط تقاطع $(C_f)$ مع المستقيم $(D_m):y=x+m$ الموازي لـ $(\Delta)$؛ $(D_0)=(\Delta)$ و $(D_1)=(T)$:
\begin{itemize}
\item $m<0$: $(D_m)$ تحت $(\Delta)$، يقطع $(C_f)$ في نقطة وحيدة (على $]e;+\infty[$): حل وحيد؛
\item $0\leqslant m<1$: لا حلول؛
\item $m=1$: حل وحيد هو $1$ (نقطة التماس)؛
\item $m>1$: حلان متمايزان في $]0;e[$.
\end{itemize}
\end{hal}

\section{مسائل بمستوى البكالوريا -- شعبة رياضيات}

%====================================================== المسألة 1
\begin{tamrin}{مسألة بمستوى البكالوريا -- شعبة رياضيات}
\textbf{الجزء الأول:} نعتبر الدالة $g$ المعرفة على $]0;+\infty[$ بـ: $g(x)=x^2-1+\ln x$.
\QQ{1} ادرس اتجاه تغير الدالة $g$ على $]0;+\infty[$.
\QQ{2} احسب $g(1)$ ثم استنتج إشارة $g(x)$ حسب قيم $x$.

\textbf{الجزء الثاني:} نعتبر الدالة $f$ المعرفة على $]0;+\infty[$ بـ: $f(x)=x-1-\dfrac{\ln x}{x}$، و $(C_f)$ تمثيلها البياني في معلم متعامد ومتجانس $(O;\vec i,\vec j)$.
\QQ{1}
\begin{sq}
\qa احسب $\displaystyle\lim_{x\to0^+}f(x)$ و $\displaystyle\lim_{x\to+\infty}f(x)$، ثم فسّر النتيجة الأولى هندسيا.
\qb بيّن أن المستقيم $(\Delta)$ ذا المعادلة $y=x-1$ مستقيم مقارب مائل للمنحنى $(C_f)$، ثم ادرس الوضع النسبي لـ $(C_f)$ و $(\Delta)$.
\end{sq}
\QQ{2}
\begin{sq}
\qa بيّن أنه من أجل كل $x$ من $]0;+\infty[$: $f'(x)=\dfrac{g(x)}{x^2}$.
\qb استنتج اتجاه تغير الدالة $f$ ثم شكّل جدول تغيراتها.
\end{sq}
\QQ{3} بيّن أن $(C_f)$ يقبل مماسا وحيدا $(T)$ موازيا لـ $(\Delta)$ يطلب تعيين معادلة له.
\QQ{4} لتكن النقطة $A(0;-1)$.
\begin{sq}
\qa تحقق أن $A\in(\Delta)$.
\qb بيّن أنه يوجد مماس وحيد $(T')$ للمنحنى $(C_f)$ يشمل النقطة $A$، ويمس $(C_f)$ في النقطة ذات الفاصلة $\sqrt e$، ثم اكتب معادلة له.
\end{sq}
\QQ{5} أنشئ $(\Delta)$ و $(T)$ و $(T')$ ثم المنحنى $(C_f)$.
\QQ{6} $m$ وسيط حقيقي. ناقش بيانيا، حسب قيم $m$، عدد حلول المعادلة: $f(x)=mx-1$.
\end{tamrin}

\begin{hal}
\textbf{الجزء الأول:}
\QQ{1} الدالة $g$ قابلة للاشتقاق على $]0;+\infty[$ (مجموع دوال قابلة للاشتقاق)، ومن أجل كل $x>0$:
\[g'(x)=2x+\frac1x .\]
بما أن $x>0$ فإن $2x>0$ و $\frac1x>0$، إذن $g'(x)>0$.
\par\bxt{الدالة $g$ متزايدة تماما على $]0;+\infty[$}
\QQ{2} $g(1)=1^2-1+\ln1=0$ إذن \bx{g(1)=0}.

بما أن $g$ متزايدة تماما على $]0;+\infty[$ و $g(1)=0$ فإن:
\begin{itemize}
\item إذا كان $0<x<1$ فإن $g(x)<g(1)$ أي $g(x)<0$.
\item إذا كان $x>1$ فإن $g(x)>g(1)$ أي $g(x)>0$.
\end{itemize}
\begin{TAB}
\tkzTabInit[lgt=2,espcl=3]{$x$/1,$g(x)$/1}{$0$,$1$,$+\infty$}
\tkzTabLine{d,-,z,+,}
\end{TAB}

\textbf{الجزء الثاني:}
\QQ{1}
\pa عند $0^+$: لدينا $\displaystyle\lim_{x\to0^+}\ln x=-\infty$ و $\displaystyle\lim_{x\to0^+}\frac1x=+\infty$، إذن $\displaystyle\lim_{x\to0^+}\frac{\ln x}{x}=\lim_{x\to0^+}\ln x\times\frac1x=-\infty$، ومنه $\displaystyle\lim_{x\to0^+}\left(-\frac{\ln x}{x}\right)=+\infty$. وبما أن $\displaystyle\lim_{x\to0^+}(x-1)=-1$:
\[\bx{\lim_{x\to0^+}f(x)=+\infty}\]
التفسير: المنحنى $(C_f)$ يقبل مستقيما مقاربا عموديا معادلته $x=0$ (محور التراتيب).

عند $+\infty$: نعلم أن $\displaystyle\lim_{x\to+\infty}\frac{\ln x}{x}=0$ (التزايد المقارن) و $\displaystyle\lim_{x\to+\infty}(x-1)=+\infty$، إذن:
\[\bx{\lim_{x\to+\infty}f(x)=+\infty}\]
\pb لدينا $f(x)-(x-1)=-\dfrac{\ln x}{x}$ و $\displaystyle\lim_{x\to+\infty}\left(-\frac{\ln x}{x}\right)=0$، إذن $\displaystyle\lim_{x\to+\infty}\big[f(x)-(x-1)\big]=0$.
\par\bxt{$(\Delta): y=x-1$ مستقيم مقارب مائل لـ $(C_f)$ بجوار $+\infty$}

\textbf{الوضع النسبي:} إشارة الفرق $f(x)-(x-1)=-\dfrac{\ln x}{x}$ هي إشارة $-\ln x$ لأن $x>0$:
\begin{itemize}
\item إذا كان $0<x<1$: $\ln x<0$ إذن $f(x)-(x-1)>0$: \textbf{$(C_f)$ فوق $(\Delta)$}.
\item إذا كان $x=1$: $f(1)-0=0$: \textbf{$(C_f)$ يقطع $(\Delta)$ في النقطة $B(1;0)$}.
\item إذا كان $x>1$: $\ln x>0$ إذن $f(x)-(x-1)<0$: \textbf{$(C_f)$ تحت $(\Delta)$}.
\end{itemize}
\QQ{2}
\pa $f$ قابلة للاشتقاق على $]0;+\infty[$، ومشتقة $\dfrac{\ln x}{x}$ هي $\dfrac{\frac1x\cdot x-\ln x\cdot1}{x^2}=\dfrac{1-\ln x}{x^2}$، إذن:
\[f'(x)=1-\frac{1-\ln x}{x^2}=\frac{x^2-1+\ln x}{x^2}=\bx{\frac{g(x)}{x^2}}\]
\pb بما أن $x^2>0$ فإن إشارة $f'(x)$ هي إشارة $g(x)$: $f$ متناقصة تماما على $]0;1]$ ومتزايدة تماما على $[1;+\infty[$، مع $f(1)=1-1-\dfrac{\ln1}{1}=0$.
\begin{TAB}
\tkzTabInit[lgt=2,espcl=3]{$x$/1,$f'(x)$/1,$f(x)$/2}{$0$,$1$,$+\infty$}
\tkzTabLine{d,-,z,+,}
\tkzTabVar{D+/$+\infty$,-/$0$,+/$+\infty$}
\end{TAB}
نلاحظ أن $f(x)\geqslant0$ من أجل كل $x>0$ (القيمة الحدية الصغرى هي $0$).
\QQ{3} يكون المماس عند النقطة ذات الفاصلة $x_0$ موازيا لـ $(\Delta)$ إذا وفقط إذا كان $f'(x_0)=1$ (معامل توجيه $(\Delta)$ هو $1$):
\[f'(x_0)=1\iff \frac{x_0^2-1+\ln x_0}{x_0^2}=1\iff x_0^2-1+\ln x_0=x_0^2\iff \ln x_0=1\iff x_0=e .\]
الحل وحيد، إذن يوجد مماس وحيد $(T)$ موازٍ لـ $(\Delta)$ عند النقطة ذات الفاصلة $e$.

$f(e)=e-1-\dfrac{\ln e}{e}=e-1-\dfrac1e$، ومعادلة $(T)$: $y=f'(e)(x-e)+f(e)=(x-e)+e-1-\dfrac1e$، أي:
\[\bx{(T):\ y=x-1-\frac1e}\]
\QQ{4}
\pa من أجل $x=0$: $y=0-1=-1$، إذن $A(0;-1)\in(\Delta)$. \ok
\pb معادلة المماس عند النقطة ذات الفاصلة $a>0$ هي: $y=f'(a)(x-a)+f(a)$. يشمل هذا المماس النقطة $A(0;-1)$ إذا وفقط إذا:
\[-1=-a\,f'(a)+f(a).\]
نحسب: $a\,f'(a)=a-\dfrac{1-\ln a}{a}$، إذن:
\[f(a)-a\,f'(a)=a-1-\frac{\ln a}{a}-a+\frac{1-\ln a}{a}=-1+\frac{1-2\ln a}{a}.\]
ومنه: $-1=f(a)-af'(a)\iff \dfrac{1-2\ln a}{a}=0\iff \ln a=\dfrac12\iff a=\sqrt e$.

الحل وحيد، إذن يوجد مماس وحيد $(T')$ يشمل $A$، وهو المماس عند النقطة ذات الفاصلة $\sqrt e$.

حساب معامل توجيهه: $g(\sqrt e)=e-1+\dfrac12=e-\dfrac12$، إذن $f'(\sqrt e)=\dfrac{e-\frac12}{e}=1-\dfrac1{2e}$.

وبما أن $(T')$ يشمل $A(0;-1)$ فإن ترتيبه عند المبدأ هو $-1$:
\[\bx{(T'):\ y=\Big(1-\frac1{2e}\Big)x-1}\qquad\big(1-\tfrac1{2e}\approx0{,}82\big)\]
تحقق: $f(\sqrt e)=\sqrt e-1-\dfrac{1}{2\sqrt e}$ و $\Big(1-\dfrac1{2e}\Big)\sqrt e-1=\sqrt e-\dfrac{1}{2\sqrt e}-1$. \ok
\QQ{5} بعض القيم: $f(0{,}3)\approx3{,}31$، $f(0{,}5)\approx0{,}89$، $f(2)\approx0{,}65$، $f(e)\approx1{,}35$، $f(3)\approx1{,}63$.
\begin{GR}{-1}{6}{-2}{6}{1.1}
\Ticks{1,2,3,4,5}{1,2,3,4,5}
\draw[cThm,thick,domain=-1:6.5] plot(\x,{\x-1});
\node[cThm,above left] at (5,4){$(\Delta)$};
\draw[cRem!80!black,thick,dashed,domain=-1:6.5] plot(\x,{\x-1-1/e});
\draw[cRem!80!black] (4.2,2.832)--(4.6,1.7) node[below]{$(T)$};
\draw[cMeth,thick,domain=-1:6.5] plot(\x,{(1-1/(2*e))*\x-1});
\draw[cMeth] (5.5,3.488)--(5.75,2.5) node[below]{$(T')$};
\draw[cDef,very thick,samples=200,domain=0.13:6.2] plot(\x,{\x-1-ln(\x)/\x});
\node[cDef] at (0.65,4.5){$(C_f)$};
\fill (0,-1) circle(1.5pt) node[left]{$A$};
\fill (1,0) circle(1.5pt) node[below]{$B$};
\fill ({exp(0.5)},{exp(0.5)-1-0.5/exp(0.5)}) circle(1.5pt);
\fill ({e},{e-1-1/e}) circle(1.5pt);
\end{GR}
\QQ{6} حلول المعادلة $f(x)=mx-1$ هي فواصل نقط تقاطع $(C_f)$ مع المستقيم $(D_m):\ y=mx-1$، وهو مستقيم يشمل النقطة الثابتة $A(0;-1)$ ومعامل توجيهه $m$ (مستقيم يدور حول $A$).

المستقيمات الخاصة: $(T')$ معامل توجيهه $m_0=1-\dfrac{1}{2e}$، و $(\Delta)$ معامل توجيهه $1$ (و $(\Delta)$ يقطع $(C_f)$ في $B$ فقط).
\begin{GR}{-1}{6}{-2}{6}{1.1}
\Ticks{1,2,3,4,5}{1,2,3,4,5}
\draw[cDef,very thick,samples=200,domain=0.13:6.2] plot(\x,{\x-1-ln(\x)/\x});
\draw[cThm,thick,domain=-1:6.5] plot(\x,{\x-1});
\node[cThm,above left] at (5,4){$m=1$};
\draw[cMeth,thick,domain=-1:6.5] plot(\x,{(1-1/(2*e))*\x-1});
\node[cMeth] at (5.2,2.9){$m=m_0$};
\draw[gray,dashed,domain=-1:6.5] plot(\x,{0.5*\x-1});
\node[gray,below] at (5.3,1.6){$m<m_0$};
\draw[cEx,dashed,thick,domain=-1:6.5] plot(\x,{0.9*\x-1});
\node[cEx] at (3.3,3.4){$m_0<m<1$};
\draw[cSol,dashed,thick,domain=-1:6.5] plot(\x,{2*\x-1});
\node[cSol,left] at (2.5,4){$m>1$};
\draw[gray,dashed,domain=-1:6.5] plot(\x,{-1.5*\x-1});
\node[gray] at (-0.4,1.3){$m<0$};
\fill (0,-1) circle(2pt) node[below right]{$A$};
\end{GR}
\textbf{التبرير التحليلي:} من أجل $x>0$:
\[f(x)=mx-1\iff x-\frac{\ln x}{x}=mx\iff 1-\frac{\ln x}{x^2}=m .\]
نضع $k(x)=1-\dfrac{\ln x}{x^2}$، فنجد $k'(x)=-\dfrac{1-2\ln x}{x^3}$، إذن $k$ متناقصة تماما على $]0;\sqrt e]$ ومتزايدة تماما على $[\sqrt e;+\infty[$، مع $k(\sqrt e)=1-\dfrac{1}{2e}=m_0$، $\displaystyle\lim_{x\to0^+}k(x)=+\infty$ و $\displaystyle\lim_{x\to+\infty}k(x)=1$ (بقيم أصغر من $1$ لما $x>1$).
\begin{TAB}
\tkzTabInit[lgt=2,espcl=3]{$x$/1,$k'(x)$/1,$k(x)$/2}{$0$,$\sqrt e$,$+\infty$}
\tkzTabLine{d,-,z,+,}
\tkzTabVar{D+/$+\infty$,-/$m_0$,+/$1$}
\end{TAB}
\textbf{النتيجة:}
\begin{itemize}
\item إذا كان $m<1-\dfrac1{2e}$: المعادلة لا تقبل حلولا.
\item إذا كان $m=1-\dfrac1{2e}$: المعادلة تقبل حلا وحيدا هو $\sqrt e$ (المماس $(T')$).
\item إذا كان $1-\dfrac1{2e}<m<1$: المعادلة تقبل حلين متمايزين (أحدهما في $]0;\sqrt e[$ والآخر في $]\sqrt e;+\infty[$).
\item إذا كان $m\geqslant1$: المعادلة تقبل حلا وحيدا (في $]0;1]$).
\end{itemize}
\end{hal}

%====================================================== المسألة 2
\begin{tamrin}{مسألة بمستوى البكالوريا -- شعبة رياضيات}
$m$ وسيط حقيقي. نعتبر عائلة الدوال $f_m$ المعرفة على $]0;+\infty[$ بـ: $f_m(x)=\dfrac{1+m\ln x}{x}$، و $(C_m)$ تمثيلها البياني في معلم متعامد ومتجانس $(O;\vec i,\vec j)$.
\QQ{1} بيّن أن جميع المنحنيات $(C_m)$ تشمل نقطة ثابتة $A$ يطلب تعيين إحداثييها، وأنها لا تشترك في أي نقطة أخرى.
\QQ{2} احسب، حسب قيم $m$، $\displaystyle\lim_{x\to0^+}f_m(x)$ و $\displaystyle\lim_{x\to+\infty}f_m(x)$.
\QQ{3}
\begin{sq}
\qa بيّن أنه من أجل كل $x>0$: $f_m'(x)=\dfrac{m-1-m\ln x}{x^2}$.
\qb ادرس، حسب قيم $m$، اتجاه تغير الدالة $f_m$ (ميّز الحالات: $m=0$، $m>0$، $m<0$).
\end{sq}
\QQ{4} من أجل $m\neq0$، نسمي $S_m$ النقطة من $(C_m)$ التي تقبل فيها $f_m$ قيمة حدية.
\begin{sq}
\qa عيّن إحداثيي $S_m$ بدلالة $m$.
\qb بيّن أنه لما يتغير $m$ في $\R^*$ فإن النقطة $S_m$ تمسح المنحنى $(\Gamma)$ الممثل للدالة $h$ المعرفة على $]0;e[\cup]e;+\infty[$ بـ: $h(x)=\dfrac{1}{x(1-\ln x)}$.
\end{sq}
\QQ{5}
\begin{sq}
\qa عيّن قيمة $m$ التي يكون من أجلها مماس $(C_m)$ عند النقطة $A$ موازيا للمستقيم ذي المعادلة $y=x$.
\qb بيّن أنه من أجل $m\neq0$، المنحنى $(C_m)$ يقطع حامل محور الفواصل في نقطة وحيدة $B_m$ يطلب تعيين فاصلتها.
\end{sq}
\QQ{6}
\begin{sq}
\qa ادرس الوضع النسبي للمنحنيين $(C_1)$ و $(C_{-1})$.
\qb شكّل جدولي تغيرات $f_1$ و $f_{-1}$ ثم أنشئ $(C_1)$ و $(C_{-1})$ في نفس المعلم.
\end{sq}
\QQ{7} $k$ وسيط حقيقي. ناقش بيانيا، حسب قيم $k$، عدد حلول المعادلة: $1+\ln x=kx$.
\end{tamrin}

\begin{hal}
\QQ{1} ليكن $m$ و $m'$ عددين حقيقيين متمايزين. من أجل $x>0$:
\[f_m(x)=f_{m'}(x)\iff 1+m\ln x=1+m'\ln x\iff (m-m')\ln x=0\iff x=1\]
(لأن $m-m'\neq0$). و $f_m(1)=\dfrac{1+0}{1}=1$ مهما كان $m$.
\par\bxt{كل المنحنيات $(C_m)$ تشمل النقطة $A(1;1)$ ولا تشترك في غيرها}
\QQ{2} \textbf{عند $+\infty$:} $f_m(x)=\dfrac1x+m\cdot\dfrac{\ln x}{x}$، و $\displaystyle\lim_{x\to+\infty}\frac1x=0$، $\displaystyle\lim_{x\to+\infty}\frac{\ln x}{x}=0$، إذن مهما كان $m$: \bx{\lim_{x\to+\infty}f_m(x)=0}، ومنه حامل محور الفواصل مستقيم مقارب أفقي لكل $(C_m)$.

\textbf{عند $0^+$:} $\displaystyle\lim_{x\to0^+}\frac1x=+\infty$ و $\displaystyle\lim_{x\to0^+}\frac{\ln x}{x}=-\infty$.
\begin{itemize}
\item إذا كان $m=0$: $f_0(x)=\dfrac1x$ و $\displaystyle\lim_{x\to0^+}f_0(x)=+\infty$.
\item إذا كان $m\neq0$: نكتب $f_m(x)=\dfrac1x\big(1+m\ln x\big)$. لدينا $\displaystyle\lim_{x\to0^+}m\ln x=-\infty$ إذا كان $m>0$، و $+\infty$ إذا كان $m<0$. إذن:
\[\bx{m>0:\ \lim_{x\to0^+}f_m(x)=-\infty}\qquad \bx{m<0:\ \lim_{x\to0^+}f_m(x)=+\infty}\]
\end{itemize}
في كل الحالات: المستقيم ذو المعادلة $x=0$ مقارب عمودي لـ $(C_m)$.
\QQ{3}
\pa $f_m=\dfrac uv$ مع $u(x)=1+m\ln x$، $u'(x)=\dfrac mx$، $v(x)=x$، $v'(x)=1$:
\[f_m'(x)=\frac{\frac mx\cdot x-(1+m\ln x)\cdot1}{x^2}=\bx{\frac{m-1-m\ln x}{x^2}}\]
\pb إشارة $f_m'(x)$ هي إشارة $N(x)=m-1-m\ln x$ لأن $x^2>0$.

\textbf{الحالة $m=0$:} $N(x)=-1<0$، إذن $f_0$ متناقصة تماما على $]0;+\infty[$.

\textbf{الحالة $m>0$:} $N(x)>0\iff m\ln x<m-1\iff \ln x<\dfrac{m-1}{m}\iff x<e^{1-\frac1m}$ (القسمة على $m>0$ لا تغير الاتجاه). نضع $x_m=e^{1-\frac1m}$:
\begin{TAB}
\tkzTabInit[lgt=2.2,espcl=3]{$x$/1,$f_m'(x)$/1,$f_m(x)$/2}{$0$,$x_m$,$+\infty$}
\tkzTabLine{d,+,z,-,}
\tkzTabVar{D-/$-\infty$,+/$f_m(x_m)$,-/$0$}
\end{TAB}
\textbf{الحالة $m<0$:} $N(x)>0\iff -m\ln x>1-m\iff \ln x>\dfrac{1-m}{-m}=1-\dfrac1m\iff x>x_m$ (القسمة على $-m>0$):
\begin{TAB}
\tkzTabInit[lgt=2.2,espcl=3]{$x$/1,$f_m'(x)$/1,$f_m(x)$/2}{$0$,$x_m$,$+\infty$}
\tkzTabLine{d,-,z,+,}
\tkzTabVar{D+/$+\infty$,-/$f_m(x_m)$,+/$0$}
\end{TAB}
\textbf{الخلاصة:} إذا كان $m>0$: $f_m$ متزايدة تماما على $]0;x_m]$ ومتناقصة تماما على $[x_m;+\infty[$ (قيمة حدية عظمى). إذا كان $m<0$: $f_m$ متناقصة تماما على $]0;x_m]$ ومتزايدة تماما على $[x_m;+\infty[$ (قيمة حدية صغرى).
\QQ{4}
\pa فاصلة $S_m$ هي $x_m=e^{1-\frac1m}$، و $\ln x_m=1-\dfrac1m$ إذن:
\[f_m(x_m)=\frac{1+m\big(1-\frac1m\big)}{x_m}=\frac{m}{x_m}=m\,e^{\frac1m-1}.\]
\[\bx{S_m\Big(e^{1-\frac1m}\ ;\ m\,e^{\frac1m-1}\Big)}\]
\pb نضع $X=x_m$ و $Y=f_m(x_m)$. لدينا $\ln X=1-\dfrac1m$، ومنه $\dfrac1m=1-\ln X$. بما أن $\dfrac1m\neq0$ فإن $\ln X\neq1$ أي $X\neq e$، و $m=\dfrac{1}{1-\ln X}$. إذن:
\[Y=\frac{m}{X}=\frac{1}{X(1-\ln X)}=h(X)\]
أي $S_m\in(\Gamma)$. \textbf{العكس:} إذا كان $X\in]0;e[\cup]e;+\infty[$ فإن $m=\dfrac{1}{1-\ln X}$ عدد حقيقي غير معدوم يحقق $x_m=X$، إذن كل نقطة من $(\Gamma)$ هي نقطة $S_m$.
\par\bxt{مجموعة النقط $S_m$ هي المنحنى $(\Gamma)$ كاملا}
\QQ{5}
\pa معامل توجيه المماس عند $A$ هو $f_m'(1)=\dfrac{m-1-0}{1}=m-1$. يكون موازيا للمستقيم $y=x$ إذا وفقط إذا $m-1=1$، أي \bx{m=2}. (معادلته عندئذ $y=x$.)
\pb من أجل $m\neq0$: $f_m(x)=0\iff1+m\ln x=0\iff\ln x=-\dfrac1m\iff x=e^{-\frac1m}$. حل وحيد، إذن \bx{B_m\big(e^{-\frac1m};0\big)}.
\QQ{6}
\pa $f_1(x)=\dfrac{1+\ln x}{x}$ و $f_{-1}(x)=\dfrac{1-\ln x}{x}$، إذن $f_1(x)-f_{-1}(x)=\dfrac{2\ln x}{x}$، وإشارته إشارة $\ln x$:
\begin{itemize}
\item على $]0;1[$: $(C_1)$ تحت $(C_{-1})$.
\item على $]1;+\infty[$: $(C_1)$ فوق $(C_{-1})$.
\item يتقاطعان في $A(1;1)$.
\end{itemize}
\pb \textbf{الدالة $f_1$} ($m=1>0$): $x_1=e^0=1$، $f_1(1)=1$، $f_1'(x)=-\dfrac{\ln x}{x^2}$، $B_1\big(e^{-1};0\big)$.
\begin{TAB}
\tkzTabInit[lgt=2,espcl=3]{$x$/1,$f_1'(x)$/1,$f_1(x)$/2}{$0$,$1$,$+\infty$}
\tkzTabLine{d,+,z,-,}
\tkzTabVar{D-/$-\infty$,+/$1$,-/$0$}
\end{TAB}
\textbf{الدالة $f_{-1}$} ($m=-1<0$): $x_{-1}=e^{2}$، $f_{-1}(e^2)=-e^{-2}\approx-0{,}14$، $f_{-1}'(x)=\dfrac{\ln x-2}{x^2}$، $B_{-1}(e;0)$.
\begin{TAB}
\tkzTabInit[lgt=2,espcl=3]{$x$/1,$f_{-1}'(x)$/1,$f_{-1}(x)$/2}{$0$,$e^2$,$+\infty$}
\tkzTabLine{d,-,z,+,}
\tkzTabVar{D+/$+\infty$,-/$-e^{-2}$,+/$0$}
\end{TAB}
\begin{GR}{-0.5}{9}{-2}{4}{1.3}
\Ticks{1,2,3,4,5,6,7,8}{-1,1,2,3}
\draw[cDef,very thick,samples=250,domain=0.15:9] plot(\x,{(1+ln(\x))/\x});
\node[cDef] at (2.5,1.15){$(C_1)$};
\draw[cThm,very thick,samples=250,domain=0.25:9] plot(\x,{(1-ln(\x))/\x});
\node[cThm] at (0.9,3.5){$(C_{-1})$};
\fill (1,1) circle(1.6pt) node[above]{$A$};
\fill ({exp(-1)},0) circle(1.4pt);
\fill ({e},0) circle(1.4pt) node[above right]{$B_{-1}$};
\fill ({exp(2)},{-exp(-2)}) circle(1.4pt);
\draw[densely dotted] ({exp(2)},0)--({exp(2)},{-exp(-2)});
\node[below,font=\scriptsize] at ({exp(2)},-0.15){$e^2$};
\end{GR}
\QQ{7} من أجل $x>0$: $1+\ln x=kx\iff\dfrac{1+\ln x}{x}=k\iff f_1(x)=k$.

حلول المعادلة هي فواصل نقط تقاطع $(C_1)$ مع المستقيم الأفقي $y=k$. من جدول تغيرات $f_1$:
\begin{itemize}
\item إذا كان $k>1$: لا توجد حلول.
\item إذا كان $k=1$: حل وحيد هو $x=1$.
\item إذا كان $0<k<1$: حلان متمايزان، أحدهما في $\big]e^{-1};1\big[$ والآخر في $]1;+\infty[$.
\item إذا كان $k\leqslant0$: حل وحيد في $\big]0;e^{-1}\big]$.
\end{itemize}
\end{hal}

%====================================================== المسألة 3
\begin{tamrin}{مسألة بمستوى البكالوريا -- شعبة رياضيات}
\textbf{الجزء الأول:} $g$ الدالة المعرفة على $]0;+\infty[$ بـ: $g(x)=x-2\ln x$.
\QQ{1} ادرس اتجاه تغير الدالة $g$.
\QQ{2} استنتج أنه من أجل كل $x>0$: $g(x)>0$.

\textbf{الجزء الثاني:} $f$ الدالة المعرفة على $]0;+\infty[$ بـ: $f(x)=x-(\ln x)^2$، و $(C_f)$ تمثيلها البياني في معلم متعامد ومتجانس.
\QQ{1}
\begin{sq}
\qa احسب $\displaystyle\lim_{x\to0^+}f(x)$ وفسّر النتيجة هندسيا.
\qb احسب $\displaystyle\lim_{x\to+\infty}f(x)$ و $\displaystyle\lim_{x\to+\infty}\frac{f(x)}{x}$ و $\displaystyle\lim_{x\to+\infty}\big[f(x)-x\big]$، ثم فسّر النتائج هندسيا.
\end{sq}
\QQ{2} بيّن أن $f'(x)=\dfrac{g(x)}{x}$ ثم شكّل جدول تغيرات الدالة $f$.
\QQ{3}
\begin{sq}
\qa ادرس الوضع النسبي لـ $(C_f)$ والمستقيم $(\Delta):\ y=x$.
\qb بيّن أن $(\Delta)$ مماس لـ $(C_f)$ في نقطة يطلب تعيينها.
\end{sq}
\QQ{4} بيّن أن المعادلة $f(x)=0$ تقبل حلا وحيدا $x_0$ وأن $0{,}49<x_0<0{,}5$.
\QQ{5} بيّن أن $(C_f)$ يقبل نقطة انعطاف $I$ يطلب تعيين إحداثييها.
\QQ{6}
\begin{sq}
\qa اكتب معادلة المماس $(T)$ للمنحنى $(C_f)$ عند النقطة ذات الفاصلة $e^2$، وتحقق أنه يشمل المبدأ $O$.
\qb أنشئ $(\Delta)$ و $(T)$ و $(C_f)$.
\end{sq}
\QQ{7} $m$ وسيط حقيقي. ناقش بيانيا، حسب قيم $m$، عدد حلول المعادلة: $x-(\ln x)^2=mx$.
\end{tamrin}

\begin{hal}
\textbf{الجزء الأول:}
\QQ{1} $g'(x)=1-\dfrac2x=\dfrac{x-2}{x}$، وإشارتها إشارة $x-2$ (لأن $x>0$): $g$ متناقصة تماما على $]0;2]$ ومتزايدة تماما على $[2;+\infty[$.
\begin{TAB}
\tkzTabInit[lgt=2,espcl=3]{$x$/1,$g'(x)$/1,$g(x)$/2}{$0$,$2$,$+\infty$}
\tkzTabLine{d,-,z,+,}
\tkzTabVar{D+/$+\infty$,-/$2-2\ln2$,+/$+\infty$}
\end{TAB}
(النهايات: $\displaystyle\lim_{x\to0^+}(-2\ln x)=+\infty$؛ و $g(x)=x\Big(1-2\dfrac{\ln x}{x}\Big)$ مع $\dfrac{\ln x}{x}\to0$ إذن $\displaystyle\lim_{x\to+\infty}g(x)=+\infty$.)
\QQ{2} القيمة الحدية الصغرى لـ $g$ هي $g(2)=2-2\ln2=2(1-\ln2)$. بما أن $\ln2<1$ (لأن $2<e$) فإن $g(2)>0$ ($\approx0{,}61$). إذن من أجل كل $x>0$: $g(x)\geqslant g(2)>0$، أي \bx{g(x)>0}.

\textbf{الجزء الثاني:}
\QQ{1}
\pa $\displaystyle\lim_{x\to0^+}\ln x=-\infty$ إذن $\displaystyle\lim_{x\to0^+}(\ln x)^2=+\infty$، ومنه \bx{\lim_{x\to0^+}f(x)=-\infty}. المستقيم $x=0$ مقارب عمودي لـ $(C_f)$.
\pb نكتب $f(x)=x\Big(1-\dfrac{(\ln x)^2}{x}\Big)$ مع $\dfrac{(\ln x)^2}{x}=4\Big(\dfrac{\ln\sqrt x}{\sqrt x}\Big)^2$ (لأن $\ln x=2\ln\sqrt x$ و $x=(\sqrt x)^2$). لما $x\to+\infty$ فإن $\sqrt x\to+\infty$ و $\dfrac{\ln\sqrt x}{\sqrt x}\to0$، إذن $\displaystyle\lim_{x\to+\infty}\frac{(\ln x)^2}{x}=0$. ومنه:
\begin{gather*}\bx{\lim_{x\to+\infty}f(x)=+\infty}\qquad \bx{\lim_{x\to+\infty}\frac{f(x)}{x}=\lim_{x\to+\infty}\Big(1-\frac{(\ln x)^2}{x}\Big)=1}\\ \bx{\lim_{x\to+\infty}\big[f(x)-x\big]=\lim_{x\to+\infty}\big(-(\ln x)^2\big)=-\infty}\end{gather*}
التفسير: $(C_f)$ يقبل بجوار $+\infty$ فرعا قطعيا (مكافئا) اتجاهه المستقيم $(\Delta):\ y=x$.
\QQ{2} $\big((\ln x)^2\big)'=2\cdot\dfrac1x\cdot\ln x$، إذن $f'(x)=1-\dfrac{2\ln x}{x}=\dfrac{x-2\ln x}{x}=\bx{\dfrac{g(x)}{x}}$.

بما أن $g(x)>0$ و $x>0$ فإن $f'(x)>0$: $f$ متزايدة تماما على $]0;+\infty[$.
\begin{TAB}
\tkzTabInit[lgt=2,espcl=4]{$x$/1,$f'(x)$/1,$f(x)$/2}{$0$,$+\infty$}
\tkzTabLine{d,+,}
\tkzTabVar{D-/$-\infty$,+/$+\infty$}
\end{TAB}
\QQ{3}
\pa $f(x)-x=-(\ln x)^2\leqslant0$، وتنعدم فقط من أجل $\ln x=0$ أي $x=1$:
\par $(C_f)$ تحت $(\Delta)$ على $]0;1[\cup]1;+\infty[$، ويشتركان في النقطة $(1;1)$.
\pb $f(1)=1$ و $f'(1)=\dfrac{g(1)}{1}=1-0=1$، إذن المماس عند النقطة ذات الفاصلة $1$ معادلته $y=1\cdot(x-1)+1=x$، أي هو $(\Delta)$.
\par\bxt{$(\Delta)$ مماس لـ $(C_f)$ في النقطة $(1;1)$}
\QQ{4} $f$ مستمرة ومتزايدة تماما على $]0;+\infty[$ وتأخذ قيمها في $]-\infty;+\infty[$، إذن المعادلة $f(x)=0$ تقبل حلا وحيدا $x_0$ (مبرهنة القيم المتوسطة مع الرتابة).

$f(0{,}49)=0{,}49-(\ln0{,}49)^2\approx-0{,}019<0$ و $f(0{,}5)=0{,}5-(\ln0{,}5)^2\approx0{,}020>0$، إذن \bx{0{,}49<x_0<0{,}5}.
\QQ{5} $f''(x)=\Big(1-\dfrac{2\ln x}{x}\Big)'=-2\cdot\dfrac{1-\ln x}{x^2}=\dfrac{2(\ln x-1)}{x^2}$.

إشارة $f''(x)$ هي إشارة $\ln x-1$: سالبة على $]0;e[$ وموجبة على $]e;+\infty[$. $f''$ تنعدم عند $e$ مغيرة إشارتها، إذن $(C_f)$ يقبل نقطة انعطاف $I\big(e;f(e)\big)$ مع $f(e)=e-1$:
\[\bx{I(e\,;\,e-1)}\]
\QQ{6}
\pa $f(e^2)=e^2-(2)^2=e^2-4$، $f'(e^2)=1-\dfrac{2\cdot2}{e^2}=1-\dfrac4{e^2}$.
\[(T):\ y=\Big(1-\frac4{e^2}\Big)(x-e^2)+e^2-4=\Big(1-\frac4{e^2}\Big)x-e^2+4+e^2-4\]
\[\bx{(T):\ y=\Big(1-\frac4{e^2}\Big)x}\qquad\Big(1-\frac4{e^2}\approx0{,}46\Big)\]
ترتيبه عند المبدأ معدوم، إذن $(T)$ يشمل $O(0;0)$. \ok
\pb بعض القيم: $f(0{,}3)\approx-1{,}15$، $f(2)\approx1{,}52$، $f(4)\approx2{,}08$، $f(e^2)\approx3{,}39$، $f(10)\approx4{,}70$.
\begin{GR}{-1}{11}{-3}{7}{0.85}
\Ticks{1,2,3,4,5,6,7,8,9,10}{-2,-1,1,2,3,4,5,6}
\draw[cThm,thick,domain=-1:11.3] plot(\x,{\x});
\node[cThm,above left] at (5,5){$(\Delta)$};
\draw[cMeth,thick,domain=-1:11.3] plot(\x,{(1-4/exp(2))*\x});
\node[cMeth,below] at (10,4.5){$(T)$};
\draw[cDef,very thick,samples=250,domain=0.12:11.3] plot(\x,{\x-ln(\x)*ln(\x)});
\node[cDef] at (9.5,3.6){$(C_f)$};
\fill (1,1) circle(1.6pt);
\fill ({e},{e-1}) circle(1.6pt) node[above left]{$I$};
\fill ({exp(2)},{exp(2)-4}) circle(1.6pt);
\fill (0.495,0) circle(1.4pt) node[below right]{$x_0$};
\end{GR}
\QQ{7} حلول المعادلة $f(x)=mx$ هي فواصل نقط تقاطع $(C_f)$ مع المستقيم $(D_m):\ y=mx$ الذي يدور حول المبدأ $O$.

المستقيمان المميزان المارّان بـ $O$ والمماسان لـ $(C_f)$ هما: $(T)$ بمعامل توجيه $m_1=1-\dfrac4{e^2}$ و $(\Delta)$ بمعامل توجيه $1$.
\begin{GR}{-1}{11}{-3}{7}{0.85}
\Ticks{2,4,6,8,10}{-2,2,4,6}
\draw[cDef,very thick,samples=250,domain=0.12:11.3] plot(\x,{\x-ln(\x)*ln(\x)});
\draw[cThm,thick,domain=-1:11.3] plot(\x,{\x});
\node[cThm,left] at (5.5,5.5){$m=1$};
\draw[cMeth,thick,domain=-1:11.3] plot(\x,{(1-4/exp(2))*\x});
\node[cMeth,below] at (10.3,4.7){$m=m_1$};
\draw[cEx,dashed,thick,domain=-1:11.3] plot(\x,{0.7*\x});
\node[cEx,above left] at (8,5.6){$m_1<m<1$};
\draw[gray,dashed,domain=-1:11.3] plot(\x,{0.2*\x});
\node[gray,below] at (10,2){$m<m_1$};
\draw[cSol,dashed,thick,domain=-1:11.3] plot(\x,{1.6*\x});
\node[cSol,left] at (3,4.8){$m>1$};
\draw[gray,dashed,domain=-1:11.3] plot(\x,{-0.5*\x});
\fill (0,0) circle(2pt);
\end{GR}
\textbf{التبرير:} من أجل $x>0$: $f(x)=mx\iff 1-\dfrac{(\ln x)^2}{x}=m$. نضع $\varphi(x)=1-\dfrac{(\ln x)^2}{x}$:
\[\varphi'(x)=-\frac{2\ln x\cdot\frac1x\cdot x-(\ln x)^2}{x^2}=\frac{\ln x\,(\ln x-2)}{x^2}.\]
$\varphi'$ تنعدم عند $1$ و $e^2$، $\varphi(1)=1$، $\varphi(e^2)=1-\dfrac4{e^2}=m_1$، $\displaystyle\lim_{x\to0^+}\varphi(x)=-\infty$، $\displaystyle\lim_{x\to+\infty}\varphi(x)=1$.
\begin{TAB}
\tkzTabInit[lgt=2,espcl=2.6]{$x$/1,$\varphi'(x)$/1,$\varphi(x)$/2}{$0$,$1$,$e^2$,$+\infty$}
\tkzTabLine{d,+,z,-,z,+,}
\tkzTabVar{D-/$-\infty$,+/$1$,-/$m_1$,+/$1$}
\end{TAB}
\textbf{النتيجة:}
\begin{itemize}
\item إذا كان $m>1$: لا توجد حلول.
\item إذا كان $m=1$: حل وحيد هو $x=1$ (التماس مع $(\Delta)$).
\item إذا كان $1-\dfrac4{e^2}<m<1$: ثلاثة حلول متمايزة، في $]0;1[$ و $]1;e^2[$ و $]e^2;+\infty[$.
\item إذا كان $m=1-\dfrac4{e^2}$: حلان، أحدهما في $]0;1[$ والآخر $e^2$ (التماس مع $(T)$).
\item إذا كان $m<1-\dfrac4{e^2}$: حل وحيد في $]0;1[$.
\end{itemize}
\end{hal}

%====================================================== المسألة 4
\begin{tamrin}{مسألة بمستوى البكالوريا -- شعبة رياضيات}
$f$ الدالة المعرفة على $D=]-\infty;-1[\cup]1;+\infty[$ بـ: $f(x)=x+\ln\Big(\dfrac{x-1}{x+1}\Big)$، و $(C_f)$ تمثيلها البياني في معلم متعامد ومتجانس $(O;\vec i,\vec j)$.
\QQ{1} بيّن أن $f$ فردية، وفسّر ذلك هندسيا.
\QQ{2} احسب نهايات $f$ عند $+\infty$ و $-\infty$ وعند $1$ بقيم أكبر و $-1$ بقيم أصغر، وفسّر النتائج هندسيا.
\QQ{3} بيّن أنه من أجل كل $x$ من $D$: $f'(x)=\dfrac{x^2+1}{x^2-1}$، ثم شكّل جدول تغيرات $f$.
\QQ{4}
\begin{sq}
\qa بيّن أن المستقيم $(\Delta):\ y=x$ مقارب مائل لـ $(C_f)$، ثم ادرس وضعيتهما النسبية.
\qb هل يقبل $(C_f)$ مماسا موازيا لـ $(\Delta)$؟ علّل.
\qc عيّن فواصل النقط التي يكون فيها المماس موازيا للمستقيم ذي المعادلة $y=2x$، واكتب معادلة المماس عند النقطة ذات الفاصلة الموجبة.
\end{sq}
\QQ{5} بيّن أن المعادلة $f(x)=0$ تقبل في $D$ حلين بالضبط $\alpha$ و $-\alpha$ حيث $1{,}5<\alpha<1{,}6$.
\QQ{6} أنشئ $(\Delta)$ و $(C_f)$.
\QQ{7} $m$ وسيط حقيقي.
\begin{sq}
\qa ناقش، حسب قيم $m$، عدد حلول المعادلة $f(x)=x+m$، ثم عيّن الحل عند وجوده.
\qb نعتبر الدالة $\varphi$ المعرفة على $D$ بـ: $\varphi(x)=\dfrac1x\ln\Big(\dfrac{x-1}{x+1}\Big)$. بيّن أن $\varphi$ زوجية، ثم ادرس اتجاه تغيرها على $]1;+\infty[$.
\qc استنتج، حسب قيم $m$، عدد حلول المعادلة $f(x)=mx$، وفسّر ذلك بيانيا.
\end{sq}
\end{tamrin}

\begin{hal}
\QQ{1} من أجل كل $x\in D$ فإن $-x\in D$ (لأن $x<-1\iff-x>1$). و:
\[f(-x)=-x+\ln\Big(\frac{-x-1}{-x+1}\Big)=-x+\ln\Big(\frac{x+1}{x-1}\Big)=-x-\ln\Big(\frac{x-1}{x+1}\Big)=-f(x).\]
\bxt{$f$ فردية}، إذن المبدأ $O$ مركز تناظر لـ $(C_f)$. يكفي دراستها على $]1;+\infty[$.
\QQ{2} عند $+\infty$: $\dfrac{x-1}{x+1}\to1$ إذن $\ln\Big(\dfrac{x-1}{x+1}\Big)\to\ln1=0$، ومنه \bx{\lim_{x\to+\infty}f(x)=+\infty}، وبالتناظر \bx{\lim_{x\to-\infty}f(x)=-\infty}.

عند $1^+$: $x-1\to0^+$ و $x+1\to2$، إذن $\dfrac{x-1}{x+1}\to0^+$ و $\ln\Big(\dfrac{x-1}{x+1}\Big)\to-\infty$، ومنه \bx{\lim_{x\to1^+}f(x)=-\infty}.

وبالتناظر ($f(x)=-f(-x)$ و $-x\to1^+$ لما $x\to-1^-$): \bx{\lim_{x\to-1^-}f(x)=+\infty}.

التفسير: المستقيمان $x=1$ و $x=-1$ مقاربان عموديان لـ $(C_f)$.
\QQ{3} نكتب $f(x)=x+\ln|x-1|-\ln|x+1|$ (لأن $\dfrac{x-1}{x+1}>0$ على $D$)، إذن:
\[f'(x)=1+\frac1{x-1}-\frac1{x+1}=1+\frac{2}{x^2-1}=\bx{\frac{x^2+1}{x^2-1}}\]
على $D$: $x^2>1$ إذن $x^2-1>0$، و $x^2+1>0$، ومنه $f'(x)>0$: $f$ متزايدة تماما على كل من المجالين.
\begin{TAB}
\tkzTabInit[lgt=2,espcl=2.8]{$x$/1,$f'(x)$/1,$f(x)$/2}{$-\infty$,$-1$,$1$,$+\infty$}
\tkzTabLine{,+,d,h,d,+,}
\tkzTabVar{-/$-\infty$,+DH/$+\infty$,D-/$-\infty$,+/$+\infty$}
\end{TAB}
\QQ{4}
\pa $f(x)-x=\ln\Big(\dfrac{x-1}{x+1}\Big)\to0$ عند $+\infty$ وعند $-\infty$، إذن \bxt{$(\Delta): y=x$ مقارب مائل لـ $(C_f)$} بجوار $\pm\infty$.

الوضعية: نقارن $\dfrac{x-1}{x+1}$ بـ $1$: $\dfrac{x-1}{x+1}-1=\dfrac{-2}{x+1}$.
\begin{itemize}
\item على $]1;+\infty[$: $\dfrac{-2}{x+1}<0$ إذن $0<\dfrac{x-1}{x+1}<1$ و $f(x)-x<0$: $(C_f)$ تحت $(\Delta)$.
\item على $]-\infty;-1[$: $x+1<0$ إذن $\dfrac{-2}{x+1}>0$ و $\dfrac{x-1}{x+1}>1$، $f(x)-x>0$: $(C_f)$ فوق $(\Delta)$.
\end{itemize}
\pb مماس موازٍ لـ $(\Delta)$ يعني $f'(x)=1$، لكن $f'(x)=1+\dfrac{2}{x^2-1}>1$ على $D$. إذن \textbf{لا يوجد} مماس موازٍ لـ $(\Delta)$.
\pc $f'(x)=2\iff x^2+1=2x^2-2\iff x^2=3\iff x=\sqrt3$ أو $x=-\sqrt3$ (وكلاهما في $D$).

عند $\sqrt3$: $\dfrac{\sqrt3-1}{\sqrt3+1}=\dfrac{(\sqrt3-1)^2}{3-1}=\dfrac{4-2\sqrt3}{2}=2-\sqrt3$، إذن $f(\sqrt3)=\sqrt3+\ln(2-\sqrt3)\approx0{,}42$.
\[y=2(x-\sqrt3)+\sqrt3+\ln(2-\sqrt3)\quad\Longrightarrow\quad \bx{y=2x-\sqrt3+\ln(2-\sqrt3)}\]
\QQ{5} على $]1;+\infty[$: $f$ مستمرة ومتزايدة تماما وتأخذ قيمها في $]-\infty;+\infty[$، إذن المعادلة $f(x)=0$ تقبل حلا وحيدا $\alpha$ في هذا المجال.
\par $f(1{,}5)=1{,}5+\ln0{,}2\approx-0{,}11<0$ و $f(1{,}6)=1{,}6+\ln\dfrac{0{,}6}{2{,}6}\approx0{,}13>0$، إذن \bx{1{,}5<\alpha<1{,}6}.

وبما أن $f$ فردية: $f(-\alpha)=-f(\alpha)=0$، و $-\alpha$ هو الحل الوحيد في $]-\infty;-1[$ (لأن $f$ متزايدة تماما عليه). إذن للمعادلة حلان بالضبط $\alpha$ و $-\alpha$.
\QQ{6}
\begin{GR}{-5}{5}{-5}{5}{0.8}
\Ticks{-4,-3,-2,2,3,4}{-4,-3,-2,-1,1,2,3,4}
\draw[cThm,thick,domain=-5:5.3] plot(\x,{\x});
\node[cThm,above left] at (4,4){$(\Delta)$};
\draw[gray,dashed] (1,-5)--(1,5) (-1,-5)--(-1,5);
\draw[cDef,very thick,samples=250,domain=1.002:5.3] plot(\x,{\x+ln((\x-1)/(\x+1))});
\draw[cDef,very thick,samples=250,domain=-5.3:-1.002] plot(\x,{\x+ln((\x-1)/(\x+1))});
\node[cDef] at (3.5,1.9){$(C_f)$};
\fill (1.5434,0) circle(1.5pt) node[below right]{$\alpha$};
\fill (-1.5434,0) circle(1.5pt) node[above left]{$-\alpha$};
\draw[cMeth,thick,domain=0.8:3.2] plot(\x,{2*\x-1.7321-1.3170});
\fill (1.7321,0.4151) circle(1.5pt);
\end{GR}
\QQ{7}
\pa $f(x)=x+m\iff\ln\Big(\dfrac{x-1}{x+1}\Big)=m\iff\dfrac{x-1}{x+1}=e^m$.
\par ومن أجل $x\in D$: $x-1=e^m(x+1)\iff x(1-e^m)=1+e^m$.
\begin{itemize}
\item إذا كان $m=0$: نحصل على $0=2$، مستحيل: \textbf{لا توجد حلول} (المستقيم هو $(\Delta)$ المقارب).
\item إذا كان $m\neq0$: حل وحيد $x=\dfrac{1+e^m}{1-e^m}$، ويتحقق أنه من $D$: $\big|1+e^m\big|>\big|1-e^m\big|$ لأن $e^m>0$، إذن $|x|>1$.
\end{itemize}
\[\bx{m\neq0:\ x=\frac{1+e^m}{1-e^m}}\]
(إذا كان $m<0$ فالحل في $]1;+\infty[$، وإذا كان $m>0$ فالحل في $]-\infty;-1[$، وهذا يوافق الوضعية النسبية في السؤال 4.)
\pb $D$ متناظر بالنسبة إلى $0$، و $\varphi(-x)=\dfrac{1}{-x}\ln\Big(\dfrac{x+1}{x-1}\Big)=\dfrac{-1}{-x}\ln\Big(\dfrac{x-1}{x+1}\Big)=\varphi(x)$: \bxt{$\varphi$ زوجية}.

على $]1;+\infty[$ نضع $L(x)=\ln\Big(\dfrac{x-1}{x+1}\Big)$ فيكون $L'(x)=\dfrac{2}{x^2-1}$ و:
\[\varphi'(x)=\frac{L'(x)\cdot x-L(x)}{x^2}=\frac{1}{x^2}\left[\frac{2x}{x^2-1}-\ln\Big(\frac{x-1}{x+1}\Big)\right].\]
على $]1;+\infty[$: $\dfrac{2x}{x^2-1}>0$ و $\ln\Big(\dfrac{x-1}{x+1}\Big)<0$ (لأن $0<\dfrac{x-1}{x+1}<1$)، إذن $\varphi'(x)>0$: \bxt{$\varphi$ متزايدة تماما على $]1;+\infty[$}.

النهايات: $\displaystyle\lim_{x\to1^+}\varphi(x)=\frac11\times(-\infty)=-\infty$ و $\displaystyle\lim_{x\to+\infty}\varphi(x)=0\times0=0$. وبالزوجية $\varphi$ متناقصة تماما على $]-\infty;-1[$ من $0$ إلى $-\infty$.
\begin{TAB}
\tkzTabInit[lgt=2,espcl=2.8]{$x$/1,$\varphi(x)$/2}{$-\infty$,$-1$,$1$,$+\infty$}
\tkzTabVar{+/$0$,-DH/$-\infty$,D-/$-\infty$,+/$0$}
\end{TAB}
\pc من أجل $x\in D$: $f(x)=mx\iff x+L(x)=mx\iff\dfrac{L(x)}{x}=m-1\iff\varphi(x)=m-1$.
\begin{itemize}
\item إذا كان $m-1<0$ أي $m<1$: للمعادلة $\varphi(x)=m-1$ حل وحيد في $]1;+\infty[$ وحل وحيد في $]-\infty;-1[$ (متناظران): \textbf{حلان}.
\item إذا كان $m\geqslant1$: $m-1\geqslant0$ و $\varphi(x)<0$ دوما: \textbf{لا توجد حلول}.
\end{itemize}
\textbf{التفسير البياني:} المستقيم $(D_m):\ y=mx$ يدور حول المبدأ $O$ (مركز تناظر $(C_f)$). إذا كان معامل توجيهه أصغر تماما من معامل توجيه المقارب $(\Delta)$ فإنه يقطع كل فرع من $(C_f)$ في نقطة واحدة، والنقطتان متناظرتان بالنسبة إلى $O$؛ وإذا كان $m\geqslant1$ فإنه لا يقطع $(C_f)$.
\begin{GR}{-5}{5}{-5}{5}{0.8}
\Ticks{-4,-2,2,4}{-4,-2,2,4}
\draw[cDef,very thick,samples=250,domain=1.002:5.3] plot(\x,{\x+ln((\x-1)/(\x+1))});
\draw[cDef,very thick,samples=250,domain=-5.3:-1.002] plot(\x,{\x+ln((\x-1)/(\x+1))});
\draw[cThm,thick,domain=-5:5.3] plot(\x,{\x});
\node[cThm,left] at (4.3,4.3){$m=1$};
\draw[cEx,dashed,thick,domain=-5.3:5.3] plot(\x,{0.5*\x});
\node[cEx,below] at (4.5,2.25){$m<1$};
\draw[cEx,dashed,thick,domain=-5.3:5.3] plot(\x,{-0.8*\x});
\node[cEx,above right] at (-4,3.2){$m<1$};
\draw[cSol,dashed,thick,domain=-5.3:5.3] plot(\x,{2*\x});
\node[cSol,left] at (2,4){$m>1$};
\fill (0,0) circle(2pt);
\end{GR}
\end{hal}

%====================================================== المسألة 5
\begin{tamrin}{مسألة بمستوى البكالوريا -- شعبة رياضيات}
\textbf{الجزء الأول:} $h$ الدالة المعرفة على $\R$ بـ: $h(x)=e^x-x-1$.
\QQ{1} ادرس اتجاه تغير $h$، ثم استنتج أنه من أجل كل عدد حقيقي $x$: $e^x-x\geqslant1$.

\textbf{الجزء الثاني:} $f$ الدالة المعرفة على $\R$ بـ: $f(x)=\ln\big(e^x-x\big)$، و $(C_f)$ تمثيلها البياني في معلم متعامد ومتجانس.
\QQ{1} برّر أن $f$ معرفة على $\R$، وأن $f(x)\geqslant0$ من أجل كل $x$.
\QQ{2}
\begin{sq}
\qa بيّن أنه من أجل كل $x<0$: $f(x)=\ln(-x)+\ln\Big(1-\dfrac{e^x}{x}\Big)$، ثم احسب $\displaystyle\lim_{x\to-\infty}f(x)$ و $\displaystyle\lim_{x\to-\infty}\frac{f(x)}{x}$ وفسّر النتيجة.
\qb بيّن أنه من أجل كل $x$: $f(x)=x+\ln\big(1-xe^{-x}\big)$، ثم استنتج $\displaystyle\lim_{x\to+\infty}f(x)$ وأن المستقيم $(\Delta):\ y=x$ مقارب مائل لـ $(C_f)$ بجوار $+\infty$.
\qc ادرس الوضع النسبي لـ $(C_f)$ و $(\Delta)$.
\end{sq}
\QQ{3} احسب $f'(x)$ ثم شكّل جدول تغيرات $f$.
\QQ{4} بيّن أن $(C_f)$ يقبل مماسا $(T)$ موازيا لـ $(\Delta)$ يطلب تعيين معادلة له.
\QQ{5} بيّن أن المعادلة $f(x)=1$ تقبل حلين $\alpha$ و $\beta$ حيث $-2{,}7<\alpha<-2{,}6$ و $1{,}4<\beta<1{,}5$.
\QQ{6} أنشئ $(\Delta)$ و $(T)$ و $(C_f)$، ثم ناقش بيانيا، حسب قيم الوسيط الحقيقي $m$، عدد حلول المعادلة $e^x-x=e^m$.

\textbf{الجزء الثالث:} نعتبر الدالة $\psi$ المعرفة على المجال $I=[1;2]$ بـ: $\psi(x)=\ln(x+e)$.
\QQ{1} بيّن أن $\beta$ هو الحل الوحيد في $I$ للمعادلة $\psi(x)=x$ (أي $\beta$ نقطة صامدة لـ $\psi$).
\QQ{2} بيّن أنه من أجل كل $x\in I$: $\psi(x)\in I$.
\QQ{3} بيّن أنه من أجل كل $x\in I$: $\big|\psi(x)-\beta\big|\leqslant\dfrac{1}{e+1}\,|x-\beta|$.
\par(يمكن دراسة الدالة $u$ المعرفة على $I$ بـ: $u(x)=\psi(x)-\dfrac{x}{e+1}$.)
\QQ{4} باستعمال $1{,}4<\beta<1{,}5$ وتزايد $\psi$، عيّن حصرا للعدد $\beta$ سعته أصغر من $0{,}03$.
\end{tamrin}

\begin{hal}
\textbf{الجزء الأول:}
\QQ{1} $h'(x)=e^x-1$. $h'(x)>0\iff e^x>1\iff x>0$: $h$ متناقصة تماما على $]-\infty;0]$ ومتزايدة تماما على $[0;+\infty[$، فتقبل قيمة حدية صغرى $h(0)=1-0-1=0$.
\begin{TAB}
\tkzTabInit[lgt=2,espcl=3]{$x$/1,$h'(x)$/1,$h(x)$/2}{$-\infty$,$0$,$+\infty$}
\tkzTabLine{,-,z,+,}
\tkzTabVar{+/$+\infty$,-/$0$,+/$+\infty$}
\end{TAB}
إذن من أجل كل $x$: $h(x)\geqslant0$ أي $e^x-x-1\geqslant0$، ومنه \bx{e^x-x\geqslant1}.

\textbf{الجزء الثاني:}
\QQ{1} $e^x-x\geqslant1>0$ من أجل كل $x$، إذن $\ln(e^x-x)$ معرفة على $\R$. وبما أن $\ln$ متزايدة تماما: $f(x)=\ln(e^x-x)\geqslant\ln1=0$. \bx{f(x)\geqslant0}، والمساواة فقط عند $x=0$.
\QQ{2}
\pa من أجل $x<0$: $e^x-x=(-x)\Big(1-\dfrac{e^x}{x}\Big)$ مع $-x>0$ و $1-\dfrac{e^x}{x}>0$ (لأن $\dfrac{e^x}{x}<0$)، إذن:
\[f(x)=\ln(-x)+\ln\Big(1-\frac{e^x}{x}\Big).\]
لما $x\to-\infty$: $e^x\to0$ إذن $\dfrac{e^x}{x}\to0$ و $\ln\Big(1-\dfrac{e^x}{x}\Big)\to0$، و $\ln(-x)\to+\infty$. ومنه \bx{\lim_{x\to-\infty}f(x)=+\infty}.
\[\frac{f(x)}{x}=\frac{\ln(-x)}{x}+\frac1x\ln\Big(1-\frac{e^x}{x}\Big)=-\frac{\ln(-x)}{-x}+\frac1x\ln\Big(1-\frac{e^x}{x}\Big).\]
بوضع $X=-x\to+\infty$: $\dfrac{\ln X}{X}\to0$، والحد الثاني $0\times0=0$، إذن \bx{\lim_{x\to-\infty}\frac{f(x)}{x}=0}.

التفسير: $(C_f)$ يقبل بجوار $-\infty$ فرعا قطعيا (مكافئا) اتجاهه حامل محور الفواصل.
\pb $e^x-x=e^x\big(1-xe^{-x}\big)$ مع $1-xe^{-x}=e^{-x}(e^x-x)>0$، إذن:
\[f(x)=\ln e^x+\ln\big(1-xe^{-x}\big)=x+\ln\big(1-xe^{-x}\big).\]
لما $x\to+\infty$: $xe^{-x}=\dfrac{x}{e^x}\to0$ (التزايد المقارن)، إذن $\ln(1-xe^{-x})\to0$، ومنه \bx{\lim_{x\to+\infty}f(x)=+\infty} و $\displaystyle\lim_{x\to+\infty}\big[f(x)-x\big]=0$:
\par\bxt{$(\Delta): y=x$ مقارب مائل لـ $(C_f)$ بجوار $+\infty$}
\pc $f(x)-x=\ln\big(1-xe^{-x}\big)$، وإشارته تتعلق بمقارنة $1-xe^{-x}$ بـ $1$، أي بإشارة $-xe^{-x}$ وهي إشارة $-x$:
\begin{itemize}
\item إذا كان $x<0$: $1-xe^{-x}>1$ و $f(x)-x>0$: $(C_f)$ فوق $(\Delta)$.
\item إذا كان $x=0$: $(C_f)$ و $(\Delta)$ يتقاطعان في $O$ (النقطة $O$ هي النقطة الوحيدة التي تحقق $f(x)=x$).
\item إذا كان $x>0$: $0<1-xe^{-x}<1$ و $f(x)-x<0$: $(C_f)$ تحت $(\Delta)$.
\end{itemize}
\QQ{3} $f=\ln u$ مع $u(x)=e^x-x>0$، إذن $f'(x)=\dfrac{u'(x)}{u(x)}=\bx{\dfrac{e^x-1}{e^x-x}}$.
\par إشارة $f'(x)$ هي إشارة $e^x-1$ (لأن $e^x-x>0$): سالبة على $]-\infty;0[$ وموجبة على $]0;+\infty[$.
\begin{TAB}
\tkzTabInit[lgt=2,espcl=3]{$x$/1,$f'(x)$/1,$f(x)$/2}{$-\infty$,$0$,$+\infty$}
\tkzTabLine{,-,z,+,}
\tkzTabVar{+/$+\infty$,-/$0$,+/$+\infty$}
\end{TAB}
\QQ{4} $f'(x)=1\iff e^x-1=e^x-x\iff x=1$. إذن يوجد مماس وحيد $(T)$ موازٍ لـ $(\Delta)$ عند النقطة ذات الفاصلة $1$، مع $f(1)=\ln(e-1)\approx0{,}54$:
\[\bx{(T):\ y=x-1+\ln(e-1)}\]
\QQ{5} $f(x)=1\iff\ln(e^x-x)=1\iff e^x-x=e$.
\begin{itemize}
\item على $]-\infty;0]$: $f$ مستمرة ومتناقصة تماما من $+\infty$ إلى $0$، و $1\in]0;+\infty[$، إذن حل وحيد $\alpha$. $f(-2{,}7)-1$ و $f(-2{,}6)-1$ لهما إشارتا $e^{-2{,}7}+2{,}7-e\approx0{,}049>0$ و $e^{-2{,}6}+2{,}6-e\approx-0{,}044<0$، إذن \bx{-2{,}7<\alpha<-2{,}6}.
\item على $[0;+\infty[$: $f$ مستمرة ومتزايدة تماما من $0$ إلى $+\infty$، إذن حل وحيد $\beta$. $e^{1{,}4}-1{,}4-e\approx-0{,}063<0$ و $e^{1{,}5}-1{,}5-e\approx0{,}263>0$، إذن \bx{1{,}4<\beta<1{,}5}.
\end{itemize}
\QQ{6} بعض القيم: $f(-3)\approx1{,}12$، $f(-2)\approx0{,}76$، $f(-1)\approx0{,}31$، $f(1)\approx0{,}54$، $f(2)\approx1{,}68$، $f(3)\approx2{,}84$.
\begin{GR}{-5}{4}{-2}{4}{1.2}
\Ticks{-4,-3,-2,-1,1,2,3}{-1,1,2,3}
\draw[cThm,thick,domain=-2:4.3] plot(\x,{\x});
\node[cThm,above left] at (3,3){$(\Delta)$};
\draw[cMeth,thick,dashed,domain=-1.5:4.3] plot(\x,{\x-1+ln(e-1)});
\node[cMeth,below right] at (3.4,1.94){$(T)$};
\draw[cDef,very thick,samples=250,domain=-5:4.3] plot(\x,{ln(exp(\x)-\x)});
\node[cDef] at (-4,2){$(C_f)$};
\draw[gray,densely dotted] (-5,1)--(4,1);
\fill (-2.6475,1) circle(1.5pt) node[above]{$\alpha$};
\fill (1.4204,1) circle(1.5pt) node[above left]{$\beta$};
\fill (1,{ln(e-1)}) circle(1.5pt);
\end{GR}
\textbf{المناقشة:} $e^x-x=e^m\iff\ln(e^x-x)=m\iff f(x)=m$. الحلول هي فواصل نقط تقاطع $(C_f)$ مع المستقيم الأفقي $y=m$:
\begin{itemize}
\item إذا كان $m<0$: لا توجد حلول.
\item إذا كان $m=0$: حل وحيد هو $x=0$.
\item إذا كان $m>0$: حلان متمايزان، أحدهما سالب تماما والآخر موجب تماما.
\end{itemize}

\textbf{الجزء الثالث:}
\QQ{1} من أجل $x\in I$: $\psi(x)=x\iff\ln(x+e)=x\iff x+e=e^x\iff e^x-x=e\iff f(x)=1$.

وبما أن $\beta\in]1{,}4;1{,}5[\subset I$ و $f(\beta)=1$ فإن $\psi(\beta)=\beta$.

\textbf{الوحدانية:} نضع $k(x)=\psi(x)-x$، $k'(x)=\dfrac1{x+e}-1<0$ على $I$ (لأن $x+e>1$)، فـ $k$ متناقصة تماما على $I$ وتنعدم مرة واحدة على الأكثر. إذن \bxt{$\beta$ هو الحل الوحيد للمعادلة $\psi(x)=x$ في $I$}.
\QQ{2} $\psi'(x)=\dfrac{1}{x+e}>0$، إذن $\psi$ متزايدة تماما على $I$، ومن أجل $1\leqslant x\leqslant2$: $\psi(1)\leqslant\psi(x)\leqslant\psi(2)$.
\par $\psi(1)=\ln(1+e)\approx1{,}31\geqslant1$ و $\psi(2)=\ln(2+e)\approx1{,}55\leqslant2$، إذن \bx{\psi(x)\in I}.
\QQ{3} $u(x)=\psi(x)-\dfrac{x}{e+1}$، $u'(x)=\dfrac{1}{x+e}-\dfrac{1}{e+1}$. من أجل $x\geqslant1$: $x+e\geqslant1+e$ إذن $\dfrac1{x+e}\leqslant\dfrac1{e+1}$، أي $u'(x)\leqslant0$: $u$ متناقصة على $I$.
\begin{itemize}
\item إذا كان $\beta\leqslant x\leqslant2$: $u(x)\leqslant u(\beta)$ أي $\psi(x)-\dfrac{x}{e+1}\leqslant\beta-\dfrac{\beta}{e+1}$، ومنه $\psi(x)-\beta\leqslant\dfrac{x-\beta}{e+1}$. وبما أن $\psi$ متزايدة: $\psi(x)\geqslant\psi(\beta)=\beta$، إذن $0\leqslant\psi(x)-\beta\leqslant\dfrac{x-\beta}{e+1}$.
\item إذا كان $1\leqslant x\leqslant\beta$: $u(x)\geqslant u(\beta)$ ومنه $\psi(x)-\beta\geqslant\dfrac{x-\beta}{e+1}$، أي $\beta-\psi(x)\leqslant\dfrac{\beta-x}{e+1}$، و $\beta-\psi(x)\geqslant0$ لأن $\psi(x)\leqslant\psi(\beta)$.
\end{itemize}
في الحالتين: \bx{\big|\psi(x)-\beta\big|\leqslant\dfrac{1}{e+1}\,|x-\beta|}.
\QQ{4} $\psi$ متزايدة تماما و $1{,}4<\beta<1{,}5$، إذن $\psi(1{,}4)<\psi(\beta)<\psi(1{,}5)$، و $\psi(\beta)=\beta$:
\[\ln(1{,}4+e)<\beta<\ln(1{,}5+e)\quad\Longrightarrow\quad\bx{1{,}415<\beta<1{,}440}\]
($\psi(1{,}4)\approx1{,}4154$ و $\psi(1{,}5)\approx1{,}4394$)، وسعة هذا الحصر $0{,}025<0{,}03$.

\textbf{ملاحظة:} يتوافق هذا مع نتيجة السؤال 3: $|x-\beta|<0{,}1$ من أجل $x\in[1{,}4;1{,}5]$ يعطي $|\psi(x)-\beta|<\dfrac{0{,}1}{e+1}\approx0{,}027$.
\end{hal}

\end{document}
